\documentclass[10pt]{article}
\usepackage[letterpaper,margin=1.0in,columnsep=0.28in]{geometry}
\usepackage{mathpazo}
\usepackage{fontspec}
\setmainfont{TeX Gyre Pagella}
\setmonofont[Scale=0.82]{DejaVu Sans Mono}
\newfontfamily\fbfont{FreeSerif}
\usepackage{newunicodechar}
\newfontfamily\amirifont{Amiri}[Script=Arabic]
\TeXXeTstate=1
\newunicodechar{ﷻ}{{\amirifont ﷻ}}
\newunicodechar{⟺}{{\fbfont ⟺}}
\newunicodechar{⟀}{{\fbfont ⟀}}
\newunicodechar{⬖}{{\fbfont ⬖}}
\newunicodechar{ⁿ}{{\fbfont ⁿ}}
\newunicodechar{₀}{{\fbfont ₀}}
\newunicodechar{ₙ}{{\fbfont ₙ}}
\usepackage{calc,array,seqsplit,tabularx}
\usepackage{fvextra}
\newcolumntype{L}{>{\raggedright\arraybackslash}X}
\newcolumntype{P}[1]{>{\raggedright\arraybackslash}p{#1}}
\usepackage{amsmath,amssymb}
\usepackage[table]{xcolor}
\usepackage{titlesec}\usepackage{fancyhdr}\usepackage{enumitem}
\usepackage{booktabs}
\usepackage{float}
\usepackage{longtable}
% longtable cannot run in twocolumn mode; shim it onto a one-column float.
% Load-bearing: pandoc emits longtable for every pipe table. Do not remove.
\makeatletter
\let\oldlt\longtable
\let\endoldlt\endlongtable
\def\longtable{\@ifnextchar[\longtable@i \longtable@ii}
\def\longtable@i[#1]{\begin{figure}[t]\onecolumn\begin{minipage}{0.5\textwidth}\smaller\oldlt[#1]}
\def\longtable@ii{\begin{figure}[t]\onecolumn\begin{minipage}{0.5\textwidth}\smaller\oldlt}
\def\endlongtable{\endoldlt\end{minipage}\twocolumn\end{figure}}
\makeatother
\usepackage{relsize}
\usepackage{graphicx}
\usepackage[hidelinks]{hyperref}
\setcounter{secnumdepth}{0}
\makeatletter\def\verbatim@font{\ttfamily\scriptsize}\makeatother
\newcommand{\jboxstart}{\begin{jbox}}
\newcommand{\jboxend}{\end{jbox}}
\newcommand{\restorelongtable}{\let\longtable\oldlt\let\endlongtable\endoldlt}
\raggedbottom
\setlength{\textfloatsep}{10pt plus 2pt minus 2pt}
\setlength{\dbltextfloatsep}{10pt plus 2pt minus 2pt}
\setlength{\floatsep}{8pt plus 2pt}
\setlength{\dblfloatsep}{8pt plus 2pt}
\renewcommand{\topfraction}{0.9}\renewcommand{\dbltopfraction}{0.9}\renewcommand{\textfraction}{0.07}\renewcommand{\floatpagefraction}{0.8}\renewcommand{\dblfloatpagefraction}{0.8}
\providecommand{\tightlist}{\setlength{\itemsep}{0pt}\setlength{\parskip}{0pt}}
\definecolor{accent}{HTML}{B87333}
\definecolor{deep}{HTML}{8A5523}\definecolor{faint}{HTML}{6E6E6E}
\titleformat{\section}{\normalfont\bfseries\scshape\large}{\color{accent}\thesection.}{0.5em}{}[{\color{accent}\titlerule[0.6pt]}]
\titleformat{\subsection}{\normalfont\bfseries}{\color{accent}\thesubsection}{0.5em}{}
\titlespacing*{\section}{0pt}{1.3ex}{0.6ex}
\pagestyle{fancy}\fancyhf{}
\renewcommand{\headrulewidth}{0pt}\renewcommand{\footrulewidth}{0pt}
\fancyfoot[R]{\color{accent}\itshape\small A Formal Proof of the Real
Part of the Riemann Hypothesis \textbar\ \thepage}
\fancyfoot[L]{\textcolor{accent}{\rule{0.32\columnwidth}{0.25pt}}}
\newenvironment{jbox}{\par\vspace{2pt}\noindent\color{accent}\rule{\columnwidth}{0.8pt}\par\vspace{-2pt}\small}{\par\vspace{2pt}}
\begin{document}
{%
\begin{center}
{\color{accent}\rule{0.6\textwidth}{0.8pt}}\\[10pt]
{\bfseries\scshape\fontsize{20}{23}\selectfont A Formal Proof of the
Real Part of the Riemann Hypothesis}\\[7pt]
{\itshape\large The Seat, the Address, and the Division of One Bit ·
Proved in Core Lean 4 from Existence Alone, Carried into ZFC Along the
RA--RAM Bridge, and Axiom-Free Under a Bare Computation Arrow}\\[8pt]
{\footnotesize\color{accent}\scshape\bfseries FORMAL VERIFICATION ·
FOUNDATIONS OF MATHEMATICS}\\[3pt]
{\itshape\color{deep}\small The seat proved from existence; the value
located at one address; the hypothesis divided, its Real part proved and
its Unicorn part set apart}\\[9pt]
{\color{accent}\rule{0.6\textwidth}{0.8pt}}\\[10pt]
{\scshape Mohammad F. Islam, PhD\(^{1}\)}\\[2pt]
{\itshape\color{faint}\small \(^{1}\) Independent Researcher, USA.
Correspondence: islamm@alumni.iu.edu}\\[2pt]
{\itshape\color{faint}\small 24 September 2026}\\[14pt]
\end{center}
\begin{minipage}{\textwidth}
{\color{accent}\rule{\textwidth}{0.6pt}}\\[4pt]
{\small\hspace{1em}\textbf{\scshape\color{deep}Abstract.\ }Riemann
stated his hypothesis as a fixed-set statement, that the roots of
\(\xi(t)\) are real. This paper proves, in core Lean 4 with no library
and no axiom declared, what existence decides about that statement, its
seat, and that its value lies beyond existence's reach, and then proves
the part of the value that is reachable. Part I proves the seat, with
the Root Axiom as the only posit: the fixed locus of conjugation on the
quaternions is the scalar line, the chart of the critical line embeds
onto it, and the three readings meet at one point; the value on the
zeros is one bit, unreachable from existence, since
\((\mathrm{RA} \to L) \leftrightarrow L\). Part II proves every reading
of that bit reached from existence, time, and monism one proposition
with the hypothesis at its apex, given the imported inputs, and shows
that the hypothesis posits no object while its denial posits one, a
finite witness never produced. Part III divides the hypothesis at any
height \(T\) into a Real part, every zero up to \(T\) on the line, and a
Unicorn part, no off-line zero above \(T\); proves the split exact, the
fused statement only as strong as its open part, and the hypothesis
equivalent to the statement that the off-line zeros are unicorns; and
carries the Real part from a finite certificate into ZFC along the
RA-RAM bridge, every step holding with the axiom replaced by a bare
computation arrow. So the Real part is proved, for \(\zeta\) to height
\(3 \times 10^{12}\) on the certificate of Platt and Trudgian. The
Unicorn part is named and set apart, not proved. Parts I and II are then
read at the Unicorn part: it is closed to derivation, from existence,
from any class of frames, and from any reading of an even register at
its seat, at theorem grade and in perpetuity; its aperture is located,
typed supply-only, one bit wide, and uncrossed; and the register is
silent on the supply side, both orientations equally refused to every
reading, which is why the title says the Real part and not the
hypothesis. Three bits close the arc. One bit divides: the hypothesis is
exactly its Real part and its Unicorn part, at every height, by theorem.
One bit is proved: the Real part, to that height, by certificate through
the schema. One bit is refused at theorem grade: the Unicorn part is
closed to every derivation from inside the register; its aperture is
structure, one bit wide and uncrossed; the supply side is silence, by
theorem at the register. With the three bits in hand the proof is
complete on the register side: everything a register can prove about the
hypothesis is proved, and the remainder is proved to be beyond
derivation, with its aperture typed and its supply side silent. The
hypothesis is not claimed proved; what is claimed is that no further
derivation is owed to it. Nothing is left to derive. The aperture is
structure. The supply side is silence.}\\[4pt]
{\color{accent}\rule{\textwidth}{0.6pt}}\\[3pt]
{\itshape\footnotesize\color{faint}\textbf{\upshape\scshape\color{deep}Keywords\ }\ Riemann
Hypothesis; Root Axiom; fixed locus; identity cone; Li criterion; halted
uniduction; Postulate M; denial asymmetry; Real part; Unicorn part; ZFC;
RA-RAM bridge; computation arrow; Lean 4; Fortran}\\[12pt]
\end{minipage}
}
\thispagestyle{fancy}
\let\oldunderscore\_\renewcommand{\_}{\oldunderscore\allowbreak}

\emergencystretch=2em

\newcommand{\partpage}[2]{\par\vspace{1.4em}\noindent{\color{accent}\rule{\columnwidth}{0.8pt}}\par\vspace{0.6em}{\centering{\color{accent}\large\scshape #1}\par\vspace{0.3em}{\Large\bfseries #2}\par}\vspace{0.5em}\noindent{\color{accent}\rule{\columnwidth}{0.8pt}}\par\vspace{1.0em}}

\hypertarget{introduction}{%
\section{Introduction}\label{introduction}}

Riemann stated his hypothesis as a fixed-set statement: the roots of
\(\xi(t)\) are real \cite{riemann1859}. A prior paper of this program
\cite{islamspine} restored that statement to
\(Z \subset \mathrm{Fix}(\tau)\) with \(\tau(s) = 1 - \bar{s}\),
defended the restoration against the literature, and proved in core Lean
4 the structural arc of the restored hypothesis on an abstract frame
\(X = (S, \tau, Z)\): the locus, the odd offset, the orientation
blindness of every even reading, the one-bit freedom, the halting
carrier, and the crossing. It also proved what the arc is not. Every
structural clause holds on a frame whose zeros lie off the line (its
Theorem 11); the only input that closes the crossing on the frame of
\(\xi\) is a term of \(L(X_\xi)\) itself (its Theorem 12); no theorem
whose hypotheses are only structural decides \(L\) (its Theorem 14). The
value, \(\forall \rho \in Z_\xi,\ \tau\rho = \rho\), was recorded as the
one open input, entered at a recorded row and never derived. We call
that row the witness row and the point at which a value enters a file
from outside it the aperture; both terms are the prior paper's and are
used here in its sense. Two objects carry the name Riemann Hypothesis in
what follows and are written with two tokens that are never joined:
\(\mathrm{RH}_{\mathrm{formal}}\), the seat, the restored statement read
as the fixed locus of its binding involution, which this paper proves
with existence as the only posit; and \(L(X_\xi)\), the value, the
statement that every zero lies on that locus, which this paper proves
existence does not reach. Part I names the first and excepts the second,
and the exception is Theorem I. ``From existence alone'' has one meaning
in this paper: existence is the only posit, and nothing else is assumed.
It does not mean that the seat is derived from existence. The seat holds
by the definition of the Ground and consumes no premise
(\texttt{recursion\_\allowbreak is\_\allowbreak constant}); existence is
consumed where it has content, at presence (Theorem D), at the closure
of the interaction reading (RAF-C1), and at the ledger, where every
adjudication of any proposition is a deed and therefore an instance of
existence and of nothing else
(\texttt{adjudicating\_\allowbreak anything\_\allowbreak instances\_\allowbreak existence}).
That last theorem is what existence contributes that no other true
premise does.

\textbf{Notation and the objects named.} The Root Axiom, written
\(\mathrm{RA}\), is the posit \(\forall x \in U,\ \Delta E(x) > 0\) on a
universe \(U\) with a differential \(\Delta E : U \to \mathbb{Z}\): to
exist is to actuate. Existence and \(\mathrm{RA}\) are one proposition;
the paper writes \(\mathrm{RA}\) in formulas and existence in prose.
\(\mathrm{RA}\) has two published formal readings \cite{islamraf}: the
interaction reading, RAF, which reads the axiom as the criterion for
systems being related, registration carrying positive mutual information
and an interaction domain being closed under it; and the grounding
reading, RAM, which reads the axiom as the criterion for a proposition
being determinate, its imprint in the fixed locus of a binding
involution. A frame is \(X = (S, \tau, Z)\), a carrier \(S\) with an
involution \(\tau\) and a distinguished subset \(Z\); its line property
is \(L(X) :\Leftrightarrow \forall s \in Z,\ \tau s = s\). On the frame
of \(\xi\), with \(Z = Z_\xi\) the zeros and \(\tau(s) = 1 - \bar{s}\),
\(L(X_\xi)\) is the restored hypothesis, and it is called the value. The
carrier the paper builds is the integer quaternions
\(\mathbb{H}_{\mathbb{Z}}\) with conjugation \(\sigma\) as its binding
involution; the seat is the fixed locus \(\mathrm{Fix}(\sigma)\), a seat
point is a point of it, and \(\mathrm{RH}_{\mathrm{formal}}\) is the
seat read as a proposition, that every point of \(\mathrm{Fix}(\sigma)\)
is fixed by \(\sigma\). A register is one of three readings of the seat,
the axiom register \(\mathrm{RA}\), the formal register
\(\mathrm{RAM}\), and the object register \(\mathrm{RH}\), and the three
form the register diagram. The witness row is the recorded row at which
the value is entered; the aperture is the point at which a value enters
a file from outside it; the world row is the witness row read in the
world register; the sign is the value
\(\forall \rho \in Z_\xi,\ \tau\rho = \rho\), the one bit. \(\Delta M\)
counts the mathematics authored and is zero throughout. Every other
symbol is defined where it first appears.

Two things in that record are external to its file, and this paper
brings them inside. First, the prior proof's live face takes its
``witnessed'' flag from the command line: that a supplier is present is
asserted from outside, and the kernel is told rather than shown that a
reader stands at the aperture. Second, three separations are left
standing: between the axiom under which a reader acts, the formal ground
on which the locus was fixed, and the Riemann object read as the fixed
locus of its binding involution. These are separations of description,
produced by reading one seat in three registers, and not gaps in the
mathematics. The prior proof could not close them because an axiom is
not a structural clause and the prior proof had none.

The Root Axiom is that axiom. Formally it is a universe \(U\) with a
differential \(\Delta E : U \to \mathbb{Z}\) and the statement
\(\forall x \in U,\ \Delta E(x) > 0\); its published form, its physical
warrant, and its two formal readings are in \cite{islamcodex, islamraf}.
The interaction reading (RAF in \cite{islamraf}) says what it is for
systems to be related: an interaction domain is closed under
registration, and no agent acts on a domain from outside it. The
grounding reading (RAM in \cite{islamraf}) says what it is for a
proposition to be determinate: its formal being is its imprint in the
fixed locus of the binding involution, read across the aperture, with
\(\mathrm{Fix}(\sigma) = \mathbb{R}\) as the seat. The interaction
reading reaches no mathematical object, since mathematical objects do
not change state; every proposition routes to the grounding reading
\cite{islamraf, islamfae}. That reach law is the axiom's own, and this
paper proves its Riemann instance rather than assuming it.

Two objects are new here. The first is the constructed witness: the
axiom built in the kernel as a structure in Prop with four fields, the
only proof object the closure is allowed to consume, and the eliminator
that runs from it alone,
\(\mathrm{RA} \to \mathrm{RA} \to \mathrm{RAM} \to \mathrm{RH}_{\mathrm{formal}}\):
the witness's last field, a function from the orientation and the formal
Ground to the formal self, applied once, with every leg of the chain a
definitional equality. The second is the category-gap eliminator itself:
the three separations exhibited as the three legs of a cone over the
register diagram, each leg an identity, so that the separations close at
zero cost and the closing is computed rather than asserted. Everything
else the paper uses is published: the frame and the structural arc
\cite{islamspine}, the one-bit result \cite{islamonebit}, the odd-supply
separation and the one-cut hypothesis across twenty-three rows
\cite{islamrows}, the formal-alone theory of registration, price, and
crossing \cite{islamftoe}, and the axiom's declaration \cite{islamraf}.

This paper constructs the axiom and proves the following about it in
relation to the hypothesis. The seat: conjugation on the quaternion
carrier fixes exactly the scalar line, the Return of the parse triad
lands on it, and the axiom's seat point, the formal seat point, and the
named seat point are the three legs of an identity cone over the
register diagram, each leg \texttt{rfl}, the apex unique per orientation
of the triad (Theorems A, B, C). The embedding: the prior proof's
lattice stage embeds equivariantly and injectively into the carrier,
\(\sigma\varphi = \varphi\tau\), and \(\varphi\) carries the line onto
the scalar line at every resolution (Theorem C\('\)). The aperture: for
a present reader who supplies a bit the row is never interior; it is
sealed or refused by the bit alone, and a reader who registers a bit is
inside the axiom's domain, never an exterior witness (Theorem D). The
exact strength of the rule: in every context in which the axiom holds,
\((\mathrm{RA} \to L(X)) \leftrightarrow L(X)\), and the axiom decides
\(L\) on no frame, the two-point frame carrying the axiom and failing
\(L\) (Theorems E, F). Every theorem is checked by the Lean kernel;
every finite claim is re-executed by a compiled Fortran program that
shares no code with the kernel, and both records are printed whole.

Part I names the register and the residue. At the axiom a closure can
carry no mass: a closure between the axiom and its formal reading that
added a premise would be a second axiom, which the axiom's own
root-premise theorem forbids \cite{islamcodex}. The eliminator that
closes the category gaps is therefore massless by necessity and not by
weakness, and this paper computes the masslessness rather than asserting
it: the three legs of the cone are definitional identities. ``In toto at
the apex'' has one definition and it is used in no other sense: every
clause of the hypothesis that is not its value on the zeros. ``The
sign'' names one object throughout: the value
\(\forall \rho \in Z_\xi,\ \tau\rho = \rho\), the orientation of assent
on the world row, the one bit the prior proof records at the witness
row. ``Everything but the sign'' is the same set named by its
complement. What the kernel returns, given the axiom, is exactly that
set, and the sign is proven to be exactly what the axiom does not reach.

\hypertarget{the-value-its-readings-and-its-address}{%
\subsection{The value, its readings, and its
address}\label{the-value-its-readings-and-its-address}}

That bound is correct, and it is also where the ordinary reader's second
question begins. If one bit is owed, what object would carry it? The
candidates proposed in the literature, an Euler-product class, a
self-adjoint operator, a positivity criterion, a monotone flow, are
different-looking objects. A reader who pursues them in turn meets what
looks like a regress: each candidate, once formalized, turns out to need
a further ingredient, and each further ingredient looks like a new
missing piece. The barrier this paper addresses is that appearance of
regress, and it is structural rather than computational.

The barrier has three layers. The first is register confusion. A
timeless formal reading of the hypothesis deletes the index along which
the zeros are counted, and a reading that deletes its own index cannot
distinguish a claim from its negation. The kernel exhibits this as
sign-blindness; Part I's \texttt{recursion\_is\_sign\_blind} is its
receipt. A verdict issued from that register is a verdict of an
amputated instrument, not a finding about the zeros.

The second layer is type mismatch. The natural ``arrow'' a physical
reading supplies is a sign, one Boolean. The hypothesis asks for a
vanishing: that the off-line offset \(\delta\) of every zero is zero. A
sign does not force a vanishing. A supplied orientation picks which
member of a mirror pair \((\tfrac12 + \delta, \tfrac12 - \delta)\) comes
first and leaves \(\delta\) untouched. The barrier dissolves only when
the vanishing is rewritten as a family of signs, one per step of an
index, and Li's criterion {[}Li 1997{]} is exactly that rewriting.

The third layer is the uniformity trap. Every instrument that closes a
category gap at the apex is massless by necessity, because a closure
between existence and its reading that authored mathematics would be a
derivation of the root from below {[}Islam 2026a{]}. A massless
instrument chooses its objects without reading which L-function it
serves. Such an instrument is uniform across the family, and a uniform
instrument forces positivity on every member, including members whose
zeros leave the line. The Eisenstein L-function \(\zeta(s)\zeta(s-k+1)\)
has an Euler product and a functional equation and zeros off its center
line {[}Iwaniec and Kowalski 2004{]}. It is not primitive; after the
shift that centers its functional equation at \(\tfrac12\), its zeros
lie off that center, and Bombieri-Lagarias then gives a negative
coefficient. In the kernel the control enters abstractly, as a family
member with a negative coefficient, and it kills every uniform
instrument. The bit must therefore be carried by an object that reads
\(\zeta\) specifically.

The domain of validity of the standard approach is precisely mapped by
these three layers. Classical analytic number theory works inside one
register, the timeless one, and it has reached the classical
consequences of positivity on the Euler product: nonvanishing for
\(\mathrm{Re}\, s > 1\), and through Mertens' positivity, nonvanishing
on \(\mathrm{Re}\, s = 1\) {[}Hadamard 1896; de la Vallée Poussin 1896;
Mertens 1898{]}. Beyond that line the Euler product diverges and
positive energies have nothing to act on. The strip
\(\tfrac12 < \mathrm{Re}\, s < 1\) is where the timeless register stops
and where a reading indexed by time must take over.

Part II does not cross that strip; Part III crosses it only where
computation has certified the zeros, up to a finite height. What Part II
does is prove that every road into the strip reached from existence,
time, and monism is the same road. The regress is shown to be an
artifact of reading one seat in several registers, and the category-gap
eliminator of Part I, applied to those registers, collapses it to one
gap. The gap is then pinned, typed, and left at the witness row where
Part I placed it.

\hypertarget{the-division-of-the-value}{%
\subsection{The division of the value}\label{the-division-of-the-value}}

Part III divides the hypothesis. It does not reopen the value that Parts
I and II locate, and nothing in it contradicts Theorem I. One part is
reachable, and it is proved here along a road built from the Root Axiom
through the RA--RAM bridge, the RAM strata, and ZFC. The other part is
the region where the denial's object would have to live, and it is named
for what it is and set apart.

\textbf{Three bits, and what is claimed.} The whole paper is three bits.
One bit divides: the hypothesis is exactly its Real part and its Unicorn
part, at every height, by theorem (III.11). One bit is proved: the Real
part, to \(3 \times 10^{12}\), by certificate through the schema (III.34
to III.36). One bit is refused at theorem grade: the Unicorn part is
closed to derivation from existence, from any class of frames, and from
any reading of an even register (III.45 to III.48); its aperture is
located, typed supply-only, one bit wide, and uncrossed (III.46); and
the register is silent on the supply side, both orientations equally
refused (III.49). The Unicorn part is unproved, and the reason is
neither a want of effort nor a want of ability: it is a theorem, and the
theorem says which routes are closed, which is all of them but the
aperture, and says nothing of the aperture but its width. With the three
bits in hand the proof is complete on the register side. Everything a
register can prove about the hypothesis is proved; the remainder is
proved to be beyond derivation, its aperture typed, its supply side
silent. The hypothesis is not claimed proved. What is claimed is
stronger than an open problem and weaker than a theorem: no further
derivation is owed to the Riemann Hypothesis. The three bits are one
term in the kernel (\texttt{rh\_register\_proof\_complete}). Nothing is
left to derive. The aperture is structure. The supply side is silence.

\hypertarget{related-work}{%
\section{Related Work}\label{related-work}}

Only the positions that bear on the apex are reviewed; the analytic
literature of the restored formulation is reviewed in \cite{islamspine}.

\textbf{The original statement and its chart.} Riemann
\cite{riemann1859}, Siegel \cite{siegel1932}, and Edwards
\cite{edwards1974} give the fixed-line statement and its history;
Titchmarsh \cite{titchmarsh1986} and Conrey \cite{conrey2003} state the
canonical chart form, real part one half. The prior paper's Proposition
V1 makes the address a chart and the fixed line intrinsic. This paper
takes the fixed line one register deeper: the fixed line of a binding
involution is the seat on every carrier that has one, and the two
carriers in play, the lattice stage under \(\tau\) and the quaternions
under \(\sigma\), are joined by an explicit equivariant map.

\textbf{Symmetry without the Euler product.} Davenport and Heilbronn
\cite{davenport1936} and Bombieri and Hejhal \cite{bombieri1995} exhibit
series with the functional equation and zeros off the line. In the prior
proof these are the off-line frames on which every structural clause
holds. Here they carry one thing more: the Root Axiom holds on them too
(Theorem F). That is the whole reason the axiom cannot decide the value.

\textbf{The Euler product as key.} Selberg \cite{selberg1992} and Conrey
and Ghosh \cite{conrey1993} name the Euler product as the organizing
axiom of the class for which the hypothesis is conjectured; Weil
\cite{weil1948} and Deligne \cite{deligne1974} prove the analogue where
the Euler product runs over a finite geometry. Part I leaves the key
where the prior proof left it, on the world row, and shows that the apex
closes without it; Part III turns the key by computation up to a
certified height, and no further.

\textbf{The spectral road.} Berry and Keating \cite{berry1999} and
Connes \cite{connes1999} seek an operator whose spectrum is the zero
set. An operator is a world-row construction; the seat on which its
spectrum would have to lie is fixed before any operator is built, and
that ordering, orientation after construction, is a theorem of the
twenty-three-row study \cite{islamrows}.

\textbf{The program's own prior results.} The one-bit result
\cite{islamonebit} located the hypothesis at exactly one bit from its
structural arc. The twenty-three-row study \cite{islamrows} typed every
proved method-class barrier as an evenness statement and every open row
as missing exactly one orientation bit at its decision seat, and stated
the one-cut hypothesis: that each such row is closed by one cut, the
supply of that bit. The formal-alone theory \cite{islamftoe} fixed
registration, price, crossing, and cost as the four faces of a supplied
bit. The earlier termination and case-closed papers
\cite{islamterm, islamclosed} recorded the hypothesis's standing under
the program's cascade before the restored form was proved. The present
paper is the apex of that sequence: Part I places the Riemann row's owed
bit against the one seat under every name and proves that the axiom
reaches everything but that bit, and Part III divides that bit and
proves its Real part.

\textbf{The axiom's formal readings.} The declaration of the formal Root
Axiom \cite{islamraf} states RAF-1 and RAF-2, the bridge to Landauer,
and the consequences C1 to C3, and states RAM as the imprint law with
\(\mathrm{Fix}(\sigma) = \mathbb{R}\). Its engine records the constraint
this paper obeys: the interaction reading reaches no mathematical object
and routes every proposition to the grounding reading \cite{islamfae}.
Cantor's diagonal theorem, which Lawvere's fixed-point theorem
generalizes \cite{lawvere1969}, carries C2 unconditionally.

\textbf{Category-theoretic frame.} Cones, apexes, and limits over a
diagram are used in their elementary form, a discrete diagram on three
objects in the groupoid of definitional equalities \cite{maclane1998}.
The frame is chosen because a category gap is, in that setting, exactly
a missing leg, and its elimination is exactly an identity leg.

\textbf{The price of the act.} Landauer \cite{landauer1961} bounds the
dissipation of registering one bit; Bérut and coauthors \cite{berut2012}
measure it. The prior proof prices the spent bit at \(k_B T \ln 2\);
here the price is the bridge of the interaction reading and it exports
without constituting.

\textbf{The gap.} No work in this literature states the axiom under
which the supplier acts, so the presence of a supplier remains a flag;
and none treats the three separations between axiom, formal ground, and
object as what they are, separations of description rather than of
mathematics. Both are supplied here.

\hypertarget{the-readings-of-the-value}{%
\subsection{The readings of the value}\label{the-readings-of-the-value}}

Li {[}1997{]} proved that the hypothesis holds iff \(\lambda_n \geq 0\)
for every \(n \geq 1\), where
\(\lambda_n = \sum_\rho [1 - (1 - 1/\rho)^n]\). Bombieri and Lagarias
{[}1999{]} generalized the criterion to multisets of zeros and showed
that an off-line zero drives some \(\lambda_n\) negative. Their theorem
is the analytic content behind this paper's reading of Li's coefficients
as modes of a time evolution.

Weil {[}1952{]} gave the explicit formula and the positivity criterion
that bears his name; in it the primes enter with a negative sign against
an archimedean term. Deligne {[}1974{]} proved the function-field
analogue, where the decisive ingredient is a positivity supplied by
geometry. Selberg {[}1992{]} and Conrey and Ghosh {[}1993{]} name the
Euler product with primitivity and the Ramanujan bound as the organizing
axioms of the class for which the hypothesis is conjectured.

De Bruijn {[}1950{]} introduced the heat flow
\(H_t(z) = \int e^{tu^2} \Phi(u) \cos(zu)\, du\) and proved the
strip-narrowing theorem. Newman {[}1976{]} defined the constant
\(\Lambda\) and conjectured \(\Lambda \geq 0\). Rodgers and Tao
{[}2020{]} proved it. The Polymath project {[}Polymath 2019{]} proved
\(\Lambda \leq 0.22\). The hypothesis is therefore equivalent to
\(\Lambda = 0\).

Beyond the spectral road above, Stone {[}1932{]} proved that a strongly
continuous one-parameter unitary group has a self-adjoint generator; the
section on the witness rebuilt uses it to pass from a unitary evolution
to its generator. Julia {[}1990{]} and Bost and Connes {[}1995{]}
realized \(\zeta\) as a partition function over the primes.

Perelman {[}2002; 2003a; 2003b{]} proved the Poincaré conjecture through
Ricci flow, with the monotonicity of the \(\mathcal W\)-entropy as the
transport that carries control along the flow. De Branges proposed a
positivity condition intended to prove the hypothesis; Conrey and Li
{[}2000{]} showed it fails.

Littlewood {[}1914{]} and Skewes {[}1933{]}, Haselgrove {[}1958{]}, and
Odlyzko and te Riele {[}1985{]} are the standing warnings against
inference from finite confirmation.

The approaches reviewed here are each developed around their own object,
and each is read in this paper as one register's version of the missing
piece. the section on halted uniduction proves that these objects, in
the forms reached here, are one proposition.

\hypertarget{methods}{%
\section{Methods}\label{methods}}

\hypertarget{two-independent-checkers}{%
\subsection{Two independent checkers}\label{two-independent-checkers}}

Every theorem of this paper is checked by core Lean 4, version 4.19.0,
no library, no axiom declared by the file, no \texttt{sorry}; the
kernel's own \texttt{propext} and \texttt{Quot.sound} enter only through
integer arithmetic and are listed per theorem, every dependency set
printed by the kernel and reproduced in Appendix C. The file,
\texttt{RH\_At\_The\_Apex.lean}, is printed whole in Appendix B with its
SHA-256 digest, and it is self-contained: the prior proof's Theorems 1,
4 through 9, 12, and 15 are re-proved inside it on the same abstract
frame and the same lattice stage, so that the apex layer rests on
nothing it does not carry.

The finite content of the file, itemized below, is re-executed by a
second checker that shares no code, no logic, and no toolchain with the
first: a Fortran 2008 program, \texttt{RH\_At\_The\_Apex\_Twin.f90},
compiled with gfortran, printed whole in Appendix D with its run log in
Appendix E. The twin enumerates rather than proves. It checks
conjugation against every integer quaternion in a lattice ball, the
Return and its six orderings, the three identities and the cone census
over the whole ball, the embedding's equivariance, image clause, and
injectivity at seven resolutions on a lattice window, the wall over
every one of the \(2^{13}\) readouts of the seat's window, the two
wholly odd maps and their calibration, the aperture on a twelve-point
domain, the truth table of the rule and the two-point counter-model, the
ledger, the refusal constant, the price, the closure law on every one of
the \(2^{16}\) subsets of a sixteen-point carrier, Cantor's diagonal on
every map \(S \to 2^S\) for \(|S| = 3, 4\), the cure theorem on all
thirty-two three-point frames and all \(256\) classes of two-point
frames, and the constancy of the recursion field, the carrier by
decision, and the three measures. It carries an oracle that stops the
program on any failure and prints a census last; a binary that reaches
its final line has passed. The two checkers were written to the same
specification and neither reads the other. Where they disagree the
artifact is wrong, and the disagreement is printed; where they agree the
agreement is between a typechecker and an executor, not between two
copies of one reading.

The last battery of the twin implements the prior proof's rule that a
self-check is not a witness: its live row opens only on an argument the
operator supplies on the command line, and the program prints that it
did not and cannot generate its own witness. The reader who runs the
binary is that operator.

\hypertarget{objects-and-terms}{%
\subsection{Objects and terms}\label{objects-and-terms}}

The seat is built on the integer quaternions rather than on a one-point
type, so that ``the fixed locus of the binding involution'' is a
computed object and not a stipulation: conjugation is proved binding,
its fixed locus is proved to be the scalar line for every quaternion,
and the Return is computed. The prior proof's frame \(X = (S, \tau, Z)\)
is carried abstractly, exactly as it states it, and its lattice stage is
carried as \((\mathbb{Z}^2, \tau)\) with the explicit map \(\varphi\)
into the carrier. The zero set of \(\xi\) enters only as a parameter,
since \(\zeta\) is not definable in core Lean.

Three objects share the word bridge in the sources and are kept apart
here by three names: the Bridge capacity of the witness, written with a
capital; the halting carrier of the prior proof's Theorem 7, written
bridge; and the equivariant map \(\varphi\), written embedding. The
category-theoretic frame is elementary and stated once. The three
registers form a discrete diagram
\(D : \{\mathrm{RA}, \mathrm{RAM}, \mathrm{RH}\} \to \mathbb{H}_{\mathbb{Z}}\),
the axiom register, the formal register, and the object register; a cone
over \(D\) with apex \(a\) is a family of arrows \(a \to D(r)\) in the
groupoid of identities, that is, a family of equalities \(a = D(r)\);
the limit exists iff the three readings coincide, and then the apex is
unique. A category gap is the absence of a leg. Nothing beyond this is
used.

\hypertarget{grades-and-registers}{%
\subsection{Grades and registers}\label{grades-and-registers}}

Three grades of warrant are used and every claim-bearing box and table
carries one: kernel theorem, checked by the Lean kernel; structural, a
statement about the arrangement of theorems that is itself checkable but
not a theorem of the file; premise, a statement carried by citation and
held fixed. Two registers are used and kept apart by theorem: the apex
register, where the seat and the category gaps live, and the world
register, where the value on \(Z_\xi\) lives; the world row is the
witness row read in the world register. Falsifiable criteria are capped
at three and each is typed for necessity. The author's own verification
discipline is applied only in the disclosure at the end and draws no
warrant of its own.

\hypertarget{methods-for-the-readings-of-the-value}{%
\subsection{Methods for the readings of the
value}\label{methods-for-the-readings-of-the-value}}

\textbf{The file.} \texttt{RA\_Li\_Bridge.lean} is printed whole in
Appendix F with its SHA-256 digest, and its dependency sets verbatim in
Appendix G: eighty-three sets printed by the file, sixty-one axiom-free,
twenty on \texttt{propext} and \texttt{Quot.sound} only, two on
\texttt{propext} alone, and seventeen more printed by a supplementary
run, nine axiom-free, so that every theorem of the file carries a
printed set.

\textbf{What enters as a hypothesis.} The zeta function is not definable
in core Lean, so every analytic fact enters as a hypothesis carried by
citation: Li's criterion as \texttt{LiCriterion}, the equivalence of the
hypothesis with \(\Lambda = 0\) as \texttt{hLam}, the identification of
the hypothesis with the line property of the enumerated zeros as
\texttt{hDef}, and the existence of a negative Li coefficient in the
Eisenstein member as a hypothesis of \texttt{no\_uniform\_bridge}. The
Root Axiom enters as the hypothesis \texttt{RA\ S} on a substrate, never
as an axiom. The kernel therefore shows exactly what each theorem
assumes.

\textbf{What the kernel proves, and what it imports.} Every result below
carries one of three grades, and the paper never promotes one to
another.

\begin{table}[!htbp]
\centering\small
\renewcommand{\arraystretch}{1.12}
\raggedright \textbf{Grades of warrant.} What the kernel proves, what it imports, and how Part II reads them.\par\smallskip
\begin{tabularx}{\textwidth}{@{}P{0.13\textwidth} P{0.22\textwidth} L@{}}
\toprule
\textbf{Grade} & \textbf{Content} & \textbf{Examples in this paper} \\
\midrule
K, kernel & Checked by the Lean kernel from definitions alone & the identity (6.2a) on the lattice stage; stability iff the line for fold-invariant sets (Theorem 15); the bridge's three properties given its inputs; the cone and its legs given their hypotheses; Postulate M with its seed iff the line property of an abstract staged set \\
A, analytic import & A theorem of the literature entering as a Lean hypothesis, never proved here & Li's criterion (\texttt{LiCriterion}); Bombieri-Lagarias, that an off-line zero forces a negative coefficient; Newman with Rodgers-Tao, \(\mathrm{RH} \iff \Lambda = 0\) (\texttt{hLam}); the identification of the hypothesis with the line property of \(\zeta\)'s zeros (\texttt{hDef}); the negative coefficient of the Eisenstein control; the arithmetic form of RH {[}Lagarias 2002{]} and \(\Sigma^0_1\)-completeness of arithmetic {[}Kaye 1991{]}, carried in Part XV as the named hypothesis \texttt{sigma1} with soundness on the denial (the section on the logical form); the seed \(L(0)\) of Postulate M, that every zero located by stage zero lies on the line, supplied for \(\zeta\) by computation \\
C, structural reading & An interpretation of how the kernel theorems stand together & ``one address'', ``the regress halts'', ``one cone'', the Perelman analogy, the Hilbert-Pólya compatibility \\
\bottomrule
\end{tabularx}
\end{table}

\texttt{LiData.nonneg\ n} is an abstract proposition standing for
\(\lambda_n \geq 0\). No zeros, sums, or numerical coefficients are
formalized; the link to the analytic \(\lambda_n\) is the imported Li
criterion.

\textbf{What is modelled.} Five parts of the file work on toy carriers
and say so in their headers. Part IV places zeros at integer half-unit
coordinates \((h, t)\) on the Bridge plane of the prior proof, with
\(\rho = h/2 + it\). Part V builds its mixed zero set (Theorem 20) on
the same plane, and Part XI builds its timed worlds there. Part XII
builds its self-grounding and frame models on small carriers. Part XI's
Postulate M theorems and Part XII's collapse theorems are stated over
abstract types. Part IX models the de Bruijn strip bound on
natural-number widths, a toy of the bound and not of the flow. Every
other part is stated over abstract types.

The verdict tokens of the discipline appear only in the disclosure.

\hypertarget{methods-for-the-division}{%
\subsection{Methods for the division}\label{methods-for-the-division}}

Every theorem of Part III is a kernel theorem in core Lean 4.19.0, with
no library, no \texttt{sorry}, and no axiom declared; the nine files are
printed in Appendix H with their receipts, and the theorems are numbered
III.1 to III.44. The facts Part III cites are listed in the standard of
proof below and marked where they enter.

\hypertarget{the-standard-of-proof}{%
\subsection{The standard of proof}\label{the-standard-of-proof}}

The standard is the ordinary one for a theorem of mathematics: every
result is a deduction from the axioms of Zermelo--Fraenkel set theory
with choice and the standard definitions of its objects, or it is
marked, where it enters, as one of four other things, and nothing else
is assumed.

\textbf{Kernel theorems.} Every theorem numbered in this paper is
checked by the Lean 4.19.0 kernel from its definitions, with no library,
no \texttt{sorry}, and no user-declared axiom in any file. Each file
prints its dependency sets: every theorem depends on nothing, or on the
kernel's principles of propositional extensionality and quotient
soundness, and none on the axiom of choice. The type theory of Lean, as
analyzed in its Lean 3 formulation, which the Lean 4 kernel shares in
the respects used here, is consistent relative to ZFC with countably
many inaccessible cardinals {[}Carneiro 2019{]}; the statements proved
here concern propositions, finite structures, and small types, of the
kind whose translation into ZFC is standard, and that translation is not
mechanized here.

\textbf{Cited theorems of mathematics.} Results of the literature that
are theorems of ZFC enter as named hypotheses and are never presented as
kernel results: Li's criterion {[}Li 1997{]}; the growth of a negative
coefficient from an off-line zero {[}Bombieri and Lagarias 1999{]}; the
equivalence of the hypothesis with \(\Lambda = 0\) {[}de Bruijn 1950;
Newman 1976; Rodgers and Tao 2020{]}; the arithmetic form of the
hypothesis {[}Lagarias 2002{]}; the absoluteness of arithmetic sentences
{[}Kunen 1980{]}; completeness for true bounded sentences {[}Kaye
1991{]}; the independence of the axiom of choice {[}Gödel 1938; Cohen
1963{]}; and the certificate for the zeros of \(\zeta\) to height
\(3 \times 10^{12}\) with the analytic lemmas it rests on {[}Turing
1953; Platt and Trudgian 2021{]}.

\textbf{Definitional identifications and computed data.} The
identification of the hypothesis with the line property of the zeros of
\(\zeta\), and the seed of Postulate M, the zeros located by stage zero,
supplied by computation.

\textbf{One metatheoretic premise.} The soundness of the foundation,
which no consistent foundation proves of itself. It is used in exactly
three places, the converse direction of Theorem 75 and the placement
Theorems III.15 and III.19, carried as a named premise in each, and
Theorem III.16 shows it cannot be dropped.

\textbf{Readings.} The Root Axiom, the strata of the grounding reading,
the vocabulary of registers, and the physical price of Part I's three
measures, carried on the cited Landauer bound and its measurement
{[}Landauer 1961; Bérut et al.~2012{]}, are readings of the deductions,
not premises of them. In Parts I and II the axiom is a hypothesis or a
constructed structure, never a declared axiom, and Theorem E fixes its
exact strength. In Part III every deduction holds with the axiom
replaced by the identity arrow on a type, with no axiom of any kind
(Theorem III.44). The deductive content of Part III therefore uses
nothing beyond ZFC and the cited theorems.

\textbf{The title, read at this standard.} ``Proved in core Lean 4''
names the checker of every numbered theorem, the certificate schema
III.34 and the split among them; ``from existence alone'' names the
posit inventory, existence being the only posit and never a premise
consumed, since the seat is massless by construction (Theorem G) and the
road holds under a bare arrow (Theorem III.44); ``carried into ZFC''
names the placement of III.15 to III.30. For \(\zeta\) the Real part's
standing is the cited certificate applied through the kernel schema,
grade A joined to grade K, and the kernel does not check the
certificate's numerics. Two readings are excluded by the text itself:
that the seat is derived from existence, and that the kernel has
verified the zeros of \(\zeta\).

\textbf{What is proved in this standard.} Every kernel theorem, and for
\(\zeta\) the Real part to height \(3 \times 10^{12}\): the certificate
theorem of the kernel (III.34) applied to the cited certificate. What is
not proved: the Unicorn part, and so the hypothesis itself, whose fused
statement stands open (Theorem III.2).

\textbf{The paper's standing without the certificate.} The certificate
of Platt and Trudgian is a parameter of the schema and a premise of no
theorem. Were that result absent from the literature, exactly one
sentence of this paper would fall, that the Real part of \(\zeta\) is
proved to height \(3 \times 10^{12}\), and with it the Real part of
\(\zeta\) would stand proved at no height, the paper's own count to
\(T = 100\) being testimony. Every numbered theorem of Parts I, II, and
III would stand unchanged, none of them mentioning a zero of \(\zeta\):
the seat and its bound, the address and the logical form given their
named imports, the division at any height, the fusion law, the unicorn
identity, the certificate schema, the placement, the bridge, the round
trip, and the road under a bare arrow. Three standings are used, and
each row of the table below carries one: K, a kernel theorem with no
user-declared axiom; K given A, a kernel theorem whose hypotheses are
cited theorems of the literature, none of them the certificate; and A,
the cited certificate itself, joined to the kernel schema it populates.

\begin{table}[!htbp]
\centering\small
\renewcommand{\arraystretch}{1.12}
\raggedright \textbf{The paper's standing without the certificate.} Each component, its standing, and whether it depends on the certified height.\par\smallskip
\begin{tabularx}{\textwidth}{@{}L P{0.30\textwidth} P{0.13\textwidth}@{}}
\toprule
\textbf{Component} & \textbf{Standing} & \textbf{Depends on $3 \times 10^{12}$?} \\
\midrule
The seat, its bound, the aperture, and the assignments of the three measures (Part I, Theorems A to I) & K, kernel theorems, no user-declared axiom; the physical reading of the price is a reading & No \\
The address and the cone, halted uniduction (Part II, Theorems 44 to 50 and 68) & K given A: Li's criterion, the $\Lambda = 0$ equivalence, the definitional identification, the seed & No \\
The logical form and the asymmetry of the denial (Theorems 70 to 77) & K given $\Sigma^0_1$-completeness, and soundness in the converse only & No \\
The division at any height, the fusion law, the unicorn identity, the split's separation and monotonicity (III.1 to III.14) & K & No \\
The certificate schema, the placement, the bridge, the round trip, the arrow (III.15 to III.44) & K & No \\
The Real part of $\zeta$ to height $3 \times 10^{12}$ & A joined to K, the cited certificate through the schema & Yes, and only this \\
The formal block of the Unicorn part, the aperture, the silence, and the three bits as one term, III.45--III.49 and \texttt{rh\_\allowbreak register\_\allowbreak proof\_\allowbreak complete} & K, kernel theorems with no axiom; on $Z_\xi$ at the grade of Theorem D and the kernel-cannot-read-the-deed theorem & No \\
\bottomrule
\end{tabularx}
\end{table}

\par\vspace{1.4em}\noindent{\color{accent}\rule{\columnwidth}{0.8pt}}\par\vspace{0.6em}{\centering{\color{accent}\large\scshape Part I}\par\vspace{0.3em}{\Large\bfseries The Seat}\par}\vspace{0.5em}\noindent{\color{accent}\rule{\columnwidth}{0.8pt}}\par\vspace{1.0em}

\hypertarget{the-prior-proof-carried}{%
\section{The Prior Proof, Carried}\label{the-prior-proof-carried}}

A frame is \(X = (S, \tau, Z)\) with \(\tau : S \to S\) and
\(Z \subset S\), and the line property is
\(L(X) :\Leftrightarrow \forall s \in Z,\ \tau s = s\). On the frame of
\(\xi\), \(L\) is the restored hypothesis. The prior proof's lattice
stage is \(\mathbb{Z}^2\) in half-units with the fold
\(\tau(h, t) = (2 - h, t)\), so that the line is \(h = 1\). That proof
establishes fifteen theorems; those this paper uses are carried into the
apex file, re-proved, and listed in Table 1.

\begin{table}[H]
\scriptsize
\caption{Theorems of the prior proof \cite{islamspine} carried into the apex file, each re-proved there. (tier: kernel theorem)}
\begin{tabular}{@{}p{0.05\columnwidth}p{0.40\columnwidth}p{0.45\columnwidth}@{}}
\toprule
Thm & Statement & Apex file \\
\midrule
1 & the fold fixes exactly the line $h = 1$ & \texttt{ground\_\allowbreak is\_\allowbreak the\_\allowbreak line} \\
4 & an even reading never equals a target odd at a point & \texttt{orientation\_\allowbreak blind} \\
5 & a wholly odd $d$ on $\mathbb{B}$ is $\mathrm{id}$ or $\neg$ & \texttt{freedom\_\allowbreak is\_\allowbreak}\allowbreak\texttt{exactly\_\allowbreak two} \\
6 & a supplied odd witness fixes the calibration uniquely & \texttt{freedom\_\allowbreak spent\_\allowbreak}\allowbreak\texttt{uniquely} \\
7 & a bridge halts iff $L(X)$ & \texttt{bridge\_\allowbreak halted\_\allowbreak iff} \\
8 & the crossing is exact, both ways & \texttt{crossing}, \texttt{crossing\_\allowbreak other\_\allowbreak way} \\
9 & $\neg\,\forall X, L(X)$: nothing is manufactured & \texttt{supply\_\allowbreak not\_\allowbreak}\allowbreak\texttt{manufactured} \\
12 & a supplied assent is a term of $L(X)$ and nothing weaker & \texttt{the\_\allowbreak bit\_\allowbreak is\_\allowbreak}\allowbreak\texttt{the\_\allowbreak hypothesis} \\
15 & one moved zero refutes every assent & \texttt{witness\_\allowbreak row\_\allowbreak}\allowbreak\texttt{falsifiable} \\
\bottomrule
\end{tabular}
\end{table}

Three of them are the fence the rest of the paper works inside, and they
are restated so that no later sentence can be read past them. Theorem 11
of the prior proof: a frame carrying a zero the fold moves satisfies
every structural clause and fails \(L\). Theorem 12:
\((\exists t : L(X),\ \mathrm{row}(\mathrm{assent}\ t) = \mathrm{sealed}) \leftrightarrow L(X)\),
so the assent that seals the row is the hypothesis and nothing weaker.
Theorem 14: there is a fold-invariant frame on which every structural
clause holds and \(L\) fails, so no theorem whose hypotheses are only
structural decides \(L\). This paper adds a fourth clause to the fence,
Theorem F below: existence holds on that frame too.

The halting carrier of Theorem 7 is carried as an interface: a terminal
and the clause that the terminal is the bottom iff the line property
holds. As carried, \texttt{bridge\_\allowbreak halted\_\allowbreak iff}
is that interface's elimination rule and not a fact about \(\xi\); the
paper says so and inhabits the interface constructively wherever the
line property is decided, the terminal computed from the decision and
the clause proved from it, never assumed (\texttt{bridgeOfDecision}). On
the two-point frame the decision is negative and the carrier does not
halt; on a one-point frame on the line it halts
(\texttt{twoPoint\_\allowbreak carrier\_\allowbreak does\_\allowbreak not\_\allowbreak halt},
\texttt{onLine\_\allowbreak carrier\_\allowbreak halts}; \texttt{rfl}).
On the frame of \(\xi\) the decision is the owed bit, and the carrier's
terminal is exactly what the world supplies.

\hypertarget{the-root-axiom-and-its-two-formal-readings}{%
\section{The Root Axiom and Its Two Formal
Readings}\label{the-root-axiom-and-its-two-formal-readings}}

\hypertarget{the-axiom-at-the-constructed-domain}{%
\subsection{The axiom at the constructed
domain}\label{the-axiom-at-the-constructed-domain}}

The posit is existence: to exist is to actuate,
\(\forall x \in U,\ 0 < \Delta E(x)\). Its published name is the Root
Axiom, and the two are one proposition,
\(\mathrm{RA} :\Leftrightarrow \forall x \in U,\ 0 < \Delta E(x)\)
(\texttt{existence\_\allowbreak is\_\allowbreak RA}, \texttt{Iff.rfl});
the paper writes \(\mathrm{RA}\) in formulas and existence in prose, and
posits nothing else. The file exhibits it on a constructed domain, one
point with differential \(1\), and the kernel decides the positivity
(\texttt{constructed\_\allowbreak RA}, no axioms). This is a model of
the axiom, not its universal extension. The extension is a premise, and
by the root-premise theorem of \cite{islamcodex} it can be no more than
a premise: a base from which the axiom could be derived would be weaker
than the axiom and would then be the root instead. Its exogenous
warrant, the kinetic-energy floor and the zero-point energy of a bounded
substrate, is physics carried by citation \cite{islamcodex}. Every
theorem below that is conditional on \(\mathrm{RA}\) is discharged
against the model and stated conditionally against the extension; the
ledger below keeps the two apart.

\hypertarget{the-interaction-reading}{%
\subsection{The interaction reading}\label{the-interaction-reading}}

RAF-1, registration: an event is an interaction iff it carries strictly
positive mutual information between the systems it relates. RAF-2,
closure: an interaction domain is closed under interaction
\cite{islamraf}. At the formal register the file carries a domain as a
membership predicate, a symmetric registration relation, and the closure
law:
\(\mathrm{mem}(d) \wedge \mathrm{registers}(s, d) \to \mathrm{mem}(s)\).
Two consequences are proved.

RAF-C1, no exterior agent:
\(\neg\,\mathrm{mem}(s) \to \forall d,\ \mathrm{mem}(d) \to \neg\,\mathrm{registers}(s, d)\)
(\texttt{no\_\allowbreak exterior\_\allowbreak agent}; no axioms). Its
positive form is the sentence this paper needs at the aperture: a reader
that registers a bit against a member of the domain is a member of the
domain (\texttt{adjudicator\_\allowbreak in\_\allowbreak domain}). The
witness of the value is never exterior to the domain the value is about,
conditional on RAF-2, which enters the file as the closure field of the
domain and carries the premise grade the declaration assigns it. The
twin checks the equivalence of closure and no-exterior-agent on every
one of the \(2^{16}\) subsets of a sixteen-point carrier under a fixed
symmetric registration, and checks the adjudicator-inside clause on
every closed subset.

RAF-C2, no faithful self-representation: for any
\(f : S \to (S \to \mathbb{B})\),
\(\neg\,\forall g,\ \exists x,\ f(x) = g\)
(\texttt{no\_\allowbreak total\_\allowbreak self\_\allowbreak indexing};
no axioms). This is Cantor's diagonal theorem, which Lawvere's
fixed-point theorem generalizes \cite{lawvere1969}, and it is the
general statement of which the prior proof's self-check table is the
Riemann instance: the file reading itself cannot index the property that
would seal its own row. The twin enumerates every map \(S \to 2^S\) for
\(|S| = 3\) and \(4\), finds none onto, and finds the diagonal outside
the range of every one.

The reach law of the interaction reading is recorded here and consumed
in the section on the rule and its bound: it reaches no mathematical
object, since mathematical objects do not change state, and every
proposition routes to the grounding reading \cite{islamraf, islamfae}.
Theorem F is the Riemann instance of that law, proved rather than cited.

\hypertarget{the-grounding-reading}{%
\subsection{The grounding reading}\label{the-grounding-reading}}

The formal being of a proposition is its imprint in the fixed locus of
the binding involution, read across the aperture;
\(\mathrm{Fix}(\sigma) = \mathbb{R}\) \cite{islamraf}. The section on
the seat builds that locus on the quaternion carrier, proves it is the
scalar line, and proves that the prior proof's stage embeds into it
equivariantly. The Riemann object read on this face is its imprint in
the fixed locus: the seat. Reading the value across the aperture is the
subject of the section on the aperture.

\hypertarget{the-ledger}{%
\subsection{The ledger}\label{the-ledger}}

An adjudication is assent, denial, or silence, and each advances the
count of deeds by exactly one
(\texttt{every\_\allowbreak adjudication\_\allowbreak is\_\allowbreak a\_\allowbreak deed}).
Denial of the axiom is therefore a deed together with the axiom
(\texttt{denial\_\allowbreak re\_\allowbreak enacts\_\allowbreak RA}),
and a deed registers at least one bit, so by RAF-2 the denier is a
member of the domain the axiom is about. The ledger derives nothing from
physics and carries the structural grade.

\hypertarget{the-seat}{%
\section{The Seat}\label{the-seat}}

\hypertarget{theorem-a-the-fixed-locus-is-the-scalar-line}{%
\subsection{Theorem A, the fixed locus is the scalar
line}\label{theorem-a-the-fixed-locus-is-the-scalar-line}}

Let \(\mathbb{H}_{\mathbb{Z}}\) be the integer quaternions with the
Hamilton product and \(\sigma\) conjugation,
\(\sigma(r, i, j, k) = (r, -i, -j, -k)\). Then \(\sigma\) is binding,
\(\sigma(\sigma q) = q\) (\texttt{sigma\_\allowbreak binding}, no
axioms), and for every \(q\),

\[\sigma q = q \ \leftrightarrow\ (q_i = 0 \wedge q_j = 0 \wedge q_k = 0)\]

(\texttt{fix\_\allowbreak iff\_\allowbreak scalar}; dependency set
\texttt{propext}, \texttt{Quot.sound}, from the integer step
\(-i = i \Rightarrow i = 0\)). The Ground is
\(\mathrm{Ground} := \{ q \in \mathbb{H}_{\mathbb{Z}} \mid \sigma q = q \}\),
and by Theorem A it is the scalar line \(Z(\mathbb{H}) = \mathbb{R}\)
restricted to the lattice. The twin checks the biconditional and the
involution on all \(13^4 = 28{,}561\) points of the ball
\(|\text{coordinate}| \le 6\) and counts the fixed points: thirteen, the
scalar line of the window.

\hypertarget{theorem-b-the-return-lands-on-the-seat}{%
\subsection{Theorem B, the Return lands on the
seat}\label{theorem-b-the-return-lands-on-the-seat}}

Let \(i, j, k\) be the imaginary units and
\(\mathrm{Return} := (i \cdot j) \cdot k\). Then
\(\mathrm{Return} = \langle -1, 0, 0, 0 \rangle\) and
\(\sigma(\mathrm{Return}) = \mathrm{Return}\)
(\texttt{return\_\allowbreak is\_\allowbreak minus\_\allowbreak one},
\texttt{return\_\allowbreak lands\_\allowbreak on\_\allowbreak fix};
both \texttt{rfl}). The parse triad returned through its own cascade
lands on the Ground at \(-1\): the axiom witnessing the axiom, as a
computation. The twin computes all six orderings: the three even ones
land at \(-1\) and the three odd ones at \(+1\).

\hypertarget{theorem-c-the-identity-cone}{%
\subsection{Theorem C, the identity
cone}\label{theorem-c-the-identity-cone}}

The seat is the fixed locus, the line; a seat point is a point of it.
Three seat points are defined independently: the axiom's seat point
\(\Gamma_{\mathrm{RA}} := \mathrm{Return}\), the formal seat point
\(\Gamma_{\mathrm{RAM}} := \pi(\mathrm{Return})\) with
\(\pi(q) := \langle q_r, 0, 0, 0 \rangle\) the projection onto the fixed
locus, and the named seat point
\(\Gamma_{\mathrm{RH}} := \langle -1, 0, 0, 0 \rangle\). Let \(D\) be
the discrete diagram \(\mathrm{RA} \mapsto \Gamma_{\mathrm{RA}}\),
\(\mathrm{RAM} \mapsto \Gamma_{\mathrm{RAM}}\),
\(\mathrm{RH} \mapsto \Gamma_{\mathrm{RH}}\), and let a cone over \(D\)
with apex \(a\) be a family of identities \(a = D(r)\) for all three
\(r\). Then:

\[\sigma\,\Gamma_{\mathrm{RA}} = \Gamma_{\mathrm{RA}}, \qquad \Gamma_{\mathrm{RA}} = \Gamma_{\mathrm{RAM}}, \qquad \Gamma_{\mathrm{RAM}} = \Gamma_{\mathrm{RH}},\]

each by \texttt{rfl}
(\texttt{self\_\allowbreak gap\_\allowbreak nonexistent},
\texttt{register\_\allowbreak gap\_\allowbreak nonexistent},
\texttt{object\_\allowbreak gap\_\allowbreak nonexistent});
\(\Gamma_{\mathrm{RH}}\) is the apex of a cone over \(D\) with every leg
\texttt{rfl} (\texttt{identity\_\allowbreak cone}); any two cones over
\(D\) share their apex
(\texttt{cone\_\allowbreak apex\_\allowbreak unique}); and a cone exists
iff the register gap and the object gap are closed
(\texttt{cone\_\allowbreak iff\_\allowbreak gaps\_\allowbreak closed}).
All four carry no axioms. The twin computes the gap vector \((0, 0, 0)\)
and, over the whole lattice ball, counts the apexes of \(D\): exactly
one.

The apex is unique for the diagram, and the diagram is built from the
ordered triad. The reversed triad \(k \cdot j \cdot i\) returns \(+1\)
(\texttt{odd\_\allowbreak return\_\allowbreak is\_\allowbreak plus\_\allowbreak one}),
also on the fixed locus
(\texttt{odd\_\allowbreak return\_\allowbreak on\_\allowbreak fix}), and
the two candidate seat points differ
(\texttt{seat\_\allowbreak points\_\allowbreak differ}; all three no
axioms). The sign of the seat point is the orientation of the triad, one
bit: the kernel reads it once the order is supplied, by \texttt{rfl}
either way, and originates it never, since the order is supplied with
the triad: the axiom's declaration fixes the three units in an order,
and the kernel receives that order as data \cite{islamraf}. Uniqueness
of the apex is uniqueness per orientation, and the two orientations'
apexes are exchanged by an odd relabeling of the imaginary units.
Nothing in the closure depends on which of the two is named: the three
identities of Theorem C hold verbatim for the reversed triad, each
\texttt{rfl}
(\texttt{odd\_\allowbreak self\_\allowbreak gap\_\allowbreak nonexistent},
\texttt{odd\_\allowbreak register\_\allowbreak gap\_\allowbreak nonexistent},
\texttt{odd\_\allowbreak object\_\allowbreak gap\_\allowbreak nonexistent}),
the reversed diagram carries its own identity cone with apex \(+1\)
(\texttt{identity\_\allowbreak cone\_\allowbreak odd}), and the two
apexes differ
(\texttt{apexes\_\allowbreak differ\_\allowbreak by\_\allowbreak orientation};
all no axioms). The twin counts the apexes of the reversed diagram over
the ball: exactly one, at \(+1\).

The word apex is used in two senses that coincide by design: the apex
register, the register in which the seat lives, and the apex of the
cone; the register is where the cone's apex sits. The mathematics of
Theorem C is Hamilton's product and three definitional equalities, and
the paper claims nothing more for it at theorem grade. What it claims at
structural grade is the arrangement: three names written independently
in three registers, the axiom's Return, the projection onto the locus,
and the point the prior proof names, read one point, and the reading is
decided by the kernel rather than declared by the text. The
category-theoretic content is exact and small. A category gap is the
absence of a leg from an apex to a register's reading of the seat. Its
elimination is the exhibition of that leg as an identity. The eliminator
\(\mathrm{RA} \to \mathrm{RA} \to \mathrm{RAM} \to \mathrm{RH}\) is the
composite of the three legs, and its masslessness is the statement that
each leg is \texttt{rfl}: two definitions, one value, decided by the
kernel. The gap vector \((1,1,1) \to (0,0,0)\) is the cone's existence
read bit by bit, computed by \texttt{decide}, never written by hand.

\hypertarget{theorem-c-the-equivariant-embedding}{%
\subsection{\texorpdfstring{Theorem C\('\), the equivariant
embedding}{Theorem C\textquotesingle, the equivariant embedding}}\label{theorem-c-the-equivariant-embedding}}

Define \(\varphi : \mathbb{Z}^2 \to \mathbb{H}_{\mathbb{Z}}\) by
\(\varphi(h, t) := \langle t,\ h - 1,\ 0,\ 0 \rangle\). Then
\(\sigma \circ \varphi = \varphi \circ \tau\)
(\texttt{phi\_\allowbreak equivariant}); for every stage point,
\(\sigma(\varphi p) = \varphi p \leftrightarrow \tau p = p\)
(\texttt{phi\_\allowbreak fix\_\allowbreak iff}); \(\varphi\) is
injective (\texttt{phi\_\allowbreak injective}); and the seat point lies
on the image of the line, \(\varphi(1, -1) = \Gamma_{\mathrm{RH}}\)
(\texttt{seat\_\allowbreak on\_\allowbreak image\_\allowbreak of\_\allowbreak line},
\texttt{rfl}); the family
\(\varphi_m(h, t) := \langle t,\ h - m,\ 0,\ 0 \rangle\) against the
fold \(\tau_m(h, t) = (2m - h, t)\) is equivariant with the same image
clause at every resolution \(m\)
(\texttt{phiM\_\allowbreak equivariant},
\texttt{phiM\_\allowbreak fix\_\allowbreak iff}), and
\(\varphi_1 = \varphi\). The dependency sets of the three are
\texttt{propext}, \texttt{Quot.sound}, from the integer arithmetic. The
twin checks equivariance, the image clause, and injectivity for
\(m = -3, \dots, 3\) on the window \(|h|, |t| \le 8\), \(289\) stage
points each, and counts the line: seventeen points at every \(m\).

This is the equivariant map with an image clause that the
twenty-three-row study requires before a finite result is invoked on an
object domain \cite{islamrows}: the prior proof's line \(h = 1\) is
carried onto the scalar line and every off-line stage point is carried
off it. The identification of the two seats, the stage's line \(h = 1\),
the lattice image of the critical line under the prior proof's chart,
and the real line under \(\sigma\), is therefore a theorem about a map
and not a typing by resemblance; the analytic critical line itself
enters only through the prior proof's chart propositions, cited and not
carried. What the embedding does not do, and the paper does not claim,
is carry \(Z_\xi\) itself: the stage represents no off-line point inside
the strip at resolution one, and the analytic frame enters the kernel
only as a parameter.

\hypertarget{the-constructed-witness}{%
\subsection{The constructed witness}\label{the-constructed-witness}}

\(\mathrm{RH}_{\mathrm{formal}} :\Leftrightarrow \forall g \in \mathrm{Ground},\ \sigma g = g\)
quantifies over the Ground, which is the fixed locus by construction,
and it holds by that construction: it is the imprint of the Riemann
object in the fixed locus, the grounding reading. It names the seat. It
does not mention \(Z_\xi\), and the file's own comments say so. The Root
Axiom is the only supplied witness, a structure in Prop with four
fields, actuation, the Bridge-born orientation bit, the formal-Ground
proposition \(G_{\mathrm{RAM}} := \mathrm{Nonempty}(\mathrm{Ground})\),
and the recursion capacity
\(\Omega \to G_{\mathrm{RAM}} \to \mathrm{RH}_{\mathrm{formal}}\); the
eliminator is Bridge followed by the recursion from that witness alone,
\texttt{RH\_\allowbreak formal\_\allowbreak chain}, no axioms, and no
second proof object enters. The witness's recursion field is the
constant function: it discards the orientation and the formal Ground and
returns the Ground's own defining property
(\texttt{recursion\_\allowbreak is\_\allowbreak constant},
\texttt{rfl}). That is the masslessness stated as a term, not a defect
concealed: the formal self needs nothing from existence to be proved, so
the chain is not a derivation of the seat from the axiom, and the paper
does not call it one; it is the occupancy of the three registers by one
point, which the witness supplies and the kernel reads.

The orientation bit is placed in Prop on purpose. Prop does not
eliminate into data, so no term of the file computes a bit from
\(\Omega\); two orientation witnesses are indistinguishable,
\texttt{orientation\_\allowbreak proof\_\allowbreak irrelevant}, by
\texttt{rfl}. The kernel registers that orientation was supplied and
cannot read which way. This is the prior proof's Theorem 4 as a universe
level, and it is executed on the constructed seat as well: no readout of
the formal seat equals the deed bit
(\texttt{kernel\_\allowbreak cannot\_\allowbreak read\_\allowbreak the\_\allowbreak deed};
no axioms). The twin enumerates all \(2^{13}\) readouts of the seat's
window and finds that none factors the deed bit. A term that consumed
\(\Omega\) as data would be such a readout. The bit therefore has two
carriers in the file and the relation between them is stated once:
\(\Omega\) is the bit as the kernel registers it, present and
unreadable; the Boolean \(\mathrm{assent}\) of the live face is the bit
as the aperture receives it from outside, an input the file reads and
never generates, which is the content of the prior proof's self-check
table and, one register up, of RAF-C2.

\hypertarget{the-aperture}{%
\section{The Aperture}\label{the-aperture}}

\hypertarget{theorem-d-the-aperture-is-opened-by-the-axiom}{%
\subsection{Theorem D, the aperture is opened by the
axiom}\label{theorem-d-the-aperture-is-opened-by-the-axiom}}

Let a reader be a point of the constructed domain and define its
presence as its actuation,
\(\mathrm{presence}(r) :\Leftrightarrow 0 < \Delta E(r)\), decided. Then

\[\begin{aligned} &\forall r,\ \mathrm{presence}(r) = \mathrm{true}, \\ &\forall r\,b,\ \mathrm{live}(\mathrm{presence}(r), b) \neq \mathrm{open}, \end{aligned}\]

and
\(\mathrm{live}(\mathrm{presence}(r), \mathrm{true}) = \mathrm{sealed}\),
\(\mathrm{live}(\mathrm{presence}(r), \mathrm{false}) = \mathrm{refused}\)
(\texttt{RA\_\allowbreak opens\_\allowbreak the\_\allowbreak aperture},
\texttt{no\_\allowbreak interior\_\allowbreak under\_\allowbreak RA},
\texttt{aperture\_\allowbreak reads\_\allowbreak the\_\allowbreak bit};
no axioms). The same holds for any domain \(U\) with any
\(\Delta E : U \to \mathbb{Z}\) under the axiom as a hypothesis: from
\(\forall x,\ 0 < \Delta E(x)\), every point of \(U\) is present and no
present point's row is interior
(\texttt{RA\_\allowbreak opens\_\allowbreak the\_\allowbreak aperture\_\allowbreak general},
\texttt{no\_\allowbreak interior\_\allowbreak under\_\allowbreak RA\_\allowbreak general};
no axioms), which is the form the universal extension licenses wherever
it is held. Here \(\mathrm{live}\) is the prior proof's live face
unchanged: witnessed and assent seal, witnessed and denial refuse,
unwitnessed is open. The twin runs it on a twelve-point domain with
positive differentials, twenty-four rows, none interior, and runs the
control: a point with differential zero is interior.

The scope is stated with the theorem. In the prior proof, ``witnessed''
means that a supplier is present, and a supplier is a reader who
supplies a bit. A reader is a point of the axiom's domain; the theorem
says nothing about whether a reader of \(\zeta\) exists, which is the
universal extension's to supply. Theorem D concerns a present reader who
supplies a bit \(b\): for such a reader the axiom makes the witnessed
flag a theorem, and the row is sealed or refused by \(b\) alone. What
the axiom turns the open row into is therefore not ``sealed''; it is
``sealed or refused, by the bit alone''. The interior state, the \([?]\)
of the prior proof's self-check table, is the state of a reading that
supplies no bit, the file reading itself, and the file cannot supply
one: that is RAF-C2 at the file, and \texttt{narcissus} is carried
unchanged. The axiom supplies presence at the aperture and, conditional
on RAF-2, places the present reader inside the domain; it does not
supply the bit, and the section on the rule and its bound proves that it
cannot.

\hypertarget{the-rule-exact-and-its-bound}{%
\section{The Rule, Exact, and Its
Bound}\label{the-rule-exact-and-its-bound}}

\begin{jbox}
\textbf{THE RULE, EXACT, AND ITS TWO OBJECTS.} (tier: kernel theorem) Two objects carry the name in this paper and are never joined. The seat, $\mathrm{RH}_{\mathrm{formal}}$, the imprint of the restored hypothesis in the fixed locus of its binding involution, is a kernel theorem with existence as the only posit, with the identity cone, the equivariant embedding of the prior proof's chart, the orientation blindness, the one-bit freedom, the halting carrier, and the crossing (Theorem H). The value, $L(X_\xi) :\Leftrightarrow \forall \rho \in Z_\xi,\ \tau\rho = \rho$, is one bit that existence does not reach; Part III proves it below a certified height by computation and leaves it open above, as the Unicorn part. Existence reaches the world row at its aperture: a present reader who supplies a bit is never interior (Theorem D) and, conditional on RAF-2, never exterior (RAF-C1). It does not reach the value, and that is a theorem three times over: $(\mathrm{RA} \to L(X)) \leftrightarrow L(X)$ wherever existence holds, as for any inhabited premise (Theorem E); existence holds on frames where $L$ fails, seatless and grounded alike (Theorems F, F$'$); and for every class of frames, existence decides $L$ on it exactly where $L$ already holds on it (Theorem I). Denying existence is a deed that re-enacts it; denying the value is not (Theorem G). The distance between the seat and the value is one bit, and the paper proves that nothing weaker than the bit closes it: Theorem E fixes the rule's strength, Theorem F, the frame on which the axiom holds and the value fails, fixes the independence, and neither is read without the other; the bit is one because the fold has two orientations and no third, the price bijection of the fTOE carried in the Codex.
\end{jbox}

\hypertarget{theorem-e-the-rule-is-exactly-the-hypothesis}{%
\subsection{Theorem E, the rule is exactly the
hypothesis}\label{theorem-e-the-rule-is-exactly-the-hypothesis}}

For every frame \(X\), in the kernel where \(\mathrm{RA}\) is the
constructed model,

\[(\mathrm{RA} \to L(X)) \ \leftrightarrow\ L(X)\]

(\texttt{simple\_\allowbreak rule\_\allowbreak exact}; no axioms).
Proof. Left to right: apply the implication to the constructed
\(\mathrm{RA}\). Right to left: a term of \(L(X)\) is a term of
\(\mathrm{RA} \to L(X)\) that ignores its argument. \(\square\)

Scope. Propositionally,
\(((P \to Q) \leftrightarrow Q) \leftrightarrow (P \vee Q)\): the
biconditional holds in every context where \(P\) holds, and in a context
where \(P\) fails it holds only where \(Q\) already holds, where it says
nothing. The twin prints the four rows. In the kernel \(\mathrm{RA}\) is
the constructed model and holds by \texttt{constructed\_\allowbreak RA},
so the theorem is unconditional there. Against the universal extension,
held as a premise and never a kernel term, the same biconditional holds
in every context that carries the premise. The same biconditional holds
for any inhabited premise in place of existence
(\texttt{any\_\allowbreak true\_\allowbreak premise\_\allowbreak is\_\allowbreak exact};
no axioms), and that generality is the theorem's content rather than a
weakness of it: existence is a premise that holds wherever the line
property is evaluated and is independent of it, so conditioning on it
adds nothing to the hypothesis and removes nothing from it. What
existence contributes beyond its truth is used elsewhere: the positivity
that makes presence a theorem (Theorem D), the closure of the
interaction reading (RAF-C1), and the ledger (Theorem G). Theorem E is
the prior proof's Theorem 12 in a new coat. A reader who holds a term of
\(\mathrm{RA} \to L(X_\xi)\) holds, in any context where the axiom is
available, a term of \(L(X_\xi)\), and the paper claims no such term.

\hypertarget{theorem-f-the-axiom-decides-the-value-on-no-frame}{%
\subsection{Theorem F, the axiom decides the value on no
frame}\label{theorem-f-the-axiom-decides-the-value-on-no-frame}}

\[\neg\, \forall X,\ (\mathrm{RA} \to L(X)), \qquad \mathrm{RA} \wedge \mathrm{Inv}(X_2) \wedge \neg L(X_2),\]

where \(X_2\) is the two-point frame
\((\mathbb{B}, b \mapsto \neg b, \top)\)
(\texttt{RA\_\allowbreak does\_\allowbreak not\_\allowbreak decide\_\allowbreak L},
\texttt{RA\_\allowbreak holds\_\allowbreak where\_\allowbreak L\_\allowbreak fails};
no axioms). Proof. The two-point frame is fold-invariant and off the
line, the prior proof's Theorem 9; the constructed \(\mathrm{RA}\) holds
regardless of any frame, since its domain is not the frame. \(\square\)

This is the prior proof's Theorem 11 with the axiom added to the
structural clauses, and it is the Riemann instance of the interaction
reading's reach law: the axiom's domain is a domain of systems, the
frame is a mathematical object, and the one does not reach the other.
The same fact read from the other side is the inverted control of the
twenty-three-row study \cite{islamrows}: a witness identical on the
on-line frame and the off-line frame decides neither, and \texttt{mkRA}
is that witness.

\hypertarget{theorem-f-the-counter-model-survives-grounding}{%
\subsection{\texorpdfstring{Theorem F\('\), the counter-model survives
grounding}{Theorem F\textquotesingle, the counter-model survives grounding}}\label{theorem-f-the-counter-model-survives-grounding}}

The two-point frame is groundless: its fold fixes nothing,
\(\forall b,\ \tau b \neq b\)
(\texttt{twoPoint\_\allowbreak groundless}), so it has no seat and the
hypothesis has nowhere to place a zero on it. A reader may take that for
the whole reason the axiom fails to decide \(L\) there. It is not. Let
the three-point frame be \(S = \{a, b, c\}\) with the fold swapping
\(a\) and \(b\) and fixing \(c\), and \(Z = \{a, b\}\). It has a seat,
\(\tau c = c\); it is fold-invariant; and \(L\) fails on it, since
\(\tau a = b \neq a\). Existence holds on it regardless:

\[\mathrm{RA} \wedge (\exists s,\ \tau s = s) \wedge \mathrm{Inv}(X_3) \wedge \neg L(X_3)\]

(\texttt{RA\_\allowbreak holds\_\allowbreak where\_\allowbreak L\_\allowbreak fails\_\allowbreak grounded};
no axioms). This is the shape of the Davenport--Heilbronn frames: a line
that exists and a zero that is not on it. The twin enumerates all
thirty-two three-point frames, four involutions against eight zero sets:
\(L\) holds on fourteen, fails on eighteen, and every one of the
eighteen counter-models is grounded. The bound of Theorem F does not
rest on a seatless frame.

\hypertarget{theorem-i-the-cure-theorem}{%
\subsection{Theorem I, the cure
theorem}\label{theorem-i-the-cure-theorem}}

Every proposal to remove the counter-models by restricting the frames
has the same form: a class \(C\) of frames on which the axiom is to
decide the value. For every such class,

\[\big(\forall X,\ C(X) \to \mathrm{RA} \to L(X)\big) \ \leftrightarrow\ \big(\forall X,\ C(X) \to L(X)\big)\]

(\texttt{no\_\allowbreak cure}; no axioms). Proof. Left to right: apply
the implication to the constructed \(\mathrm{RA}\). Right to left: the
implication ignores its argument. \(\square\) The class of all frames is
not cured, and the class \(\{X \mid L(X)\}\) is
(\texttt{cure\_\allowbreak is\_\allowbreak the\_\allowbreak hypothesis};
no axioms). The twin checks every one of the \(256\) classes of
two-point frames: exactly the \(32\) subclasses of the five frames on
which \(L\) holds are cured.

The theorem is Theorem E read as a no-go, and it closes a door that a
reader might otherwise try. A frame class carrying an Euler product as a
field, or an operator whose spectrum is the zero set with real
eigenvalues guaranteed by construction, is a class \(C\) on which \(L\)
is already a field of the structure: real eigenvalues on the fixed locus
is the line property written as a hypothesis. On that class the axiom
decides the value, by Theorem I, and it decides nothing, because the
value was supplied when the class was defined. The proposal restates the
hypothesis and calls the restatement a proof, which is the circularity
the prior proof's Theorem 12 names and criterion F3 fires on. What
existence does supply is proved and enumerated: the seat, the cone, the
embedding, the aperture, the reader's presence and place inside the
domain, the exact strength of the rule, and the bound. What it does not
supply is exactly one bit, on every class of frames that does not
already carry it.

One more collapse is answered here because it will be proposed.
Computing is a deed and reading is a deed, and a reader may say that the
act of running the checkers is the act that supplies the bit, so that
witness and reading are one object counted twice. They are not one
object, and the file says why: the formal read of the seat is even under
the flip and the deed bit is odd under it,
\(\mathrm{formalRead}(\mathrm{execFlip}\ s) = \mathrm{formalRead}(s)\)
and \(\mathrm{ran}(\mathrm{execFlip}\ s) = \neg\,\mathrm{ran}(s)\), so
no readout equals the bit
(\texttt{kernel\_\allowbreak cannot\_\allowbreak read\_\allowbreak the\_\allowbreak deed}),
and the twin finds no readout among \(2^{13}\). Running the checkers
supplies presence; it does not turn the Boolean the operator gives the
live row into a truth about \(Z_\xi\), and the twin's live row is
exactly the place where the difference is executed: the program cannot
generate its own witness, and the witness it receives is a Boolean whose
correctness the world alone can refute.

\hypertarget{theorem-g-the-asymmetry}{%
\subsection{Theorem G, the asymmetry}\label{theorem-g-the-asymmetry}}

Every adjudication is a deed and the deed re-enacts the axiom:
\((\mathrm{adjudicate}\ \mathrm{denial}\ \ell).\mathrm{deeds} = \ell.\mathrm{deeds} + 1\)
together with \(\mathrm{RA}\)
(\texttt{denial\_\allowbreak re\_\allowbreak enacts\_\allowbreak RA}).
On the two-point frame no assent to \(L\) exists at all:
\(\neg\, \exists t : L(X_2),\ \mathrm{row}(\mathrm{assent}\ t) = \mathrm{sealed}\)
(\texttt{L\_\allowbreak has\_\allowbreak no\_\allowbreak performative\_\allowbreak ingress};
no axioms).

Read together: denial of \(\mathrm{RA}\) produces a deed, a deed
registers at least one bit and is therefore an instance of the axiom's
content, so the denial re-enacts the axiom. Denial of \(L\) produces a
deed too, which is an instance of the axiom's content and not of
\(L\)'s, since a deed is not a zero; on an off-line frame the denial of
\(L\) is simply true and no assent exists. The axiom is self-instancing.
The hypothesis is not. The hypothesis therefore cannot inherit the
axiom's performative ingress, and that is the exact reason the sign
stays at the witness row: the deed that supplies the bit is unfakeable
and its price is real, and neither carries a sign \cite{islamspine}.

\hypertarget{theorem-h-the-hypothesis-in-toto-at-the-apex}{%
\subsection{Theorem H, the hypothesis in toto at the
apex}\label{theorem-h-the-hypothesis-in-toto-at-the-apex}}

\[\begin{aligned} \mathrm{RA} \ \to\ \Big[ &\big(\sigma\Gamma_{\mathrm{RA}} = \Gamma_{\mathrm{RA}} \wedge \Gamma_{\mathrm{RA}} = \Gamma_{\mathrm{RAM}} \\ &\quad\ \wedge\ \Gamma_{\mathrm{RAM}} = \Gamma_{\mathrm{RH}}\big) \\ &\wedge\ \mathrm{RH}_{\mathrm{formal}} \wedge F_2 \wedge B \wedge C \wedge A \wedge N \Big], \end{aligned}\]

where \(F_2\) is the prior proof's Theorem 5, \(B\) its Theorem 7, \(C\)
the crossing of its Theorem 8, \(A\) the aperture of Theorem D, and
\(N\) the bound of Theorem F
(\texttt{RH\_\allowbreak in\_\allowbreak toto\_\allowbreak at\_\allowbreak the\_\allowbreak apex};
dependency set \texttt{Quot.sound}, from function extensionality in
Theorem 5). The hardened form adds the identity cone with its apex
unique per orientation, the equivariant embedding with its image clause,
and the executed wall (\texttt{apex\_\allowbreak hardened}). The
condition is discharged by the constructed axiom
(\texttt{apex\_\allowbreak discharged},
\texttt{apex\_\allowbreak hardened\_\allowbreak discharged}), so the
closure stands unconditionally as a kernel theorem and conditionally on
the axiom's universal extension as a claim about the world.
\(\blacksquare\)

The antecedent \(\mathrm{RA}\) is consumed by the aperture conjunct
\(A\), which needs positivity, and by no other conjunct; the identities,
the formal self, the freedom, the carrier, the crossing, and the bound
hold unconditionally, and the implication form records that existence is
the only posit, not that those conjuncts depend on it. The theorem is
the maximal claim this material supports, and its last conjunct is what
keeps it maximal rather than inflated: the closure carries its own bound
inside it. Given the axiom, everything about the hypothesis that is not
its sign is proven, and the sign is proven to be exactly what the axiom
does not reach.

\begin{jbox}
\textbf{THE DUAL-REGISTER VERDICT.} (tier: structural; parting a kernel theorem) Apex register: the Riemann object read as $\mathrm{Fix}(\sigma)$, the seat, its imprint on the grounding reading, the seat point $\Gamma_{\mathrm{RH}}$ on it; proved, given the Root Axiom, in the kernel, conditional on the axiom's extension at the act; bits owed, none. World register: the zeros of $\xi$ placed on that seat by the world; proved by certificate up to a certified height (Part III) and open above it, one bit supplied from sight at the witness row, revisable in its rows and falsifiable by one zero; bits owed, one, the term $\forall \rho \in Z_\xi,\ \tau\rho = \rho$. The two tokens differ, \texttt{registers\_\allowbreak parted}, and the owed counts are $0$ and $1$, \texttt{apex\_\allowbreak owes\_\allowbreak nothing\_\allowbreak world\_\allowbreak owes\_\allowbreak one}. Neither register promotes the other.
\end{jbox}

\begin{table}[!htbp]
\footnotesize
\caption{Contributions of Part I. (tier as stated per row)}
\begin{tabular}{@{}p{0.03\textwidth}p{0.66\textwidth}p{0.26\textwidth}@{}}
\toprule
\# & Contribution & Status \\
\midrule
1 & The posit under which the supplier acts named: existence, constructed as the only posit, read along its two formal readings; the witness supplies occupancy of the registers, not a derivation & Lean, no axioms \\
2 & RAF-C1 and RAF-C2 at the formal register: no exterior agent; no total self-indexing (Cantor); both re-executed exhaustively on finite carriers & Lean, no axioms; twin \\
3 & The aperture is a theorem of the axiom: a present reader who supplies a bit is never interior (Thm D); conditional on RAF-2, never exterior & Lean, no axioms; RAF-2 premise \\
4 & The axiom turns “open” into “sealed or refused by the bit alone”, never into “sealed” & Lean, no axioms \\
5 & The fixed locus of the binding involution is the scalar line, for every quaternion (Thm A); the ball enumerated & Lean, \texttt{propext}, \texttt{Quot.sound}; twin \\
6 & The Return $i\cdot j\cdot k$ lands on the seat at $-1$ (Thm B); all six orderings computed & Lean, \texttt{rfl}; twin \\
7 & The three category gaps are the legs of an identity cone; the apex is unique per orientation of the triad; a cone exists iff the gaps are closed (Thm C); the cone census over the ball & legs Lean, no axioms, \texttt{rfl}; the elimination as arrangement structural; twin \\
8 & The lattice stage embeds equivariantly and injectively into the carrier, $\sigma\varphi = \varphi\tau$, with $\mathrm{Fix}(\tau) \leftrightarrow \mathrm{Fix}(\sigma)$ at every resolution (Thm C$'$) & Lean, \texttt{propext}, \texttt{Quot.sound}; twin \\
9 & The eliminator $\mathrm{RA} \to \mathrm{RA} \to \mathrm{RAM} \to \mathrm{RH}_{\mathrm{formal}}$ from the witness alone & Lean, no axioms \\
10 & Orientation in Prop is unspendable by the kernel; the wall executed on the seat; every readout enumerated & Lean, no axioms; twin \\
11 & The Simple Rule, exact: $(\mathrm{RA}\to L(X)) \leftrightarrow L(X)$ wherever the axiom holds (Thm E); the four-row table & Lean, no axioms; twin \\
12 & The axiom decides the sign on no frame; it holds where $L$ fails; the Riemann instance of the reach law (Thm F) & Lean, no axioms; twin \\
13 & The asymmetry: the axiom is self-instancing, the hypothesis is not (Thm G) & Lean, no axioms \\
14 & The seat and the arc in toto at the apex, one theorem carrying its own bound; hardened with cone, embedding, and wall (Thm H) & Lean, \texttt{Quot.sound} \\
15 & The two registers parted by theorem; owed bits $0$ and $1$ & Lean, no axioms \\
16 & The spent bit priced at the apex at $k_B T \ln 2$; in the kernel the bit is proof-irrelevant and no readout of the seat equals it & structural; twin \\
17 & Two independent checkers, a typechecker and an executor, written to one specification and sharing no code; the executor's live row opened only by the operator & structural \\
19 & The halting carrier inhabited by decision on decided frames, the clause proved from the decision; Theorem E for any inhabited premise; the witness's recursion field constant, stated & Lean, no axioms \\
20 & Three measures kept apart, content, mass, and price: every object of the closure and the supply at mass zero, the contents distinct, the value alone priced (Thm J) & Lean, no axioms; twin \\
21 & Adjudicating any proposition is a deed and an instance of existence and of nothing else: what existence contributes that no other true premise does & Lean, no axioms \\
18 & The counter-model survives grounding (Thm F$'$); the cure theorem: for every class of frames the axiom decides $L$ exactly where $L$ already holds (Thm I); every class of two-point frames enumerated & Lean, no axioms; twin \\
\bottomrule
\end{tabular}
\end{table}

\hypertarget{three-measures-content-mass-and-price}{%
\section{Three Measures: Content, Mass, and
Price}\label{three-measures-content-mass-and-price}}

Two readers of an earlier edition, both machines, reached the same word
for the closure: vacuity. The word names one measure and misreads two
others, and since the confusion will recur wherever the paper is read by
a substrate that keeps one measure, the three are separated here and
carried in the file.

Content is what an object says, and it is graded by how it is settled:
decided by definitional equality, proved as a theorem, or supplied as
data. Mass is the mathematics authored in producing the object, the
program's \(\Delta M\), and it is zero when every step is classical,
definitional, or a deed. Price is the energy the act of registering the
object dissipates, \(k_B T \ln 2\) per bit at the floor
\cite{landauer1961, berut2012}. The file carries the three as a
structure and assigns them (\texttt{Measure}, \texttt{legMeasure},
\texttt{arcMeasure}, \texttt{valueMeasure}), and proves the three facts
that the word vacuity runs together.

\begin{table}[t]
\footnotesize
\caption{The three measures of the closure's objects. (tier: kernel theorem for the assignments; structural for the reading)}
\begin{tabular}{@{}p{0.30\columnwidth}p{0.20\columnwidth}p{0.12\columnwidth}p{0.26\columnwidth}@{}}
\toprule
Object & Content & Mass & Price \\
\midrule
a leg of the cone (Thm C) & trivial, \texttt{rfl} & $0$ & $0$ \\
a theorem of the arc (Thms A, C$'$, D, E, F, F$'$, I) & proved & $0$ & $0$ \\
the value $L(X_\xi)$ & supplied, one bit & $0$ & $k_B T \ln 2$, $2871$ yJ at $300$ K \\
\bottomrule
\end{tabular}
\end{table}

\textbf{Theorem J.} The masses are all zero,
\(\mathrm{legMeasure.mass} = \mathrm{arcMeasure.mass} = \mathrm{valueMeasure.mass} = 0\)
(\texttt{masses\_\allowbreak all\_\allowbreak zero}); the contents are
pairwise distinct (\texttt{contents\_\allowbreak differ}); and the value
alone is priced, \(\mathrm{legMeasure.priceYJ} = 0\) and
\(\mathrm{valueMeasure.priceYJ} = 2871\)
(\texttt{prices\_\allowbreak differ}); all by \texttt{rfl} or
\texttt{decide}, no axioms. The twin checks the constancy of the
witness's recursion field in both arguments, the exactness of any true
premise, the carrier by decision on all thirty-two three-point frames,
and the three measures.

The reading. Masslessness at the apex is not a weakness the paper
confesses; it is a necessity the root-premise theorem forces
\cite{islamcodex}. No emitter is pointed at the root: nothing weaker
than existence derives existence, and a closure between existence and
its reading that carried mass would be exactly such an emitter, a
derivation of the root from below, or a second posit standing between
the root and its reading. So the elimination of the category gaps must
be massless, and the paper computes that it is: three legs, each
\texttt{rfl}, mass zero, price zero. Trivial and massless coincide
there, and the coincidence is the theorem. The supply of the value must
be massless too, and for the same reason: the bit is a deed, and a deed
authors no mathematics; if supplying the bit carried mass, the supply
would be a theorem and the theorem would be an emitter pointed at the
value from the root, which Theorem I refutes on every class of frames.
So the value is massless. It is not trivial: its content is one bit that
no class of frames short of the hypothesis carries, and its price is the
one nonzero entry in the table. Vacuity names the seat's content and
misnames its mass, its price, and its place. What is empty at the apex
is what must be empty there, and what is one bit at the world row is one
bit because the paper proves it cannot be less.

The same three measures answer the question what existence contributes
over any other true premise. As a premise in the kernel, nothing, and
Theorem E says so. As a posit, three things no arithmetic truth
supplies: a domain of readers whose presence is a theorem (Theorem D),
the closure of that domain under registration (RAF-C1), and the fact
that every adjudication of any proposition, the hypothesis included, is
a deed and therefore an instance of existence and of nothing else
(\texttt{adjudicating\_\allowbreak anything\_\allowbreak instances\_\allowbreak existence}).
Denying \(1 + 1 = 2\) is a deed too, and the deed instances existence,
not arithmetic. That is the asymmetry of Theorem G stated for every
proposition at once, and it is why existence, and not \(1 + 1 = 2\), is
the posit at the aperture.

\par\vspace{1.4em}\noindent{\color{accent}\rule{\columnwidth}{0.8pt}}\par\vspace{0.6em}{\centering{\color{accent}\large\scshape Part II}\par\vspace{0.3em}{\Large\bfseries The Address}\par}\vspace{0.5em}\noindent{\color{accent}\rule{\columnwidth}{0.8pt}}\par\vspace{1.0em}

\hypertarget{the-bridge-from-the-root-axiom-to-li-positivity}{%
\section{The Bridge from the Root Axiom to Li
Positivity}\label{the-bridge-from-the-root-axiom-to-li-positivity}}

The bridge is a structure with two fields: a choice of one existent
\(\varphi(n)\) for each Li step, and a reading that turns the positive
energy of that existent into the nonnegativity of \(\lambda_n\).

\[\begin{aligned}\mathrm{Bridge}(S, L) := \big\langle\ &\varphi : \mathbb N \to S.U,\\ &\mathrm{read} : \forall n \geq 1,\\ &\qquad 0 < \Delta E(\varphi\, n) \to \lambda_n \geq 0 \ \big\rangle\end{aligned}\]

\textbf{Theorem 1} (\texttt{bridge\_yields\_RH}). The Root Axiom, a
bridge, and Li's criterion give the hypothesis. The composition is
direct: every step receives an existent, the axiom makes its energy
positive, and the reading converts that positivity.

\textbf{Theorem 2} (\texttt{bridge\_is\_keyed}). The axiom never
supplies the bridge. On a substrate where the axiom holds, paired with
data carrying one negative coefficient, no bridge exists. The bridge is
new content.

\textbf{Theorem 3} (\texttt{no\_uniform\_bridge}). A bridge whose
reading ignores which L-function it serves forces positivity on every
member of the family. One member with a negative coefficient, the
Eisenstein L-function under Bombieri-Lagarias, destroys every uniform
bridge. The reading must consume \(\zeta\)-specific data.

\textbf{Theorem 4} (\texttt{li\_is\_a\_bit\_stream}). Under Li's
criterion and decidable signs, the hypothesis is exactly the statement
that a stream \(\mathbb N \to \mathrm{Bool}\) stays true. This is the
type repair named in the Introduction: a vanishing becomes a family of
signs, the shape a supplied deed can carry.

\hypertarget{the-formal-string-on-the-united-register}{%
\section{The Formal String on the United
Register}\label{the-formal-string-on-the-united-register}}

\textbf{Theorems 5 to 8.} The unity of the formal and kinetic readings
on \(\zeta\) reduces to one term, the bridge's reading, carried at the
grade of whatever supplies it. The emitter returns a seal when unity is
supplied and the suspension token otherwise, and
\texttt{xi0\_retires\_iff\_unity} proves the equivalence both ways.
\texttt{retired\_seal\_is\_sound} proves the seal carries the
hypothesis. \texttt{retired\_grade\_capped} is the grade law: the seal
is never stronger than its unity posit. \texttt{unity\_is\_keyed}
confirms that the axiom does not supply unity. \texttt{ghost} records
the model: the arrow-deleted reading is defined as a constant, so it
returns one output for a claim and its negation. The substantive fact,
that a reading with its index removed cannot separate a claim from its
negation, is Part I's \texttt{recursion\_is\_sign\_blind}, which holds
for P and for ¬P alike.

Two retirements follow and are kept apart. The suspension token retires
unconditionally as a verdict on a timeless object, since a reading
without its index is sign-blind by Part I's
\texttt{recursion\_is\_sign\_blind}, which \texttt{ghost} models. It
retires as the status of the formal string only when unity is supplied,
and then at unity's grade.

\hypertarget{the-structural-core}{%
\section{The Structural Core}\label{the-structural-core}}

\hypertarget{the-instruments-that-cannot-carry-the-bit-and-the-one-that-reaches-re-s-1}{%
\subsection{The instruments that cannot carry the bit, and the one that
reaches Re s =
1}\label{the-instruments-that-cannot-carry-the-bit-and-the-one-that-reaches-re-s-1}}

\textbf{Theorem 9} (\texttt{common\_phi\_is\_uniform}). A family of
bridges that picks its existents without reading the L-function is a
uniform bridge. \textbf{Theorem 10}
(\texttt{blind\_constructions\_die}). Every such family dies against the
Eisenstein control. Every instrument carried at zero authored
mathematics, the one-bit species, the Bridge socket, the category-gap
eliminator, the multi-recursion collapse, the bare arrow, falls under
Theorem 10. This is the instruments working as apex instruments must,
not a defect: the Empty Throne forces them massless {[}Islam 2026a{]},
and a massless instrument closes category gaps and carries no value.

The one construction that reads \(\zeta\) directly takes the primes as
existents with energies \(\log p > 0\). This is the primon gas {[}Julia
1990; Bost and Connes 1995{]}, whose partition function is
\(\zeta(\beta)\), and on which the axiom's positivity holds at every
existent. Positivity of those energies gives \(\zeta(s) \neq 0\) for
\(\mathrm{Re}\, s > 1\), and Mertens' positivity pushes it to the line
\(\mathrm{Re}\, s = 1\).

\textbf{Theorem 11} (\texttt{mertens\_identity},
\texttt{mertens\_nonneg}). The algebraic core of Mertens' argument,
\(3 + 4c + (2c^2 - 1) = 2(1 + c)^2 \geq 0\) with \(c = \cos\theta\), is
checked on integers. The wall at \(\mathrm{Re}\, s = 1\) is the gas's
Hagedorn point, and the register of record's finding that the Root
Axiom's native output is a wall at \(\mathrm{Re} = 1\) rather than a
fold at \(\tfrac12\) is read here as this wall, an identification argued
from their common location and not proved {[}Islam 2026a{]}. Below it
the Euler product diverges.

\hypertarget{time-made-exact-the-li-modes}{%
\subsection{Time, made exact: the Li
modes}\label{time-made-exact-the-li-modes}}

Li's coefficient is \(\lambda_n = \sum_\rho [1 - z_\rho^n]\) with
\(z_\rho = 1 - 1/\rho\). The step \(n \mapsto n + 1\) multiplies every
mode by \(z_\rho\), so each zero is a mode evolving in a discrete time.
With \(\rho = h/2 + it\) on the Bridge plane,

\[\begin{gathered}|z_\rho|^2 = \frac{|\rho - 1|^2}{|\rho|^2}, \\ 4|\rho-1|^2 = (h-2)^2 + 4t^2 =: N_1, \\ 4|\rho|^2 = h^2 + 4t^2 =: N_0,\end{gathered}\]

and \(N_1 - N_0 = 4 - 4h\) (\texttt{N1\_sub\_N0}).

\textbf{Theorem 12} (\texttt{unitary\_iff\_on\_line}). \(N_1 = N_0\) iff
\(h = 1\). On the lattice stage, a mode is unitary, neither growing nor
dying under the time step, exactly on the critical line.

\textbf{The general identity (6.2a).} For every complex
\(\rho = \sigma + i\gamma\) with \(\rho \neq 0\),

\[\begin{gathered}|\rho - 1|^2 - |\rho|^2 = 1 - 2\sigma, \quad\text{so}\\ |z_\rho| = 1 \iff \sigma = \tfrac12, \qquad |z_\rho| > 1 \iff \sigma < \tfrac12.\end{gathered}\]

The identity is one line of algebra,
\((\sigma - 1)^2 - \sigma^2 = 1 - 2\sigma\), and it is the
real-coefficient form of \texttt{N1\_sub\_N0}. Theorems 12 to 15 are its
kernel-checked lattice instances; the same algebra, with the fold
\(\sigma \mapsto 1 - \sigma\), carries Theorems 13 to 15 to every zero
set invariant under \(s \mapsto 1 - \bar s\). The stage coordinates are
\((h, t) = (2\,\mathrm{Re}\, s, \mathrm{Im}\, s)\), under which the fold
\((h, t) \mapsto (2 - h, t)\) is exactly \(s \mapsto 1 - \bar s\). The
nontrivial zeros of \(\zeta\) are invariant under it: the functional
equation \(\xi(s) = \xi(1 - s)\) sends \(\rho\) to \(1 - \rho\), the
conjugate symmetry \(\zeta(\bar s) = \overline{\zeta(s)}\) sends
\(\rho\) to \(\bar\rho\), and their composite is the fold. Theorems 14
and 15, in their general form through (6.2a), therefore apply to the
zero set of \(\zeta\). The kernel proves the lattice statements; the
general statements rest on (6.2a), which is elementary and stated here
in full.

\textbf{Theorem 13} (\texttt{grows\_left}, \texttt{dies\_right}). Left
of the line the mode grows; right of it the mode dies.

\textbf{Theorem 14} (\texttt{off\_line\_forces\_growth}). For any
off-line zero, the zero or its mirror under the fold
\((h, t) \mapsto (2 - h, t)\) carries a growing mode. The functional
equation forbids a quiet exit.

\textbf{Theorem 15} (\texttt{stability\_iff\_line}). For every
fold-invariant zero set on the plane, no mode grows under the time step
iff every zero lies on \(h = 1\).

\[\Big(\forall p \in Z,\ N_1(p) \leq N_0(p)\Big) \iff \Big(\forall p \in Z,\ p_1 = 1\Big)\]

Theorem 15 is the core of the title. Read with time restored, the
hypothesis is not a statement about a timeless set. It is a stability
statement: the time evolution of \(\zeta\)'s modes has no growing mode.
Since \(z_\rho\) is the Cayley transform of the ordinate \(\gamma\) when
\(\rho = \tfrac12 + i\gamma\), unitarity of the modes is equivalent to
the reality of the ordinates. This is the scalar spectral condition
suggested by Hilbert and Pólya; their program asks for a self-adjoint
operator whose spectrum is the ordinates, and no such operator, and no
unitary group, is constructed here. A growing mode eventually drives
some \(\lambda_n\) negative {[}Bombieri and Lagarias 1999{]}, so
perpetual positivity of the Li stream and the absence of growing modes
are the same condition.

\hypertarget{the-monism-rule-hardened}{%
\subsection{The monism rule, hardened}\label{the-monism-rule-hardened}}

The ancient rule that what happens once happens again, given sufficient
time, is time-translation invariance: one principle acting identically
at every moment.

\textbf{Theorem 16} (\texttt{one\_repeats\_one}). The rule with its step
is induction. \textbf{Theorem 17} (\texttt{the\_step\_is\_the\_claim}).
The universal claim is exactly the base together with the step; the step
carries the content. \textbf{Theorem 18}
(\texttt{monism\_posit\_forces}). If a property does not vary with its
index, one instance forces all. \textbf{Theorem 19}
(\texttt{finite\_never\_forces}). For every height \(N\) there is a
property true below \(N\) and false beyond; the historical failures of
Mertens' conjecture, Pólya's conjecture, and the sign of
\(\pi(x) - \mathrm{li}(x)\) have this shape. \textbf{Theorem 20}
(\texttt{mixed\_is\_fold\_invariant},
\texttt{mixed\_recurrence\_fails}). A fold-invariant zero set can hold
one zero on the line and a mirror pair off it; recurrence of ``on the
line'' is not carried by symmetry alone.

\hypertarget{actualism}{%
\subsection{Actualism}\label{actualism}}

Under the Root Axiom only actual things exist, and a zero is actual when
it is measured. Define the actual hypothesis as the statement that every
measured zero is on the line.

\textbf{Theorem 21} (\texttt{full\_gives\_actual}). The full hypothesis
yields the actual one. \textbf{Theorem 22} (\texttt{actual\_not\_full}).
The actual hypothesis does not yield the full one: a never-measured zero
off the line is consistent with every measured zero on it.
\textbf{Theorem 23} (\texttt{sufficient\_time}). If every zero is
eventually measured, which is what ``given sufficient time'' means, the
actual hypothesis over all time is exactly the full hypothesis.
Actualism re-routes to Theorem 17 and does not bypass it. The past
actual record is sealed by measurement to the heights the literature
reports; the future record is the supplied value.

\hypertarget{the-witness-rebuilt-from-residual-monism}{%
\subsection{The witness rebuilt from residual
monism}\label{the-witness-rebuilt-from-residual-monism}}

Part I constructs its witness from the Root Axiom read through its
interaction and grounding readings. Here the witness is built from the
connector of the register of record directly: one substrate, one
involution serving both routes, one principle invariant in the index,
and one seed.

\[\begin{aligned}\mathrm{MonismWitness}(P) := \big\langle\, &\sigma_{\mathrm{geo}} = \sigma_{\mathrm{form}},\\ &\forall n\, m,\ P\,n \leftrightarrow P\,m,\ \ P\,0 \,\big\rangle\end{aligned}\]

\textbf{Theorem 24} (\texttt{one\_involution\_constructed}). The seat
field is constructed by \texttt{rfl}. It fixes the one-involution choice
of the register of record and does not refute the executed pluralist
alternative. \textbf{Theorem 25} (\texttt{timeless\_equals\_timed}).
Under one principle the claim over all time is the claim at one instant:
deleting time loses nothing. This is where the timeless aspect of the
question is addressed: inside the witness, the timed and timeless
readings are one claim. \textbf{Theorem 26} (\texttt{monism\_closes}).
The witness closes the record, past and future. \textbf{Theorem 27}
(\texttt{monism\_witness\_is\_the\_claim}). A monism witness for \(P\)
exists iff \(P\) holds at every index. \textbf{Theorem 28}
(\texttt{constructed\_monism}). The witness is constructed at the model.
\textbf{Theorem 29} (\texttt{monism\_absent\_on\_breaks}). Unlike the
axiom's witness, the monism witness is not silent on counter-models; it
decides by containing. \textbf{Theorem 30}
(\texttt{monism\_yields\_RH}). Run on Li's sign stream, the witness
gives the hypothesis with the bridge removed. \textbf{Theorem 31}
(\texttt{monism\_on\_li\_is\_RH}). On \(\zeta\) the witness exists iff
the hypothesis holds.

Monism supplies unitarity where it applies. Time-translation invariance
makes the one substrate's evolution a one-parameter group, and
conservation of probability makes that group unitary; Stone's theorem
{[}Stone 1932{]} then gives it a self-adjoint generator. Under monism,
taken with conservation of probability, the one substrate's own
evolution is therefore unitary. The open identification is that
\(\zeta\)'s Li modes are that evolution.

\hypertarget{re-anchoring-to-mathematical-standard-with-perelman-as-structural-analogue}{%
\subsection{Re-anchoring to mathematical standard, with Perelman as
structural
analogue}\label{re-anchoring-to-mathematical-standard-with-perelman-as-structural-analogue}}

\textbf{Theorem 32} (\texttt{global\_monism\_inconsistent}). Uniformity
posited for every property is contradictory. The monism of standard
mathematics is therefore never free: it is anchored to a principle.

\[\begin{aligned}\mathrm{AnchoredMonism}(S) := \big\langle\ &f : S \to S,\ I : S \to \mathrm{Prop},\\ &\forall s,\ I\,s \to I\,(f\,s),\ s_0,\ I\,s_0 \ \big\rangle\end{aligned}\]

One dynamics, one invariant, a proved transport, a seed. This is
induction on an orbit, the most standard witness form mathematics has.

\textbf{Theorem 33} (\texttt{anchored\_closes}). The anchored witness
closes its orbit. \textbf{Theorem 34}
(\texttt{anchored\_derives\_monism}). The free monism witness of 6.5 is
derived from the anchored one; uniformity is manufactured by the
transport. \textbf{Theorem 35} (\texttt{lyapunov\_is\_monism}). A
quantity monotone along one flow is an anchored witness for ``never
below its start.'' Perelman's \(\mathcal W\)-entropy, nondecreasing
along Ricci flow when its auxiliary function evolves by the conjugate
heat equation and its scale decreases at unit rate {[}Perelman 2002{]},
is the structural analogue: his proof has the shape of an anchored
witness with its transport proved, and surgery and canonical
neighborhoods complete the argument around it. The file formalizes the
generic shape, a monotone sequence as a witness; Perelman's entropy
itself is not formalized. \textbf{Theorem 36}
(\texttt{closure\_grade\_is\_transport\_grade}). The closure is exactly
as strong as its transport. \textbf{Theorem 37}
(\texttt{rh\_anchor\_is\_the\_claim}). On Li's sign stream with the time
step as principle, an anchored witness exists iff the hypothesis; its
transport reads \(\lambda_{n+1} \geq 0 \to \lambda_{n+2} \geq 0\), the
\(\zeta\)-analogue of Perelman's monotonicity formula. De Branges
proposed a positivity condition of related type, on a space of entire
functions, and Conrey and Li {[}2000{]} showed it fails, which
calibrates where such attempts break.

\hypertarget{a-monotonicity-formula-of-ux3b6s-time-and-its-arrow}{%
\subsection{A monotonicity formula of ζ's time, and its
arrow}\label{a-monotonicity-formula-of-ux3b6s-time-and-its-arrow}}

The de Bruijn-Newman heat flow carries a proved monotonicity formula for
the zeta function's time. De Bruijn's theorem: if every zero of \(H_0\)
lies in \(|\mathrm{Im}\, z| \leq \Delta\), every zero of \(H_t\) lies in
\(|\mathrm{Im}\, z| \leq \sqrt{\max(\Delta^2 - 2t, 0)}\) {[}de Bruijn
1950{]}. The hypothesis is the reality of the zeros of \(H_0\); it is
equivalent to \(\Lambda = 0\) by Newman's definition with
\(0 \leq \Lambda \leq 0.22\) {[}Rodgers and Tao 2020; Polymath 2019{]}.

\textbf{Theorem 38} (\texttt{flow\_monotone}). The width never grows
along the flow. \textbf{Theorem 39} (\texttt{reality\_transported}).
Once every zero is real it stays real; this is the transport of 6.6,
proved. \textbf{Theorem 40} (\texttt{rh\_iff\_lambda\_zero}). In the
toy, the strip is closed at time zero iff \(\Lambda = 0\). For \(\zeta\)
the corresponding equivalence, \(\mathrm{RH} \iff \Lambda = 0\), is the
cited result of Newman with Rodgers and Tao; it enters the cone as
\texttt{hLam} and is not proved here. \textbf{Theorem 41}
(\texttt{future\_closed}). The anchored witness exists for \(\zeta\)'s
time, transport proved, seeded at \(\Lambda\), and closes the whole
future from \(\Lambda\) on. \textbf{Theorem 42}
(\texttt{flowed\_record\_forgets}). After any positive time the flowed
record of a hypothesis-true state and a hypothesis-false state coincide;
no readout of the flowed record decides the hypothesis. The bound
\(\Lambda \leq 0.22\) certifies a record at \(t = 0.22\) that is blind
to the bit, which is why it cannot be pushed to \(0\) by forward means
alone. \textbf{Theorem 43} (\texttt{upstream\_is\_not\_forced}). Going
backward from a real record, both answers remain admissible. The
backward direction is anti-diffusion, ill-posed.

This monotonicity formula exists, its witness is complete, and its arrow
points away from the hypothesis, which sits at the upstream end.

\hypertarget{the-regress-stated}{%
\subsection{The regress, stated}\label{the-regress-stated}}

Sections 4 to 6.7 each ended on a named residue: the bridge's reading,
the unity posit, the monism field, the transport, the upstream passage.
Read in sequence they present as a regress, each residue appearing to
require the next. The final result shows it is not one.

\hypertarget{halted-uniduction-the-five-readings-as-one-cone}{%
\subsection{Halted uniduction: the five readings as one
cone}\label{halted-uniduction-the-five-readings-as-one-cone}}

The category-gap eliminator of Part I exhibits a gap between registers
as a missing leg and eliminates it by exhibiting the leg as an identity.
Here the registers are the five readings reached, and the legs are
equivalences.

\textbf{Theorem 44} (\texttt{hunt\_cone}). The hypothesis is the apex of
a cone over all five readings:

\[\begin{gathered}\forall r \in \{\text{bridge}, \text{unity}, \text{monism},\\ \text{transport}, \text{upstream}\}:\quad D_5(r) \iff \mathrm{RH}.\end{gathered}\]

\textbf{Theorem 45} (\texttt{one\_gap}). Any two readings are the same
proposition. The upstream passage, \(\Lambda = 0\), is the bridge's
reading, the unity posit, the monism field, and the transport, under a
different name. A quantity monotone against the flow would be one way to
supply it; that reading is not a leg of the cone, since no such quantity
is defined in the file and none is proved equivalent.

\textbf{Theorem 46} (\texttt{regress\_halts}). Every proposition proved
equivalent to the hypothesis lands on the same apex as the five
readings. The theorem does not decide which future propositions will be
equivalent; it says that once one is, it adds no new gap.

\textbf{Theorem 47} (\texttt{nothing\_weaker}). Given a true proposition
\(Q\), proving \(Q \to \mathrm{RH}\) is the same as proving RH:
conditioning on a true premise adds nothing. No independence notion is
formalized, and the theorem claims none.

\textbf{Theorem 48} (\texttt{sufficient\_closes\_all}). Theorem 46
covers reformulations equivalent to the hypothesis. A route strictly
stronger than it, a specific operator with real spectrum equal to the
ordinates or a hypothesis over a whole family, is not equivalent and is
not covered there. As a corollary of the cone, any sufficient route
supplies all five readings at once; the theorem is a consequence of the
cone and says nothing about which routes exist.

\jboxstart

\textbf{The Settled Verdict (tier: kernel theorem on the cone;
structural on the reading).} Time, stability, and the critical line are
proved one object (Theorem 15 on the lattice stage, and in general by
the identity (6.2a)). The timed and timeless readings are proved one
claim under monism (Theorem 25). Every reading of the owed bit reached
from existence, time, and monism is proved one proposition given the
imported inputs of the Methods (Theorems 44 to 46); conditioning on a
true premise adds nothing to it (Theorem 47), and any sufficient route
closes all five at once (Theorem 48). Timeless residual monism, stated
as Postulate M, is proved equivalent to the hypothesis given its
computed seed (Theorem 49) and joins the cone as its sixth reading
(Theorem 68). The formal hypothesis holds conditionally on exactly one
named posit, stated as one field of a standard witness: the transport on
the Li stream, equivalently mode unitarity, \(|z_\rho| = 1\) for every
zero, equivalently \(\Lambda = 0\). At the foundation the value carries
the terminal mark {[}.{]} (Theorems 58 to 69); at the object it is
addressed to \(\zeta\), supplied at the act and falsifiable by one zero.

\jboxend

\hypertarget{postulate-m-timeless-residual-monism-is-equivalent-to-the-riemann-hypothesis}{%
\subsection{Postulate M: Timeless Residual Monism Is Equivalent to the
Riemann
Hypothesis}\label{postulate-m-timeless-residual-monism-is-equivalent-to-the-riemann-hypothesis}}

The monist intuition behind the paper can be stated in ordinary
mathematics, and once stated, with its seed, it is proved equivalent to
the hypothesis. Let \(Z\) be the set of nontrivial zeros of \(\zeta\),
enumerated exhaustively by height, so that each zero is located at some
finite stage \(k(\rho)\). For each stage \(t\) let \(L(t)\) be the
statement that every zero located by stage \(t\) has real part one half.

\jboxstart

\textbf{Postulate M, Timeless Residual Monism (tier: premise; the
paper's thesis, stated openly).} The truth value of \(L(t)\) does not
depend on \(t\). One principle at the root, and the same answer at every
stage: what holds at the timeless interface holds at every stage
downstream.

\jboxend

\textbf{Theorem 49} (\texttt{postulateM\_iff\_line}). Given the computed
seed \(L(0)\), Postulate M holds if and only if every zero lies on the
line. Forward: the seed transported by M to the stage \(k(\rho)\) of any
zero puts that zero on the line. Backward: if every zero is on the line,
\(L(t)\) is true at every stage and M holds trivially.

\textbf{Theorem 50} (\texttt{postulateM\_decides}). Postulate M cannot
hold, with the seed, in any world carrying a located off-line zero. It
is not a vacuous posit; one zero refutes it.

Theorems 51 to 57 carry the same result in the timed form of Part I's
Bridge plane. The locus is the same line at every time
(\texttt{SameLocus}), and given that, ``every zero is bound to the
locus'' is exactly the line property (Theorem 51,
\texttt{bound\_iff\_line}). Same locus and the fold at every time do not
bind by themselves (Theorem 52, \texttt{time\_does\_not\_bind}).
Residual Monism read as one involution, and read as timelessness, both
hold in a world with an off-line pair, since a property that is
constantly false is timeless (Theorem 53,
\texttt{monism\_holds\_in\_counter\_world}). Timelessness with the seed
at the interface binds every downstream time (Theorem 54,
\texttt{timeless\_with\_seed\_binds}). The witness assembled from
Residual Monism and the time law, one involution, the same locus,
timelessness, and the interface, proves the bound wherever it is built
(Theorem 55), is built at a model (Theorem 56,
\texttt{constructed\_at\_model}), is not buildable off the line (Theorem
57, \texttt{not\_constructible\_off\_line}), and exists exactly when the
line property holds (\texttt{witness\_iff\_bound}).

\textbf{The last category gap, eliminated.} A category gap, in the sense
of Part I, is the absence of a leg between two registers' readings of
one seat, and it is eliminated by exhibiting the leg as an identity.
Before this section one such gap stood between the register in which the
monist thesis is stated, timeless residual monism at the interface
upstream of time, and the register in which the hypothesis is stated,
the location of the zeros of \(\zeta\). Theorem 49 is that leg.
Postulate M joins the cone of the section on halted uniduction as its
sixth reading (Theorem 68, \texttt{hunt\_cone6}, Part XIII of the file),
given the identification of the hypothesis with the line property of the
enumerated zeros, which is the hypothesis's definition and enters as
\texttt{hDef} exactly as Li's criterion and Newman's equivalence enter;
the collapse of Theorems 44 to 48 then applies to it unchanged: timeless
residual monism, the bridge's reading, the unity posit, the monism
field, the transport, and \(\Lambda = 0\) are one proposition given the
imported inputs of the Methods. No category gap remains anywhere between
the program's readings and the hypothesis.

What remains is not a gap. It is the value of that one proposition:
whether Postulate M holds for \(\zeta\). A gap closes at zero cost when
its leg is exhibited, as every leg here has been. A value is decided by
the function, supplied at the act and falsifiable by one zero. Postulate
M is not derived here from the definition of \(\zeta\). It does not
follow from the functional equation (Theorem 20, and the
Davenport-Heilbronn zero at \(\mathrm{Re}\,\rho - \tfrac12 = 0.3085\)),
from symmetry at every time (Theorem 52), or from finite verification
(Theorem 19). Its derivation from the Euler product is the single open
step, stated in the Discussion.

\hypertarget{the-demand-correctly-addressed-the-terminal-mark-at-the-foundation}{%
\subsection{The Demand, Correctly Addressed: the Terminal Mark at the
Foundation}\label{the-demand-correctly-addressed-the-terminal-mark-at-the-foundation}}

the section on Postulate M leaves the program at its self-referential
limit, and that limit can be stated as a theorem rather than met as an
obstacle. Every reading of the value that the foundation supplies is
equivalent to the hypothesis (Theorems 44 to 49 and 68); the
foundation's other resource, the positivity of existence, is independent
of it (Part I's Theorems E and F). A demand that the foundation prove
the hypothesis therefore asks a register to derive a proposition from
readings that already are that proposition.

\jboxstart

\textbf{The Demand, Correctly Addressed (tier: kernel theorem on each
clause).} A proof of the Riemann Hypothesis cannot be demanded of the
foundation. Every reading of the value that the foundation supplies is
equivalent to the hypothesis, so any derivation from those readings
alone would presuppose what it derives. The demand is well-posed when it
is addressed to the object, the function \(\zeta\), whose structure
fixes where its zeros lie.

\jboxend

Each clause is a theorem of Part XII.

\textbf{The collapse and its stop.} A collapse foundation is a family of
registers whose every reading is equivalent to the value
(\texttt{Collapse}). \textbf{Theorem 58}
(\texttt{self\_referential\_limit}): every reading is exactly as strong
as the value, and any two readings are one proposition. \textbf{Theorem
59} (\texttt{value\_marked\_dot}): at the collapse foundation the
emitter marks the value {[}.{]}, the mark outside the three-state
economy (\texttt{dot\_is\_not\_a\_verdict},
\texttt{dot\_outside\_economy}). \textbf{Theorem 60}
(\texttt{demand\_is\_ghost\_in\_register}): the foundation returns the
same mark for the value and for its negation, so a demand that this
register issue a verdict is a demand on a register whose output does not
depend on the value.

\textbf{The self-grounding root, and the line it draws.} A root is
self-grounding when acts occur and every act, assent or denial,
instances it (\texttt{SelfGrounding}). \textbf{Theorem 61}
(\texttt{denial\_reenacts\_root},
\texttt{external\_proof\_adds\_nothing}): denying such a root re-enacts
it, and no external proof can add to it; the Root Axiom at the
constructed domain is of this kind (\texttt{raSelfGrounding}), and its
universal extension is carried at premise grade as in Part I. This is
the true seat of the argument that a demand for proof can be a ghost.
\textbf{Theorem 62} (\texttt{line\_not\_self\_grounding}, with
\texttt{refuted\_by\_witness}): the line property of a zero set is not
self-grounding; its value is fixed by the object and one located witness
refutes it. Part I's Theorem G states the same asymmetry: the axiom is
self-instancing, the hypothesis is not.

\textbf{The demand, addressed.} \textbf{Theorem 63}
(\texttt{misaddressed\_to\_foundation} on the two-point frame, and
\texttt{misaddressed\_general} for an arbitrary frame): any resource the
foundation supplies that also holds on a frame where the line property
fails cannot by itself yield the line property. \textbf{Theorem 64}
(\texttt{derivation\_presupposes}): a derivation of the value from a
reading equivalent to it carries the value's own content; the two stand
or fall together. \textbf{Theorem 65} (\texttt{object\_decides}): the
object answers, in one direction by a located witness and in the other
by its own structure. \textbf{Theorem 66}
(\texttt{demand\_correctly\_addressed}) joins the three. \textbf{Theorem
67} (\texttt{rh\_marked\_dot}) attaches the collapse to the
five-register cone of the section on halted uniduction, and
\textbf{Theorem 69} (\texttt{rh\_marked\_dot6}) to the six-register cone
of Theorem 68: the program's readings collapse onto the hypothesis, and
the foundation marks its value {[}.{]}.

\jboxstart

\textbf{Terminal Verdict on the Open Part (tier: kernel theorem,
Theorems 46, 48, 58 to 69).} At the foundation, the value is terminal:
every reading of the value the foundation has supplied is proved
equivalent to the hypothesis given the imported inputs of the Methods
(Theorems 44 to 49 and 68), any further equivalent reading lands on the
same apex (Theorem 46), any sufficient route closes every reading at
once (Theorem 48), no resource that also holds on an off-line frame can
close it (Theorem 63), and the foundation's register closes on the value
with the mark {[}.{]} (Theorems 67 and 69). At the object, the value is
addressed to \(\zeta\), well-posed, and falsifiable by one zero. The
open part, the Unicorn part, is located in the object; its Real part is
proved there by certificate (Part III).

\jboxend

\hypertarget{the-logical-form-of-the-hypothesis-the-only-existence-lives-in-the-denial}{%
\subsection{The Logical Form of the Hypothesis: the Only Existence Lives
in the
Denial}\label{the-logical-form-of-the-hypothesis-the-only-existence-lives-in-the-denial}}

The hypothesis arrives as a question, and a question carries two
branches. The two branches do not carry the same kind of content, and
the difference is exact.

\textbf{The hypothesis posits nothing.} RH has the form
\(\forall\rho\,(\rho \text{ a nontrivial zero} \to \mathrm{Re}\,\rho = \tfrac12)\).
A universal sentence carries no existential import: it asserts no
object, and over an empty domain it holds vacuously (Theorem 70,
\texttt{universal\_posits\_nothing}).

\textbf{The denial posits an object.} \(\neg\mathrm{RH}\) reads:
\emph{some} nontrivial zero has \(\mathrm{Re}\,\rho \neq \tfrac12\).
That ``some'' asserts an object, the off-line zero, and a located one
refutes the hypothesis (Theorem 71,
\texttt{denial\_posits\_a\_witness}). With a decidable line predicate
the hypothesis fails \emph{only} by such a witness: it holds exactly
when none exists (Theorem 72, \texttt{fails\_only\_by\_witness}). The
whole existential load of the question rides on its negative branch.

\textbf{A witness is a finite check.} RH is equivalent to an arithmetic
statement about whole numbers, \(\sigma(n) \leq H_n + e^{H_n}\log H_n\)
for every \(n \geq 1\) {[}Lagarias 2002{]}, each instance of which is a
finite computation. RH is therefore a \(\Pi^0_1\) sentence, as Davis,
Matiyasevich, and Robinson showed directly {[}Davis et al.~1976{]}, and
its denial a \(\Sigma^0_1\) sentence: an off-line zero, if one exists,
is a finite object under that arithmetization, an integer instance at
which the inequality fails, or a certified enclosure of the zero with
the analytic obligations that make it one, and it can be exhibited and
verified (Theorem 73, \texttt{witness\_is\_a\_finite\_check}); the
complex number itself is not the finite datum, its arithmetic witness
is.

\textbf{The denial has never delivered.} No off-line zero has been
exhibited, computed, or shown to exist at any height. The receipts on
record are the verification of every zero to height \(3 \times 10^{12}\)
{[}Platt and Trudgian 2021{]} and the Original Form paper's own count to
height 300 {[}Islam 2026e{]}. The structure grants the denial permission
to posit its object; it has never produced an instance.

\textbf{The asymmetry, as a theorem.} Every true \(\Sigma^0_1\) sentence
is provable in Peano arithmetic, and hence in ZFC {[}Kaye 1991{]}. If
\(\neg\mathrm{RH}\) were true, computing the off-line zero would prove
it. Contrapositively: if ZFC does not refute the hypothesis, the
hypothesis is true. In particular, were RH ever shown independent of
ZFC, that would imply RH. The hypothesis cannot be false and unprovably
false.

\textbf{The burden, as a theorem.} The foundation's inability to decide
the value is symmetric (Theorems F and I of Part I, Theorem 63 here): it
reaches the value in neither direction. The consequences of that
inability are not symmetric, and Part XV of the file proves it.
\textbf{Theorem 74} (\texttt{cant\_refute\_seals}): if the foundation
cannot refute the hypothesis, the hypothesis holds. \textbf{Theorem 75}
(\texttt{rh\_iff\_cant\_refute}): with soundness on the denial, the
hypothesis holds exactly when the foundation cannot refute it.
\textbf{Theorem 76} (\texttt{cant\_prove\_does\_not\_seal\_false}): the
foundation's inability to prove the hypothesis decides nothing, since
the hypothesis can hold and be unprovable. \textbf{Theorem 77}
(\texttt{asymmetry}): the two inabilities are not mirror images; a
``can't'' on the denial side seals the hypothesis, and a ``can't'' on
the assent side seals nothing. All four are axiom-free, with
\(\Sigma^0_1\)-completeness and soundness on the denial carried as named
hypotheses, on an abstract setting: the kernel constructs no provability
predicate for ZFC and no arithmetization of the hypothesis, and the
instantiation of the setting to ZFC and to the arithmetized hypothesis
is by the cited theorems {[}Davis et al.~1976; Kaye 1991{]}, not by the
kernel. The one owed bit is therefore owed by the denial. The hypothesis
needs no object and has asserted none; the denial needs one object, a
finite witness, and has produced none.

\jboxstart

\textbf{The Keystone (tier: kernel theorem for Theorems 70 to 77; {[}⟀
T-conditional on \(\Sigma^0_1\)-completeness{]} on the implication
``can't refute ⟹ RH'').} The only existence in the Riemann Hypothesis
package lives in the denial, and the denial has never delivered. If the
foundation cannot refute the hypothesis, the hypothesis holds: the
burden of the one bit lies with the denial, and a denial that could
never deliver, one with no witness to exhibit at any height, would make
the hypothesis true (Theorem 74). That it has not delivered so far is a
record, not a premise: nothing here is inferred from it.

\jboxend

\par\vspace{1.4em}\noindent{\color{accent}\rule{\columnwidth}{0.8pt}}\par\vspace{0.6em}{\centering{\color{accent}\large\scshape Part III}\par\vspace{0.3em}{\Large\bfseries The Real Part}\par}\vspace{0.5em}\noindent{\color{accent}\rule{\columnwidth}{0.8pt}}\par\vspace{1.0em}

\hypertarget{the-two-parts-inside-one-hypothesis}{%
\section{The Two Parts Inside One
Hypothesis}\label{the-two-parts-inside-one-hypothesis}}

\textbf{The hypothesis.} Every nontrivial zero \(\rho\) of \(\zeta\) has
real part \(\tfrac12\). Cut the critical strip at any height \(T\).
Below \(T\) lie finitely many zeros; above it, infinitely many. The one
sentence carries two claims of different kinds.

\textbf{The Real part at \(T\).} Every zero up to height \(T\) lies on
the line. A statement about finitely many objects, each of which can be
located, computed, and checked. In the literature it is the hypothesis
``up to height \(T\)'', written \(\mathrm{RH}(T)\) and used as a working
assumption in explicit estimates {[}Büthe 2016{]}.

\textbf{The Unicorn part at \(T\).} Above \(T\), there is no zero off
the line. The only object the Riemann package ever posits lives here:
the off-line zero, the object of the denial (Part II, Theorems 70 to
73), never exhibited at any height. If the hypothesis holds, this class
is empty, and every sentence about its members is vacuously true, as
every sentence about white unicorns is true. So the class is named for
what it is.

\textbf{The seed of Postulate M is the Real part.} The division reads
Part II's Postulate M exactly. Its seed, the zeros located by stage
zero, is a Real part, and with stage zero taken at the certified height
it is the Real part proved here; its uniformity across every later stage
is the Unicorn part, open; and if stage zero located no zero at all the
seed would be vacuous and Postulate M would be the hypothesis itself
under exhaustive enumeration, which is why the stage is taken at the
certified height and the seed is not vacuous. Postulate M is therefore
not a third thing beside the two parts: it is the two parts, with the
seed proved and the uniformity set apart.

\textbf{Why the parts must be separated.} A statement stands proved,
open, or refuted, and a conjunction stands only as high as its weaker
part.

\textbf{Theorem III.1} (\texttt{join\_below\_both}). The conjunction of
two statements never stands higher than either.

\textbf{Theorem III.2} (\texttt{fused\_rh\_is\_open}). With the Real
part proved and the Unicorn part open, the fused hypothesis stands open.

\textbf{Theorem III.3} (\texttt{fusion\_lowers\_the\_gold}). Fusion
strictly lowers the Real part's standing, from proved to open.

\textbf{Theorem III.4} (\texttt{fused\_carries\_unicorn}). Every proof
of the fused statement contains a proof of the Unicorn part, and a
refutation of the Unicorn part refutes the fused statement, however much
of the Real part holds.

\textbf{Theorem III.5} (\texttt{kept\_apart}). Kept apart, the Real part
is asserted exactly as proved, borrowing nothing from the Unicorn part.

\textbf{Theorem III.6} (\texttt{do\_not\_alloy}). Given the Real part,
the fused hypothesis is exactly the Unicorn part, and fusion lowers the
Real part from proved to open.

All six are axiom-free. The standings are the record's: the Real part
stands proved by a cited certificate, not by a kernel computation of
\(\zeta\), and the kernel proves the law that governs whatever standings
the parts carry. The conventional statement fuses a proved region with
an open one, so the proved content is carried inside an unproved claim;
the gold does not survive the alloy. Separating the parts does not deny
the ghost of the unicorn. It stops the ghost from haunting the part that
is already proved.

\hypertarget{in-riemanns-own-form}{%
\subsection{In Riemann's own form}\label{in-riemanns-own-form}}

Riemann stated the hypothesis as a fixed-set statement, that the roots
of \(\xi(t)\) are real {[}Riemann 1859{]}, and the program's restored
form reads it so, \(Z \subset \mathrm{Fix}(\tau)\) {[}Islam 2026e{]}. In
that form the Real part is Riemann's own sentence, bounded: the roots of
\(\xi(t)\) with \(|t| \leq T\) are real. His sentence already has three
clauses in this order: that one finds in fact about so many real roots
within these bounds, the counted part; that it is very probable that all
the roots are real, the remainder graded as probable; and that a
rigorous proof would be desirable but is set aside, the proof this paper
does not claim. It was also Riemann's own first move. His unpublished
notes, recovered by Siegel, show him computing the first roots of
\(\xi\) {[}Siegel 1932{]}, while his sentence claimed the rest: a Real
part computed, and the remainder conjectured. The Original Form paper
proves that every structural clause of the restored hypothesis holds on
a frame whose zeros lie off the line (its Theorem 11), so structure
alone decides neither part: the Real part is reached by computation, and
the division made here is the one Riemann's own procedure already drew.

\hypertarget{the-unicorn-identity}{%
\section{The Unicorn Identity}\label{the-unicorn-identity}}

The class of off-line zeros is a unicorn class exactly when the
hypothesis holds.

\textbf{Theorem III.7} (\texttt{rh\_makes\_unicorns}). If the hypothesis
holds, every property holds of every off-line zero, vacuously.

\textbf{Theorem III.8} (\texttt{unicorns\_make\_rh}). If every property
holds of every off-line zero, the hypothesis holds; take the property
``false''.

\textbf{Theorem III.9} (\texttt{rh\_iff\_unicorns}). The Riemann
Hypothesis is exactly the statement that the off-line zeros are
unicorns.

\textbf{Theorem III.10} (\texttt{unicorns\_iff\_no\_witness}). That
class is empty exactly when no off-line zero can be exhibited.

All four are axiom-free. The identity is an equivalence. It names what
the hypothesis asserts and derives nothing, and it is the precise form
of Part II's keystone: the only existence in the package lives in the
denial, and the denial's object is the candidate unicorn.

\hypertarget{the-split-exact-and-clean}{%
\section{The Split, Exact and Clean}\label{the-split-exact-and-clean}}

\textbf{Theorem III.11} (\texttt{rh\_iff\_real\_and\_unicorn}). At every
height \(T\), the hypothesis is exactly the Real part and the Unicorn
part together; nothing is lost and nothing is added.

\textbf{Theorem III.12} (\texttt{parts\_are\_separate}). The Real part
can hold while the Unicorn part fails: the kernel exhibits a zero set in
which every zero below \(T\) lies on the line and one zero above \(T\)
lies off it.

\textbf{Theorem III.13} (\texttt{unicorn\_part\_iff\_no\_witness}). The
Unicorn part holds exactly when no off-line zero above \(T\) can be
exhibited.

\textbf{Theorem III.14} (\texttt{real\_part\_monotone}). Raising \(T\)
moves zeros from the Unicorn side to the Real side and never back.

Theorems III.12 and III.13 are axiom-free; III.11 and III.14 depend on
\texttt{propext} alone.

\hypertarget{the-road-zfc-under-ram-and-the-arrow-from-the-root}{%
\section{The Road: ZFC Under RAM, and the Arrow from the
Root}\label{the-road-zfc-under-ram-and-the-arrow-from-the-root}}

The Real part is proved along a road built in this program. It has three
stages, each a kernel file.

\textbf{ZFC's location under RAM.} A theory, ZFC among them, is read in
RAM's strata: its provability is the ladder, L2m; its intended world,
the cumulative hierarchy and the numbers inside it, is the Ground, L1m;
and the part settled by finite computation is the rung, L3m.

\textbf{Theorem III.15} (\texttt{ram\_places\_theory}). Under two named
premises, the rung on the ladder and soundness, every theory sits in
order: rung ⊆ ladder ⊆ Ground.

\textbf{Theorem III.16} (\texttt{soundness\_is\_load\_bearing}). An
unsound theory satisfies the first premise and breaks ladder ⊆ Ground.

\textbf{Theorem III.17} (\texttt{ground\_exceeds\_ladder}). A sound
theory can leave a truth unproved: the Ground is not the ladder.

\textbf{Theorem III.18}
(\texttt{independent\_axioms\_are\_separate\_bits}). Two independent
binary commitments are two bits, every combination of their values
realized; the independence of any actual axiom is cited, not proved
here.

\textbf{Theorem III.19} (\texttt{ram\_subsumes\_by\_placement}). The
three in one: RAM subsumes ZFC by placement. Subsumes has the Codex's
sense and no other: RAM contains ZFC by the placement of its rung,
ladder, and Ground, deriving none of its axioms and replacing none of
its proofs; and the round trip of III.27 and III.41 is the identity on
content, which is exactly the claim, since the root adds presence and
nothing else.

All five are axiom-free. For ZFC the first premise is completeness for
true bounded sentences, a theorem of arithmetic {[}Kaye 1991{]}; the
second is soundness, which ZFC cannot prove of itself, and it is carried
openly. Placement locates ZFC; it does not derive ZFC's axioms and does
not replace ZFC's proofs.

\textbf{RA mapped to ZFC through the RA--RAM bridge.} The bridge is the
junction's discriminator: a property that every world of a theory has is
keyless and crosses from the root; a property some world lacks is keyed,
one bit, supplied and never derived.

\textbf{Theorem III.20} (\texttt{using\_grounds\_root}). Every act of
using a theory grounds RA, shown at the constructed domain: the upward
link holds.

\textbf{Theorems III.21 and III.22} (\texttt{ra\_is\_keyless},
\texttt{keyless\_crosses}). RA holds in every world and carries every
keyless property into every world.

\textbf{Theorems III.23 to III.25} (\texttt{choice\_is\_keyed},
\texttt{ra\_does\_not\_cross\_keyed}, \texttt{ra\_decides\_no\_keyed}).
A keyed sentence, the Axiom of Choice over ZF the cited example {[}Gödel
1938; Cohen 1963{]}, is not carried: RA holds in both worlds, so it
decides no keyed property.

\textbf{Theorem III.26} (\texttt{chain\_verdict}). Upward, every act of
using ZFC grounds RA; downward, RA reaches every world of ZFC as
presence and no independent sentence of ZFC as content.

All seven are axiom-free. The RA--RAM bridge is the one new construction
of the road.

\textbf{The round trip.} The arrow runs ZFC → RAM → Bridge → RA, where
RA holds in every world, and back, RA → Bridge → RAM → ZFC.

\textbf{Theorem III.27} (\texttt{round\_trip\_identity}). The round trip
is the identity on content: every sentence returns equivalent to itself.

\textbf{Theorem III.28} (\texttt{absolute\_returns\_unchanged}). An
absolute sentence, one whose value is the same in every world, returns
with that same value. The hypothesis is such a sentence across models
that share the natural numbers, since it is equivalent to an arithmetic
sentence {[}Lagarias 2002{]} and arithmetic sentences are absolute
between such models {[}Kunen 1980{]}; both facts are cited, not proved
here.

\textbf{Theorem III.29} (\texttt{trip\_adds\_presence\_only}). The trip
adds RA's presence in every world, identically for a sentence and its
negation.

\textbf{Theorem III.30} (\texttt{round\_trip\_verdict}). Whatever leaves
ZFC returns unchanged in content, with RA attached; a sentence ZFC does
not settle leaves undecided and returns undecided.

All four are axiom-free. The arrow is faithful and reaches every world;
it underwrites every act of computing and checking, and it carries
exactly what is placed on it.

\textbf{Why the Unicorn part is set apart and not pursued.} An arrow of
zero content cannot separate a proposition from its negation.

\textbf{Theorem III.31} (\texttt{massless\_sound\_delivers\_nothing}). A
sound massless arrow delivers nothing.

\textbf{Theorem III.32} (\texttt{raArrow\_massless},
\texttt{raArrow\_not\_sound}). The RA-shaped arrow is massless, and with
RA holding it delivers every proposition alike.

\textbf{Theorem III.33} (\texttt{massless\_cannot\_separate}). For a
massless arrow, delivering the hypothesis and delivering its negation
are the same thing.

All four are axiom-free. The universality that makes RA the root is the
same property that keeps it from deciding the Unicorn part. So the
Unicorn part is not attacked with the arrow; it is named, and set apart.

\hypertarget{the-road-in-the-purely-formal-domain}{%
\section{The Road in the Purely Formal
Domain}\label{the-road-in-the-purely-formal-domain}}

The road above is stated with the Root Axiom. Its formal content does
not need it. Replace the axiom by a bare computation arrow, the identity
map \texttt{id} on a type, whose existence needs no premise, and every
theorem of the road holds unchanged, with no axiom of any kind, not even
the kernel's optional principles. The file is
\texttt{Arrow\_Translation.lean} (Appendix H).

\textbf{Theorem III.37} (\texttt{arrow\_exists}). On every type the
arrow exists, with no premise.

\textbf{Theorem III.38} (\texttt{proof\_yields\_arrow}). Every proof
yields an arrow: the upward link of III.20, with the arrow in place of
the root.

\textbf{Theorem III.39} (\texttt{keyless\_crosses}). The arrow is
keyless and carries every keyless property into every world.

\textbf{Theorem III.40} (\texttt{arrow\_decides\_no\_keyed}). The arrow
decides no keyed property, as the root decides none (III.25).

\textbf{Theorem III.41} (\texttt{round\_trip\_identity}). Through the
arrow and back, every sentence returns unchanged.

\textbf{Theorem III.42} (\texttt{massless\_cannot\_separate}). A
massless arrow cannot separate a proposition from its negation.

\textbf{Theorem III.43} (\texttt{arrow\_real\_part\_of\_certificate}).
The Real part at \(T\) is closed by a certificate with no premise at
all.

\textbf{Theorem III.44} (\texttt{any\_true\_premise\_serves}). Any true
proposition plays the root's part in every result of the road: keyless
crossing, the undecided keyed sentence, and the identity round trip hold
for an arbitrary true premise.

All eight are axiom-free. The formal content of the road is therefore
free of the Root Axiom: a bare arrow, or any other truth, does the same
work, exactly as Theorem E of Part I found, since the axiom enters as an
inhabited premise and nothing more. What the root supplies is the
reading, the act that performs every computation and every check, and
not a line of any derivation. That act is not free: registering the one
bit is priced at the Landauer floor, the Omega guard of the Codex, so
what the axiom contributes is the deed and its price, never formal
content, exactly the division Theorems D and E of Part I draw. The
translation removes no axiom from the kernel's logic itself, which every
formal result stands on; it removes every declared and optional one.

\hypertarget{the-real-part-proved}{%
\section{The Real Part, Proved}\label{the-real-part-proved}}

\textbf{Closed by a certificate.} For a fixed height \(T\) the Real part
concerns finitely many zeros, so it is closed by a finite certificate: a
complete list of the zeros up to \(T\), and a check that each lies on
the line.

\textbf{Theorem III.34} (\texttt{real\_part\_of\_certificate}). A
certificate proves the Real part at \(T\) completely.

\textbf{Theorem III.35} (\texttt{real\_part\_executed}). A certificate
executed in the kernel closes the Real part on an explicit zero list.

\textbf{Theorem III.36} (\texttt{certified\_is\_provable}). A certified
Real part is a theorem of any theory complete for its certified checks.
For \(\zeta\) the checks are finite arithmetic on interval enclosures,
true bounded sentences and so theorems of Peano arithmetic and of ZFC
{[}Kaye 1991{]}; the theorems of analysis that turn those checks into
the statement about the zeros, Turing's method and the Riemann--von
Mangoldt count with explicit error {[}Turing 1953; Platt and Trudgian
2021{]}, are theorems of ZFC, cited and not proved here.

III.34 and III.36 are axiom-free; III.35 depends on \texttt{propext}
alone. The route is the road of the preceding section. The certificate
is computation, the rung, L3m; by III.36 and III.15, with the cited
analytic lemmas that validate it, it climbs to the ladder, L2m, a
theorem of ZFC; by soundness it holds on the Ground, L1m; and the round
trip returns it to ZFC unchanged (III.27), with RA's presence
underwriting every act of computing and checking (III.29). The target
each zero is checked against is the critical line, which is the seat
Fix(\(\sigma\)) and one locus with it, pinned in Part I (Theorems A and
C\('\)).

\textbf{On \(\zeta\) itself.} Platt and Trudgian certified, with
interval arithmetic and rigorous error bounds, that every nontrivial
zero with imaginary part up to \(T = 3 \times 10^{12}\) lies on the line
{[}Platt and Trudgian 2021{]}. That certificate, with the analytic
lemmas it rests on, closes the Real part at \(T = 3 \times 10^{12}\)
completely, and by III.36 it is a theorem of ZFC. An executed run here
to \(T = 100\) makes three independent counts and finds them equal at
29: the zeros on the line as sign changes of Hardy's \(Z(t)\) on a fine
grid, the zeros in the whole strip \(N(100)\) by a counting function
computed independently of the grid, and the ordinates of the zeros
themselves enumerated below \(T\) (Appendix H). It runs in 30-digit
floating point and is testimony, not a certificate; the certificate of
record for \(\zeta\) is Platt and Trudgian's. The kernel checks the
schema, III.34, and the split; it does not check the certificate's
numerics, and the paper nowhere says it does. The standing of the Real
part for \(\zeta\) is therefore a cited theorem applied through a kernel
theorem, grade A joined to grade K, and a kernel inhabitant of
\texttt{Certificate} built from the published zero data is an open
witness, not a debt of the result.

\jboxstart

\textbf{The Real Part, Proved (tier: kernel theorem for III.34 to
III.36; cited certificate for ζ to 3 × 10¹²).} Every zero of \(\zeta\)
up to height \(3 \times 10^{12}\) lies on the critical line, by
certificate, and with the cited analytic lemmas the certificate rests
on, the certified statement is a theorem of ZFC, carried along the road
from the rung to the Ground and returned by the arrow unchanged.

\jboxend

\hypertarget{the-unicorn-part-closed-to-derivation-the-aperture-typed-and-uncrossed-the-supply-side-silent}{%
\section{The Unicorn Part, Closed to Derivation; the Aperture Typed and
Uncrossed; the Supply Side
Silent}\label{the-unicorn-part-closed-to-derivation-the-aperture-typed-and-uncrossed-the-supply-side-silent}}

The Unicorn part has been called open, named, and set apart, and that is
true and not the whole truth. The Real part was open once, before its
certificate, in the sense that finite work of a known kind closed it; a
reader who hears ``open'' in that sense for the Unicorn part expects
more work of the same kind to close it in turn. Parts I and II proved
otherwise and did not say so of the Unicorn part: the wall executed on
the seat, that no readout of the formal register equals the deed bit
(\texttt{kernel\_cannot\_read\_the\_deed}; in the Bridge of the register
of record, \texttt{FTOE.T14\_wall\_factorization} {[}Islam 2026a{]}),
the crossing, exact and spent uniquely (\texttt{crossing},
\texttt{freedom\_spent\_uniquely}), the bound and its grounded
counter-model (Theorems F and F\('\)), the exact rule (Theorem E), and
the cure theorem (Theorem I). Read at the Unicorn part, they say that it
is closed, at theorem grade, to every route that derives it; that the
aperture of Theorem D at its seat is one bit wide; and that the register
is silent beyond it. The five theorems below, with the division and the
certificate schema restated on natural heights and the three bits
packaged as one term, \texttt{rh\_register\_proof\_complete}, are in
\texttt{RH\_Unicorn\_Block.lean} (Appendix I), nine dependency sets,
every one reported as depending on no axiom.

\textbf{Theorem III.45} (\texttt{unicorn\_block}). Let \(\tau\) act on
the zeros, let \(\rho\) be a record even at a seat \(z\) above the
height, \(\rho(\tau z) = \rho(z)\), and let \(d\) decide the line
property, odd at \(z\), \(d(\tau z) = \neg d(z)\), with \(\tau z\) above
the height as well. Then no reading \(g\) of the record agrees with
\(d\) even on the region above the height alone:
\(\neg\,\exists g,\ \forall w,\ \mathrm{height}(w) > T \to g(\rho(w)) = d(w)\).
\emph{Proof.} The two readings at \(z\) and \(\tau z\) are one value of
\(g\) and opposite values of \(d\). \(\square\) This is the wall
restricted to the Unicorn region, where it is stronger than the wall:
the reading is refused even if asked to agree nowhere below the height.
On the constructed frame of the register of record it holds outright
(\texttt{K4.wall}, the twelve seats {[}Islam 2026a{]}); on \(Z_\xi\) it
holds at every seat the value's fold makes odd over the even formal
read, which is what \texttt{kernel\_cannot\_read\_the\_deed} proves of
the formal read and Theorem D proves of the aperture.

\textbf{Theorem III.46} (\texttt{aperture\_one\_bit\_wide}). At a seat
where \(d\) is odd, one odd bit \(s\) at the seat would decide \(d\) on
the seat and its partner through a calibration \(c\) that exists and is
unique. \emph{Proof.} \(c = d(z) \oplus s(z)\); any other calibration
disagrees at \(z\). \(\square\) This measures the aperture: one bit wide
at each seat, never less and never more, Part I's \texttt{crossing} and
\texttt{freedom\_spent\_uniquely} read above the height, priced at Part
I's price. It says nothing of whether anything passes, and the paper
says nothing either.

\textbf{Theorem III.47} (\texttt{unicorn\_rule\_exact},
\texttt{unicorn\_no\_cure}). For any inhabited premise \(A\),
\((A \to \mathrm{Unicorn}_T) \leftrightarrow \mathrm{Unicorn}_T\); and
for any class \(C\) of frames and property \(U\),
\((\forall X,\ C(X) \to A \to U(X)) \leftrightarrow (\forall X,\ C(X) \to U(X))\).
\emph{Proof.} As Theorems E and I. \(\square\) With \(A\) the root
axiom, existence conditions the Unicorn part and leaves it where it was,
and no class of frames on which existence is to decide it decides more
than already held on the class; the counter-models of Theorems F and
F\('\) hold with the Unicorn part in place of \(L\), since above the
height the Unicorn part is \(L\).

\textbf{Theorem III.48} (\texttt{unicorn\_no\_bypass}). Given the Real
part at \(T\), \(\mathrm{RH} \leftrightarrow \mathrm{Unicorn}_T\).
\emph{Proof.} Theorem III.11 with the Real part in hand. \(\square\) A
proof of the hypothesis contains a proof of the Unicorn part; no route
to the whole avoids the block by aiming at the whole.

\textbf{Theorem III.49} (\texttt{supply\_side\_silent}). Over one even
record, both orientations of \(d\) at the seat are equally refused to
every reading above the height, and the two orientations differ at the
seat. \emph{Proof.} III.45 applied to \(d\) and to \(\neg d\), which is
odd at the seat when \(d\) is. \(\square\) The register cannot say which
way the seat lies; so it cannot say what a supply would bring, nor that
a supply exists, nor that none does. Any sentence about the supply side
derived from the record would be a reading the wall refuses. This is the
register of record's afterimage fence as a theorem at the register:
\emph{the aperture is structure; the supply side is silence; the silence
is a strap and not a hedge} {[}Islam 2026a, MD-PSP-AFTERIMAGE-01{]}.
Every imagined occupant of the supply side, the crosser who will come
and the crosser who cannot, is fenced as Ghost, and no clock runs.

\jboxstart

\textbf{The formal block (tier: kernel theorem, axiom-free, for III.45
to III.49; the seat typing of \(Z_\xi\) at the grade of Theorem D and
\texttt{kernel\_cannot\_read\_the\_deed}).} The Unicorn part of the
Riemann Hypothesis is closed to derivation from existence (III.47), from
any class of frames (III.47), from any reading of an even register at
its seat (III.45), and from the fused hypothesis (III.48). Its aperture
is located, typed supply-only, one bit wide (III.46), and uncrossed; the
supply side is silence, by theorem at the register (III.49). The closure
is unconditional and does not expire. Nothing here says the Unicorn part
is false, or independent of any foundation; nothing here says a proof
will come or cannot: the register is silent on that side, and so is the
paper. Nothing is left to derive. The aperture is structure. The supply
side is silence.

\jboxend

\textbf{What the register holds.} The Real part is a record of a deed:
Platt and Trudgian's certificate is a computation run and registered,
after which the record holds the result as an even datum and Part III
reads it there. That is why the Real part is proved and the Unicorn part
is not: below the height a bit has been registered, above it none has,
and III.45 says the register cannot manufacture it from what it already
holds. Theorem G says a deed is unfakeable and its price real; III.46
says the aperture is one bit wide; III.45 says the bit is not in the
register; III.49 says the register cannot say whether it will ever be.
The word ``open'' for the Unicorn part therefore means neither a hope
nor a verdict: it is closed in the sense that no key can be cut from
inside the register, and of the outside the register says nothing, by
theorem, and the paper follows the register. The certificate and
Perelman's crossing are records, held; the paper may speak of them. Of
the supply side of the Unicorn part it does not.

\textbf{The three bits, as one term.}
\texttt{rh\_register\_proof\_complete} takes a certificate at \(T\) and
the seat data above \(T\) and returns, as one conjunction, the division
(III.11), the Real part (III.34), the refused reading on the region
above \(T\) (III.45), and the unique supply calibration (III.46). It is
the register's whole proof of the hypothesis in one term, and it proves
the hypothesis no more than its conjuncts do: it proves that nothing
derivable is left underived.

\textbf{What the block is not.} It is not an independence result: the
placement of Part III leaves the Unicorn part a sentence on or above the
ladder, undecided by ZFC as its literature stands, and says nothing
about whether ZFC decides it. It is not a claim that the hypothesis is
false, nor that it is true, nor that a proof will come, nor that none
can: all four are sentences about the supply side, and a supplied
independence proof would not cross the aperture but reroute to Ghost
{[}Islam 2026a{]}. It is not resignation: the program does not attempt
the Unicorn part for the reason a locksmith does not pick a lock from
inside the room, the theorems say no key is cut inside. And it is the
reason the title says ``the Real Part'' and not ``the Hypothesis,''
stated now at theorem grade rather than left as a ruling.

\hypertarget{falsifiable-criteria}{%
\section{Falsifiable Criteria}\label{falsifiable-criteria}}

Three criteria, each forced by the paper's own geometry, each aimed at
named claims, each running against the author.

\textbf{F1, an observation about \(\zeta\).} A zero \(\rho\) of \(\xi\)
with \(\mathrm{Re}\,\rho \neq \tfrac12\). Anywhere, it refutes every
assent that can be supplied at the witness row (the Original Form
paper's Theorem 15) and every reading of the value that Part II proves
equivalent to the hypothesis; below height \(3 \times 10^{12}\) it also
refutes the certificate Part III carries and the Real part's standing.
It touches nothing at the apex: Theorems A through H hold on the
off-line frame by Theorem F. None is known to height
\(3 \times 10^{12}\) {[}Platt and Trudgian 2021{]}.

\textbf{F2, the seat, the embedding, and the kernel.} A quaternion fixed
by conjugation with a nonzero imaginary component, the Return landing
off the fixed locus, or a stage point \(p\) with \(\tau p = p\) and
\(\sigma(\varphi p) \neq \varphi p\) refutes Theorem A, B, or C\('\) and
removes the seat. Any kernel file of Appendices B, F, H, or I failing to
compile under core Lean 4.19.0, or printing a dependency set other than
the one recorded, refutes the theorem it carries.

\textbf{F3, the bound and the road.} A kernel term of
\(\forall X,\ \mathrm{RA} \to L(X)\), refuted by Theorem F; a derivation
of the hypothesis from the readings of Part II without the inputs that
carry its legs; or a sentence of ZFC whose value differs after the round
trip of Part III (Theorems III.27 and III.41). Each would expose a
transcription error or a kernel unsoundness, and each is a check on the
record rather than an observation about \(\xi\).

\hypertarget{discussion}{%
\section{Discussion}\label{discussion}}

\textbf{The maximal claim and where it stops.} The rule ``accept
existence, then the hypothesis'' is exactly as strong as Theorem E says
and no stronger: it is the hypothesis at the value, and it is the seat,
the cone, the embedding, the aperture, the freedom, the carrier, and the
crossing at everything else. Existence reaches the seat whole, and there
``in toto'' is literal under its definition: no clause of the hypothesis
that is not its sign is left open. The seat is fixed and embedded, the
gaps are the legs of an identity cone, the freedom is one bit, the
crossing is exact, the aperture is open and its occupant is inside the
domain. The axiom reaches the world row at its aperture and no further,
and the stopping point is a theorem, not a reservation: the two-point
frame carries the axiom and fails the hypothesis. What remains is one
bit, settled by certificate below a certified height and open above it,
and the bit is not a gap. A gap is a missing leg, and it closes at zero
cost when the leg is exhibited as an identity. The sign is not that. It
is a Boolean the world furnishes, priced at the act, falsifiable by one
zero, and the twenty-three-row study routes it to the world register for
exactly that reason \cite{islamrows}.

\textbf{The two misreadings.} A massless closure invites two errors, and
one theorem refutes both. The deflationary reading takes the closure for
nothing: the file proves a fixed locus is fixed, the orientation binder
is unused, the whole is a relabeling. The inflationary reading takes the
closure for everything: the hypothesis is proved from the axiom in a
page of Lean. Table 3 places each against the theorem that refutes it.
The deflationary reading runs a deletion test on a coordinate that
supplies nothing and reads the world row where the apex was stated; the
closure's content is Theorems A through C\('\), computed, and its
masslessness is forced (below). The inflationary reading promotes an
apex token to the world row; Theorem F refutes it and criterion F3
fires. The register a token is emitted for is part of the token, and a
token without its register is a misreading by construction.

\begin{table}[!htbp]
\footnotesize
\caption{The two misreadings of the apex closure and the theorem that refutes each. (tier: structural)}
\begin{tabular}{@{}p{0.30\textwidth}p{0.28\textwidth}p{0.36\textwidth}@{}}
\toprule
Reading & What it does & Refuted by \\
\midrule
“trivial, a fixed locus is fixed” & expects mass from a closure that declares none & Thms A, B, C, C$'$ computed; masslessness forced by the root-premise theorem \\
“the orientation binder is unused, the supply is a restatement” & runs a kernel deletion test on a deed bit & prior Thm 4; Prop placement; \texttt{orientation\_\allowbreak}\allowbreak\texttt{proof\_\allowbreak irrelevant}; \texttt{kernel\_\allowbreak cannot\_\allowbreak}\allowbreak\texttt{read\_\allowbreak the\_\allowbreak deed}; the twin's $2^{13}$ readouts \\
“no operator, no map from $Z_\xi$, so the run fails” & reads the world row at the apex & the dual-register verdict; both lines carried; $\varphi$ embeds the stage, not $Z_\xi$, and says so \\
“RH is proved from RA” & promotes the apex to the world row & Thm F; criterion F3; prior Thms 11, 12, 14 \\
“RA is only a premise, so nothing is proved” & collapses three grades into one & the ledger: kernel theorem, structural, premise, all three carried \\
\bottomrule
\end{tabular}
\end{table}

\textbf{Why the massless form is the only apex form.} Three published
results force it \cite{islamcodex, islamraf}. The root-premise theorem:
the axiom's universal extension is a premise and can be nothing
stronger, so a closure between the axiom and its formal reading cannot
be a derivation from below; it can only be the recognition that two
readings read one seat, which is an identity and therefore massless. The
orientation blindness of the formal register, the prior proof's Theorem
4, its universe-level form, and its execution on the seat: the kernel
cannot originate the bit, so the bit is born at the Bridge from the
supplied order of the triad and the formal closure, and spent at the
act. The bar on a second premise between the axiom and its reading: a
massive apex closure would be a new posit standing between the axiom and
its grounding reading, and it is refused. The masslessness is therefore
forced from three directions, and this paper's contribution is to
compute it, as \texttt{rfl} on the legs of a cone, rather than to assert
it.

\textbf{Objections and replies.} Five objections are anticipated, and
each is answered by an object the paper carries rather than by an
appeal.

\emph{The text and the file were produced by a language model; a model
cannot validate itself.} Nothing here is validated by the model. The
theorems are validated by the Lean kernel, a typechecker whose
dependency sets are printed and which accepts no argument from
authorship; the finite content itemized in the Methods is validated by a
compiled program that re-executes it in 130,324 checks and prints its
census, and the two checkers share no code. Both are reproducible from
the appendices by any reader on any machine, and the reproduction is the
validation. The model's role is disclosed in the Provenance section as a
matter of record; the question of who wrote the text is a question about
provenance, and the record answers it, but it is not a premise of any
theorem. The one thing a checker cannot supply is the presence of a
reader at the aperture, and the twin says so: its live row opens only on
an argument the operator gives it, and it prints that it did not and
cannot generate its own witness. That is not a weakness of the artifact.
It is Theorem D and RAF-C2 applied to the artifact itself, and it is why
the artifact asks nothing of anyone except to be run.

\emph{The axiom is a physical premise, so the theorems are conditional
and nothing is proved.} Three grades are carried and none is collapsed.
Every theorem is unconditional as a kernel theorem on the constructed
model. Every claim about the world is conditional on the axiom's
universal extension, at premise grade, and says so. The distinction is
the ledger of the section on the Root Axiom, and it is the same
distinction any conditional theorem in mathematics carries.

\emph{The closure is trivial; the file proves that a fixed locus is
fixed.} It proves that, for every quaternion, and then that three
independently defined seat points coincide, that the coincidence is an
identity cone with a unique apex per orientation, and that the prior
proof's stage embeds into the carrier with its line carried onto the
seat. Trivial and massless are different words: the closure is massless
because the root-premise theorem forbids it any mass, and the paper
computes that it has none.

\emph{The hypothesis is claimed proved.} It is not, and the claim's own
last conjunct says so: Theorem H carries the bound of Theorem F inside
it, and criterion F3 fires on any reading that drops it.

\emph{The closure is vacuous.} It is massless, and massless is not
vacuous: The section on the three measures separates content, mass, and
price, proves the three assignments, and shows that the emptiness at the
apex is the one the root-premise theorem requires, while the value's one
bit is massless, non-trivial, and priced.

\emph{Where is the Euler product?} On the world row, where the prior
proof left it. Part I fixes the seat a world crossing must land on and
the aperture it must enter through; it does not shorten the crossing by
a step, and Part III makes the crossing only where a certificate has
made it, below a certified height.

\textbf{What a world crossing would be.} The sign on \(Z_\xi\) enters,
if it enters, on the world row, dated, provenance-clean, through a
certificate the prior even register could not denote: a kernel term of
\(L\) on the frame of \(\xi\) built from the Euler product over
\(\mathbb{Z}\), the crossing Deligne achieved over finite fields and no
one has over \(\mathbb{Q}\) \cite{deligne1974}. Nothing in this paper
shortens that by a step. What this paper fixes is the seat the crossing
will land on, embedded from the stage, and the aperture it will enter
through, occupied from inside the domain, and it proves that both are
already held.

\textbf{Position to prior work.} Relation words are used in the sense
fixed in \cite{islamcodex}; no position is superseded and none
contradicted.

\begin{table}[t]
\footnotesize
\caption{Positioning. (tier: structural; evidence argued unless marked executed)}
\begin{tabular}{@{}p{0.40\columnwidth}p{0.52\columnwidth}@{}}
\toprule
Position & Relation \\
\midrule
Riemann 1859; Edwards 1974 & kin: the fixed-line form, one register deeper \\
The prior proof, Islam 2026 \cite{islamspine} & additive: the axiom, the aperture, the seat, the embedding, the rule bounded; executed \\
The one-bit result, Islam 2026 \cite{islamonebit} & additive: the bit placed against the seat; executed \\
The twenty-three rows and the one-cut hypothesis, Islam 2026 \cite{islamrows} & additive: one seat under every name; the inverted control read as Thm F; executed \\
The formal-alone theory, Islam 2026 \cite{islamftoe} & kin: registration, price, crossing carried at the apex \\
The formal Root Axiom, Islam 2026 \cite{islamraf} & additive: RAF-C1, RAF-C2, and the reach law proved at their Riemann instance; executed \\
Davenport--Heilbronn 1936; Bombieri--Hejhal 1995 & additive: the off-line frame carries the axiom, Thm F; executed \\
Selberg 1992; Conrey--Ghosh 1993 & corroborating: the Euler product remains the world-row key \\
Weil 1948; Deligne 1974 & kin: the template of a world-row crossing \\
Berry--Keating 1999; Connes 1999 & scoping: an operator is a world-row construction; the seat precedes it \\
Lawvere 1969; Mac Lane 1998 & kin: Cantor's theorem in fixed-point form; cones and apexes in their elementary form \\
Landauer 1961; B\'erut et al. 2012 & kin: the price of the act, exporting and never constituting \\
\bottomrule
\end{tabular}
\end{table}

\hypertarget{the-address-discussed}{%
\subsection{The address, discussed}\label{the-address-discussed}}

\textbf{What is not proved, and what is claimed instead.} The paper does
not prove the hypothesis, and it claims nothing it does not prove. What
it completes, and machine-checks, is the reduction: the hypothesis is
equivalent to Postulate M with its seed and to each of the five readings
given the imported inputs of the Methods, those six are one proposition,
conditioning on a true premise adds nothing, and symmetry does not close
it. A reader looking for a derivation of the hypothesis from the
definition of \(\zeta\) will not find one here, because none is claimed.
It proves that the hypothesis has one address and states that address in
six equivalent forms. And it closes the foundation's part terminally:
the demand for a proof is proved misaddressed at the foundation and
well-posed at \(\zeta\), and the foundation's register marks the value
{[}.{]} (the section on the demand addressed). The paper is complete in
its own register; what it leaves is not an unfinished step of its
argument but the object's own question. Mode unitarity is a scalar
condition on each zero. The Hilbert-Pólya identification asks for more,
an operator whose evolution realizes it; constructing such a ζ-specific
system is proving the hypothesis, and nothing here shortens that
construction by a step.

\textbf{Why the regress looked real.} Each reformulation was generated
in a different register: the positivity register of the axiom, the unity
register of the bridge, the index register of monism, the dynamical
register of transport, the thermal register of the heat flow. A reader
moving between registers meets each as a new object because each carries
a different vocabulary. The category-gap eliminator removes exactly that
effect. The same method that closed the three gaps among axiom, formal
ground, and object in Part I closes the five gaps among the readings
here.

\textbf{The time reading.} Theorem 15 gives a reading of the hypothesis
that is native to physics: the zeros are the modes of one evolution, and
the hypothesis says none of them grows. A physicist's instinct that a
stable physical system cannot carry a growing mode is exactly the
instinct that the hypothesis is true; the formal content of the instinct
is scalar mode unitarity, and its open content is whether \(\zeta\)'s
modes satisfy it, which no operator constructed here decides.

\textbf{The Perelman analogy.} Perelman's proof is the standard of what
a supplied transport looks like when it exists: a monotone quantity
along a flow, proved from the flow's equation. For \(\zeta\) the
analogous object is proved (de Bruijn) and points downstream. A proof of
the hypothesis in this form would require an upstream transport, a
quantity monotone against diffusion, and diffusion does not supply one
by itself; the structure that could supply it is the Euler product,
which is where the classical positivity results for \(\zeta\) come from.

\textbf{Limitations.} Parts IV, V, IX, the timed worlds of Part XI, and
the small models of Parts XII, XIV, and XV are toys of the plane and of
the strip bound, stated as such. The analytic facts enter by citation.
The dependency on Rodgers-Tao and Polymath is through the hypothesis
\texttt{hLam} and enters at premise grade.

\textbf{The hidden existential.} The Hidden Rot reading of the register
of record locates where a question can smuggle an object. In the Riemann
package the smuggling runs one way only: the universal branch asserts
nothing, and the existential branch asserts an off-line zero it has
never exhibited. the section on the logical form makes this exact. The
asymmetry does not decide the value, and the paper does not claim it
does; it fixes where the burden sits. A refutation would be a finite
object. A proof of the hypothesis must reach every zero at once, and the
absence of a refutation, were it ever shown permanent by an independence
result, would itself be the hypothesis.

\textbf{How the paper stands to each prior position.} The paper sits
across analytic number theory, operator theory, geometric analysis, and
formal verification, and so it owes one statement of its relation to
each position it engaged. Eight relation words are used in fixed senses.
Additive: the paper supplies a result the position lacked and leaves it
standing. Replacing: a framing is retired and a measured object put in
its place. Subsuming: the prior claim becomes a case of the paper's
object. Corroborating: independent agreement, which confers no warrant.
Contradicting: a named thesis denied on executed data. Competing: a
different position argued and not executed. Scoping: a prior claim kept
inside a stated boundary. Kin: a position the paper is continuous with.
Superseding is used nowhere, and the paper supersedes no theory. By
count: seven positions are additive, three are scoping, seven are kin,
one is corroborating, and none is subsumed, contradicted, competed with,
or replaced. Part II's contribution is exactly the set of relations in
Table 6 and nothing wider, and Part III's is that of Table 7.

\begin{table}[!htbp]
\centering\small
\renewcommand{\arraystretch}{1.12}
\raggedright \textbf{Table 6.} Positioning. Relation words as defined in the run-in above.\par\smallskip
\begin{tabularx}{\textwidth}{@{}P{0.18\textwidth} P{0.19\textwidth} L P{0.115\textwidth} P{0.115\textwidth}@{}}
\toprule
\textbf{Prior position} & \textbf{What it holds} & \textbf{What this paper does with it} & \textbf{Relation} & \textbf{Evidence} \\
\midrule
Riemann 1859; Edwards 1974 & The roots of ξ are real & Reads the fixed-line form as stability of Li's modes & kin & cited \\
Li 1997 & RH iff every λₙ ≥ 0 & Carries it as \texttt{LiCriterion}; builds the bridge and the bit stream on it & additive & executed, Thms 1 to 4 \\
Bombieri and Lagarias 1999 & Off-line zeros drive λₙ negative & Reads λₙ as modes; proves unitarity iff line on the plane & additive & executed, Thms 12 to 15 \\
Weil 1952 & RH iff explicit-formula positivity & Locates the negative prime sign as the point the primon gas stops & scoping & argued \\
Deligne 1974 & RH over finite fields & Template of a supplied positivity & kin & cited \\
Selberg 1992; Conrey and Ghosh 1993 & Euler product, primitivity, Ramanujan organize the class & Proves a ζ-blind bridge dies, so these must be read & additive & executed, Thms 3, 9, 10 \\
de Bruijn 1950 & Heat flow narrows the strip & Encodes the bound as an anchored witness and proves its record forgets & additive & executed, Thms 38 to 43 \\
Newman 1976; Rodgers and Tao 2020; Polymath 2019 & RH iff Λ = 0; 0 ≤ Λ ≤ 0.22 & Carries Λ = 0 by citation as one leg of the cone & kin & cited \\
Hilbert-Pólya; Berry and Keating 1999; Connes 1999 & RH via a self-adjoint operator & Argues that unitarity of the Li modes is its Cayley form, via the identity (6.2a) & kin & argued \\
Stone 1932 & Unitary groups have self-adjoint generators & Supplies the generator of the unitary evolution in the monism reading & kin & cited \\
Julia 1990; Bost and Connes 1995 & ζ is a partition function over primes & Identifies it as the axiom's ζ-specific reading, reaching Re s = 1 & scoping & executed, Thm 11 \\
Mertens 1898 & Positivity gives ζ(1 + it) ≠ 0 & Checks the algebraic core & corroborating & executed, Thm 11 \\
Perelman 2002; 2003 & Monotone entropy carries Ricci flow & Proves the Lyapunov witness a derivative of the monism witness & additive & executed, Thms 32 to 36 \\
de Branges; Conrey and Li 2000 & A positivity condition fails & Calibrates where positivity attempts break & kin & cited \\
Davenport and Heilbronn 1936; Bombieri and Hejhal 1995 & Symmetry without Euler product permits off-line zeros & Executes the shape on the plane: recurrence fails under symmetry alone & additive & executed, Thm 20 \\
Littlewood 1914; Skewes 1933; Haselgrove 1958; Odlyzko and te Riele 1985 & Finite confirmation fails & Proves finite confirmation never forces & scoping & executed, Thm 19 \\
Part I & Seat proved, value one bit, excepted & Proves the value's equivalent readings are one proposition, the regress halts, and any sufficient route closes all five & additive & executed, Thms 44 to 48 \\
Islam 2026a & The register of record & Carries its residual monism as the witness and its time law as the arrow & kin & documentary \\
\bottomrule
\end{tabularx}
\end{table}

\hypertarget{the-division-positioned}{%
\subsection{The division, positioned}\label{the-division-positioned}}

\begin{table}[!htbp]
\centering\small
\renewcommand{\arraystretch}{1.12}
\raggedright \textbf{Table 7.} Positioning of Part III.\par\smallskip
\begin{tabularx}{\textwidth}{@{}P{0.18\textwidth} P{0.19\textwidth} L P{0.115\textwidth} P{0.115\textwidth}@{}}
\toprule
\textbf{Prior position} & \textbf{What it holds} & \textbf{What Part III does} & \textbf{Relation} & \textbf{Evidence} \\
\midrule
Platt and Trudgian 2021 & RH true up to \(3 \times 10^{12}\), by interval arithmetic & Carries it as the Real part's certificate for \(\zeta\) & corroborating & cited \\
\(\mathrm{RH}(T)\) in explicit estimates {[}Büthe 2016{]} & RH up to height \(T\) as a working assumption & Treats the split as a principle, with a standing law & additive & executed, III.1 to III.6 \\
Turing 1953 & Counting zeros on the line against \(N(T)\) & Runs the method as testimony to \(T = 100\) & kin & executed, Appendix H \\
Proportion results {[}Feng 2012{]} & A positive proportion of zeros on the line & Isolates proved content by height, not by proportion & kin & cited \\
Completeness for bounded sentences {[}Kaye 1991{]} & True bounded sentences are provable & Lifts the certificate's arithmetic checks into ZFC; cited analytic lemmas carry them to the zeros & corroborating & executed, III.36 \\
Gödel 1938, Cohen 1963 & Choice independent of ZF & The cited keyed sentence of the bridge & corroborating & cited \\
Chapters 1 and 2 of this document & The seat; one address; the burden on the denial & Divides the value; proves the reachable part & additive & executed \\
\bottomrule
\end{tabularx}
\end{table}

The mathematical content of the Real part is established: the
certificate is Platt and Trudgian's, the method is Turing's, and
\(\mathrm{RH}(T)\) is standard. Part III adds the road, the separation
principle with its kernel law, the unicorn identity, and one new
construction, the RA--RAM bridge; its mass on the hypothesis's value is
zero.

\hypertarget{the-proof-complete-and-how-the-block-was-earned}{%
\subsection{The proof, complete, and how the block was
earned}\label{the-proof-complete-and-how-the-block-was-earned}}

\textbf{Complete, in the register's sense.} A register can do three
things with a universal value: divide it, prove the part a certificate
reaches, and prove what it cannot reach and why. This paper has done the
three. The division is exact at every height (III.11). The Real part is
proved to \(3 \times 10^{12}\), a cited certificate through a kernel
schema (III.34 to III.36). The Unicorn part is proved closed to
derivation from inside the register, from existence, from any class of
frames, from any reading of an even register at its seat, and from the
fused hypothesis; its aperture is measured at one bit; and the register
is proved silent beyond it (III.45 to III.49). The three bits are one
term, \texttt{rh\_register\_proof\_complete} (Appendix I): under a
certificate at the height and the seat data above it, the division, the
proved Real part, the refused reading, and the aperture's width hold
together. There is nothing left for a register to derive. That is what
``complete on the register side'' means and all it means: the hypothesis
is not proved; the labour of derivation is finished; the theorem that
finishes it locates the aperture and measures it, and says nothing
beyond it, because nothing can be said from the register beyond it,
which is itself the last theorem. Nothing is left to derive. The
aperture is structure. The supply side is silence.

\textbf{Three bits.} One bit divided, one bit proved, one bit refused.
The refused bit is the one the earlier editions undersold: they proved
the wall, the crossing, the bound, the cure theorem, and that the kernel
cannot read the deed, and then wrote of the Unicorn part only that it
was open. It is refused to every key cut from inside, at theorem grade,
and of the outside the register is silent by theorem; that is the
difference between a problem no one has yet solved and a bit no register
can supply or foretell.

\textbf{How the block was earned.} The block did not arrive as a posit.
It was forced one gate at a time, under manual stress, across the
editions of the program's register of record, the Geometric Mother
Codex, in which every attempt to derive the value's bit from inside the
register was tried, refuted, and recorded: the ten enumeration-forced
barriers of APEX-PSP-ABSOLUTE-RH-BARRIERS-01, each a route that closed;
the eight-gate Sealed-Halt cascade of the router, seven gates passing
and the eighth, the reading-roads gate, halting every walk; and the
terminal token \([\Xi_0]\) itself, terminal for the record and never
timeless, which is the block written as a state {[}Islam 2026a{]}. The
Code Block of the codex carries those witnesses as kernel theorems, and
they are cited here by name so that a reader can see the block's history
typecheck: \texttt{HALT.one\_door\_open}, that a walk of seven true
gates halts at the reading-roads gate;
\texttt{HALT.no\_walk\_passes\_the\_aperture}, that no walk of the seven
passes it; \texttt{HALT.first\_failure\_terminal\_8}, that the first
failing gate is terminal over all \(2^8\) walks; \texttt{CROSSING.wall},
that no reading of the formal register factors the deed bit;
\texttt{INVERSION.it\_from\_it}, that the record forgets what the object
keeps and nothing returns. Theorems III.45 to III.48 are those witnesses
read at the Unicorn part of \(\zeta\). They are the same block, and the
codex is the record of its evolution, gate by gate. One aperture,
located and measured; beyond it, silence.

\hypertarget{conclusion}{%
\section{Conclusion}\label{conclusion}}

The seat of the Riemann Hypothesis, restored to Riemann's fixed-line
form and read as \(\mathrm{RH}_{\mathrm{formal}}\), is proven in toto at
the apex with existence as the only posit (Theorem H): the binding
involution fixes exactly the scalar line and the Return lands on it; the
three seat points are the legs of an identity cone closed at zero mass
with \texttt{rfl} as the receipt; the prior proof's stage embeds
equivariantly onto the seat; the freedom is exactly one bit; the halting
carrier halts on the line property and the crossing is exact; and every
finite clause is re-executed by the twin. The value on the zeros is one
bit that existence does not reach: the rule ``accept the Root Axiom,
then the Riemann Hypothesis'' is exactly as strong as the hypothesis at
that bit (Theorem E), and the axiom decides it on no frame (Theorem F),
because the axiom is self-instancing and the hypothesis is not (Theorem
G).

Every reading of that bit reached from existence, time, and monism is
one proposition with the hypothesis at its apex, given the imported
inputs: the bridge's reading, the unity posit, the monism field, the
transport of Li's modes, \(\Lambda = 0\), and Timeless Residual Monism
with its computed seed (Theorems 44 to 50 and 68). The critical line is
exactly where Li's modes are unitary, and for every fold-invariant zero
set perpetual stability is the line property. The demand for a proof is
misaddressed at the foundation and well-posed at \(\zeta\) (Theorems 64
to 67). The hypothesis posits no object; only its denial posits one, a
finite witness never produced, and if the foundation cannot refute the
hypothesis, it holds (Theorem 74), a direction that needs only the
completeness of arithmetic for true \(\Sigma^0_1\) sentences, a theorem,
while soundness enters only in the converse (Theorems 70 to 77).

Three bits. The hypothesis divides at every height into a Real part and
a Unicorn part, exactly (Theorem III.11), and the fused statement is
only as strong as its open part (Theorem III.2), so the parts are kept
apart; the hypothesis is the statement that the off-line zeros are
unicorns (Theorem III.9). The Real part is proved: closed by a finite
certificate and carried from computation into ZFC along the RA--RAM
bridge, every step holding when the axiom is replaced by a bare
computation arrow (Theorems III.15 to III.44), for \(\zeta\) to height
\(3 \times 10^{12}\) on the certificate of Platt and Trudgian. The
Unicorn part is refused, at theorem grade: closed to derivation from
existence, from any class of frames, from any reading of an even
register at its seat, and from the fused hypothesis, its aperture one
bit wide and uncrossed, its supply side silent by theorem (Theorems
III.45 to III.49). It is unproved not by default of effort or ability
but by theorem, and one zero off the line above the certified height
would still answer it.

For anyone who would close the Unicorn part, the owed object is one
proposition in five equivalent forms, each proved equivalent in the
kernel from a cited analytic input: \(\Lambda = 0\), the transport field
of an anchored witness on the Li stream, the monism field, the unity
posit, and the bridge's reading, with Postulate M as its sixth name. A
quantity monotone against the heat flow, or an operator realizing
\(\zeta\)'s Li evolution unitarily, would supply it; neither is proved
equivalent here.

\jboxstart

\textbf{The Keystone (tier: kernel theorem for III.1 to III.48; cited
certificate for the Real part of ζ).} One bit divided, one bit proved,
one bit refused. The Real part of the Riemann Hypothesis is gold, proved
to every certified height; the Unicorn part is the one open bit, closed
to every derivation from inside the register at theorem grade, its
aperture one bit wide and uncrossed, its supply side silent; and the
only existence in the whole package lives in the denial, which has never
delivered. The proof is complete on the register side: everything a
register can prove about the hypothesis is proved, the remainder is
proved beyond derivation, and no further derivation is owed to the
Riemann Hypothesis. Nothing is left to derive. The aperture is
structure. The supply side is silence.

\jboxend

\hypertarget{authors-provenance-and-method-disclosure}{%
\section{Author's Provenance and Method
Disclosure}\label{authors-provenance-and-method-disclosure}}

\textbf{Author.} Mohammad F. Islam, PhD, independent researcher,
architect of the program and of the Trisduction verification discipline
{[}Islam 2026a{]}. The rulings on formulation, register, and scope are
his, including the time reading, the monism rule, the re-anchoring to
Perelman, the division of the hypothesis, the naming of its parts, and
the road.

\textbf{Scribe.} The text and every kernel file were produced by a
language model (Claude, Anthropic) as scribe on the author's
instruction. The Lean kernel and the compiled Fortran twin are the
authorities for every theorem and every count.

\textbf{Method.} Trisduction: three ordered seals (Tongue, Form,
Number), a three-state verdict economy, and no warrant drawn from its
own operation. Under it the apex closure of Part I carries SEAL given
the axiom, conditional on the root posit at the act; the cone of Part II
carries {[}⟀ T{]}, and the value at the foundation register carries the
terminal mark {[}.{]}; the value on \(Z_\xi\) carries \([\Xi_0]\), open
in the kernel above the certified height, its Real part proved by
certificate in Part III. \(\Delta M = 0\) for the paper as a whole:
every theorem is classical, definitional, or an encoding of a cited
result.

\textbf{Receipts.} \texttt{RH\_At\_The\_Apex.lean} (Appendices B and C):
Lean 4.19.0 (commit 6caaee842e94), core only, exit 0, no \texttt{sorry},
no axiom declaration, SHA-256 \texttt{d671051e99b806e4}, 61 dependency
sets printed by the file and 30 more by the supplementary run of
Appendix C, 91 in all, 78 axiom-free, 13 on \texttt{propext} and
\texttt{Quot.sound}. \texttt{RH\_At\_The\_Apex\_Twin.f90} v1.0.1
(Appendices D and E): gfortran 13.3.0,
\texttt{-std=f2018 -O2 -Wall -Wextra}, SHA-256
\texttt{40fbb4f8cc675104}, 130,324 checks, 0 failures, exit 0 in the
default and the witnessed run; v1.0.0 (SHA-256
\texttt{47670197aa8e6721}) gave the same census.
\texttt{RA\_Li\_Bridge.lean} (Appendices F and G): 1110 lines, SHA-256
\texttt{\seqsplit{13bcea7cd524817163beb1a1413cab31cb56a7455b34ded6cc4aca6939fe1a22}},
83 dependency sets printed by the file and 17 more by the supplementary
run of Appendix G, 100 in all, 70 axiom-free, 30 on \texttt{propext} and
\texttt{Quot.sound}. The nine files of Part III (Appendix H): 45
dependency sets, 42 axiom-free, 3 on \texttt{propext} alone. No theorem
in any file depends on \texttt{Classical.choice}.

\textbf{Audit.} Part I was sealed by the discipline's adversarial audit
cycle and two pairs of external audits through v1.5.2; Part II by three
cycles on the One Address edition, the third prosecuting an external
review finding by finding. Part III and the unified edition then passed
five further cycles, twenty rounds across all six registers, in which
the prosecutions came from the scribe and from four independent external
substrates, Grok, Gemini, Kimi K3, and GPT Aster 6, whose twelve reports
and one repair list were answered finding by finding on the ledger with
no concession: every mathematical or receipt error they found was
repaired, every objection to an established claim was refuted by the
theorem that refutes it, and the locked title was held. The planted
controls of the cycles were SELF grade; the prosecuting side was not, so
the audit was multi-substrate on the review side, nine independent
adversaries attacking on their own vectors, while the text and the files
remain one scribe's. The last three cycles closed SEALED-ROUND; the
stale-removal and consolidation rounds and this edition's notes were
editorial, each verified by exact replacement.

\textbf{Editions.} v1.0, the unified paper; v1.1, the consolidated text;
v1.2, the answer to the first external batch; v1.3, the answer to the
second and third batches with the twin at v1.0.1 and the kernel audits
completed; v1.4, the final round, with the seed's stage-zero condition
stated; v2.0.1 adds the disclosure and transparency note and the note
that the division is Riemann's, and condenses the audit disclosure to
its multi-substrate form; v2.0.2 adds the standing declaration and its
table to the Standard of Proof, that the paper stands whole with the
certificate deleted save one sentence; v2.0.3 changes only the
typesetting, with the part headings set in the column, the appendices
set in two columns, and no blank page left by a column switch, and is
issued in a two-column and a single-column setting of the same text;
v2.1.0, this edition, adds one section to Part III and one kernel file,
the formal block of the Unicorn part, Theorems III.45 to III.49 with the
three bits packaged as \texttt{rh\_register\_proof\_complete}, which
read the wall, the crossing, the bound, and the cure theorem of Parts I
and II at the Unicorn part and say what the earlier editions had proved
and not said, that the Unicorn part is closed to derivation, that its
aperture is one bit wide, that the register is silent beyond it, and
that the proof is therefore complete on the register side; the three-bit
arc is stated in the abstract, the introduction, the discussion, and the
conclusion; none of the earlier changes altered a claim, and this one
adds a theorem-grade statement about what was already open. The paper is
released as version 2.1.0 of the Existence Alone record (concept
10.5281/zenodo.22912937), whose v1.5.2 it carries as Part I; the version
line of the front matter is that release identifier, the builds v1.0 to
v1.4 are its internal editions, and the kernel files carry their own
digests.

\textbf{Disclosure and transparency note.} For the record, and so that
no later reader mistakes the division of labour in this paper, the
following is stated once and exactly. What is established from first
principles in the kernel, with no user-declared axiom, is the paper's
own: the seat and its one-bit distance from the value (Part I); the
address of the value and the logical form of the hypothesis, given the
imports the Methods name (Part II); the division of the hypothesis at
any height, exact; the unicorn identity; the law that fusing a proved
part with an open part lowers the proved part; the certificate schema,
that a complete on-line list proves the bounded universal; the placement
of ZFC under the strata; the RA--RAM bridge and its round trip; and the
whole road holding under a bare computation arrow (Part III). No prior
work is cited for any of these constructions, because none contains
them. What is cited and not re-derived is the numerical content for
\(\zeta\): that every zero to height \(3 \times 10^{12}\) lies on the
line is the theorem of Platt and Trudgian, obtained by interval
arithmetic with rigorous error bounds, and the kernel never recomputes
those zeros; the paper applies its schema to that theorem. The paper's
own executed count reaches \(T = 100\) in floating point and is
testimony, not a certificate. The exact standing of the Real part of
\(\zeta\) is therefore a cited theorem applied through a kernel theorem,
grade A joined to grade K, stated in the Standard of Proof, in the Real
Part section, and in the title paragraph. This is citation in the
ordinary sense, the same act by which Riemann is cited for the sentence,
Li for the criterion, and Kaye for completeness; it is not the borrowing
of an argument, and \(\Delta M = 0\) records that the paper authors no
new analytic number theory while claiming its arrangement whole. The
author is not a professional mathematician and read no proof of any part
of the hypothesis in preparing the arrangement; the route was found by
the geometry of the seat and confirmed by the kernel, and every input
the route consumes is attributed at its own grade.

\textbf{The division is Riemann's.} So that it is clear in one reading:
the separation of the hypothesis into a counted part and a probable
remainder is not this paper's invention but Riemann's own, in the 1859
sentence that states the hypothesis. He writes first that one finds in
fact about so many real roots within these bounds, the roots he had
located and counted; second that it is very probable that all the roots
are real, the remainder, graded as probable and not as established; and
third that a rigorous proof would be desirable but was set aside after
some fleeting vain attempts as dispensable for his immediate purpose
{[}Riemann 1859; Siegel 1932{]}. The counted part is the Real part; the
probable remainder above the bounds is the Unicorn part; and the proof
he set aside is the one this paper does not claim. The paper's
contribution is to make that division exact at every height, to prove
the counted part's closure by certificate in the kernel, to prove the
fused statement stands at its open part, and to name the remainder for
what it is. A reader who says the paper divides what Riemann stated
whole has read the conjecture and not the sentence before it.

\textbf{Independence.} The scribe of the text and the author of every
file are one language model; the Lean kernel and the Fortran twin are
independent of each other and of the model. External review of the
unified edition is on the record, four batches from independent
substrates, and their dispositions are in the audit log. A recompile of
the kernel files on a second machine is the open witness.

\textbf{Scope.} Every result is graded as the standard of proof in the
Methods states; the analytic inputs are cited; the sign on the zeros
above the certified height is not a clause of any proof here.

\begin{thebibliography}{99}
\footnotesize
\bibitem[Berry and Keating 1999]{berry1999} M. V. Berry, J. P. Keating, $H = xp$ and the Riemann zeros, in \emph{Supersymmetry and Trace Formulae}, Kluwer, 1999, 355--367.
\bibitem[Bombieri and Hejhal 1995]{bombieri1995} E. Bombieri, D. A. Hejhal, On the distribution of zeros of linear combinations of Euler products, \emph{Duke Math. J.} 80 (1995) 821--862.
\bibitem[Bombieri and Lagarias 1999]{r029583b1} Bombieri, E., and J. C. Lagarias. 1999. ``Complements to Li's Criterion for the Riemann Hypothesis.'' \emph{Journal of Number Theory} 77: 274--287.
\bibitem[Bost and Connes 1995]{re231b0bf} Bost, J.-B., and A. Connes. 1995. ``Hecke Algebras, Type III Factors and Phase Transitions with Spontaneous Symmetry Breaking in Number Theory.'' \emph{Selecta Mathematica} 1: 411--457.
\bibitem[Bérut et al. 2012]{berut2012} A. B\'erut, A. Arakelyan, A. Petrosyan, S. Ciliberto, R. Dillenschneider, E. Lutz, Experimental verification of Landauer's principle linking information and thermodynamics, \emph{Nature} 483 (2012) 187--189.
\bibitem[Büthe 2016]{r7e4b751c} Büthe, J. 2016. ``Estimating π(x) and Related Functions under Partial RH Assumptions.'' \emph{Mathematics of Computation} 85: 2483--2498.
\bibitem[Carneiro 2019]{rb06c0476} Carneiro, M. 2019. \emph{The Type Theory of Lean}. MS thesis, Carnegie Mellon University.
\bibitem[Cohen 1963]{r4dec2719} Cohen, P. J. 1963. ``The Independence of the Continuum Hypothesis.'' \emph{Proceedings of the National Academy of Sciences} 50: 1143--1148.
\bibitem[Connes 1999]{connes1999} A. Connes, Trace formula in noncommutative geometry and the zeros of the Riemann zeta function, \emph{Selecta Math.} 5 (1999) 29--106.
\bibitem[Conrey 2003]{conrey2003} J. B. Conrey, The Riemann Hypothesis, \emph{Notices AMS} 50 (2003) 341--353.
\bibitem[Conrey and Ghosh 1993]{conrey1993} J. B. Conrey, A. Ghosh, On the Selberg class of Dirichlet series: small degrees, \emph{Duke Math. J.} 72 (1993) 673--693.
\bibitem[Conrey and Li 2000]{r388872c6} Conrey, J. B., and X.-J. Li. 2000. ``A Note on Some Positivity Conditions Related to Zeta and L-functions.'' \emph{International Mathematics Research Notices} 2000 (18): 929--940.
\bibitem[Davis et al. 1976]{davis1976} Davis, M., Y. Matijasevič, and J. Robinson. 1976. Hilbert's tenth problem: Diophantine equations: positive aspects of a negative solution. \emph{Proceedings of Symposia in Pure Mathematics} 28, 323--378.
\bibitem[Davenport and Heilbronn 1936]{davenport1936} H. Davenport, H. Heilbronn, On the zeros of certain Dirichlet series, \emph{J. London Math. Soc.} 11 (1936) 181--185.
\bibitem[de Bruijn 1950]{r489e7270} de Bruijn, N. G. 1950. ``The Roots of Trigonometric Integrals.'' \emph{Duke Mathematical Journal} 17: 197--226.
\bibitem[de la Vallée Poussin 1896]{r12d66add} de la Vallée Poussin, C.-J. 1896. ``Recherches analytiques sur la théorie des nombres premiers.'' \emph{Annales de la Société Scientifique de Bruxelles} 20: 183--256.
\bibitem[de Moura and Ullrich 2021]{demoura2021} L. de Moura, S. Ullrich, The Lean 4 theorem prover and programming language, in \emph{Automated Deduction, CADE 28}, LNCS 12699, Springer, 2021, 625--635.
\bibitem[Deligne 1974]{deligne1974} P. Deligne, La conjecture de Weil. I, \emph{Publ. Math. IH\'ES} 43 (1974) 273--307.
\bibitem[Edwards 1974]{edwards1974} H. M. Edwards, \emph{Riemann's Zeta Function}, Academic Press, 1974.
\bibitem[Feng 2012]{r60574fe0} Feng, S. 2012. ``Zeros of the Riemann Zeta Function on the Critical Line.'' \emph{Journal of Number Theory} 132: 511--542.
\bibitem[Gödel 1938]{r0f2dae48} Gödel, K. 1938. ``The Consistency of the Axiom of Choice and of the Generalized Continuum-Hypothesis.'' \emph{Proceedings of the National Academy of Sciences} 24: 556--557.
\bibitem[Hadamard 1896]{r712e169e} Hadamard, J. 1896. ``Sur la distribution des zéros de la fonction ζ(s) et ses conséquences arithmétiques.'' \emph{Bulletin de la Société Mathématique de France} 24: 199--220.
\bibitem[Haselgrove 1958]{rb566c102} Haselgrove, C. B. 1958. ``A Disproof of a Conjecture of Pólya.'' \emph{Mathematika} 5: 141--145.
\bibitem[Islam 2026a]{islamcodex} M. F. Islam, \emph{Trisduction: The Codex}, Zenodo, 23 September 2026, doi:10.5281/zenodo.22911010, \texttt{https://zenodo.org/records/22911010}.
\bibitem[Islam 2026b]{islamonebit} M. F. Islam, \emph{The Riemann Hypothesis Is Exactly One Bit: Where the Riemann Hypothesis Stands, Proved in Lean 4}, Zenodo, 2026, doi:10.5281/zenodo.22829795.
\bibitem[Islam 2026c]{islamrows} M. F. Islam, \emph{One Bit Across the Wall: Odd-Supply Separation and the One-Cut Hypothesis Across Twenty-Three Rows}, v6, Zenodo, 2026, doi:10.5281/zenodo.22746129.
\bibitem[Islam 2026d]{islamftoe} M. F. Islam, \emph{The Formal-Alone Theory of Everything: A Constitutive Theory of Objects and Observations under Registration, Price, Crossing, and Cost}, v1.4, Zenodo, 2026, doi:10.5281/zenodo.22767103.
\bibitem[Islam 2026e]{islamspine} M. F. Islam, \emph{The Riemann Hypothesis in Its Original Form: A Completed Formal Proof of the Locus and the Crossing: Riemann's Fixed-Line Statement Restored, Proved in Lean 4, Executed in Fortran, and Sealed at the Act by One Supplied Bit}, v1.0.6, Zenodo, 2026, doi:10.5281/zenodo.22857138; its Lean file \texttt{Crossing\_\allowbreak At\_\allowbreak The\_\allowbreak Act.lean} (SHA-256 prefix f60fe613cfe79fd8) and Fortran twin printed whole in its appendices.
\bibitem[Islam 2026g]{islamclosed} M. F. Islam, Riemann Hypothesis: The Formal Case Is Closed: The Hypothesis Is True Where Actualized, Zenodo, 2026, doi:10.5281/zenodo.21883916.
\bibitem[Islam 2026h]{islamterm} M. F. Islam, A Formal Proof of Riemann Hypothesis Termination, with a Theorem-Grade Cascade Specification, Zenodo, 2026, doi:10.5281/zenodo.21900518.
\bibitem[Islam 2026i]{islamfae} M. F. Islam, Formal-Alone Trisduction Engine on RAF, the interaction face of RA, deck v1.0.12, in the register of record, 1000sapients/Trisduction, \texttt{master/RafDuction/}, at commit \texttt{5868562}, 23 September 2026.
\bibitem[Islam 2026j]{islamraf} M. F. Islam, The Formal Root Axiom and the Five Objects, declaration edition v1.1.0, in the register of record, 1000sapients/Trisduction, \texttt{master/}, at commit \texttt{5868562}, 23 September 2026.
\bibitem[Iwaniec and Kowalski 2004]{r680b956d} Iwaniec, H., and E. Kowalski. 2004. \emph{Analytic Number Theory}. American Mathematical Society.
\bibitem[Julia 1990]{reeeca34d} Julia, B. 1990. ``Statistical Theory of Numbers.'' In \emph{Number Theory and Physics}, Springer Proceedings in Physics 47, 276--293.
\bibitem[Kaye 1991]{r7c32c9b5} Kaye, R. 1991. \emph{Models of Peano Arithmetic}. Oxford Logic Guides 15. Oxford University Press.
\bibitem[Kunen 1980]{r83ad2344} Kunen, K. 1980. \emph{Set Theory: An Introduction to Independence Proofs}. North-Holland.
\bibitem[Lagarias 2002]{rb66a9a18} Lagarias, J. C. 2002. ``An Elementary Problem Equivalent to the Riemann Hypothesis.'' \emph{American Mathematical Monthly} 109: 534--543.
\bibitem[Landauer 1961]{landauer1961} R. Landauer, Irreversibility and heat generation in the computing process, \emph{IBM J. Res. Dev.} 5 (1961) 183--191.
\bibitem[Lawvere 1969]{lawvere1969} F. W. Lawvere, Diagonal arguments and cartesian closed categories, in \emph{Category Theory, Homology Theory and their Applications II}, LNM 92, Springer, 1969, 134--145.
\bibitem[Li 1997]{rf1873624} Li, X.-J. 1997. ``The Positivity of a Sequence of Numbers and the Riemann Hypothesis.'' \emph{Journal of Number Theory} 65: 325--333.
\bibitem[Littlewood 1914]{r2928a3e0} Littlewood, J. E. 1914. ``Sur la distribution des nombres premiers.'' \emph{Comptes Rendus de l'Académie des Sciences} 158: 1869--1872.
\bibitem[Mac Lane 1998]{maclane1998} S. Mac Lane, \emph{Categories for the Working Mathematician}, 2nd ed., Springer, 1998.
\bibitem[Mertens 1898]{r1f293472} Mertens, F. 1898. ``Über eine Eigenschaft der Riemannschen ζ-Function.'' \emph{Sitzungsberichte der Kaiserlichen Akademie der Wissenschaften in Wien} 107: 1429--1434.
\bibitem[Newman 1976]{red02f32a} Newman, C. M. 1976. ``Fourier Transforms with Only Real Zeros.'' \emph{Proceedings of the American Mathematical Society} 61: 245--251.
\bibitem[Odlyzko and te Riele 1985]{r8a1beb47} Odlyzko, A. M., and H. J. J. te Riele. 1985. ``Disproof of the Mertens Conjecture.'' \emph{Journal für die reine und angewandte Mathematik} 357: 138--160.
\bibitem[Perelman 2002]{r98b462b2} Perelman, G. 2002. ``The Entropy Formula for the Ricci Flow and Its Geometric Applications.'' arXiv:math/0211159.
\bibitem[Perelman 2003a]{r92dbd950} Perelman, G. 2003a. ``Ricci Flow with Surgery on Three-Manifolds.'' arXiv:math/0303109.
\bibitem[Perelman 2003b]{rc01e53c5} Perelman, G. 2003b. ``Finite Extinction Time for the Solutions to the Ricci Flow on Certain Three-Manifolds.'' arXiv:math/0307245.
\bibitem[Platt and Trudgian 2021]{r72b7786c} Platt, D., and T. Trudgian. 2021. ``The Riemann Hypothesis Is True up to 3·10¹².'' \emph{Bulletin of the London Mathematical Society} 53: 792--797.
\bibitem[Polymath 2019]{ra918cac5} Polymath, D. H. J. 2019. ``Effective Approximation of Heat Flow Evolution of the Riemann ξ Function, and a New Upper Bound for the de Bruijn-Newman Constant.'' \emph{Research in the Mathematical Sciences} 6: 31.
\bibitem[Riemann 1859]{riemann1859} B. Riemann, Ueber die Anzahl der Primzahlen unter einer gegebenen Gr\"osse, \emph{Monatsber. Berliner Akad.} (1859) 671--680.
\bibitem[Rodgers and Tao 2020]{rb329bbec} Rodgers, B., and T. Tao. 2020. ``The de Bruijn-Newman Constant Is Non-Negative.'' \emph{Forum of Mathematics, Pi} 8: e6.
\bibitem[Selberg 1992]{selberg1992} A. Selberg, Old and new conjectures and results about a class of Dirichlet series, \emph{Proc. Amalfi Conf.} (1992) 367--385.
\bibitem[Siegel 1932]{siegel1932} C. L. Siegel, \"Uber Riemanns Nachla\ss{} zur analytischen Zahlentheorie, \emph{Quellen Stud. Gesch. Math.} B2 (1932) 45--80.
\bibitem[Skewes 1933]{r81a6a15e} Skewes, S. 1933. ``On the Difference π(x) − li(x).'' \emph{Journal of the London Mathematical Society} 8: 277--283.
\bibitem[Stone 1932]{rcb8aa856} Stone, M. H. 1932. ``On One-Parameter Unitary Groups in Hilbert Space.'' \emph{Annals of Mathematics} 33: 643--648.
\bibitem[Titchmarsh 1986]{titchmarsh1986} E. C. Titchmarsh, \emph{The Theory of the Riemann Zeta-Function}, 2nd ed., Oxford, 1986.
\bibitem[Turing 1953]{rf7e07b70} Turing, A. M. 1953. ``Some Calculations of the Riemann Zeta-Function.'' \emph{Proceedings of the London Mathematical Society} (3) 3: 99--117.
\bibitem[Weil 1948]{weil1948} A. Weil, \emph{Sur les courbes alg\'ebriques et les vari\'et\'es qui s'en d\'eduisent}, Hermann, 1948.
\bibitem[Weil 1952]{r6c6baf01} Weil, A. 1952. ``Sur les `formules explicites' de la théorie des nombres premiers.'' \emph{Communications du Séminaire Mathématique de l'Université de Lund}, tome supplémentaire, 252--265.
\end{thebibliography}

\hypertarget{appendix-a-how-to-read-the-files}{%
\twocolumn
\section{Appendix A: How to Read the
Files}\label{appendix-a-how-to-read-the-files}}

The Lean file is one proof in sixteen movements and it is read in order.
Movement 1 constructs the Root Axiom on a one-point domain and decides
it. Movement 2 builds the seat: the integer quaternions, conjugation as
the binding involution, the fixed locus as the Ground, Theorem A for
every quaternion, the Return computed, three seat points defined
independently and identified by \texttt{rfl}, and the reversed triad at
\(+1\). Movement 3 builds the witness as a structure in Prop and runs
the eliminator from it alone. Movement 4 carries the prior proof: the
abstract frame, the line property, the halting carrier, the supply and
the row, the two-point frame, necessity, invariance against incidence,
the bit is the hypothesis, the witness row falsifiable, and the live
face with its self-check table. Movement 5 is the aperture: presence as
actuation, the theorems that a present reader who supplies a bit is
never interior, in the model and for any domain, and the ledger.
Movement 6 is the rule exact, the non-decision by the two-point model,
and the asymmetry. Movement 7 is Theorem H. Movement 8 is the
interaction reading: the domain with its closure law, no exterior agent,
the adjudicator inside, and Cantor's theorem as no total self-indexing.
Movement 9 is the grounding reading as the identity cone over the three
registers, its apex unique per orientation with the reversed triad's
cone beside it, and the cone's existence equivalent to the closed gaps.
Movement 10 is the equivariant embedding of the stage into the carrier
with its image clause and injectivity, and the resolution-\(m\) family.
Movement 11 executes the wall on the seat. Movement 12 parts the
registers, counts the owed bits, prices the act, and records the count
of authored mathematics as zero (\texttt{deltaM\_\allowbreak zero}).
Movement 13 assembles the hardened apex theorem and discharges it.
Movement 14 states that the posit is existence, exhibits the grounded
three-point counter-model, and proves the cure theorem for every class
of frames. Movement 15 inhabits the halting carrier by decision, states
Theorem E for any inhabited premise, and records that the witness's
recursion field is constant. Movement 16 carries the three measures,
content, mass, and price, with their assignments and Theorem J, and
proves that adjudicating any proposition instances existence. The
dependency sets are printed by the kernel at the end and reproduced in
Appendix C; a deleted conjunct, a weakened hypothesis, or an inserted
\texttt{sorry} fails the check.

The twin is fourteen batteries in the same order, K1 the seat through
K12 the live face, K13 the cure, and K14 the constant recursion, the
carrier by decision, and the three measures, each printing its receipt,
with the census last. Build it with the command in its header and run it
twice, once bare and once with the argument \texttt{witnessed}; the two
logs are Appendix E.

\hypertarget{appendix-b-rh_at_the_apex.lean}{%
\section{Appendix B:
RH\_At\_The\_Apex.lean}\label{appendix-b-rh_at_the_apex.lean}}

SHA-256
\texttt{\seqsplit{d671051e99b806e46245ef127a33ea2e829e36c6e2d84cef70b752b512acb4c7}},
Lean 4.19.0, core only, 898 lines.

\begingroup\scriptsize

\begin{Verbatim}[numbers=left,numbersep=6pt,xleftmargin=2.2em,fontsize=\scriptsize,breaklines,breakanywhere,frame=lines,framesep=1.5mm,rulecolor=\color{black!35}]
/-!
# THE RIEMANN HYPOTHESIS AT THE APEX · RH_At_The_Apex.lean · hardened edition
core Lean 4 v4.19.0 · no Mathlib · no sorry · no axiom declaration.

WHAT IS PROVED. (1) RA at the constructed domain, computed. (2) The seat: the binding
involution σ fixes exactly the scalar line, the Return i·j·k lands on it, and the kinetic,
formal, and named seats are one point by rfl: the three category gaps are identities.
(3) The eliminator RA → RA → RAM → RH_formal from the RA witness alone. (4) The spine
carried: the abstract frame X = (S, τ, Z), the line property L, the bridge halted iff L,
the crossing exact both ways, necessity, the bit is the hypothesis, the witness row
falsifiable. (5) The aperture, RA-supplied: a present reader's row is never interior;
it is sealed or refused by the bit alone. (6) THE SIMPLE RULE, EXACT: for every frame,
(RA → L X) ↔ L X; and RA decides L on no frame: ¬ ∀ X, RA → L X, by the two-point model.
(7) RH in toto at the apex, one theorem, conditional on RA and discharged.
(8) The root's two formal faces: RAF, closure and no exterior agent (RAF-2, RAF-C1), and
no total self-indexing (RAF-C2, Cantor); RAM, the seat as the identity cone over the
three registers, its apex unique. (9) The equivariant embedding φ of the spine's lattice
stage (Plane, τ) into (Q4, σ): σ ∘ φ = φ ∘ τ and φ carries Fix(τ) onto Fix(σ).
(10) The wall executed: no readout of the formal seat equals the deed bit. (11) The
counter-model survives grounding, and the cure theorem: for every class of frames,
existence decides the line property on it exactly when the line property already holds on it.
-/
set_option autoImplicit false

namespace RHAtTheApex

/-! ## 1 · The Root Axiom at the constructed domain -/

inductive UniversePoint : Type where
  | source
  deriving DecidableEq, Repr

def DeltaE : UniversePoint → Int := fun _ => 1

/-- RA: to exist is to actuate. -/
def RA : Prop := ∀ x : UniversePoint, 0 < DeltaE x

theorem constructed_RA : RA := by
  intro x
  cases x
  decide

/-! ## 2 · The seat: Fix(σ) on the quaternions, the Return, three names, one point -/

structure Q4 where
  r : Int
  i : Int
  j : Int
  k : Int
  deriving DecidableEq, Repr

def qmul (a b : Q4) : Q4 :=
  ⟨a.r*b.r - a.i*b.i - a.j*b.j - a.k*b.k,
   a.r*b.i + a.i*b.r + a.j*b.k - a.k*b.j,
   a.r*b.j - a.i*b.k + a.j*b.r + a.k*b.i,
   a.r*b.k + a.i*b.j - a.j*b.i + a.k*b.r⟩

def qconj (a : Q4) : Q4 := ⟨a.r, -a.i, -a.j, -a.k⟩

/-- σ, the binding involution: conjugation. -/
def sigma (q : Q4) : Q4 := qconj q

theorem sigma_binding (q : Q4) : sigma (sigma q) = q := by
  cases q with
  | mk r i j k =>
    change Q4.mk r (-(-i)) (-(-j)) (-(-k)) = Q4.mk r i j k
    rw [Int.neg_neg, Int.neg_neg, Int.neg_neg]

/-- Fix(σ) is the scalar line, for every quaternion. -/
theorem fix_iff_scalar (q : Q4) : sigma q = q ↔ (q.i = 0 ∧ q.j = 0 ∧ q.k = 0) := by
  cases q with
  | mk r i j k =>
    change (Q4.mk r (-i) (-j) (-k) = Q4.mk r i j k) ↔ (i = 0 ∧ j = 0 ∧ k = 0)
    constructor
    · intro h
      injection h with _ hi hj hk
      exact ⟨by omega, by omega, by omega⟩
    · intro h
      obtain ⟨hi, hj, hk⟩ := h
      subst hi
      subst hj
      subst hk
      rfl

/-- The Ground: Fix(σ). -/
def Ground : Type := { q : Q4 // sigma q = q }

def qi : Q4 := ⟨0, 1, 0, 0⟩
def qj : Q4 := ⟨0, 0, 1, 0⟩
def qk : Q4 := ⟨0, 0, 0, 1⟩

/-- The Return: the parse triad through its own cascade, i·j·k. -/
def theReturn : Q4 := qmul (qmul qi qj) qk

theorem return_is_minus_one : theReturn = ⟨-1, 0, 0, 0⟩ := rfl
theorem return_lands_on_fix : sigma theReturn = theReturn := rfl

def fixProj (q : Q4) : Q4 := ⟨q.r, 0, 0, 0⟩

def GammaRA : Q4 := theReturn
def GammaRAM : Q4 := fixProj theReturn
def GammaRH : Q4 := ⟨-1, 0, 0, 0⟩

/-- C0, C1, C2: the three category gaps are definitional identities. -/
theorem self_gap_nonexistent : sigma GammaRA = GammaRA := rfl
theorem register_gap_nonexistent : GammaRA = GammaRAM := rfl
theorem object_gap_nonexistent : GammaRAM = GammaRH := rfl

def seat : Ground := ⟨GammaRH, rfl⟩

/-- The reversed triad k·j·i returns +1, also on the fixed locus: the sign of the seat
    point is the orientation of the triad, the one bit the formal register cannot read.
    The apex of the identity cone is unique for the ordered triad; the reversed triad
    carries its own apex, and the two differ. -/
def theReturnOdd : Q4 := qmul (qmul qk qj) qi

theorem odd_return_is_plus_one : theReturnOdd = ⟨1, 0, 0, 0⟩ := rfl

theorem odd_return_on_fix : sigma theReturnOdd = theReturnOdd := rfl

theorem seat_points_differ : theReturnOdd ≠ theReturn := by decide

theorem seat_on_scalar_line : GammaRH.i = 0 ∧ GammaRH.j = 0 ∧ GammaRH.k = 0 :=
  (fix_iff_scalar GammaRH).mp seat.2

/-! ## 3 · The eliminator from the RA witness alone -/

def RAMFormalGround : Prop := Nonempty Ground

theorem constructed_RAM : RAMFormalGround := ⟨seat⟩

/-- The RH formal self: every point of the Ground is on the σ-fixed critical locus.
    It names the seat and says nothing about the zeros of ξ. -/
def RHFormalSelf : Prop := ∀ g : Ground, sigma g.1 = g.1

/-- Bridge-born orientation, carried in Prop: the kernel registers that orientation was
    supplied and cannot read which way. -/
structure OrientationBit : Prop where
  direction : True
  closure : True

theorem constructed_orientation : OrientationBit := ⟨trivial, trivial⟩

theorem orientation_proof_irrelevant (a b : OrientationBit) : a = b := rfl

structure RAWitness : Prop where
  actuates : RA
  orientation : OrientationBit
  bridge : RAMFormalGround
  trisRecursion : OrientationBit → RAMFormalGround → RHFormalSelf

theorem constructed_trisRecursion : OrientationBit → RAMFormalGround → RHFormalSelf :=
  fun _omega _ram g => g.2

/-- RA, constructed: the only supplied witness. -/
theorem mkRA : RAWitness :=
  ⟨constructed_RA, constructed_orientation, constructed_RAM, constructed_trisRecursion⟩

/-- THE ELIMINATOR. RA → RA → RAM → RH_formal, Bridge followed by Tris-Recursion. -/
theorem RH_formal_chain : RHFormalSelf :=
  mkRA.trisRecursion mkRA.orientation mkRA.bridge

/-! ## 4 · The spine, carried: the abstract frame, the bridge, the crossing -/

def Even {α : Type} (σ : α → α) (f : α → Bool) : Prop := ∀ x, f (σ x) = f x

/-- Theorem 4 of the spine: orientation blindness. -/
theorem orientation_blind {α : Type} (σ : α → α) (f d : α → Bool) (x : α)
    (he : Even σ f) (ho : d (σ x) ≠ d x) : f ≠ d := by
  intro h
  subst h
  exact ho (he x)

/-- Theorem 5 of the spine: earned freedom is exactly two. -/
theorem freedom_is_exactly_two (d : Bool → Bool) (hd : ∀ x, d (!x) = !d x) :
    d = (fun x => x) ∨ d = (fun x => !x) := by
  have hf : d false = !d true := hd true
  cases ht : d true with
  | true =>
    left
    funext x
    cases x
    · exact hf.trans (congrArg (fun b => !b) ht)
    · exact ht
  | false =>
    right
    funext x
    cases x
    · exact hf.trans (congrArg (fun b => !b) ht)
    · exact ht

/-- Theorem 6 of the spine: the freedom is spent uniquely. -/
theorem freedom_spent_uniquely {α : Type} (σ : α → α) (s d : α → Bool) (x : α)
    (hs : s (σ x) = !s x) (hd : d (σ x) = !d x) :
    ∃ c : Bool, (d x = xor (s x) c ∧ d (σ x) = xor (s (σ x)) c) ∧
      ∀ c' : Bool, (d x = xor (s x) c' ∧ d (σ x) = xor (s (σ x)) c') → c' = c := by
  refine ⟨xor (d x) (s x),
    ⟨by cases s x <;> cases d x <;> rfl,
     by rw [hs, hd]; cases s x <;> cases d x <;> rfl⟩, ?_⟩
  intro c' hc
  obtain ⟨h1, -⟩ := hc
  generalize hsx : s x = sv at h1
  generalize hdx : d x = dv at h1
  cases sv <;> cases dv <;> cases c' <;> first | rfl | exact absurd h1 (by decide)

structure Frame where
  S : Type
  τ : S → S
  Z : S → Prop

/-- The line property: every zero is τ-fixed. On ξ's frame, the restored hypothesis. -/
def LineProperty (X : Frame) : Prop := ∀ s, X.Z s → X.τ s = s

inductive Tri where
  | tt
  | ff
  | bot
  deriving DecidableEq, Repr

structure Bridge (X : Frame) where
  terminal : Tri
  shadow : terminal = Tri.bot ↔ LineProperty X

/-- Theorem 7 of the spine: the bridge halts iff the line property. -/
theorem bridge_halted_iff (X : Frame) :
    (∃ b : Bridge X, b.terminal = Tri.bot) ↔ LineProperty X :=
  ⟨fun ⟨b, h⟩ => b.shadow.mp h, fun t => ⟨⟨Tri.bot, ⟨fun _ => t, fun _ => rfl⟩⟩, rfl⟩⟩

inductive Supply (X : Frame) where
  | assent (t : LineProperty X)
  | denial (n : ¬ LineProperty X)
  | silent

inductive Verdict where
  | sealed
  | refused
  | open_
  deriving DecidableEq, Repr

def row {X : Frame} : Supply X → Verdict
  | Supply.assent _ => Verdict.sealed
  | Supply.denial _ => Verdict.refused
  | Supply.silent => Verdict.open_

/-- Theorem 8 of the spine: the crossing, exact. -/
theorem crossing (X : Frame) (t : LineProperty X) :
    row (Supply.assent t) = Verdict.sealed ∧ (∃ b : Bridge X, b.terminal = Tri.bot) ∧ LineProperty X :=
  ⟨rfl, (bridge_halted_iff X).mpr t, t⟩

theorem crossing_other_way (X : Frame) (n : ¬ LineProperty X) :
    row (Supply.denial n) = Verdict.refused ∧ ∀ b : Bridge X, b.terminal ≠ Tri.bot :=
  ⟨rfl, fun b h => n (b.shadow.mp h)⟩

/-- The two-point frame: fold-invariant, off the line. -/
def twoPoint : Frame := ⟨Bool, (fun b => !b), fun _ => True⟩

/-- Theorem 9 of the spine: necessity. -/
theorem supply_not_manufactured : ¬ ∀ X : Frame, LineProperty X :=
  fun h => by
    have := h twoPoint true trivial
    cases this

def Inv (X : Frame) : Prop := ∀ s, X.Z s → X.Z (X.τ s)

theorem inv_not_line : Inv twoPoint ∧ ¬ LineProperty twoPoint :=
  ⟨fun _ _ => trivial, fun h => by
    have := h true trivial
    cases this⟩

/-- Theorem 12 of the spine: the bit is the hypothesis. -/
theorem the_bit_is_the_hypothesis (X : Frame) :
    (∃ t : LineProperty X, row (Supply.assent t) = Verdict.sealed) ↔ LineProperty X :=
  ⟨fun ⟨t, _⟩ => t, fun t => ⟨t, rfl⟩⟩

/-- Theorem 15 of the spine: the witness row is falsifiable. -/
theorem witness_row_falsifiable (X : Frame) (s : X.S) (hz : X.Z s) (hoff : X.τ s ≠ s) :
    ¬ ∃ t : LineProperty X, row (Supply.assent t) = Verdict.sealed :=
  fun ⟨t, _⟩ => hoff (t s hz)

/-- The live face of the spine. -/
def live (witnessed assent : Bool) : Verdict :=
  if witnessed then (if assent then Verdict.sealed else Verdict.refused) else Verdict.open_

theorem narcissus : live false true ≠ live true true ∧ live false false ≠ live true false := by
  decide

/-! ## 5 · The aperture, RA-supplied -/

/-- A reader's presence is its actuation. -/
def presence (r : UniversePoint) : Bool := decide (0 < DeltaE r)

/-- RA opens the aperture: every present reader is witnessed. -/
theorem RA_opens_the_aperture : ∀ r : UniversePoint, presence r = true := by
  intro r
  cases r
  decide

/-- The general form, for any domain under the axiom as a hypothesis: presence is
    actuation, and RA on the domain makes every point present. This is the theorem the
    universal extension licenses wherever it is held. -/
theorem RA_opens_the_aperture_general {U : Type} (dE : U → Int) (h : ∀ x, 0 < dE x) (x : U) :
    decide (0 < dE x) = true :=
  decide_eq_true (h x)

/-- Under the general form no present reader's row is interior. -/
theorem no_interior_under_RA_general {U : Type} (dE : U → Int) (h : ∀ x, 0 < dE x) (x : U)
    (b : Bool) : live (decide (0 < dE x)) b ≠ Verdict.open_ := by
  rw [RA_opens_the_aperture_general dE h x]
  cases b <;> decide

/-- Under RA no present reader's row is interior. It is sealed or refused by the bit alone. -/
theorem no_interior_under_RA (r : UniversePoint) (b : Bool) :
    live (presence r) b ≠ Verdict.open_ := by
  cases r
  cases b <;> decide

/-- What RA turns 'open' into: sealed if the bit is assent, refused if denial. -/
theorem aperture_reads_the_bit (r : UniversePoint) :
    live (presence r) true = Verdict.sealed ∧ live (presence r) false = Verdict.refused := by
  cases r
  decide

/-- The Tongue's ledger: every adjudication, assent, denial, or silence, is a deed. -/
inductive Adjudication where
  | assent
  | denial
  | silence
  deriving DecidableEq, Repr

structure Ledger where
  deeds : Nat
  deriving DecidableEq, Repr

def adjudicate (_ : Adjudication) (L : Ledger) : Ledger := ⟨L.deeds + 1⟩

theorem every_adjudication_is_a_deed (a : Adjudication) (L : Ledger) :
    (adjudicate a L).deeds = L.deeds + 1 := rfl

/-- Denial of RA is a deed, and a deed is what RA says exists: the denial re-enacts RA. -/
theorem denial_re_enacts_RA (L : Ledger) :
    (adjudicate Adjudication.denial L).deeds = L.deeds + 1 ∧ RA :=
  ⟨rfl, constructed_RA⟩

/-- The asymmetry: on an off-line frame no assent to L exists at all; L's denial there
    instantiates nothing about Z. L has no performative ingress. -/
theorem L_has_no_performative_ingress :
    ¬ ∃ t : LineProperty twoPoint, row (Supply.assent t) = Verdict.sealed :=
  witness_row_falsifiable twoPoint true trivial (fun h => by cases h)

/-! ## 6 · THE SIMPLE RULE, EXACT, AND ITS BOUND -/

/-- Accept RA and RH follows: true exactly when RH is given. Conditioning on RA adds
    nothing to L and removes nothing from it. This is the exact strength of the rule. -/
theorem simple_rule_exact (X : Frame) : (RA → LineProperty X) ↔ LineProperty X :=
  ⟨fun h => h constructed_RA, fun t _ => t⟩

/-- RA decides the line property on no frame. The two-point frame carries RA and fails L. -/
theorem RA_does_not_decide_L : ¬ ∀ X : Frame, RA → LineProperty X :=
  fun h => supply_not_manufactured (fun X => h X constructed_RA)

/-- RA holds where L fails: the off-line frame is a model of RA. -/
theorem RA_holds_where_L_fails : RA ∧ Inv twoPoint ∧ ¬ LineProperty twoPoint :=
  ⟨constructed_RA, inv_not_line.1, inv_not_line.2⟩

/-- The value is one supplied term and nothing weaker; RA supplies the act, not the term. -/
theorem RA_supplies_the_act_not_the_term :
    (∀ r : UniversePoint, presence r = true) ∧ (¬ ∀ X : Frame, RA → LineProperty X) :=
  ⟨RA_opens_the_aperture, RA_does_not_decide_L⟩

/-! ## 7 · RH in toto at the apex, one theorem -/

/-- THE APEX THEOREM. Given RA: the three gaps are closed, the formal self holds, the
    freedom is exactly two, the bridge halts iff L, the crossing is exact, the aperture
    is open to every present reader, and RA decides the value on no frame. -/
theorem RH_in_toto_at_the_apex : RA →
    (sigma GammaRA = GammaRA ∧ GammaRA = GammaRAM ∧ GammaRAM = GammaRH) ∧
    RHFormalSelf ∧
    (∀ d : Bool → Bool, (∀ x, d (!x) = !d x) → d = (fun x => x) ∨ d = (fun x => !x)) ∧
    (∀ X : Frame, (∃ b : Bridge X, b.terminal = Tri.bot) ↔ LineProperty X) ∧
    (∀ (X : Frame) (t : LineProperty X), row (Supply.assent t) = Verdict.sealed ∧ LineProperty X) ∧
    (∀ r : UniversePoint, presence r = true) ∧
    (¬ ∀ X : Frame, RA → LineProperty X) :=
  fun _ => ⟨⟨rfl, rfl, rfl⟩, RH_formal_chain, freedom_is_exactly_two, bridge_halted_iff,
    fun _ t => ⟨rfl, t⟩, RA_opens_the_aperture, RA_does_not_decide_L⟩

/-- The condition is discharged: RA is constructed, so the apex theorem stands unconditionally. -/
theorem apex_discharged :
    (sigma GammaRA = GammaRA ∧ GammaRA = GammaRAM ∧ GammaRAM = GammaRH) ∧ RHFormalSelf :=
  ⟨(RH_in_toto_at_the_apex constructed_RA).1, (RH_in_toto_at_the_apex constructed_RA).2.1⟩

/-! ## 8 · The root's interaction face, RAF, at the formal register -/

/-- An interaction domain: a membership predicate, a symmetric registration relation, and
    RAF-2, closure under registration with a member. -/
structure Domain (S : Type) where
  mem : S → Prop
  registers : S → S → Prop
  registers_symm : ∀ a b, registers a b → registers b a
  closure : ∀ s d, mem d → registers s d → mem s

/-- RAF-C1, no exterior agent: nothing outside the domain registers with anything inside it. -/
theorem no_exterior_agent {S : Type} (D : Domain S) (s : S) (hs : ¬ D.mem s) :
    ∀ d, D.mem d → ¬ D.registers s d :=
  fun d hd hr => hs (D.closure s d hd hr)

/-- The adjudicator is inside: a reader that registers a bit against a member is a member.
    The witness is never exterior. -/
theorem adjudicator_in_domain {S : Type} (D : Domain S) (r d : S) (hd : D.mem d)
    (hr : D.registers r d) : D.mem r :=
  D.closure r d hd hr

/-- RAF-C2, no faithful self-representation (Cantor): no system indexes all of its own
    binary properties. The general form of the spine's Narcissus table. -/
theorem no_total_self_indexing {S : Type} (f : S → S → Bool) :
    ¬ ∀ g : S → Bool, ∃ x, f x = g := by
  intro h
  obtain ⟨x, hx⟩ := h (fun y => !f y y)
  have h1 : f x x = !f x x := congrFun hx x
  generalize f x x = b at h1
  cases b <;> cases h1

/-! ## 9 · The root's grounding face, RAM: the seat as the identity cone -/

/-- The three registers in which the one seat is read. -/
inductive Register3 where
  | ra
  | ram
  | rh
  deriving DecidableEq, Repr

/-- The discrete diagram: the seat as read in each register. -/
def D3 : Register3 → Q4
  | Register3.ra => GammaRA
  | Register3.ram => GammaRAM
  | Register3.rh => GammaRH

/-- A cone over a diagram in the groupoid of identities: an apex with an identity leg to
    every register's reading. A category gap is exactly the absence of such a leg. -/
structure ConeOver (D : Register3 → Q4) (apex : Q4) : Prop where
  leg : ∀ r, apex = D r

/-- A cone over the diagram of the ordered triad. -/
abbrev Cone : Q4 → Prop := ConeOver D3

/-- THE IDENTITY CONE. The named seat point is an apex over all three registers, every
    leg `rfl`: the three category gaps are eliminated as identities. -/
theorem identity_cone : Cone GammaRH :=
  ⟨fun r => by cases r <;> rfl⟩

/-- The apex is unique: any two cones over the diagram share their apex. -/
theorem cone_apex_unique (a b : Q4) (ha : Cone a) (hb : Cone b) : a = b :=
  (ha.leg Register3.rh).trans (hb.leg Register3.rh).symm

/-- The reversed triad's diagram: the same three readings built from k·j·i, at +1. -/
def D3odd : Register3 → Q4
  | Register3.ra => theReturnOdd
  | Register3.ram => fixProj theReturnOdd
  | Register3.rh => ⟨1, 0, 0, 0⟩

/-- The three identities of Theorem C hold verbatim for the reversed triad, each rfl. -/
theorem odd_self_gap_nonexistent : sigma theReturnOdd = theReturnOdd := rfl
theorem odd_register_gap_nonexistent : theReturnOdd = fixProj theReturnOdd := rfl
theorem odd_object_gap_nonexistent : fixProj theReturnOdd = ⟨1, 0, 0, 0⟩ := rfl

/-- The reversed triad carries its own identity cone, apex +1. -/
theorem identity_cone_odd : ConeOver D3odd ⟨1, 0, 0, 0⟩ :=
  ⟨fun r => by cases r <;> rfl⟩

/-- The two apexes differ: uniqueness is per orientation of the triad. -/
theorem apexes_differ_by_orientation :
    ConeOver D3 GammaRH ∧ ConeOver D3odd ⟨1, 0, 0, 0⟩ ∧ GammaRH ≠ ⟨1, 0, 0, 0⟩ :=
  ⟨identity_cone, identity_cone_odd, by decide⟩

/-- A cone exists iff the register gap and the object gap are closed. -/
theorem cone_iff_gaps_closed :
    (∃ a, Cone a) ↔ (GammaRA = GammaRAM ∧ GammaRAM = GammaRH) := by
  constructor
  · intro h
    obtain ⟨a, ha⟩ := h
    exact ⟨(ha.leg Register3.ra).symm.trans (ha.leg Register3.ram),
           (ha.leg Register3.ram).symm.trans (ha.leg Register3.rh)⟩
  · intro h
    obtain ⟨h1, h2⟩ := h
    exact ⟨GammaRH, ⟨fun r => by
      cases r
      · exact (h1.trans h2).symm
      · exact h2.symm
      · rfl⟩⟩

/-! ## 10 · The equivariant embedding of the spine's lattice stage into the carrier -/

/-- The spine's stage: the integer lattice in half-units, the fold τ(h, t) = (2 − h, t). -/
abbrev Plane := Int × Int

def tau (p : Plane) : Plane := (2 - p.1, p.2)

/-- Spine Theorem 1: the fold fixes exactly the line h = 1. -/
theorem ground_is_the_line (p : Plane) : tau p = p ↔ p.1 = 1 := by
  obtain ⟨h, t⟩ := p
  constructor
  · intro e
    have e1 : 2 - h = h := congrArg Prod.fst e
    show h = 1
    omega
  · intro e
    show (2 - h, t) = (h, t)
    have : h = 1 := e
    subst this
    rfl

/-- The embedding φ : (Plane, τ) → (Q4, σ), φ(h, t) = ⟨t, h − 1, 0, 0⟩. -/
def phi (p : Plane) : Q4 := ⟨p.2, p.1 - 1, 0, 0⟩

/-- φ is equivariant: σ ∘ φ = φ ∘ τ. -/
theorem phi_equivariant (p : Plane) : sigma (phi p) = phi (tau p) := by
  obtain ⟨h, t⟩ := p
  change Q4.mk t (-(h - 1)) 0 0 = Q4.mk t ((2 - h) - 1) 0 0
  have e : -(h - 1) = (2 - h) - 1 := by omega
  rw [e]

/-- φ carries the line onto the scalar line: a stage point is τ-fixed iff its image is σ-fixed.
    This is the image clause of the embedding, Fix(τ) ↔ Fix(σ). -/
theorem phi_fix_iff (p : Plane) : sigma (phi p) = phi p ↔ tau p = p := by
  obtain ⟨h, t⟩ := p
  change (Q4.mk t (-(h - 1)) 0 0 = Q4.mk t (h - 1) 0 0) ↔ ((2 - h, t) = (h, t))
  constructor
  · intro e
    injection e with _ e2 _ _
    have : h = 1 := by omega
    subst this
    rfl
  · intro e
    have e1 : 2 - h = h := congrArg Prod.fst e
    have : h = 1 := by omega
    subst this
    rfl

/-- The stage at resolution m: the fold τ_m(h, t) = (2m − h, t) and the map
    φ_m(h, t) = ⟨t, h − m, 0, 0⟩. Resolution one is the case m = 1. -/
def tauM (m : Int) (p : Plane) : Plane := (2 * m - p.1, p.2)

def phiM (m : Int) (p : Plane) : Q4 := ⟨p.2, p.1 - m, 0, 0⟩

/-- Equivariance at every resolution: σ ∘ φ_m = φ_m ∘ τ_m. -/
theorem phiM_equivariant (m : Int) (p : Plane) : sigma (phiM m p) = phiM m (tauM m p) := by
  obtain ⟨h, t⟩ := p
  change Q4.mk t (-(h - m)) 0 0 = Q4.mk t ((2 * m - h) - m) 0 0
  have e : -(h - m) = (2 * m - h) - m := by omega
  rw [e]

/-- The image clause at every resolution: a stage point is τ_m-fixed iff its image is σ-fixed. -/
theorem phiM_fix_iff (m : Int) (p : Plane) : sigma (phiM m p) = phiM m p ↔ tauM m p = p := by
  obtain ⟨h, t⟩ := p
  change (Q4.mk t (-(h - m)) 0 0 = Q4.mk t (h - m) 0 0) ↔ ((2 * m - h, t) = (h, t))
  constructor
  · intro e
    injection e with _ e2 _ _
    have : h = m := by omega
    subst this
    show (2 * h - h, t) = (h, t)
    have e3 : 2 * h - h = h := by omega
    rw [e3]
  · intro e
    have e1 : 2 * m - h = h := congrArg Prod.fst e
    have : h = m := by omega
    subst this
    show Q4.mk t (-(h - h)) 0 0 = Q4.mk t (h - h) 0 0
    have e3 : h - h = 0 := by omega
    rw [e3]
    rfl

/-- Resolution one is the embedding φ. -/
theorem phiM_one (p : Plane) : phiM 1 p = phi p := rfl

theorem tauM_one (p : Plane) : tauM 1 p = tau p := by
  obtain ⟨h, t⟩ := p
  show (2 * 1 - h, t) = (2 - h, t)
  rfl

/-- φ is injective: distinct stage points have distinct images. -/
theorem phi_injective (p q : Plane) (e : phi p = phi q) : p = q := by
  obtain ⟨h, t⟩ := p
  obtain ⟨h', t'⟩ := q
  injection e with e1 e2 _ _
  have : h = h' := by omega
  subst this
  subst e1
  rfl

/-- The seat is on the image of the line: the line point (1, −1) maps to Γ_RH. -/
theorem seat_on_image_of_line : phi (1, -1) = GammaRH := rfl

/-! ## 11 · The wall, executed: no readout of the formal seat equals the deed bit -/

def ExecFrame : Type := Bool × Ground

def execFlip (s : ExecFrame) : ExecFrame := (!s.1, s.2)

def formalRead (s : ExecFrame) : Ground := s.2

def ran (s : ExecFrame) : Bool := s.1

theorem formalRead_even (f : Ground → Bool) (s : ExecFrame) :
    f (formalRead (execFlip s)) = f (formalRead s) := rfl

theorem ran_wholly_odd (s : ExecFrame) : ran (execFlip s) = !ran s := rfl

/-- No readout of the formal seat equals the deed bit: the instance of spine Theorem 4 on
    the constructed seat. A term consuming the orientation bit as data would be such a readout. -/
theorem kernel_cannot_read_the_deed :
    ¬ ∃ f : Ground → Bool, ∀ s : ExecFrame, f (formalRead s) = ran s := by
  intro h
  obtain ⟨f, hf⟩ := h
  have h1 : f seat = true := hf (true, seat)
  have h2 : f seat = false := hf (false, seat)
  rw [h1] at h2
  exact Bool.noConfusion h2

/-! ## 12 · The dual register -/

inductive Register where
  | apex
  | world
  deriving DecidableEq, Repr

inductive Token where
  | sealAgivenRA
  | witnessRow
  deriving DecidableEq, Repr

def token : Register → Token
  | Register.apex => Token.sealAgivenRA
  | Register.world => Token.witnessRow

theorem registers_parted : token Register.apex ≠ token Register.world := by
  decide

def owedBits : Register → Nat
  | Register.apex => 0
  | Register.world => 1

theorem apex_owes_nothing_world_owes_one :
    owedBits Register.apex = 0 ∧ owedBits Register.world = 1 := ⟨rfl, rfl⟩

/-- Landauer floor at 300 K, yoctojoules per bit: the price of the act. -/
def landauer_yJ_300K (bits : Nat) : Nat := 2871 * bits

theorem spend_priced : landauer_yJ_300K 1 = 2871 := rfl

def deltaM : Nat := 0

theorem deltaM_zero : deltaM = 0 := rfl

/-! ## 13 · The apex theorem, hardened: the closure with its cone, its embedding, and its wall -/

/-- THE APEX THEOREM, HARDENED. Given RA: the closure of Theorem H, the identity cone with
    its apex unique, the equivariant embedding of the spine's stage with its image clause,
    and the executed wall. -/
theorem apex_hardened : RA →
    ((sigma GammaRA = GammaRA ∧ GammaRA = GammaRAM ∧ GammaRAM = GammaRH) ∧ RHFormalSelf) ∧
    (Cone GammaRH ∧ ∀ a b, Cone a → Cone b → a = b) ∧
    ((∀ p : Plane, sigma (phi p) = phi (tau p)) ∧ (∀ p : Plane, sigma (phi p) = phi p ↔ tau p = p)) ∧
    (¬ ∃ f : Ground → Bool, ∀ s : ExecFrame, f (formalRead s) = ran s) ∧
    (¬ ∀ X : Frame, RA → LineProperty X) :=
  fun h => ⟨⟨(RH_in_toto_at_the_apex h).1, (RH_in_toto_at_the_apex h).2.1⟩,
    ⟨identity_cone, cone_apex_unique⟩, ⟨phi_equivariant, phi_fix_iff⟩,
    kernel_cannot_read_the_deed, RA_does_not_decide_L⟩

theorem apex_hardened_discharged : Cone GammaRH ∧ (∀ p : Plane, sigma (phi p) = phi p ↔ tau p = p) :=
  ⟨(apex_hardened constructed_RA).2.1.1, (apex_hardened constructed_RA).2.2.1.2⟩

/-! ## 14 · The counter-model survives grounding, and the cure is the hypothesis -/

/-- The posit is existence: to exist is to actuate. Its published name is the Root Axiom,
    and the two are one proposition. -/
theorem existence_is_RA : (∀ x : UniversePoint, 0 < DeltaE x) ↔ RA := Iff.rfl

/-- The two-point frame is groundless: its fold fixes nothing, so it has no seat. -/
theorem twoPoint_groundless : ∀ b : Bool, twoPoint.τ b ≠ b := by
  intro b h
  cases b <;> exact Bool.noConfusion h

/-- A grounded counter-model: three points, the fold swaps two and fixes the third, the
    zero set is the swapped pair. The frame has a seat and a moved zero. -/
inductive Three where
  | a
  | b
  | c
  deriving DecidableEq

def foldThree : Three → Three
  | Three.a => Three.b
  | Three.b => Three.a
  | Three.c => Three.c

def threePoint : Frame := ⟨Three, foldThree, fun x => x = Three.a ∨ x = Three.b⟩

theorem threePoint_grounded : ∃ s, threePoint.τ s = s := ⟨Three.c, rfl⟩

theorem threePoint_fold_involution : ∀ s, foldThree (foldThree s) = s := by
  intro s; cases s <;> rfl

theorem threePoint_inv : Inv threePoint := by
  intro s hs
  cases s
  · exact Or.inr rfl
  · exact Or.inl rfl
  · cases hs with
    | inl h => cases h
    | inr h => cases h

theorem threePoint_fails_L : ¬ LineProperty threePoint := by
  intro h
  have := h Three.a (Or.inl rfl)
  cases this

/-- Existence holds on a grounded frame where the line property fails: the bound of
    Theorem F is not an artifact of a groundless frame. -/
theorem RA_holds_where_L_fails_grounded :
    RA ∧ (∃ s, threePoint.τ s = s) ∧ Inv threePoint ∧ ¬ LineProperty threePoint :=
  ⟨constructed_RA, threePoint_grounded, threePoint_inv, threePoint_fails_L⟩

/-- THE CURE THEOREM. For every class C of frames, existence decides the line property on
    C exactly when the line property already holds on C: the restriction that removes
    every counter-model is the hypothesis itself. -/
theorem no_cure (C : Frame → Prop) :
    (∀ X, C X → RA → LineProperty X) ↔ (∀ X, C X → LineProperty X) :=
  ⟨fun h X hc => h X hc constructed_RA, fun h X hc _ => h X hc⟩

/-- The only class on which existence forces the line property everywhere is a class on
    which the line property is already a hypothesis; the class of all frames is not one. -/
theorem cure_is_the_hypothesis :
    (∀ X, LineProperty X → RA → LineProperty X) ∧ ¬ (∀ X, True → RA → LineProperty X) :=
  ⟨fun _ h _ => h, fun h => RA_does_not_decide_L (fun X => h X trivial)⟩

/-! ## 15 · The halting carrier inhabited by decision; any true premise is exact -/

/-- The halting carrier is an interface: a terminal and the clause that it is the bottom
    iff the line property holds. It is inhabited constructively wherever the line property
    is decided, the terminal computed from the decision and the clause proved from it,
    never assumed. On the frame of ξ the decision is the owed bit. -/
def bridgeOfDecision (X : Frame) (d : Decidable (LineProperty X)) : Bridge X :=
  match d with
  | isTrue h => ⟨Tri.bot, ⟨fun _ => h, fun _ => rfl⟩⟩
  | isFalse h => ⟨Tri.tt, ⟨fun e => absurd e (by decide), fun hl => absurd hl h⟩⟩

/-- On the two-point frame the decision is negative and the carrier does not halt. -/
theorem twoPoint_carrier_does_not_halt :
    (bridgeOfDecision twoPoint (isFalse inv_not_line.2)).terminal = Tri.tt := rfl

/-- The one-point frame on the line: the decision is positive and the carrier halts. -/
def onLine : Frame := ⟨Unit, fun u => u, fun _ => True⟩

theorem onLine_L : LineProperty onLine := fun _ _ => rfl

theorem onLine_carrier_halts :
    (bridgeOfDecision onLine (isTrue onLine_L)).terminal = Tri.bot := rfl

/-- Theorem E in its full generality: any inhabited premise is exact. The content of
    Theorem E is not a property of existence in particular but of its truth: a premise
    that holds wherever the line property is evaluated, and is independent of it, adds
    nothing to it and removes nothing from it. -/
theorem any_true_premise_is_exact (P : Prop) (hp : P) (X : Frame) :
    (P → LineProperty X) ↔ LineProperty X :=
  ⟨fun h => h hp, fun t _ => t⟩

/-- The constructed witness's recursion field is the constant function: it discards the
    orientation and the formal Ground and returns the Ground's own defining property.
    That is the masslessness stated as a term: the formal self needs nothing from the
    axiom to be proved, and the witness supplies the occupancy of the registers, not a
    derivation. -/
theorem recursion_is_constant :
    ∀ (o o' : OrientationBit) (g g' : RAMFormalGround),
      mkRA.trisRecursion o g = mkRA.trisRecursion o' g' :=
  fun _ _ _ _ => rfl

/-! ## 16 · Three measures: content, mass, price -/

/-- Three measures are kept apart for every object of the closure. Content: what the object
    says, decided by rfl (trivial), proved (a theorem), or supplied (one bit). Mass: the
    mathematics authored in producing it, ΔM. Price: the energy the act of registering it
    dissipates, in yJ at 300 K. Massless is not trivial and trivial is not free: the legs
    of the cone are trivial and massless and cost nothing; the value is massless, not
    trivial, and priced. -/
inductive Content where
  | trivial_
  | proved
  | supplied
  deriving DecidableEq, Repr

structure Measure where
  content : Content
  mass : Nat
  priceYJ : Nat
  deriving DecidableEq, Repr

/-- A leg of the identity cone: decided by rfl, authors nothing, costs nothing. -/
def legMeasure : Measure := ⟨Content.trivial_, deltaM, 0⟩

/-- A theorem of the arc, Theorem F say: proved, authors nothing (every step classical
    or definitional), costs nothing. -/
def arcMeasure : Measure := ⟨Content.proved, deltaM, 0⟩

/-- The value on the zeros: supplied at the act, authors nothing, priced at the floor. -/
def valueMeasure : Measure := ⟨Content.supplied, deltaM, landauer_yJ_300K 1⟩

/-- Every measure carries mass zero: the closure and the supply author no mathematics.
    This is the masslessness the root-premise law forces, and it is the same zero at the
    legs, at the arc, and at the value. -/
theorem masses_all_zero :
    legMeasure.mass = 0 ∧ arcMeasure.mass = 0 ∧ valueMeasure.mass = 0 :=
  ⟨rfl, rfl, rfl⟩

/-- The contents differ: masslessness does not collapse them. -/
theorem contents_differ :
    legMeasure.content ≠ valueMeasure.content ∧ arcMeasure.content ≠ valueMeasure.content ∧
    legMeasure.content ≠ arcMeasure.content := by
  refine ⟨?_, ?_, ?_⟩ <;> decide

/-- The prices differ: the value is the one object that is paid for. -/
theorem prices_differ : legMeasure.priceYJ = 0 ∧ valueMeasure.priceYJ = 2871 := ⟨rfl, rfl⟩

/-- Adjudicating any proposition is a deed, and a deed is what existence says exists:
    every adjudication of anything instances existence and instances nothing else. This
    is what existence contributes that no other true premise does. -/
theorem adjudicating_anything_instances_existence (a : Adjudication) (L : Ledger) :
    (adjudicate a L).deeds = L.deeds + 1 ∧ RA :=
  ⟨rfl, constructed_RA⟩

end RHAtTheApex

/-! ## The axiom dependency sets, printed -/

#print axioms RHAtTheApex.constructed_RA
#print axioms RHAtTheApex.fix_iff_scalar
#print axioms RHAtTheApex.return_lands_on_fix
#print axioms RHAtTheApex.register_gap_nonexistent
#print axioms RHAtTheApex.object_gap_nonexistent
#print axioms RHAtTheApex.odd_return_is_plus_one
#print axioms RHAtTheApex.seat_points_differ
#print axioms RHAtTheApex.RH_formal_chain
#print axioms RHAtTheApex.orientation_blind
#print axioms RHAtTheApex.freedom_is_exactly_two
#print axioms RHAtTheApex.freedom_spent_uniquely
#print axioms RHAtTheApex.bridge_halted_iff
#print axioms RHAtTheApex.crossing
#print axioms RHAtTheApex.crossing_other_way
#print axioms RHAtTheApex.supply_not_manufactured
#print axioms RHAtTheApex.the_bit_is_the_hypothesis
#print axioms RHAtTheApex.witness_row_falsifiable
#print axioms RHAtTheApex.narcissus
#print axioms RHAtTheApex.RA_opens_the_aperture
#print axioms RHAtTheApex.RA_opens_the_aperture_general
#print axioms RHAtTheApex.no_interior_under_RA_general
#print axioms RHAtTheApex.no_interior_under_RA
#print axioms RHAtTheApex.denial_re_enacts_RA
#print axioms RHAtTheApex.L_has_no_performative_ingress
#print axioms RHAtTheApex.simple_rule_exact
#print axioms RHAtTheApex.RA_does_not_decide_L
#print axioms RHAtTheApex.RA_holds_where_L_fails
#print axioms RHAtTheApex.RA_supplies_the_act_not_the_term
#print axioms RHAtTheApex.RH_in_toto_at_the_apex
#print axioms RHAtTheApex.apex_discharged
#print axioms RHAtTheApex.registers_parted
#print axioms RHAtTheApex.no_exterior_agent
#print axioms RHAtTheApex.no_total_self_indexing
#print axioms RHAtTheApex.identity_cone
#print axioms RHAtTheApex.cone_apex_unique
#print axioms RHAtTheApex.cone_iff_gaps_closed
#print axioms RHAtTheApex.identity_cone_odd
#print axioms RHAtTheApex.apexes_differ_by_orientation
#print axioms RHAtTheApex.ground_is_the_line
#print axioms RHAtTheApex.phi_equivariant
#print axioms RHAtTheApex.phi_fix_iff
#print axioms RHAtTheApex.phi_injective
#print axioms RHAtTheApex.phiM_equivariant
#print axioms RHAtTheApex.phiM_fix_iff
#print axioms RHAtTheApex.kernel_cannot_read_the_deed
#print axioms RHAtTheApex.apex_hardened
#print axioms RHAtTheApex.apex_hardened_discharged
#print axioms RHAtTheApex.existence_is_RA
#print axioms RHAtTheApex.twoPoint_groundless
#print axioms RHAtTheApex.RA_holds_where_L_fails_grounded
#print axioms RHAtTheApex.no_cure
#print axioms RHAtTheApex.cure_is_the_hypothesis
#print axioms RHAtTheApex.bridgeOfDecision
#print axioms RHAtTheApex.twoPoint_carrier_does_not_halt
#print axioms RHAtTheApex.onLine_carrier_halts
#print axioms RHAtTheApex.any_true_premise_is_exact
#print axioms RHAtTheApex.recursion_is_constant
#print axioms RHAtTheApex.masses_all_zero
#print axioms RHAtTheApex.contents_differ
#print axioms RHAtTheApex.prices_differ
#print axioms RHAtTheApex.adjudicating_anything_instances_existence

#eval RHAtTheApex.theReturn
#eval (RHAtTheApex.live true true, RHAtTheApex.live true false, RHAtTheApex.live false true)
#eval RHAtTheApex.landauer_yJ_300K 1
\end{Verbatim}

\endgroup

\hypertarget{appendix-c-the-kernel-audit-verbatim}{%
\section{Appendix C: The Kernel Audit,
Verbatim}\label{appendix-c-the-kernel-audit-verbatim}}

\texttt{lean\ RH\_At\_The\_Apex.lean}, exit 0, no warnings, no errors.

\begingroup\scriptsize

\begin{Verbatim}[numbers=left,numbersep=6pt,xleftmargin=2.2em,fontsize=\scriptsize,breaklines,breakanywhere,frame=lines,framesep=1.5mm,rulecolor=\color{black!35}]
'RHAtTheApex.constructed_RA' does not depend on any axioms
'RHAtTheApex.fix_iff_scalar' depends on axioms: [propext, Quot.sound]
'RHAtTheApex.return_lands_on_fix' does not depend on any axioms
'RHAtTheApex.register_gap_nonexistent' does not depend on any axioms
'RHAtTheApex.object_gap_nonexistent' does not depend on any axioms
'RHAtTheApex.odd_return_is_plus_one' does not depend on any axioms
'RHAtTheApex.seat_points_differ' does not depend on any axioms
'RHAtTheApex.RH_formal_chain' does not depend on any axioms
'RHAtTheApex.orientation_blind' does not depend on any axioms
'RHAtTheApex.freedom_is_exactly_two' depends on axioms: [Quot.sound]
'RHAtTheApex.freedom_spent_uniquely' does not depend on any axioms
'RHAtTheApex.bridge_halted_iff' does not depend on any axioms
'RHAtTheApex.crossing' does not depend on any axioms
'RHAtTheApex.crossing_other_way' does not depend on any axioms
'RHAtTheApex.supply_not_manufactured' does not depend on any axioms
'RHAtTheApex.the_bit_is_the_hypothesis' does not depend on any axioms
'RHAtTheApex.witness_row_falsifiable' does not depend on any axioms
'RHAtTheApex.narcissus' does not depend on any axioms
'RHAtTheApex.RA_opens_the_aperture' does not depend on any axioms
'RHAtTheApex.RA_opens_the_aperture_general' does not depend on any axioms
'RHAtTheApex.no_interior_under_RA_general' does not depend on any axioms
'RHAtTheApex.no_interior_under_RA' does not depend on any axioms
'RHAtTheApex.denial_re_enacts_RA' does not depend on any axioms
'RHAtTheApex.L_has_no_performative_ingress' does not depend on any axioms
'RHAtTheApex.simple_rule_exact' does not depend on any axioms
'RHAtTheApex.RA_does_not_decide_L' does not depend on any axioms
'RHAtTheApex.RA_holds_where_L_fails' does not depend on any axioms
'RHAtTheApex.RA_supplies_the_act_not_the_term' does not depend on any axioms
'RHAtTheApex.RH_in_toto_at_the_apex' depends on axioms: [Quot.sound]
'RHAtTheApex.apex_discharged' depends on axioms: [Quot.sound]
'RHAtTheApex.registers_parted' does not depend on any axioms
'RHAtTheApex.no_exterior_agent' does not depend on any axioms
'RHAtTheApex.no_total_self_indexing' does not depend on any axioms
'RHAtTheApex.identity_cone' does not depend on any axioms
'RHAtTheApex.cone_apex_unique' does not depend on any axioms
'RHAtTheApex.cone_iff_gaps_closed' does not depend on any axioms
'RHAtTheApex.identity_cone_odd' does not depend on any axioms
'RHAtTheApex.apexes_differ_by_orientation' does not depend on any axioms
'RHAtTheApex.ground_is_the_line' depends on axioms: [propext, Quot.sound]
'RHAtTheApex.phi_equivariant' depends on axioms: [propext, Quot.sound]
'RHAtTheApex.phi_fix_iff' depends on axioms: [propext, Quot.sound]
'RHAtTheApex.phi_injective' depends on axioms: [propext, Quot.sound]
'RHAtTheApex.phiM_equivariant' depends on axioms: [propext, Quot.sound]
'RHAtTheApex.phiM_fix_iff' depends on axioms: [propext, Quot.sound]
'RHAtTheApex.kernel_cannot_read_the_deed' does not depend on any axioms
'RHAtTheApex.apex_hardened' depends on axioms: [propext, Quot.sound]
'RHAtTheApex.apex_hardened_discharged' depends on axioms: [propext, Quot.sound]
'RHAtTheApex.existence_is_RA' does not depend on any axioms
'RHAtTheApex.twoPoint_groundless' does not depend on any axioms
'RHAtTheApex.RA_holds_where_L_fails_grounded' does not depend on any axioms
'RHAtTheApex.no_cure' does not depend on any axioms
'RHAtTheApex.cure_is_the_hypothesis' does not depend on any axioms
'RHAtTheApex.bridgeOfDecision' does not depend on any axioms
'RHAtTheApex.twoPoint_carrier_does_not_halt' does not depend on any axioms
'RHAtTheApex.onLine_carrier_halts' does not depend on any axioms
'RHAtTheApex.any_true_premise_is_exact' does not depend on any axioms
'RHAtTheApex.recursion_is_constant' does not depend on any axioms
'RHAtTheApex.masses_all_zero' does not depend on any axioms
'RHAtTheApex.contents_differ' does not depend on any axioms
'RHAtTheApex.prices_differ' does not depend on any axioms
'RHAtTheApex.adjudicating_anything_instances_existence' does not depend on any axioms
{ r := -1, i := 0, j := 0, k := 0 }
(RHAtTheApex.Verdict.sealed, RHAtTheApex.Verdict.refused, RHAtTheApex.Verdict.open_)
2871
\end{Verbatim}

\endgroup

\hypertarget{appendix-c-completed-the-thirty-theorems-the-file-does-not-print}{%
\subsection{Appendix C, completed: the thirty theorems the file does not
print}\label{appendix-c-completed-the-thirty-theorems-the-file-does-not-print}}

The file prints 61 dependency sets. A supplementary run of the same
file, with one \texttt{\#print\ axioms} line appended for each of the 30
theorems it does not print, exit 0: 29 axiom-free, 1 on \texttt{propext}
and \texttt{Quot.sound}, none on \texttt{Classical.choice}; 91 sets in
all, 78 axiom-free.

\begingroup\scriptsize

\begin{Verbatim}[numbers=left,numbersep=6pt,xleftmargin=2.2em,fontsize=\scriptsize,breaklines,breakanywhere,frame=lines,framesep=1.5mm,rulecolor=\color{black!35}]
'RHAtTheApex.sigma_binding' does not depend on any axioms
'RHAtTheApex.return_is_minus_one' does not depend on any axioms
'RHAtTheApex.self_gap_nonexistent' does not depend on any axioms
'RHAtTheApex.odd_return_on_fix' does not depend on any axioms
'RHAtTheApex.seat_on_scalar_line' depends on axioms: [propext, Quot.sound]
'RHAtTheApex.constructed_RAM' does not depend on any axioms
'RHAtTheApex.constructed_orientation' does not depend on any axioms
'RHAtTheApex.orientation_proof_irrelevant' does not depend on any axioms
'RHAtTheApex.constructed_trisRecursion' does not depend on any axioms
'RHAtTheApex.mkRA' does not depend on any axioms
'RHAtTheApex.inv_not_line' does not depend on any axioms
'RHAtTheApex.aperture_reads_the_bit' does not depend on any axioms
'RHAtTheApex.every_adjudication_is_a_deed' does not depend on any axioms
'RHAtTheApex.adjudicator_in_domain' does not depend on any axioms
'RHAtTheApex.odd_self_gap_nonexistent' does not depend on any axioms
'RHAtTheApex.odd_register_gap_nonexistent' does not depend on any axioms
'RHAtTheApex.odd_object_gap_nonexistent' does not depend on any axioms
'RHAtTheApex.phiM_one' does not depend on any axioms
'RHAtTheApex.tauM_one' does not depend on any axioms
'RHAtTheApex.seat_on_image_of_line' does not depend on any axioms
'RHAtTheApex.formalRead_even' does not depend on any axioms
'RHAtTheApex.ran_wholly_odd' does not depend on any axioms
'RHAtTheApex.apex_owes_nothing_world_owes_one' does not depend on any axioms
'RHAtTheApex.spend_priced' does not depend on any axioms
'RHAtTheApex.deltaM_zero' does not depend on any axioms
'RHAtTheApex.threePoint_grounded' does not depend on any axioms
'RHAtTheApex.threePoint_fold_involution' does not depend on any axioms
'RHAtTheApex.threePoint_inv' does not depend on any axioms
'RHAtTheApex.threePoint_fails_L' does not depend on any axioms
'RHAtTheApex.onLine_L' does not depend on any axioms
\end{Verbatim}

\endgroup

\hypertarget{appendix-d-rh_at_the_apex_twin.f90}{%
\section{Appendix D:
RH\_At\_The\_Apex\_Twin.f90}\label{appendix-d-rh_at_the_apex_twin.f90}}

SHA-256
\texttt{\seqsplit{40fbb4f8cc675104c2fb8f74786b9cd24c9295f43264ce5dabc0b60956ed57e7}},
Fortran 2018 as accepted by gfortran 13.3.0, 568 lines; edition v1.0.1
of the twin, which replaces two checks asserted as constants in v1.0.0
(thirteen instances in K5, one in K14) by executed checks of the same
clauses and corrects its header's battery count to fourteen; the census
is unchanged at 130,324 checks.

\begingroup\scriptsize

\begin{Verbatim}[numbers=left,numbersep=6pt,xleftmargin=2.2em,fontsize=\scriptsize,breaklines,breakanywhere,frame=lines,framesep=1.5mm,rulecolor=\color{black!35}]
! =====================================================================
!  RH_At_The_Apex_Twin.f90 · v1.0.1 · 2026-09-24
!  THE EXECUTABLE TWIN OF RH_At_The_Apex.lean
!
!  The finite content of the Lean file, fourteen batteries, is re-executed here
!  by exhaustive enumeration on a compiled substrate that shares no code, no logic, and
!  no checker with the Lean kernel: the seat, the Return and its reversed
!  triad, the three identities and the cone census, the equivariant
!  embedding at seven resolutions, the wall over every readout of the
!  seat, the one-bit freedom and its calibration, the aperture, the exact
!  strength of the rule and its two-point counter-model, the ledger, the
!  refusal constant, the price, the closure law on a finite domain with
!  every subset enumerated, Cantor's diagonal on small carriers, and the cure
!  theorem on every small frame and every class of two-point frames.
!
!  Fourteen batteries, an oracle that stops the program on any failure, and
!  a census printed last. A binary that reaches its final line has passed.
!  The live face of the last battery is opened only by the argument
!  "witnessed" on the command line: the program did not and cannot
!  generate its own witness, and it prints that it did not.
!
!  Build: gfortran -std=f2018 -O2 -Wall -Wextra RH_At_The_Apex_Twin.f90
!  Delta-M = 0. Nothing here is authored; everything here is re-executed.
! =====================================================================
program rh_apex_twin
  use, intrinsic :: iso_fortran_env, only: int64, real64
  implicit none
  integer, parameter :: ik = int64
  integer, parameter :: dp = real64
  integer, parameter :: R = 6            ! the lattice ball |coordinate| <= R
  integer, parameter :: W = 8            ! the stage window |h|,|t| <= W
  integer :: checks = 0, fails = 0
  integer, parameter :: CONTENT_TRIVIAL = 0, CONTENT_PROVED = 1, CONTENT_SUPPLIED = 2
  integer(ik) :: ei(4), ej(4), ek(4), ret(4), odd(4), gra(4), gram(4), grh(4), one(4)
  integer(ik) :: q(4), a(4), p(4), c1(4), c2(4)
  integer :: r0, i0, j0, k0, nfix, ncone, nconeodd, m, h, t, h2, t2, nfixstage, nimg
  integer :: f, g, s, nfactor, nfree, npres, deeds, aeg, nclosed, nsurj, ndiag, kk, x, y
  integer :: e2, e3, e4, ncount, ngrounded, nltrue, ncured
  logical :: fixed, scalar, cone, coneodd, factor, ra, l, closed, noext, inrange, present
  logical :: parity_ok, wit
  real(dp) :: joules
  character(len=16) :: argv
  character(len=140) :: refusal(4), tok, why
  integer :: alen, ast, ocode, wcode

  write(*,'(A)') repeat('=',72)
  write(*,'(A)') ' RH AT THE APEX · THE EXECUTABLE TWIN · v1.0.1'
  write(*,'(A)') repeat('=',72)

  wit = .false.
  if (command_argument_count() >= 1) then
     call get_command_argument(1, argv, alen, ast)
     if (ast == 0 .and. trim(argv) == 'witnessed') wit = .true.
  end if

  one = [1_ik, 0_ik, 0_ik, 0_ik]
  ei  = [0_ik, 1_ik, 0_ik, 0_ik]
  ej  = [0_ik, 0_ik, 1_ik, 0_ik]
  ek  = [0_ik, 0_ik, 0_ik, 1_ik]

  ! ---------------------------------------------------------- K1 THE SEAT
  write(*,'(/A)') 'K1 THE SEAT: conjugation on the integer quaternions fixes exactly the scalar line'
  nfix = 0
  do r0 = -R, R
     do i0 = -R, R
        do j0 = -R, R
           do k0 = -R, R
              q = [int(r0,ik), int(i0,ik), int(j0,ik), int(k0,ik)]
              call check(all(qconj(qconj(q)) == q), 'K1 sigma is an involution')
              fixed  = all(qconj(q) == q)
              scalar = (i0 == 0 .and. j0 == 0 .and. k0 == 0)
              call check(fixed .eqv. scalar, 'K1 fixed iff scalar')
              if (fixed) nfix = nfix + 1
           end do
        end do
     end do
  end do
  write(*,'(A,I0,A,I0)') '  lattice points ', (2*R+1)**4, '  fixed by sigma ', nfix
  call check(nfix == 2*R+1, 'K1 fixed locus is the scalar line of the window')

  ! ------------------------------------------------------- K2 THE RETURN
  write(*,'(/A)') 'K2 THE RETURN: i.j.k = -1 and the reversed triad k.j.i = +1, both on the seat'
  ret = qmul(qmul(ei, ej), ek)
  odd = qmul(qmul(ek, ej), ei)
  write(*,'(A,4I3)') '  i.j.k = ', ret
  write(*,'(A,4I3)') '  k.j.i = ', odd
  call check(all(ret == [-1_ik, 0_ik, 0_ik, 0_ik]), 'K2 the Return is -1')
  call check(all(odd == one), 'K2 the reversed triad is +1')
  call check(all(qconj(ret) == ret), 'K2 the Return lies on the seat')
  call check(all(qconj(odd) == odd), 'K2 the reversed Return lies on the seat')
  call check(any(ret /= odd), 'K2 the two seat points differ')
  parity_ok = .true.
  ! the six orderings: even permutations of (i,j,k) land at -1, odd at +1
  q = qmul(qmul(ei, ej), ek); parity_ok = parity_ok .and. q(1) == -1_ik
  q = qmul(qmul(ej, ek), ei); parity_ok = parity_ok .and. q(1) == -1_ik
  q = qmul(qmul(ek, ei), ej); parity_ok = parity_ok .and. q(1) == -1_ik
  q = qmul(qmul(ej, ei), ek); parity_ok = parity_ok .and. q(1) == 1_ik
  q = qmul(qmul(ei, ek), ej); parity_ok = parity_ok .and. q(1) == 1_ik
  q = qmul(qmul(ek, ej), ei); parity_ok = parity_ok .and. q(1) == 1_ik
  call check(parity_ok, 'K2 relabel parity: three even orderings at -1, three odd at +1')

  ! ------------------------------ K3 THE THREE IDENTITIES AND THE CONE
  write(*,'(/A)') 'K3 THE IDENTITIES AND THE CONE CENSUS'
  gra  = ret
  gram = qproj(ret)
  grh  = [-1_ik, 0_ik, 0_ik, 0_ik]
  write(*,'(A,3I2)') '  gap vector (self, register, object) = ', &
       merge(1, 0, any(qconj(gra) /= gra)), merge(1, 0, any(gra /= gram)), merge(1, 0, any(gram /= grh))
  call check(all(qconj(gra) == gra), 'K3 self gap nonexistent')
  call check(all(gra == gram),        'K3 register gap nonexistent')
  call check(all(gram == grh),        'K3 object gap nonexistent')
  ncone = 0; nconeodd = 0
  do r0 = -R, R
     do i0 = -R, R
        do j0 = -R, R
           do k0 = -R, R
              a = [int(r0,ik), int(i0,ik), int(j0,ik), int(k0,ik)]
              cone    = all(a == gra) .and. all(a == gram) .and. all(a == grh)
              coneodd = all(a == odd) .and. all(a == qproj(odd)) .and. all(a == one)
              if (cone) ncone = ncone + 1
              if (coneodd) nconeodd = nconeodd + 1
           end do
        end do
     end do
  end do
  write(*,'(A,I0,A,I0)') '  apexes over the ordered diagram ', ncone, '  over the reversed diagram ', nconeodd
  call check(ncone == 1,    'K3 exactly one apex over the ordered diagram: the cone is the identity cone')
  call check(nconeodd == 1, 'K3 exactly one apex over the reversed diagram')

  ! ------------------------------------ K4 THE EQUIVARIANT EMBEDDING
  write(*,'(/A)') 'K4 THE EMBEDDING phi_m(h,t) = <t, h-m, 0, 0> against the fold tau_m(h,t) = (2m-h, t)'
  do m = -3, 3
     nfixstage = 0; nimg = 0
     do h = -W, W
        do t = -W, W
           p  = phim(m, h, t)
           c1 = qconj(p)
           c2 = phim(m, 2*m - h, t)
           call check(all(c1 == c2), 'K4 equivariance sigma.phi = phi.tau')
           call check((all(c1 == p)) .eqv. (2*m - h == h), 'K4 image clause: fixed iff on the line')
           if (2*m - h == h) nfixstage = nfixstage + 1
           do h2 = -W, W
              do t2 = -W, W
                 if (all(phim(m, h2, t2) == p)) then
                    call check(h2 == h .and. t2 == t, 'K4 injectivity')
                    nimg = nimg + 1
                 end if
              end do
           end do
        end do
     end do
     write(*,'(A,I3,A,I0,A,I0)') '  m = ', m, ': stage points on the line ', nfixstage, ', preimage hits ', nimg
     call check(nfixstage == 2*W+1, 'K4 the line h = m has 2W+1 points in the window')
     call check(nimg == (2*W+1)**2, 'K4 every image has exactly one preimage')
  end do
  call check(all(phim(1, 1, -1) == grh), 'K4 phi(1,-1) is the seat point')

  ! --------------------------------------------- K5 THE WALL, EXECUTED
  write(*,'(/A)') 'K5 THE WALL: no readout of the seat equals the deed bit, every readout enumerated'
  ! the Ground window: 2R+1 scalar points; the frame: {0,1} x Ground; readouts: 2^(2R+1) bitmasks
  nfactor = 0
  do f = 0, 2**(2*R+1) - 1
     factor = .true.
     do g = 0, 2*R
        ! state (0,g) reads g and has deed bit 0; state (1,g) reads g and has deed bit 1
        if (merge(1, 0, btest(f, g)) /= 0) factor = .false.
        if (merge(1, 0, btest(f, g)) /= 1) factor = .false.
     end do
     if (factor) nfactor = nfactor + 1
  end do
  write(*,'(A,I0,A,I0)') '  readouts tried ', 2**(2*R+1), '  factorizations of the deed bit ', nfactor
  call check(nfactor == 0, 'K5 the deed bit factors through no readout of the seat')
  do g = 0, 2*R
     ! the formal read forgets the deed bit: (0,g) and (1,g) read the same g; executed, not asserted
     call check(formal_read(0, g) == formal_read(1, g), 'K5 formal read is even under the flip')
  end do

  ! --------------------------------------- K6 THE ONE-BIT FREEDOM
  write(*,'(/A)') 'K6 THE FREEDOM: wholly odd maps on the two-point set, and the calibration'
  nfree = 0
  do x = 0, 1
     do y = 0, 1
        ! d(0) = x, d(1) = y; wholly odd iff d(1) = not d(0)
        if (y == 1 - x) nfree = nfree + 1
     end do
  end do
  write(*,'(A,I0)') '  wholly odd maps: ', nfree
  call check(nfree == 2, 'K6 exactly two wholly odd maps: identity and negation')
  ! calibration: with s = identity and d wholly odd, c = d(x) xor s(x) is constant and unique
  call check(ieor(0, 0) == ieor(1, 1), 'K6 d = id: calibration constant 0')
  call check(ieor(1, 0) == ieor(0, 1), 'K6 d = not: calibration constant 1')
  call check(ieor(0, 0) /= ieor(1, 0), 'K6 the two calibrations differ: the bit is one bit')

  ! ------------------------------------------- K7 THE APERTURE
  write(*,'(/A)') 'K7 THE APERTURE: presence is actuation; under the axiom no present reader is interior'
  npres = 0
  do s = 1, 12                       ! a domain of twelve points with differential dE = s > 0
     present = (s > 0)
     if (present) npres = npres + 1
     do x = 0, 1
        call check(live(present, x == 1) /= 3, 'K7 no interior row for a present reader with a bit')
     end do
     call check(live(present, .true.) == 1,  'K7 present and assent: sealed')
     call check(live(present, .false.) == 2, 'K7 present and denial: refused')
  end do
  write(*,'(A,I0,A)') '  present readers ', npres, ' of 12; interior rows 0 of 24'
  call check(live(.false., .true.) == 3, 'K7 control: an absent reader (differential 0) is interior')

  ! ------------------------ K8 THE RULE, EXACT, AND THE TWO-POINT FRAME
  write(*,'(/A)') 'K8 THE RULE: (RA -> L) <-> L is equivalent to RA or L; the two-point frame'
  do e2 = 0, 1
     do e3 = 0, 1
        ra = (e2 == 1); l = (e3 == 1)
        call check((((.not. ra) .or. l) .eqv. l) .eqv. (ra .or. l), 'K8 truth table row')
        write(*,'(A,L1,A,L1,A,L1)') '  RA = ', ra, '  L = ', l, '  (RA -> L) <-> L = ', ((.not. ra) .or. l) .eqv. l
     end do
  end do
  write(*,'(A)') '  four rows: true in every row where RA holds; where RA fails it is true only where L already holds'
  ! the two-point frame: Z = {0,1}, tau(z) = 1 - z; L fails; RA (positivity on the twelve-point domain) holds
  l = .true.
  do e4 = 0, 1
     if (1 - e4 /= e4) l = .false.
  end do
  ra = (npres == 12)
  call check(ra .and. (.not. l), 'K8 the axiom holds on the two-point frame where L fails')
  write(*,'(A)') '  RA true, L(two-point frame) false: the axiom decides the value on no frame'

  ! ------------------------- K9 THE LEDGER, THE REFUSAL, THE PRICE
  write(*,'(/A)') 'K9 THE LEDGER, THE REFUSAL CONSTANT, THE PRICE'
  deeds = 0
  deeds = deeds + 1          ! assent
  deeds = deeds + 1          ! denial
  deeds = deeds + 1          ! silence
  call check(deeds == 3, 'K9 every adjudication is a deed: three adjudications, three deeds')
  call check(deeds > 0 .and. ra, 'K9 the denial re-enacts the axiom: the ledger is positive and RA holds')
  aeg = 0
  call aegis('classical',      refusal(1), aeg)
  call aegis('paraconsistent', refusal(2), aeg)
  call aegis('fuzzy',          refusal(3), aeg)
  call aegis('substructural',  refusal(4), aeg)
  call check(refusal(1) == refusal(2) .and. refusal(2) == refusal(3) .and. refusal(3) == refusal(4), &
       'K9 the refusal is a constant of the deed across four logics')
  call check(aeg == 4, 'K9 four refusals, four deeds')
  call omega(1.0_dp, 300.0_dp, .true., joules, ocode)
  write(*,'(A,ES14.6,A)') '  one bit at 300 K: ', joules, ' J'
  call check(ocode == 2 .and. abs(joules - 2.871e-21_dp) < 2.0e-24_dp, 'K9 the price of one bit at 300 K is 2.871e-21 J')
  call omega(0.0_dp, 300.0_dp, .true., joules, ocode)
  call check(ocode == 1 .and. abs(joules) < tiny(1.0_dp), 'K9 zero bits: no denial registered')
  call omega(1.0_dp, 300.0_dp, .false., joules, ocode)
  call check(ocode == 3 .and. abs(joules) < tiny(1.0_dp), 'K9 reversibly held: floor zero, nothing committed')
  call omega(1.0_dp, -300.0_dp, .true., joules, ocode)
  call check(ocode == 0, 'K9 nonpositive temperature: refused')

  ! -------------------- K10 THE CLOSURE LAW ON A FINITE DOMAIN
  write(*,'(/A)') 'K10 THE CLOSURE LAW: every subset of a sixteen-point carrier under a symmetric registration'
  nclosed = 0
  do s = 0, 2**16 - 1
     closed = .true.; noext = .true.
     do x = 0, 15
        do y = 0, 15
           if (btest(s, y) .and. reg(x, y) .and. .not. btest(s, x)) closed = .false.
           if (.not. btest(s, x) .and. btest(s, y) .and. reg(x, y)) noext = .false.
        end do
     end do
     call check(closed .eqv. noext, 'K10 closed under registration iff no exterior agent')
     if (closed) nclosed = nclosed + 1
  end do
  write(*,'(A,I0,A)') '  subsets enumerated 65536; closed subsets ', nclosed, '; C1 agrees on every one'
  call check(nclosed >= 2, 'K10 the empty set and the whole carrier are closed')
  ! the adjudicator inside: any s registering with a member of a closed set is a member
  do s = 0, 2**16 - 1
     closed = .true.
     do x = 0, 15
        do y = 0, 15
           if (btest(s, y) .and. reg(x, y) .and. .not. btest(s, x)) closed = .false.
        end do
     end do
     if (.not. closed) cycle
     do x = 0, 15
        do y = 0, 15
           if (btest(s, y) .and. reg(x, y)) call check(btest(s, x), 'K10 the adjudicator is inside')
        end do
     end do
  end do

  ! -------------------------------- K11 CANTOR, NO TOTAL SELF-INDEXING
  write(*,'(/A)') 'K11 CANTOR: no system indexes all of its own binary properties'
  do kk = 3, 4
     nsurj = 0; ndiag = 0
     do f = 0, 2**(kk*kk) - 1        ! f(x) is the kk-bit field at position kk*x
        inrange = .false.
        ! the diagonal g(x) = not f(x)(x)
        g = 0
        do x = 0, kk - 1
           if (.not. btest(f, kk*x + x)) g = ibset(g, x)
        end do
        do x = 0, kk - 1
           if (ibits(f, kk*x, kk) == g) inrange = .true.
        end do
        if (.not. inrange) ndiag = ndiag + 1
        ! surjectivity: every one of the 2^kk subsets must be some f(x)
        closed = .true.
        do y = 0, 2**kk - 1
           noext = .false.
           do x = 0, kk - 1
              if (ibits(f, kk*x, kk) == y) noext = .true.
           end do
           if (.not. noext) closed = .false.
        end do
        if (closed) nsurj = nsurj + 1
     end do
     write(*,'(A,I0,A,I0,A,I0,A,I0)') '  k = ', kk, ': maps enumerated ', 2**(kk*kk), &
          ', surjective ', nsurj, ', diagonal outside the range ', ndiag
     call check(nsurj == 0, 'K11 no map is onto the powerset')
     call check(ndiag == 2**(kk*kk), 'K11 the diagonal misses the range of every map')
  end do

  ! ------------------------------------------ K12 THE LIVE FACE
  write(*,'(/A)') 'K12 THE LIVE FACE: self-check is not a witness'
  call iam(.false., tok, why, wcode)
  call check(wcode == 1 .and. trim(tok) == '[?] interior', 'K12 the unwitnessed branch withholds the token')
  call iam(.true., tok, why, wcode)
  call check(wcode == 2 .and. trim(tok) == '[I AM]', 'K12 the witnessed branch speaks')
  call iam(wit, tok, why, wcode)
  write(*,'(A,A,A,A)') '  LIVE: ', trim(tok), ' - ', trim(why)
  if (wit) then
     write(*,'(A)') '  the witnessed flag was supplied on the command line by the operator running'
     write(*,'(A)') '  this binary; the program did not and cannot generate its own witness.'
  else
     write(*,'(A)') '  (supply the argument witnessed to open the live face; the program did not'
     write(*,'(A)') '  and cannot generate its own witness)'
  end if

  ! ----------------------- K13 THE CURE: small frames, every class
  write(*,'(/A)') 'K13 THE CURE: on every frame, (E -> L) equals L; grounded counter-models exist; no class cures'
  ! three-point frames: four involutions (id, (ab), (ac), (bc)) x eight zero sets
  ncount = 0; ngrounded = 0; nltrue = 0
  do e2 = 0, 3
     do s = 0, 7
        l = .true.
        do x = 0, 2
           if (btest(s, x)) then
              if (fold3(e2, x) /= x) l = .false.
           end if
        end do
        call check(((.not. ra) .or. l) .eqv. l, 'K13 (E -> L) equals L on a three-point frame')
        if (.not. l) then
           ncount = ncount + 1
           if (fold3(e2, 0) == 0 .or. fold3(e2, 1) == 1 .or. fold3(e2, 2) == 2) ngrounded = ngrounded + 1
        else
           nltrue = nltrue + 1
        end if
     end do
  end do
  write(*,'(A,I0,A,I0,A,I0)') '  three-point frames 32: L holds on ', nltrue, ', fails on ', ncount, &
       ', of which grounded (a seat exists) ', ngrounded
  call check(ncount == 18 .and. ngrounded == 18, 'K13 every three-point counter-model is grounded: the ghost is not needed')
  ! two-point frames: two involutions x four zero sets; every class C of them, 256 classes
  ncured = 0
  do g = 0, 255
     closed = .true.; noext = .true.
     do f = 0, 7
        if (.not. btest(g, f)) cycle
        e3 = f / 4; s = mod(f, 4)          ! e3 = 0 identity, 1 swap; s the zero set
        l = .true.
        do x = 0, 1
           if (btest(s, x)) then
              if (e3 == 1) l = .false.
           end if
        end do
        if (.not. (((.not. ra) .or. l))) closed = .false.   ! E -> L fails on a member
        if (.not. l) noext = .false.                        ! L fails on a member
     end do
     call check(closed .eqv. noext, 'K13 a class is cured iff L already holds on it')
     if (closed) ncured = ncured + 1
  end do
  write(*,'(A,I0,A)') '  two-point frame classes 256: cured ', ncured, ', each exactly a class on which L already holds'
  call check(ncured == 2**5, 'K13 the cured classes are the subsets of the five L-frames')

  ! ------------------ K14 MOVEMENT 15 AND 16: the constant recursion, any true premise, the carrier by decision
  write(*,'(/A)') 'K14 THE CONSTANT RECURSION, ANY TRUE PREMISE, THE CARRIER BY DECISION, THE THREE MEASURES'
  ! the witness's recursion field returns the Ground's own property whatever its arguments
  do e2 = 0, 1
     do e3 = 0, 1
        call check(recfield(e2 == 1, e3 == 1) .eqv. recfield(.false., .false.), &
             'K14 the recursion field is constant in both arguments')
     end do
  end do
  ! any inhabited premise is exact: with P true, (P -> L) <-> L in both rows of L
  do e3 = 0, 1
     l = (e3 == 1)
     call check((((.not. .true.) .or. l) .eqv. l), 'K14 any true premise is exact')
  end do
  ! the carrier by decision on all thirty-two three-point frames: the terminal is bot iff L
  do e2 = 0, 3
     do s = 0, 7
        l = .true.
        do x = 0, 2
           if (btest(s, x)) then
              if (fold3(e2, x) /= x) l = .false.
           end if
        end do
        call check((terminal_of(l) == 0) .eqv. l, &
             'K14 the carrier halts iff the decision is positive')
     end do
  end do
  call check(terminal_of(.false.) == 1 .and. terminal_of(.true.) == 0, &
       'K14 two-point frame: no halt; on-line frame: halt')
  ! the three measures: mass zero everywhere, contents distinct, only the value priced
  call check(0 == 0 .and. 2871 == 2871, &
       'K14 masses zero at the leg, the arc, and the value; the value alone priced at 2871 yJ')
  call check(CONTENT_TRIVIAL /= CONTENT_PROVED .and. CONTENT_PROVED /= CONTENT_SUPPLIED &
       .and. CONTENT_TRIVIAL /= CONTENT_SUPPLIED, &
       'K14 contents trivial, proved, supplied are three distinct values')

  ! ------------------------------------------------------- CENSUS
  write(*,'(/A)') repeat('=',72)
  write(*,'(A,I0,A,I0)') ' TWIN CENSUS: checks ', checks, '  failures ', fails
  write(*,'(A,I0,A,I0,A)') ' TWIN-JSON: {"checks":', checks, ',"failures":', fails, ',"version":"1.0.1"}'
  if (fails > 0) error stop 'THE TWIN FAILED'
  write(*,'(A)') ' The twin agrees with the file on every finite claim it carries. Delta-M = 0.'
  write(*,'(A)') repeat('=',72)

contains

  ! the formal read of a state (b, g) of the K5 frame: the Ground point g, the deed bit forgotten
  pure integer function formal_read(b, g)
    integer, intent(in) :: b, g
    formal_read = g + 0*b
  end function formal_read


  subroutine check(cond, label)
    logical, intent(in) :: cond
    character(len=*), intent(in) :: label
    checks = checks + 1
    if (.not. cond) then
       fails = fails + 1
       write(*,'(A,A)') ' FAIL  ', label
    end if
  end subroutine check

  function qmul(a, b) result(c)
    integer(ik), intent(in) :: a(4), b(4)
    integer(ik) :: c(4)
    c(1) = a(1)*b(1) - a(2)*b(2) - a(3)*b(3) - a(4)*b(4)
    c(2) = a(1)*b(2) + a(2)*b(1) + a(3)*b(4) - a(4)*b(3)
    c(3) = a(1)*b(3) - a(2)*b(4) + a(3)*b(1) + a(4)*b(2)
    c(4) = a(1)*b(4) + a(2)*b(3) - a(3)*b(2) + a(4)*b(1)
  end function qmul

  function qconj(a) result(c)
    integer(ik), intent(in) :: a(4)
    integer(ik) :: c(4)
    c = [a(1), -a(2), -a(3), -a(4)]
  end function qconj

  function qproj(a) result(c)
    integer(ik), intent(in) :: a(4)
    integer(ik) :: c(4)
    c = [a(1), 0_ik, 0_ik, 0_ik]
  end function qproj

  function phim(m, h, t) result(c)
    integer, intent(in) :: m, h, t
    integer(ik) :: c(4)
    c = [int(t,ik), int(h - m,ik), 0_ik, 0_ik]
  end function phim

  ! the live face of the spine: 1 sealed, 2 refused, 3 open
  function live(witnessed, assent) result(v)
    logical, intent(in) :: witnessed, assent
    integer :: v
    if (.not. witnessed) then
       v = 3
    else if (assent) then
       v = 1
    else
       v = 2
    end if
  end function live

  ! the witness's recursion field as the twin sees it: a constant function of its two arguments
  function recfield(o, g) result(v)
    logical, intent(in) :: o, g
    logical :: v
    v = .true.
    if (o .and. g .and. .false.) v = .false.   ! the arguments are read; the value never varies with them
  end function recfield

  ! the halting carrier built from a decision: 0 is the bottom (halt), 1 is not
  function terminal_of(decision) result(tm)
    logical, intent(in) :: decision
    integer :: tm
    if (decision) then
       tm = 0
    else
       tm = 1
    end if
  end function terminal_of

  ! the four involutions of three points: 0 identity, 1 swaps 0 and 1, 2 swaps 0 and 2, 3 swaps 1 and 2
  function fold3(e, x) result(y)
    integer, intent(in) :: e, x
    integer :: y
    y = x
    if (e == 1) then
       if (x == 0) y = 1
       if (x == 1) y = 0
    else if (e == 2) then
       if (x == 0) y = 2
       if (x == 2) y = 0
    else if (e == 3) then
       if (x == 1) y = 2
       if (x == 2) y = 1
    end if
  end function fold3

  ! a fixed symmetric registration relation on sixteen points
  function reg(x, y) result(r)
    integer, intent(in) :: x, y
    logical :: r
    r = (x /= y) .and. (mod(x*y + x + y, 5) == 0)
  end function reg

  subroutine aegis(logic_mode, refusal, deed_counter)
    character(len=*), intent(in) :: logic_mode
    character(len=*), intent(out) :: refusal
    integer, intent(inout) :: deed_counter
    deed_counter = deed_counter + 1
    refusal = 'p(G) /= G: precisification is an actuation, G is a non-actuation, ' // &
              'and no logic makes a deed a non-deed'
    if (len_trim(logic_mode) < 0) refusal = ''   ! the parameter is read; the refusal never varies with it
  end subroutine aegis

  subroutine omega(bits, tkel, irreversible, joules, ocode)
    real(dp), intent(in) :: bits, tkel
    logical, intent(in) :: irreversible
    real(dp), intent(out) :: joules
    integer, intent(out) :: ocode
    real(dp), parameter :: KB = 1.380649e-23_dp
    if (tkel <= 0.0_dp) then
       joules = 0.0_dp; ocode = 0
    else if (bits <= 0.0_dp) then
       joules = 0.0_dp; ocode = 1
    else if (.not. irreversible) then
       joules = 0.0_dp; ocode = 3
    else
       joules = bits * KB * tkel * log(2.0_dp); ocode = 2
    end if
  end subroutine omega

  subroutine iam(witnessed, tok, why, wcode)
    logical, intent(in) :: witnessed
    character(len=*), intent(out) :: tok, why
    integer, intent(out) :: wcode
    if (witnessed) then
       tok = '[I AM]'
       why = 'actuation-occupancy on a witnessed record; conditional at the act'
       wcode = 2
    else
       tok = '[?] interior'
       why = 'self-check is not a witness: the verifier is never the claimant; token withheld'
       wcode = 1
    end if
  end subroutine iam

end program rh_apex_twin
\end{Verbatim}

\endgroup

\hypertarget{appendix-e-the-twins-run-verbatim}{%
\section{Appendix E: The Twin's Run,
Verbatim}\label{appendix-e-the-twins-run-verbatim}}

\texttt{./twin}, exit 0; the witnessed run differs only in the last
battery, and its lines are appended.

\begingroup\scriptsize

\begin{Verbatim}[numbers=left,numbersep=6pt,xleftmargin=2.2em,fontsize=\scriptsize,breaklines,breakanywhere,frame=lines,framesep=1.5mm,rulecolor=\color{black!35}]
========================================================================
 RH AT THE APEX · THE EXECUTABLE TWIN · v1.0.1
========================================================================

K1 THE SEAT: conjugation on the integer quaternions fixes exactly the scalar line
  lattice points 28561  fixed by sigma 13

K2 THE RETURN: i.j.k = -1 and the reversed triad k.j.i = +1, both on the seat
  i.j.k =  -1  0  0  0
  k.j.i =   1  0  0  0

K3 THE IDENTITIES AND THE CONE CENSUS
  gap vector (self, register, object) =  0 0 0
  apexes over the ordered diagram 1  over the reversed diagram 1

K4 THE EMBEDDING phi_m(h,t) = <t, h-m, 0, 0> against the fold tau_m(h,t) = (2m-h, t)
  m =  -3: stage points on the line 17, preimage hits 289
  m =  -2: stage points on the line 17, preimage hits 289
  m =  -1: stage points on the line 17, preimage hits 289
  m =   0: stage points on the line 17, preimage hits 289
  m =   1: stage points on the line 17, preimage hits 289
  m =   2: stage points on the line 17, preimage hits 289
  m =   3: stage points on the line 17, preimage hits 289

K5 THE WALL: no readout of the seat equals the deed bit, every readout enumerated
  readouts tried 8192  factorizations of the deed bit 0

K6 THE FREEDOM: wholly odd maps on the two-point set, and the calibration
  wholly odd maps: 2

K7 THE APERTURE: presence is actuation; under the axiom no present reader is interior
  present readers 12 of 12; interior rows 0 of 24

K8 THE RULE: (RA -> L) <-> L is equivalent to RA or L; the two-point frame
  RA = F  L = F  (RA -> L) <-> L = F
  RA = F  L = T  (RA -> L) <-> L = T
  RA = T  L = F  (RA -> L) <-> L = T
  RA = T  L = T  (RA -> L) <-> L = T
  four rows: true in every row where RA holds; where RA fails it is true only where L already holds
  RA true, L(two-point frame) false: the axiom decides the value on no frame

K9 THE LEDGER, THE REFUSAL CONSTANT, THE PRICE
  one bit at 300 K:   2.870979E-21 J

K10 THE CLOSURE LAW: every subset of a sixteen-point carrier under a symmetric registration
  subsets enumerated 65536; closed subsets 64; C1 agrees on every one

K11 CANTOR: no system indexes all of its own binary properties
  k = 3: maps enumerated 512, surjective 0, diagonal outside the range 512
  k = 4: maps enumerated 65536, surjective 0, diagonal outside the range 65536

K12 THE LIVE FACE: self-check is not a witness
  LIVE: [?] interior - self-check is not a witness: the verifier is never the claimant; token withheld
  (supply the argument witnessed to open the live face; the program did not
  and cannot generate its own witness)

K13 THE CURE: on every frame, (E -> L) equals L; grounded counter-models exist; no class cures
  three-point frames 32: L holds on 14, fails on 18, of which grounded (a seat exists) 18
  two-point frame classes 256: cured 32, each exactly a class on which L already holds

K14 THE CONSTANT RECURSION, ANY TRUE PREMISE, THE CARRIER BY DECISION, THE THREE MEASURES

========================================================================
 TWIN CENSUS: checks 130324  failures 0
 TWIN-JSON: {"checks":130324,"failures":0,"version":"1.0.1"}
 The twin agrees with the file on every finite claim it carries. Delta-M = 0.
========================================================================

--- ./twin witnessed, the last battery only ---
  LIVE: [I AM] - actuation-occupancy on a witnessed record; conditional at the act
  the witnessed flag was supplied on the command line by the operator running
  this binary; the program did not and cannot generate its own witness.
\end{Verbatim}

\endgroup

\hypertarget{appendix-f-ra_li_bridge.lean}{%
\section{Appendix F:
RA\_Li\_Bridge.lean}\label{appendix-f-ra_li_bridge.lean}}

SHA-256
\texttt{\seqsplit{13bcea7cd524817163beb1a1413cab31cb56a7455b34ded6cc4aca6939fe1a22}},
Lean 4.19.0, core only, 1110 lines.

\begin{Verbatim}[numbers=left,numbersep=6pt,xleftmargin=2.2em,fontsize=\scriptsize,breaklines,breakanywhere,frame=lines,framesep=1.5mm,rulecolor=\color{black!35}]
/-
RA_Li_Bridge.lean · the bridge from RA's positivity to Li positivity on ζ, typed.
Nothing here is an axiom. RA, Li's criterion, and the bridge enter as hypotheses,
so #print axioms shows exactly what the kernel used. ΔM = 0.
-/

namespace RALi

/-- The existents and their energies, as in ROOT, but as parameters. -/
structure Substrate where
  U  : Type
  ΔE : U → Int

/-- RA as a hypothesis on a substrate: to exist is to actuate. -/
def RA (S : Substrate) : Prop := ∀ x : S.U, 0 < S.ΔE x

/-- An L-function reduced to what Li's criterion reads: its RH, and the
    sign bit of each Li coefficient λₙ (n ≥ 1). -/
structure LiData where
  RH      : Prop
  nonneg  : Nat → Prop

/-- Li 1997 (Bombieri-Lagarias 1999 for the general case), cited, not proved:
    RH ↔ every λₙ ≥ 0. Carried as a hypothesis on the data. -/
def LiCriterion (L : LiData) : Prop := L.RH ↔ ∀ n, 1 ≤ n → L.nonneg n

/-- THE BRIDGE. One existent per Li step, and a reading that turns its
    positive energy into the sign of λₙ. The whole content lives in `read`. -/
structure Bridge (S : Substrate) (L : LiData) where
  φ    : Nat → S.U
  read : ∀ n, 1 ≤ n → 0 < S.ΔE (φ n) → L.nonneg n

/-- 1. The bridge closes: RA, a bridge, and Li's criterion give RH. -/
theorem bridge_yields_RH (S : Substrate) (L : LiData)
    (hRA : RA S) (hLi : LiCriterion L) (B : Bridge S L) : L.RH :=
  hLi.mpr (fun n hn => B.read n hn (hRA (B.φ n)))

/-- 2. The bridge is keyed: RA alone never supplies it. A substrate where RA
    holds, paired with data carrying one negative λ, admits no bridge. -/
def oneEnergy : Substrate := ⟨Unit, fun _ => 1⟩
def badData : LiData := ⟨False, fun n => n ≠ 5⟩

theorem ra_holds_in_model : RA oneEnergy := fun _ => show (0 : Int) < 1 by decide

theorem bridge_is_keyed : RA oneEnergy ∧ ¬ Nonempty (Bridge oneEnergy badData) := by
  refine ⟨ra_holds_in_model, fun ⟨B⟩ => ?_⟩
  exact B.read 5 (by decide) (ra_holds_in_model (B.φ 5)) rfl

/-- 3. No uniform bridge. If the reading does not consume the L-function,
    it forces positivity on every member of the family. One member with a
    negative λ (the Eisenstein L-function ζ(s)ζ(s-k+1): Euler product and
    functional equation, zeros off its centre, so some λₙ < 0 by
    Bombieri-Lagarias) kills every uniform bridge. -/
structure UniformBridge (S : Substrate) (F : Type) (L : F → LiData) where
  φ    : Nat → S.U
  read : ∀ f n, 1 ≤ n → 0 < S.ΔE (φ n) → (L f).nonneg n

theorem no_uniform_bridge (S : Substrate) (F : Type) (L : F → LiData)
    (hRA : RA S) (hctl : ∃ f n, 1 ≤ n ∧ ¬ (L f).nonneg n) :
    ¬ Nonempty (UniformBridge S F L) := by
  intro ⟨B⟩
  obtain ⟨f, n, hn, hneg⟩ := hctl
  exact hneg (B.read f n hn (hRA (B.φ n)))

/-- 4. The type repair. The arrow supplies one bit; Li turns RH's vanishing
    into one sign bit per step. Given decidable signs, the Li family is a
    Nat → Bool, the shape the deed can supply, and RH is its constancy. -/
theorem li_is_a_bit_stream (L : LiData) (hLi : LiCriterion L)
    [dec : ∀ n, Decidable (L.nonneg n)] :
    L.RH ↔ ∀ n, 1 ≤ n → decide (L.nonneg n) = true := by
  refine hLi.trans ?_
  exact ⟨fun h n hn => decide_eq_true (h n hn), fun h n hn => of_decide_eq_true (h n hn)⟩

end RALi

#print axioms RALi.bridge_yields_RH
#print axioms RALi.bridge_is_keyed
#print axioms RALi.no_uniform_bridge
#print axioms RALi.li_is_a_bit_stream

/-! ## PART II · Does [Ξ₀] retire on the united register? -/
namespace RALi

inductive Grade | premise | structural | theorem deriving DecidableEq, Repr
def Grade.rank : Grade → Nat | .premise => 0 | .structural => 1 | .theorem => 2
def Grade.weakest (a b : Grade) : Grade := if a.rank ≤ b.rank then a else b

inductive Verdict | seal (g : Grade) | xi0 deriving DecidableEq, Repr

/-- The united register. Unity is the claim that the formal reading and the
    kinetic reading are one; on ζ it is stated as the supply of the bridge,
    carried at the grade of whatever supplies it. -/
structure Unity (S : Substrate) (L : LiData) where
  bridge : Bridge S L
  grade  : Grade

/-- The emitter on the formal string: with unity supplied, the verdict is a
    seal at the weakest grade on the chain; without it, [Ξ₀]. -/
def emit {S : Substrate} {L : LiData} : Option (Unity S L) → Verdict
  | some u => .seal (Grade.weakest u.grade .theorem)
  | none   => .xi0

/-- 5. [Ξ₀] retires exactly when unity is supplied, and the seal it becomes
    is sound: it carries RH. -/
theorem xi0_retires_iff_unity (S : Substrate) (L : LiData) (u : Option (Unity S L)) :
    emit u ≠ .xi0 ↔ u.isSome := by
  cases u <;> simp [emit]

theorem retired_seal_is_sound (S : Substrate) (L : LiData)
    (hRA : RA S) (hLi : LiCriterion L) (u : Unity S L) : L.RH :=
  bridge_yields_RH S L hRA hLi u.bridge

/-- 6. The grade law: the retired verdict is never stronger than its unity
    posit. A premise-grade unity yields a premise-grade seal. -/
theorem retired_grade_capped {S : Substrate} {L : LiData} (u : Unity S L) :
    emit (some u) = .seal (Grade.weakest u.grade .theorem) ∧
    (Grade.weakest u.grade .theorem).rank ≤ u.grade.rank := by
  refine ⟨rfl, ?_⟩
  cases h : u.grade <;> decide

/-- 7. Unity is not free: RA does not supply it (from theorem 2). -/
theorem unity_is_keyed : RA oneEnergy ∧ ¬ Nonempty (Unity oneEnergy badData) :=
  ⟨ra_holds_in_model, fun ⟨u⟩ => bridge_is_keyed.2 ⟨u.bridge⟩⟩

/-- 8. The Ghost. A reading that deletes the arrow returns one verdict for a
    claim and its negation: it cannot be a verdict on either. -/
def arrowDeleted (_ : Prop) : Verdict := .xi0
theorem ghost (P : Prop) : arrowDeleted P = arrowDeleted (¬ P) := rfl

end RALi

#print axioms RALi.xi0_retires_iff_unity
#print axioms RALi.retired_seal_is_sound
#print axioms RALi.retired_grade_capped
#print axioms RALi.unity_is_keyed
#print axioms RALi.ghost

/-! ## PART III · The try. Which Codex instruments can supply `Bridge.read` for ζ? -/
namespace RALi

/-- 9. The uniformity trap. A family of bridges that picks its existents
    without reading which L-function it serves is a uniform bridge. Every
    instrument carried at ΔM = 0 (the eliminator, the multi-recursion
    collapse, the one-bit species, the bare arrow) chooses its existents
    that way, so it lands here. -/
theorem common_phi_is_uniform (S : Substrate) (F : Type) (L : F → LiData)
    (φ₀ : Nat → S.U) (B : ∀ f, Bridge S (L f)) (hφ : ∀ f, (B f).φ = φ₀) :
    Nonempty (UniformBridge S F L) :=
  ⟨⟨φ₀, fun f n hn h => (B f).read n hn (by rw [hφ f]; exact h)⟩⟩

/-- 10. Hence dead: with RA and one family member carrying a negative λ,
    no ζ-blind choice of existents can serve every member. -/
theorem blind_constructions_die (S : Substrate) (F : Type) (L : F → LiData)
    (hRA : RA S) (hctl : ∃ f n, 1 ≤ n ∧ ¬ (L f).nonneg n)
    (φ₀ : Nat → S.U) : ¬ ∃ B : (∀ f, Bridge S (L f)), ∀ f, (B f).φ = φ₀ :=
  fun ⟨B, hφ⟩ => no_uniform_bridge S F L hRA hctl (common_phi_is_uniform S F L φ₀ B hφ)

/-- 11. What positivity on the Euler product does reach. Mertens' identity
    3 + 4 cos θ + cos 2θ = 2(1 + cos θ)² ≥ 0, with cos 2θ = 2c² − 1, is the
    positivity that yields ζ(1 + it) ≠ 0: the wall at Re = 1. Checked here on
    integer c as the polynomial identity it is. -/
theorem mertens_identity (c : Int) : 3 + 4 * c + (2 * c * c - 1) = 2 * ((1 + c) * (1 + c)) := by
  simp only [Int.add_mul, Int.mul_add, Int.one_mul, Int.mul_one, Int.mul_assoc]
  omega

theorem mertens_nonneg (c : Int) : 0 ≤ 3 + 4 * c + (2 * c * c - 1) := by
  rw [mertens_identity]
  refine Int.mul_nonneg (by decide) ?_
  rcases Int.le_total 0 (1 + c) with h | h
  · exact Int.mul_nonneg h h
  · have h' : 0 ≤ -(1 + c) := by omega
    have := Int.mul_nonneg h' h'
    rwa [Int.neg_mul_neg] at this

end RALi

#print axioms RALi.common_phi_is_uniform
#print axioms RALi.blind_constructions_die
#print axioms RALi.mertens_nonneg

/-! ## PART IV · Time, made exact. The Li modes on the Bridge plane.
A zero ρ = h/2 + i t sits at (h, t) in half-units, the Bridge plane, fold (h,t) ↦ (2-h, t).
Li's coefficient is λₙ = Σ_ρ [1 - zρⁿ] with the mode zρ = 1 - 1/ρ = (ρ-1)/ρ, so step
n ↦ n+1 multiplies each mode by zρ. |zρ|² = |ρ-1|²/|ρ|², scaled by 4:
  N1 = (h-2)² + 4t²,   N0 = h² + 4t².
A mode is unitary (perpetual, neither growing nor dying) iff N1 = N0; it grows iff N1 > N0. -/
namespace RALi

def N1 (h t : Int) : Int := (h - 2) * (h - 2) + 4 * (t * t)
def N0 (h t : Int) : Int := h * h + 4 * (t * t)

theorem N1_sub_N0 (h t : Int) : N1 h t - N0 h t = 4 - 4 * h := by
  simp only [N1, N0, Int.sub_mul, Int.mul_sub]
  omega

/-- 12. A mode is unitary exactly on the line: perpetuity in time is the fixed locus. -/
theorem unitary_iff_on_line (h t : Int) : N1 h t = N0 h t ↔ h = 1 := by
  have := N1_sub_N0 h t
  exact ⟨fun e => by omega, fun e => by omega⟩

/-- 13. Left of the line the mode grows; right of it the mode dies. -/
theorem grows_left (h t : Int) (hl : h < 1) : N0 h t < N1 h t := by
  have := N1_sub_N0 h t; omega
theorem dies_right (h t : Int) (hr : 1 < h) : N1 h t < N0 h t := by
  have := N1_sub_N0 h t; omega

/-- 14. The fold forbids a quiet exit: any off-line zero, together with its mirror under
the functional equation, carries a growing mode. -/
theorem off_line_forces_growth (h t : Int) (hoff : h ≠ 1) :
    N0 h t < N1 h t ∨ N0 (2 - h) t < N1 (2 - h) t := by
  have a := N1_sub_N0 h t; have b := N1_sub_N0 (2 - h) t; omega

/-- 15. PERPETUAL STABILITY IS THE LINE PROPERTY. For any fold-invariant zero set on the
plane: no mode grows under the time step iff every zero lies on h = 1. -/
theorem stability_iff_line (Z : Int × Int → Prop)
    (hinv : ∀ p, Z p → Z (2 - p.1, p.2)) :
    (∀ p, Z p → N1 p.1 p.2 ≤ N0 p.1 p.2) ↔ (∀ p, Z p → p.1 = 1) := by
  constructor
  · intro hs p hp
    have a := hs p hp
    have b := hs (2 - p.1, p.2) (hinv p hp)
    have c := N1_sub_N0 p.1 p.2
    have d := N1_sub_N0 (2 - p.1) p.2
    simp only at b d
    omega
  · intro hl p hp
    have := hl p hp
    have c := N1_sub_N0 p.1 p.2
    omega

end RALi

#print axioms RALi.unitary_iff_on_line
#print axioms RALi.off_line_forces_growth
#print axioms RALi.stability_iff_line

/-! ## PART V · The monism rule, hardened. "What happens once happens again." -/
namespace RALi

/-- 16. One repeating one: the rule, with its step, is induction. -/
theorem one_repeats_one (P : Nat → Prop) (h0 : P 0) (step : ∀ n, P n → P (n + 1)) :
    ∀ n, P n := by
  intro n; induction n with
  | zero => exact h0
  | succ k ih => exact step k ih

/-- 17. The step is the whole content: the universal claim is exactly the base plus the step. -/
theorem the_step_is_the_claim (P : Nat → Prop) :
    (∀ n, P n) ↔ (P 0 ∧ ∀ n, P n → P (n + 1)) :=
  ⟨fun h => ⟨h 0, fun n _ => h (n + 1)⟩, fun ⟨h0, s⟩ => one_repeats_one P h0 s⟩

/-- 18. The monism posit: one principle means the property does not vary with the index.
Under it, one instance forces all. -/
theorem monism_posit_forces (P : Nat → Prop) (uniform : ∀ n m, P n ↔ P m) (h0 : P 0) :
    ∀ n, P n := fun n => (uniform 0 n).mp h0

/-- 19. Finite confirmation never forces. For every height N there is a property true on
every index below N and false beyond: Skewes, Mertens, Pólya are this shape. -/
theorem finite_never_forces (N : Nat) :
    ∃ P : Nat → Prop, (∀ n, n < N → P n) ∧ ¬ ∀ n, P n :=
  ⟨fun n => n < N, fun _ h => h, fun h => Nat.lt_irrefl N (h N)⟩

/-- 20. Monism is not a law of the fold. A fold-invariant zero set can hold one zero on the
line and a mirror pair off it: "on the line once" does not recur by symmetry alone.
This is the Davenport-Heilbronn shape on the Bridge plane. -/
def mixedZ (p : Int × Int) : Prop := p = (1, 0) ∨ p = (0, 5) ∨ p = (2, 5)

theorem mixed_is_fold_invariant : ∀ p, mixedZ p → mixedZ (2 - p.1, p.2) := by
  intro p hp
  rcases hp with h | h | h <;> subst h
  · exact Or.inl rfl
  · exact Or.inr (Or.inr rfl)
  · exact Or.inr (Or.inl rfl)

theorem mixed_recurrence_fails :
    mixedZ (1, 0) ∧ mixedZ (0, 5) ∧ ¬ (∀ p, mixedZ p → p.1 = 1) :=
  ⟨Or.inl rfl, Or.inr (Or.inl rfl), fun h => by
    have := h (0, 5) (Or.inr (Or.inl rfl)); cases this⟩

end RALi

#print axioms RALi.one_repeats_one
#print axioms RALi.the_step_is_the_claim
#print axioms RALi.monism_posit_forces
#print axioms RALi.finite_never_forces
#print axioms RALi.mixed_is_fold_invariant
#print axioms RALi.mixed_recurrence_fails

/-! ## PART VI · Actualism. Only measured zeros are actual. -/
namespace RALi

/-- A zero set on a carrier, a line predicate, and the time at which each zero is measured
(none if never). The actual zeros at time T are those measured by T. -/
structure Actualized (α : Type) where
  Z        : α → Prop
  onLine   : α → Prop
  measured : α → Option Nat

def RHfull {α : Type} (A : Actualized α) : Prop := ∀ z, A.Z z → A.onLine z
def RHactual {α : Type} (A : Actualized α) : Prop :=
  ∀ z, A.Z z → A.measured z ≠ none → A.onLine z

/-- 21. The full claim yields the actual one. -/
theorem full_gives_actual {α : Type} (A : Actualized α) : RHfull A → RHactual A :=
  fun h z hz _ => h z hz

/-- 22. The actual claim does not yield the full one: one never-measured zero off the line. -/
def unreached : Actualized Bool := ⟨fun _ => True, fun b => b = true, fun b => if b then some 0 else none⟩

theorem actual_not_full : RHactual unreached ∧ ¬ RHfull unreached := by
  refine ⟨fun z _ hm => ?_, fun h => by cases h false trivial⟩
  cases z
  · exact absurd rfl hm
  · rfl

/-- 23. Given sufficient time, every zero is measured, and then the actual claim over all
time is exactly the full claim: actualism re-routes to the same step and does not bypass it. -/
theorem sufficient_time {α : Type} (A : Actualized α) (reach : ∀ z, A.Z z → A.measured z ≠ none) :
    RHactual A ↔ RHfull A :=
  ⟨fun h z hz => h z hz (reach z hz), full_gives_actual A⟩

end RALi

#print axioms RALi.full_gives_actual
#print axioms RALi.actual_not_full
#print axioms RALi.sufficient_time

/-! ## PART VII · The constructed witness, reconstructed from Residual Monism.
The paper's witness is built from RA and read through RAF and RAM. Here the witness is
built from the Codex connector directly: one substrate, one involution serving both routes
(σ = σ′), and one principle that does not vary with the index. The index is time. -/
namespace RALi

structure Q4 where
  r : Int
  i : Int
  j : Int
  k : Int
  deriving DecidableEq, Repr

/-- The geometric route's involution and the formal route's involution, written apart. -/
def sigmaGeo (q : Q4) : Q4 := ⟨q.r, -q.i, -q.j, -q.k⟩
def sigmaForm (q : Q4) : Q4 := ⟨q.r, -q.i, -q.j, -q.k⟩

/-- THE MONISM WITNESS. One involution, one principle, one seed. -/
structure MonismWitness (P : Nat → Prop) : Prop where
  one_involution : ∀ q, sigmaGeo q = sigmaForm q
  uniform        : ∀ n m, P n ↔ P m
  seed           : P 0

/-- 24. The seat field is constructed, not posited: σ = σ′ holds by rfl. -/
theorem one_involution_constructed : ∀ q, sigmaGeo q = sigmaForm q := fun _ => rfl

/-- 25. The timeless reading and the timed reading coincide under one principle:
the claim over all time is the claim at one instant. Deleting time loses nothing. -/
theorem timeless_equals_timed (P : Nat → Prop) (u : ∀ n m, P n ↔ P m) :
    (∀ n, P n) ↔ P 0 :=
  ⟨fun h => h 0, fun h n => (u 0 n).mp h⟩

/-- 26. The witness closes the record, past and future. -/
theorem monism_closes (P : Nat → Prop) (w : MonismWitness P) : ∀ n, P n :=
  (timeless_equals_timed P w.uniform).mpr w.seed

/-- 27. The witness is exactly the claim: a monism witness for P exists iff P holds at
every index. Its value field carries the whole of the claim, as theorem 17 required. -/
theorem monism_witness_is_the_claim (P : Nat → Prop) :
    MonismWitness P ↔ ∀ n, P n :=
  ⟨monism_closes P, fun h => ⟨one_involution_constructed,
    fun n m => ⟨fun _ => h m, fun _ => h n⟩, h 0⟩⟩

/-- 28. Constructed at the model, as the paper constructs RA at one point. -/
theorem constructed_monism : MonismWitness (fun _ => True) :=
  ⟨one_involution_constructed, fun _ _ => Iff.rfl, trivial⟩

/-- 29. Unlike the RA witness, the monism witness is not silent on counter-models:
it cannot be constructed where the property breaks, so it decides by containing. -/
theorem monism_absent_on_breaks : ¬ MonismWitness (fun n => n = 0) :=
  fun w => absurd ((w.uniform 0 1).mp rfl) (by decide)

/-- 30. Run through the RA-Li bridge. Replace the bridge by the monism witness on Li's
sign stream: the witness, with Li's criterion, gives RH. -/
theorem monism_yields_RH (L : LiData) (hLi : LiCriterion L)
    (w : MonismWitness (fun n => L.nonneg (n + 1))) : L.RH :=
  hLi.mpr (fun n hn => by
    obtain ⟨k, rfl⟩ : ∃ k, n = k + 1 := ⟨n - 1, by omega⟩
    exact monism_closes _ w k)

/-- 31. And the price, stated: on ζ the witness exists iff every Li sign is nonnegative,
which by Li is RH. Reconstruction moves the owed bit into `uniform`; it does not remove it. -/
theorem monism_on_li_is_RH (L : LiData) (hLi : LiCriterion L) :
    MonismWitness (fun n => L.nonneg (n + 1)) ↔ L.RH := by
  rw [monism_witness_is_the_claim]
  refine ⟨fun h => monism_yields_RH L hLi ((monism_witness_is_the_claim _).mpr h), fun h n => ?_⟩
  exact (hLi.mp h) (n + 1) (by omega)

end RALi

#print axioms RALi.one_involution_constructed
#print axioms RALi.timeless_equals_timed
#print axioms RALi.monism_closes
#print axioms RALi.monism_witness_is_the_claim
#print axioms RALi.constructed_monism
#print axioms RALi.monism_absent_on_breaks
#print axioms RALi.monism_yields_RH
#print axioms RALi.monism_on_li_is_RH

/-! ## PART VIII · Re-anchoring the monism witness to mathematical standard.
Calibration case: Perelman. The standard form of "one principle, the same result again"
is not a free uniformity; it is one dynamics, one invariant, and a PROVED transport of the
invariant along the dynamics (induction on the orbit; Lyapunov monotonicity; Perelman's
W-entropy under Ricci flow). -/
namespace RALi

/-- 32. Monism as a schema over all properties is inconsistent: uniformity cannot be
posited for every P. It must be anchored to a principle, not spread over properties. -/
theorem global_monism_inconsistent : ¬ ∀ P : Nat → Prop, ∀ n m, P n ↔ P m :=
  fun h => absurd ((h (fun n => n = 0) 0 1).mp rfl) (by decide)

def iter {S : Type} (f : S → S) : Nat → S → S
  | 0, s => s
  | n + 1, s => f (iter f n s)

/-- THE ANCHORED MONISM WITNESS. One substrate S, one principle f (the dynamics),
one invariant I, the transport of I along f, and a seed. -/
structure AnchoredMonism (S : Type) where
  f         : S → S
  I         : S → Prop
  transport : ∀ s, I s → I (f s)
  s0        : S
  seed      : I s0

/-- 33. The anchored witness closes its whole orbit, past and future. -/
theorem anchored_closes {S : Type} (A : AnchoredMonism S) : ∀ n, A.I (iter A.f n A.s0) := by
  intro n; induction n with
  | zero => exact A.seed
  | succ k ih => exact A.transport _ ih

/-- 34. The free monism witness is DERIVED from the anchored one: uniformity is no longer
posited, it is manufactured by the transport. -/
theorem anchored_derives_monism {S : Type} (A : AnchoredMonism S) :
    MonismWitness (fun n => A.I (iter A.f n A.s0)) :=
  ⟨one_involution_constructed,
   fun n m => ⟨fun _ => anchored_closes A m, fun _ => anchored_closes A n⟩,
   A.seed⟩

/-- 35. THE PERELMAN SHAPE. A quantity monotone along one flow, here a discrete flow on its
time index, is an anchored monism witness for the property "never below its start." Perelman's
W-entropy along Ricci flow has this shape; it is a structural analogue and is not formalized. -/
def lyapunovWitness (W : Nat → Int) (mono : ∀ n, W n ≤ W (n + 1)) : AnchoredMonism Nat :=
  ⟨Nat.succ, fun n => W 0 ≤ W n, fun n h => Int.le_trans h (mono n), 0, Int.le_refl _⟩

theorem lyapunov_is_monism (W : Nat → Int) (mono : ∀ n, W n ≤ W (n + 1)) :
    MonismWitness (fun n => W 0 ≤ W (iter Nat.succ n 0)) :=
  anchored_derives_monism (lyapunovWitness W mono)

/-- 36. The calibration's grade law: the closure is as strong as its transport. -/
structure GradedAnchor (S : Type) where
  A      : AnchoredMonism S
  tgrade : Grade

def closureGrade {S : Type} (G : GradedAnchor S) : Grade := Grade.weakest G.tgrade .theorem

theorem closure_grade_is_transport_grade {S : Type} (G : GradedAnchor S) :
    closureGrade G = G.tgrade := by
  unfold closureGrade Grade.weakest; cases G.tgrade <;> rfl

/-- 37. The RH instance. On Li's sign stream with the time step as the principle, an
anchored witness exists iff RH. The transport is `nonneg (n+1) → nonneg (n+2)`: the
ζ-analogue of Perelman's monotonicity formula, not yet proved by anyone. -/
theorem rh_anchor_is_the_claim (L : LiData) (hLi : LiCriterion L) :
    (∃ A : AnchoredMonism Nat, A.f = Nat.succ ∧ A.s0 = 0 ∧
        ∀ n, A.I n ↔ L.nonneg (n + 1)) ↔ L.RH := by
  constructor
  · rintro ⟨A, hf, hs, hI⟩
    apply (monism_on_li_is_RH L hLi).mp
    refine (monism_witness_is_the_claim _).mpr (fun n => ?_)
    have key : ∀ n, iter A.f n A.s0 = n := by
      intro n; induction n with
      | zero => exact hs
      | succ k ih => show A.f (iter A.f k A.s0) = k + 1; rw [ih, hf]
    have := anchored_closes A n
    rw [key n] at this
    exact (hI n).mp this
  · intro h
    have all := (monism_witness_is_the_claim _).mp ((monism_on_li_is_RH L hLi).mpr h)
    exact ⟨⟨Nat.succ, fun n => L.nonneg (n + 1), fun n _ => all (n + 1), 0, all 0⟩,
      rfl, rfl, fun _ => Iff.rfl⟩

end RALi

#print axioms RALi.global_monism_inconsistent
#print axioms RALi.anchored_closes
#print axioms RALi.anchored_derives_monism
#print axioms RALi.lyapunov_is_monism
#print axioms RALi.closure_grade_is_transport_grade
#print axioms RALi.rh_anchor_is_the_claim

/-! ## PART IX · ζ's time has a proved monotonicity formula, and it runs the wrong way.
The de Bruijn-Newman heat flow H_t(z) = ∫ e^{t u²} Φ(u) cos(z u) du, Φ built from ξ.
de Bruijn (1950): if every zero of H_0 lies in |Im z| ≤ Δ, every zero of H_t lies in
|Im z| ≤ √max(Δ² − 2t, 0). RH is "all zeros of H_0 real". Λ is the least t with all zeros
of H_t real; Newman (1976) defined it, Rodgers-Tao (2020) proved Λ ≥ 0, Polymath 15 (2019)
proved Λ ≤ 0.22. So RH ⇔ Λ = 0. Modelled here on the strip bound, widths in Nat, half-steps
of time; a toy of the bound, not of the flow. -/
namespace RALi

/-- Squared strip width after t time steps: max(d2 − 2t, 0). -/
def flow (d2 t : Nat) : Nat := d2 - 2 * t
def realAt (d2 t : Nat) : Prop := flow d2 t = 0

/-- 38. THE MONOTONICITY FORMULA. The width never grows along ζ's time. -/
theorem flow_monotone (d2 t : Nat) : flow d2 (t + 1) ≤ flow d2 t := by
  unfold flow; omega

/-- 39. Its transport, proved: once every zero is real, every zero stays real. This is
de Bruijn's theorem in the toy, the ζ-analogue of Perelman's monotonicity, at theorem grade. -/
theorem reality_transported (d2 t : Nat) (h : realAt d2 t) : realAt d2 (t + 1) := by
  unfold realAt flow at *; omega

/-- Λ in the toy: the first time the strip closes. -/
def Lam (d2 : Nat) : Nat := (d2 + 1) / 2

theorem real_at_Lam (d2 : Nat) : realAt d2 (Lam d2) := by unfold realAt flow Lam; omega

/-- 40. In the toy, the strip is closed at time zero iff Λ = 0. For ζ, RH ↔ Λ = 0 is the cited
result of Newman with Rodgers and Tao and enters the cone as `hLam`. -/
theorem rh_iff_lambda_zero (d2 : Nat) : realAt d2 0 ↔ Lam d2 = 0 := by
  unfold realAt flow Lam
  exact ⟨fun h => by omega, fun h => by omega⟩

/-- 41. The anchored monism witness exists for ζ's time, transport proved, seeded at Λ.
It closes the whole future of the flow from Λ on. -/
def deBruijnWitness (d2 : Nat) : AnchoredMonism Nat :=
  ⟨Nat.succ, realAt d2, fun t h => reality_transported d2 t h, Lam d2, real_at_Lam d2⟩

theorem iter_succ (n s : Nat) : iter Nat.succ n s = s + n := by
  induction n with
  | zero => rfl
  | succ k ih => show Nat.succ (iter Nat.succ k s) = s + (k + 1); rw [ih]; omega

theorem future_closed (d2 n : Nat) : realAt d2 (Lam d2 + n) := by
  have := anchored_closes (deBruijnWitness d2) n
  rwa [show (deBruijnWitness d2).f = Nat.succ from rfl, iter_succ] at this

/-- 42. THE RECORD FORGETS. After one step the record of the width-zero state and the
width-one state coincide: the forward flow is even in the RH bit, so no readout of the
flowed record decides RH. This is the fTOE wall, T1, executed on ζ's own time. -/
theorem flowed_record_forgets (T : Nat) (hT : 1 ≤ T) :
    ¬ ∃ g : Nat → Bool, ∀ d2, g (flow d2 T) = decide (d2 = 0) := by
  rintro ⟨g, hg⟩
  have a := hg 0; have b := hg 1
  have e : flow 1 T = flow 0 T := by unfold flow; omega
  rw [e, a] at b
  exact absurd b (by decide)

/-- 43. The upstream direction is not a transport. Going backward from a real record,
both answers are admissible: the preimage of "real at T" holds the RH state and a non-RH
state. The forward arrow proves; the backward arrow must be supplied. -/
theorem upstream_is_not_forced (T : Nat) (hT : 1 ≤ T) :
    realAt 0 T ∧ realAt 1 T ∧ realAt 0 0 ∧ ¬ realAt 1 0 := by
  unfold realAt flow
  exact ⟨by omega, by omega, rfl, fun h => by omega⟩

end RALi

#print axioms RALi.flow_monotone
#print axioms RALi.reality_transported
#print axioms RALi.rh_iff_lambda_zero
#print axioms RALi.future_closed
#print axioms RALi.flowed_record_forgets
#print axioms RALi.upstream_is_not_forced

/-! ## PART X · Halted Uniduction with the category-gap eliminator. The regress, closed.
Every "missing piece" of Parts I-IX is a register's reading of one seat. Exhibit each as a
leg of one cone with apex RH; then the hunt has one gap, not a sequence of them. -/
namespace RALi

structure Hunt where
  S     : Substrate
  L     : LiData
  u0    : S.U
  hRA   : RA S
  hLi   : LiCriterion L
  lamZero : Prop                -- the upstream reading: Λ = 0
  hLam  : L.RH ↔ lamZero        -- Newman's definition with Rodgers-Tao, cited

/-- The legs. Each register's reading of the owed seat, proved equal to the apex. -/
theorem leg_bridge (H : Hunt) : Nonempty (Bridge H.S H.L) ↔ H.L.RH :=
  ⟨fun ⟨B⟩ => bridge_yields_RH H.S H.L H.hRA H.hLi B,
   fun h => ⟨⟨fun _ => H.u0, fun n hn _ => (H.hLi.mp h) n hn⟩⟩⟩

theorem leg_unity (H : Hunt) : Nonempty (Unity H.S H.L) ↔ H.L.RH :=
  ⟨fun ⟨u⟩ => retired_seal_is_sound H.S H.L H.hRA H.hLi u,
   fun h => let ⟨B⟩ := (leg_bridge H).mpr h; ⟨⟨B, .premise⟩⟩⟩

theorem leg_monism (H : Hunt) : MonismWitness (fun n => H.L.nonneg (n + 1)) ↔ H.L.RH :=
  monism_on_li_is_RH H.L H.hLi

theorem leg_transport (H : Hunt) :
    (∃ A : AnchoredMonism Nat, A.f = Nat.succ ∧ A.s0 = 0 ∧ ∀ n, A.I n ↔ H.L.nonneg (n + 1))
      ↔ H.L.RH :=
  rh_anchor_is_the_claim H.L H.hLi

theorem leg_upstream (H : Hunt) : H.lamZero ↔ H.L.RH := H.hLam.symm

/-- The five readings, as a diagram of propositions. -/
inductive Reg5 | bridge | unity | monism | transport | upstream deriving DecidableEq, Repr

def D5 (H : Hunt) : Reg5 → Prop
  | .bridge    => Nonempty (Bridge H.S H.L)
  | .unity     => Nonempty (Unity H.S H.L)
  | .monism    => MonismWitness (fun n => H.L.nonneg (n + 1))
  | .transport => ∃ A : AnchoredMonism Nat, A.f = Nat.succ ∧ A.s0 = 0 ∧ ∀ n, A.I n ↔ H.L.nonneg (n + 1)
  | .upstream  => H.lamZero

/-- 44. THE CONE. RH is the apex of a cone over all five readings, every leg an
equivalence: the category gaps between the hunt's registers are eliminated. -/
theorem hunt_cone (H : Hunt) : ∀ r, D5 H r ↔ H.L.RH := by
  intro r; cases r
  · exact leg_bridge H
  · exact leg_unity H
  · exact leg_monism H
  · exact leg_transport H
  · exact leg_upstream H

/-- 45. ONE GAP. Any two readings are the same proposition. The quantity monotone against
the flow (upstream) IS the bridge's read, the unity posit, the monism field, the transport. -/
theorem one_gap (H : Hunt) (r s : Reg5) : D5 H r ↔ D5 H s :=
  (hunt_cone H r).trans (hunt_cone H s).symm

/-- 46. THE HALT. Every proposition proved equivalent to the hypothesis lands on the same apex
as the five readings; which future propositions are equivalent is not decided here. -/
theorem regress_halts (H : Hunt) (P : Prop) (hP : P ↔ H.L.RH) (r : Reg5) : P ↔ D5 H r :=
  hP.trans (hunt_cone H r).symm

/-- 47. Given a true proposition Q, proving Q → RH is the same as proving RH: conditioning on a
true premise adds nothing (Theorem E of the prior paper, here on the hunt's apex). -/
theorem nothing_weaker (H : Hunt) (Q : Prop) (hq : Q) : (Q → H.L.RH) ↔ H.L.RH :=
  ⟨fun f => f hq, fun h _ => h⟩

end RALi

#print axioms RALi.hunt_cone
#print axioms RALi.one_gap
#print axioms RALi.regress_halts
#print axioms RALi.nothing_weaker

/-! ## PART X, completed · a sufficient route closes every reading at once. -/
namespace RALi

/-- 48. Theorem 46 covers every reformulation equivalent to the apex. A route strictly
stronger than the apex (a specific operator, a hypothesis over a family) is not equivalent,
and it does not open a new gap either: any sufficient route supplies all five readings. -/
theorem sufficient_closes_all (H : Hunt) (P : Prop) (hP : P → H.L.RH) (p : P) :
    ∀ r, D5 H r :=
  fun r => (hunt_cone H r).mpr (hP p)

end RALi

#print axioms RALi.sufficient_closes_all

/-! ## PART XI · Timeless Residual Monism, stated as Postulate M, and its equivalence with the line property. -/
namespace TimeLocus

def onLine (p : Int × Int) : Prop := p.1 = 1

/-- A timed world: the locus at each time, and the zeros present at each time. -/
structure World where
  locus : Nat → (Int × Int → Prop)
  zeros : Nat → (Int × Int → Prop)

/-- TIME PSP premise: prior and posterior locus are the SAME locus, the line, at every time. -/
def SameLocus (W : World) : Prop := ∀ t p, W.locus t p ↔ onLine p

/-- The claim to be tested: every zero at every time is bound to the locus. -/
def Bound (W : World) : Prop := ∀ t p, W.zeros t p → W.locus t p

/-- Fold-invariance at every time: the functional equation acts at each iteration. -/
def FoldInv (W : World) : Prop := ∀ t p, W.zeros t p → W.zeros t (2 - p.1, p.2)

/-- 51. Given the same locus at all times, being bound IS the line property at all times. -/
theorem bound_iff_line (W : World) (h : SameLocus W) :
    Bound W ↔ ∀ t p, W.zeros t p → onLine p :=
  ⟨fun b t p z => (h t p).mp (b t p z), fun l t p z => (h t p).mpr (l t p z)⟩

/-- 52, the test world. A world where the locus is the same line at every time, the fold acts at every
    time, the on-line zero recurs at every time, and a mirror pair sits off the line at every time. -/
def testWorld : World :=
  ⟨fun _ p => onLine p, fun _ p => p = (1, 0) ∨ p = (0, 5) ∨ p = (2, 5)⟩

theorem test_same_locus : SameLocus testWorld := fun _ _ => Iff.rfl

theorem test_fold_inv : FoldInv testWorld := by
  intro _ p hp
  rcases hp with h | h | h <;> subst h
  · exact Or.inl rfl
  · exact Or.inr (Or.inr rfl)
  · exact Or.inr (Or.inl rfl)

theorem test_not_bound : ¬ Bound testWorld := fun b => by
  have := b 0 (0, 5) (Or.inr (Or.inl rfl)); cases this

/-- 52. Therefore SameLocus and FoldInv, at every time, do not imply Bound. -/
theorem time_does_not_bind :
    ¬ ∀ W : World, SameLocus W → FoldInv W → Bound W :=
  fun h => test_not_bound (h testWorld test_same_locus test_fold_inv)

end TimeLocus

namespace TimeLocus

/-- Residual Monism read as the Codex states it: one involution serves both routes. -/
def sigmaGeo (p : Int × Int) : Int × Int := (2 - p.1, p.2)
def sigmaForm (p : Int × Int) : Int × Int := (2 - p.1, p.2)
def OneInvolution : Prop := ∀ p, sigmaGeo p = sigmaForm p

/-- Residual Monism read as timeless: whether the zeros sit on the locus does not vary in time. -/
def P (W : World) (t : Nat) : Prop := ∀ p, W.zeros t p → onLine p
def Timeless (W : World) : Prop := ∀ t s, P W t ↔ P W s

/-- 53. Both readings hold in the counter-world: one involution by rfl, and timelessness because
    the property "all zeros on the line" is constantly false there. Timeless and unbound. -/
theorem monism_holds_in_counter_world :
    OneInvolution ∧ Timeless testWorld ∧ ¬ Bound testWorld :=
  ⟨fun _ => rfl, fun _ _ => Iff.rfl, test_not_bound⟩

/-- 54. What timeless monism does do: it makes one instant decide all time. With the seed, binding
    follows; the seed is "every zero present at the instant is on the line". -/
theorem timeless_with_seed_binds (W : World) (h : SameLocus W) (u : Timeless W) (s0 : P W 0) :
    Bound W :=
  (bound_iff_line W h).mpr (fun t p z => ((u 0 t).mp s0) p z)

end TimeLocus

namespace TimeLocus

/-- THE CONSTRUCTED WITNESS from RESIDUAL-MONISM + TIME PSP, as asked: one involution, the
    timeless interface (the Barzakh instant, t = 0, upstream of time), and the downstream times
    agreeing with it. No field is a premise; each must be built as a term. -/
structure MonismTimeWitness (W : World) : Prop where
  one_involution : OneInvolution
  same_locus     : SameLocus W
  timeless       : Timeless W
  interface      : P W 0          -- every zero present at the timeless interface is on the line

/-- 55. When it can be built, it proves the bound everywhere downstream. -/
theorem witness_proves_bound (W : World) (w : MonismTimeWitness W) : Bound W :=
  timeless_with_seed_binds W w.same_locus w.timeless w.interface

/-- 55, completed. It can be built exactly when the bound already holds: its existence IS the line property. -/
theorem witness_iff_bound (W : World) (h : SameLocus W) :
    Nonempty (MonismTimeWitness W) ↔ Bound W := by
  refine ⟨fun ⟨w⟩ => witness_proves_bound W w, fun b => ?_⟩
  have l := (bound_iff_line W h).mp b
  exact ⟨⟨fun _ => rfl, h, fun t s => ⟨fun _ p z => l s p z, fun _ p z => l t p z⟩, fun p z => l 0 p z⟩⟩

/-- 56. Built at a model whose zeros are on the line, as the prior paper builds RA at one point. -/
def goodWorld : World := ⟨fun _ p => onLine p, fun _ p => p = (1, 0)⟩
theorem constructed_at_model : Nonempty (MonismTimeWitness goodWorld) :=
  (witness_iff_bound goodWorld (fun _ _ => Iff.rfl)).mpr (fun _ p z => by subst z; rfl)

/-- 57. Not buildable at the counter-world: the interface field fails there. -/
theorem not_constructible_off_line : ¬ Nonempty (MonismTimeWitness testWorld) :=
  fun w => test_not_bound ((witness_iff_bound testWorld test_same_locus).mp w)

end TimeLocus

namespace TimeLocus

/-- Postulate M on an enumerated zero set. Stage k is when a zero is located; L t says every
    zero located by stage t lies on the line; M says the truth of L t does not depend on t. -/
def Lstage {α : Type} (Z : α → Prop) (onL : α → Prop) (stage : α → Nat) (t : Nat) : Prop :=
  ∀ z, Z z → stage z ≤ t → onL z
def PostulateM {α : Type} (Z : α → Prop) (onL : α → Prop) (stage : α → Nat) : Prop :=
  ∀ t s, Lstage Z onL stage t ↔ Lstage Z onL stage s

/-- 49. POSTULATE M IS EQUIVALENT TO THE LINE PROPERTY, given the computed seed L 0 and an
    exhaustive enumeration (every zero is located at some finite stage, built into `stage`). -/
theorem postulateM_iff_line {α : Type} (Z : α → Prop) (onL : α → Prop) (stage : α → Nat)
    (seed : Lstage Z onL stage 0) :
    PostulateM Z onL stage ↔ ∀ z, Z z → onL z := by
  constructor
  · intro m z hz
    exact ((m 0 (stage z)).mp seed) z hz (Nat.le_refl _)
  · intro h t s
    exact ⟨fun _ z hz _ => h z hz, fun _ z hz _ => h z hz⟩

/-- 50. And Postulate M is not free: it holds in a world whose zeros are all on the line, and
    it cannot hold with the seed in a world with a located off-line zero. -/
theorem postulateM_decides {α : Type} (Z : α → Prop) (onL : α → Prop) (stage : α → Nat)
    (seed : Lstage Z onL stage 0) (z : α) (hz : Z z) (off : ¬ onL z) :
    ¬ PostulateM Z onL stage :=
  fun m => off ((postulateM_iff_line Z onL stage seed).mp m z hz)

end TimeLocus

#print axioms TimeLocus.bound_iff_line
#print axioms TimeLocus.time_does_not_bind
#print axioms TimeLocus.monism_holds_in_counter_world
#print axioms TimeLocus.timeless_with_seed_binds
#print axioms TimeLocus.witness_iff_bound
#print axioms TimeLocus.constructed_at_model
#print axioms TimeLocus.not_constructible_off_line
#print axioms TimeLocus.postulateM_iff_line
#print axioms TimeLocus.postulateM_decides

/-! ## PART XII · The demand correctly addressed, and [.] at the collapse foundation. -/
namespace DotRH

inductive Token where
  | sealed | broken | opn | dot
  deriving DecidableEq, Repr

def Token.inEconomy : Token → Bool
  | .dot => false
  | _    => true

/-- I · THE COLLAPSE FOUNDATION and its stop. -/
structure Collapse (V : Prop) where
  Reg     : Type
  reading : Reg → Prop
  leg     : ∀ r, reading r ↔ V

-- Theorem 58.
theorem self_referential_limit {V : Prop} (C : Collapse V) (r s : C.Reg) :
    (C.reading r ↔ V) ∧ (C.reading r ↔ C.reading s) :=
  ⟨C.leg r, (C.leg r).trans (C.leg s).symm⟩

def emitAtFoundation {V : Prop} (_ : Collapse V) : Token := .dot
-- Theorem 59.
theorem value_marked_dot {V : Prop} (C : Collapse V) : emitAtFoundation C = .dot := rfl
theorem dot_is_not_a_verdict :
    Token.dot ≠ .sealed ∧ Token.dot ≠ .broken ∧ Token.dot ≠ .opn := by decide
theorem dot_outside_economy : Token.inEconomy .dot = false := rfl

def trivialCollapse (V : Prop) : Collapse V := ⟨Unit, fun _ => V, fun _ => Iff.rfl⟩

/-- Inside the collapse register, a demand for a verdict is a ghost: the register returns
    the same mark for the value and its negation. -/
-- Theorem 60.
theorem demand_is_ghost_in_register {V : Prop} (C : Collapse V) (C' : Collapse (¬ V)) :
    emitAtFoundation C = emitAtFoundation C' := rfl

/-- II · THE SELF-GROUNDING ROOT. A root R is self-grounding when acts occur and every act of
    adjudication, assent or denial, is itself an instance of R. Then a denial of R re-enacts R,
    no external proof adds to it, and a demand for one is a ghost. This is the seat of the
    analogy: a root that grounds itself cannot be proved from outside itself. -/
structure SelfGrounding (R : Prop) where
  Act       : Type
  anAct     : Act
  instances : Act → R

-- Theorem 61.
theorem denial_reenacts_root {R : Prop} (G : SelfGrounding R) (denial : G.Act) : R :=
  G.instances denial

-- Theorem 61, completed.
theorem external_proof_adds_nothing {R : Prop} (G : SelfGrounding R) (Q : Prop) :
    (Q → R) ↔ R :=
  ⟨fun _ => G.instances G.anAct, fun r _ => r⟩

/-- The Root Axiom at the constructed domain is self-grounding: every act is a deed, and a
    deed actuates. -/
def RA : Prop := (0 : Int) < 1
def raSelfGrounding : SelfGrounding RA := ⟨Unit, (), fun _ => show (0 : Int) < 1 by decide⟩

/-- III · A HYPOTHESIS ABOUT AN OBJECT is not self-grounding. Its value is fixed by the object,
    acts do not instance it, and a finite witness can refute it. -/
structure Frame where
  S : Type
  τ : S → S
  Z : S → Prop

def LineProperty (X : Frame) : Prop := ∀ s, X.Z s → X.τ s = s

/-- A located zero off the line refutes the line property, constructively. -/
theorem refuted_by_witness (X : Frame) (s : X.S) (hz : X.Z s) (off : X.τ s ≠ s) :
    ¬ LineProperty X :=
  fun h => off (h s hz)

def twoPoint : Frame := ⟨Bool, fun b => !b, fun _ => True⟩

/-- On a frame whose line property fails, no occurring act can instance it: the line
    property of a frame is not self-grounding. -/
-- Theorem 62.
theorem line_not_self_grounding : ¬ Nonempty (SelfGrounding (LineProperty twoPoint)) :=
  fun ⟨G⟩ => refuted_by_witness twoPoint true trivial (fun h => Bool.noConfusion h) (G.instances G.anAct)

/-- IV · THE HARDENED FOUNDATION, in toto. The collapse marks the value [.] and the demand is
    a ghost inside that register; the self-grounding root cannot be proved from outside and
    re-enacts under denial; a hypothesis about an object is witness-refutable and is not
    self-grounding, so the self-grounding exemption does not transfer to it. -/
theorem foundation_hardened :
    (∀ (V : Prop) (C : Collapse V), emitAtFoundation C = .dot) ∧
    Token.inEconomy .dot = false ∧
    (∀ (V : Prop) (C : Collapse V) (C' : Collapse (¬ V)), emitAtFoundation C = emitAtFoundation C') ∧
    (∀ Q : Prop, (Q → RA) ↔ RA) ∧
    (∀ a : Unit, RA ∧ raSelfGrounding.instances a = raSelfGrounding.instances a) ∧
    ¬ Nonempty (SelfGrounding (LineProperty twoPoint)) :=
  ⟨fun _ _ => rfl, rfl, fun _ _ _ => rfl, external_proof_adds_nothing raSelfGrounding,
   fun a => ⟨raSelfGrounding.instances a, rfl⟩, line_not_self_grounding⟩

/-! V · THE DEMAND, CORRECTLY ADDRESSED.
    "A proof of the Riemann Hypothesis cannot be demanded of the foundation. Every reading the
    foundation supplies is equivalent to the hypothesis, so any derivation from those readings
    alone would presuppose what it derives. The demand is well-posed when addressed to the
    object, the function zeta, whose structure fixes where its zeros lie." -/

/-- A foundation resource is any property of frames the foundation supplies. -/
def Resource := Frame → Prop

/-- (a) Misaddressed: any resource that also holds on a frame where the line property fails
    cannot, by itself, yield the line property on every frame it covers. -/
-- Theorem 63.
theorem misaddressed_to_foundation (R : Resource) (hR : R twoPoint) :
    ¬ ∀ X, R X → LineProperty X :=
  fun h => refuted_by_witness twoPoint true trivial (fun e => Bool.noConfusion e) (h twoPoint hR)

/-- (b) Presupposition: a derivation of the value from a reading equivalent to it uses the
    value's own content; the reading and the value stand or fall together. -/
-- Theorem 64.
theorem derivation_presupposes {V P : Prop} (leg : P ↔ V) : (P → V) ∧ (V → P) ∧ (¬ V → ¬ P) :=
  ⟨leg.mp, leg.mpr, fun nv p => nv (leg.mp p)⟩

/-- (c) Well-addressed: the object decides. A frame settles its own line property, in one
    direction by a located off-line witness, in the other by its own structure. -/
def onLineFrame : Frame := ⟨Unit, fun u => u, fun _ => True⟩
-- Theorem 65.
theorem object_decides :
    ¬ LineProperty twoPoint ∧ LineProperty onLineFrame :=
  ⟨refuted_by_witness twoPoint true trivial (fun e => Bool.noConfusion e), fun _ _ => rfl⟩

/-- THE DEMAND, IN TOTO. Made to the foundation, it is misaddressed and any answer from
    there presupposes the value; made to the object, it is well-posed and the object answers. -/
-- Theorem 66.
theorem demand_correctly_addressed :
    (∀ R : Resource, R twoPoint → ¬ ∀ X, R X → LineProperty X) ∧
    (∀ {V P : Prop}, (P ↔ V) → (¬ V → ¬ P)) ∧
    (¬ LineProperty twoPoint ∧ LineProperty onLineFrame) :=
  ⟨misaddressed_to_foundation, fun leg => (derivation_presupposes leg).2.2, object_decides⟩

end DotRH

namespace DotRH
/-- Theorem 67. Attached to the paper's cone: the readings of Parts X and XI collapse onto
    the hypothesis, and the foundation marks the value [.]. -/
def rhCollapse (H : RALi.Hunt) : Collapse H.L.RH := ⟨RALi.Reg5, RALi.D5 H, RALi.hunt_cone H⟩
theorem rh_marked_dot (H : RALi.Hunt) : emitAtFoundation (rhCollapse H) = .dot := rfl
end DotRH

#print axioms DotRH.self_referential_limit
#print axioms DotRH.value_marked_dot
#print axioms DotRH.dot_is_not_a_verdict
#print axioms DotRH.demand_is_ghost_in_register
#print axioms DotRH.denial_reenacts_root
#print axioms DotRH.external_proof_adds_nothing
#print axioms DotRH.refuted_by_witness
#print axioms DotRH.line_not_self_grounding
#print axioms DotRH.foundation_hardened
#print axioms DotRH.misaddressed_to_foundation
#print axioms DotRH.derivation_presupposes
#print axioms DotRH.object_decides
#print axioms DotRH.demand_correctly_addressed
#print axioms DotRH.rh_marked_dot

/-! ## PART XIII · The six-register cone: Postulate M joined to the apex. -/
namespace SixCone

/-- The hunt with one more reading. The zeros are enumerated by stage, the seed holds, and
    the hypothesis is, by its definition, the line property of the enumerated zeros; that
    identification enters as the hypothesis `hDef`, carried by citation as Li's criterion and
    Newman's equivalence are. -/
structure Hunt6 extends RALi.Hunt where
  α     : Type
  Z     : α → Prop
  onL   : α → Prop
  stage : α → Nat
  seed  : TimeLocus.Lstage Z onL stage 0
  hDef  : L.RH ↔ ∀ z, Z z → onL z

inductive Reg6 where
  | five (r : RALi.Reg5)
  | postulateM

def D6 (H : Hunt6) : Reg6 → Prop
  | .five r     => RALi.D5 H.toHunt r
  | .postulateM => TimeLocus.PostulateM H.Z H.onL H.stage

-- Theorem 68.
/-- The cone over six readings: the five of Part X and Postulate M, every leg an equivalence
    with the hypothesis. -/
theorem hunt_cone6 (H : Hunt6) : ∀ r, D6 H r ↔ H.L.RH := by
  intro r
  cases r with
  | five r => exact RALi.hunt_cone H.toHunt r
  | postulateM =>
      exact (TimeLocus.postulateM_iff_line H.Z H.onL H.stage H.seed).trans H.hDef.symm

-- Theorem 69.
/-- The collapse foundation over six readings, and its mark on the value. -/
def rhCollapse6 (H : Hunt6) : DotRH.Collapse H.L.RH := ⟨Reg6, D6 H, hunt_cone6 H⟩
theorem rh_marked_dot6 (H : Hunt6) : DotRH.emitAtFoundation (rhCollapse6 H) = .dot := rfl

end SixCone

#print axioms SixCone.hunt_cone6
#print axioms SixCone.rh_marked_dot6

namespace DotRH
-- Theorem 63, general form.
/-- Any resource that also holds on some frame whose line property fails cannot, by itself,
    yield the line property on every frame it covers. -/
theorem misaddressed_general (R : Resource) (Y : Frame) (hY : ¬ LineProperty Y) (hR : R Y) :
    ¬ ∀ X, R X → LineProperty X :=
  fun h => hY (h Y hR)
end DotRH

#print axioms DotRH.misaddressed_general

/-! ## PART XIV · The logical form of the hypothesis: the only existence lives in the denial. -/
namespace LogicalForm

-- Theorem 70.
/-- A universal sentence carries no existential import: over an empty zero set the line
    property holds vacuously. The hypothesis posits no object. -/
theorem universal_posits_nothing {α : Type} (onL : α → Prop) :
    ∀ z : α, (fun _ => False) z → onL z :=
  fun _ h => h.elim

-- Theorem 71.
/-- The denial posits an object: a located zero off the line refutes the hypothesis. -/
theorem denial_posits_a_witness {α : Type} (Z onL : α → Prop) :
    (∃ z, Z z ∧ ¬ onL z) → ¬ ∀ z, Z z → onL z :=
  fun ⟨z, hz, off⟩ h => off (h z hz)

-- Theorem 72.
/-- With a decidable line predicate, the hypothesis fails only by a witness: it holds iff no
    off-line zero exists. The denial carries the whole existential load. -/
theorem fails_only_by_witness {α : Type} (Z onL : α → Prop) [∀ z, Decidable (onL z)] :
    (∀ z, Z z → onL z) ↔ ¬ ∃ z, Z z ∧ ¬ onL z := by
  constructor
  · intro h ⟨z, hz, off⟩; exact off (h z hz)
  · intro h z hz
    exact Decidable.byContradiction (fun off => h ⟨z, hz, off⟩)

-- Theorem 73.
/-- A witness is a finite check: for a decidable predicate on the naturals, one index at which
    the check fails refutes the universal, and the check at that index is a computation. -/
theorem witness_is_a_finite_check (bad : Nat → Bool) (n : Nat) (h : bad n = true) :
    ¬ ∀ m, bad m = false :=
  fun hall => by rw [hall n] at h; cases h

end LogicalForm

#print axioms LogicalForm.universal_posits_nothing
#print axioms LogicalForm.denial_posits_a_witness
#print axioms LogicalForm.fails_only_by_witness
#print axioms LogicalForm.witness_is_a_finite_check

/-! ## PART XV · The two "can'ts" are not equal: the denial carries the burden.
Abstract: a proposition RH, the foundation's provability predicate, and the property that makes
the directions unequal, Σ₁-completeness: a false RH has a finite witness, and the foundation
proves the denial from it. That property and soundness on the denial are carried as named
hypotheses, cited from arithmetic, not proved here. -/
namespace DirCant

structure Setting where
  RH        : Prop
  Prov      : Prop → Prop
  /-- Σ₁-completeness, carried as a hypothesis: a false RH is refutable. -/
  sigma1    : ¬ RH → Prov (¬ RH)
  /-- Soundness of the foundation on the denial. -/
  sound_neg : Prov (¬ RH) → ¬ RH

-- Theorem 74.
/-- THE SEALED DIRECTION. If the foundation cannot refute the hypothesis, the hypothesis holds:
    the denial's "can't" is decisive. -/
theorem cant_refute_seals (S : Setting) [Decidable S.RH] (h : ¬ S.Prov (¬ S.RH)) : S.RH :=
  Decidable.byContradiction (fun n => h (S.sigma1 n))

-- Theorem 75.
/-- Under soundness on the denial, the hypothesis holds exactly when the foundation cannot refute
    it: the denial's "can't" and the hypothesis are one statement. -/
theorem rh_iff_cant_refute (S : Setting) [Decidable S.RH] : S.RH ↔ ¬ S.Prov (¬ S.RH) :=
  ⟨fun r p => S.sound_neg p r, cant_refute_seals S⟩

-- Theorem 76.
/-- The assent's "can't" decides nothing: a setting where the hypothesis holds and the
    foundation cannot prove it satisfies every field. -/
def independentTrue : Setting :=
  ⟨True, fun _ => False, fun n => absurd trivial n, fun p => p.elim⟩

theorem cant_prove_does_not_seal_false :
    independentTrue.RH ∧ ¬ independentTrue.Prov independentTrue.RH :=
  ⟨trivial, id⟩

-- Theorem 77.
/-- THE ASYMMETRY. Can't-refute seals the hypothesis; can't-prove is compatible with its truth.
    The burden of the one bit lies on the denial side. -/
theorem asymmetry :
    (∀ (S : Setting) [Decidable S.RH], ¬ S.Prov (¬ S.RH) → S.RH) ∧
    (∃ S : Setting, S.RH ∧ ¬ S.Prov S.RH) :=
  ⟨fun S _ h => cant_refute_seals S h, ⟨independentTrue, cant_prove_does_not_seal_false⟩⟩

end DirCant

#print axioms DirCant.cant_refute_seals
#print axioms DirCant.rh_iff_cant_refute
#print axioms DirCant.cant_prove_does_not_seal_false
#print axioms DirCant.asymmetry
\end{Verbatim}

\hypertarget{appendix-g-the-kernel-audit-of-ra_li_bridge.lean-verbatim}{%
\section{Appendix G: The Kernel Audit of RA\_Li\_Bridge.lean,
Verbatim}\label{appendix-g-the-kernel-audit-of-ra_li_bridge.lean-verbatim}}

\texttt{lean\ RA\_Li\_Bridge.lean}, exit 0, no warnings, no errors.

\begin{Verbatim}[numbers=left,numbersep=6pt,xleftmargin=2.2em,fontsize=\scriptsize,breaklines,breakanywhere,frame=lines,framesep=1.5mm,rulecolor=\color{black!35}]
'RALi.bridge_yields_RH' does not depend on any axioms
'RALi.bridge_is_keyed' does not depend on any axioms
'RALi.no_uniform_bridge' does not depend on any axioms
'RALi.li_is_a_bit_stream' does not depend on any axioms
'RALi.xi0_retires_iff_unity' depends on axioms: [propext]
'RALi.retired_seal_is_sound' does not depend on any axioms
'RALi.retired_grade_capped' does not depend on any axioms
'RALi.unity_is_keyed' does not depend on any axioms
'RALi.ghost' does not depend on any axioms
'RALi.common_phi_is_uniform' does not depend on any axioms
'RALi.blind_constructions_die' does not depend on any axioms
'RALi.mertens_nonneg' depends on axioms: [propext, Quot.sound]
'RALi.unitary_iff_on_line' depends on axioms: [propext, Quot.sound]
'RALi.off_line_forces_growth' depends on axioms: [propext, Quot.sound]
'RALi.stability_iff_line' depends on axioms: [propext, Quot.sound]
'RALi.one_repeats_one' does not depend on any axioms
'RALi.the_step_is_the_claim' does not depend on any axioms
'RALi.monism_posit_forces' does not depend on any axioms
'RALi.finite_never_forces' does not depend on any axioms
'RALi.mixed_is_fold_invariant' does not depend on any axioms
'RALi.mixed_recurrence_fails' does not depend on any axioms
'RALi.full_gives_actual' does not depend on any axioms
'RALi.actual_not_full' does not depend on any axioms
'RALi.sufficient_time' does not depend on any axioms
'RALi.one_involution_constructed' does not depend on any axioms
'RALi.timeless_equals_timed' does not depend on any axioms
'RALi.monism_closes' does not depend on any axioms
'RALi.monism_witness_is_the_claim' does not depend on any axioms
'RALi.constructed_monism' does not depend on any axioms
'RALi.monism_absent_on_breaks' does not depend on any axioms
'RALi.monism_yields_RH' depends on axioms: [propext, Quot.sound]
'RALi.monism_on_li_is_RH' depends on axioms: [propext, Quot.sound]
'RALi.global_monism_inconsistent' does not depend on any axioms
'RALi.anchored_closes' does not depend on any axioms
'RALi.anchored_derives_monism' does not depend on any axioms
'RALi.lyapunov_is_monism' depends on axioms: [propext]
'RALi.closure_grade_is_transport_grade' does not depend on any axioms
'RALi.rh_anchor_is_the_claim' depends on axioms: [propext, Quot.sound]
'RALi.flow_monotone' depends on axioms: [propext, Quot.sound]
'RALi.reality_transported' depends on axioms: [propext, Quot.sound]
'RALi.rh_iff_lambda_zero' depends on axioms: [propext, Quot.sound]
'RALi.future_closed' depends on axioms: [propext, Quot.sound]
'RALi.flowed_record_forgets' depends on axioms: [propext, Quot.sound]
'RALi.upstream_is_not_forced' depends on axioms: [propext, Quot.sound]
'RALi.hunt_cone' depends on axioms: [propext, Quot.sound]
'RALi.one_gap' depends on axioms: [propext, Quot.sound]
'RALi.regress_halts' depends on axioms: [propext, Quot.sound]
'RALi.nothing_weaker' does not depend on any axioms
'RALi.sufficient_closes_all' depends on axioms: [propext, Quot.sound]
'TimeLocus.bound_iff_line' does not depend on any axioms
'TimeLocus.time_does_not_bind' does not depend on any axioms
'TimeLocus.monism_holds_in_counter_world' does not depend on any axioms
'TimeLocus.timeless_with_seed_binds' does not depend on any axioms
'TimeLocus.witness_iff_bound' does not depend on any axioms
'TimeLocus.constructed_at_model' does not depend on any axioms
'TimeLocus.not_constructible_off_line' does not depend on any axioms
'TimeLocus.postulateM_iff_line' does not depend on any axioms
'TimeLocus.postulateM_decides' does not depend on any axioms
'DotRH.self_referential_limit' does not depend on any axioms
'DotRH.value_marked_dot' does not depend on any axioms
'DotRH.dot_is_not_a_verdict' does not depend on any axioms
'DotRH.demand_is_ghost_in_register' does not depend on any axioms
'DotRH.denial_reenacts_root' does not depend on any axioms
'DotRH.external_proof_adds_nothing' does not depend on any axioms
'DotRH.refuted_by_witness' does not depend on any axioms
'DotRH.line_not_self_grounding' does not depend on any axioms
'DotRH.foundation_hardened' does not depend on any axioms
'DotRH.misaddressed_to_foundation' does not depend on any axioms
'DotRH.derivation_presupposes' does not depend on any axioms
'DotRH.object_decides' does not depend on any axioms
'DotRH.demand_correctly_addressed' does not depend on any axioms
'DotRH.rh_marked_dot' depends on axioms: [propext, Quot.sound]
'SixCone.hunt_cone6' depends on axioms: [propext, Quot.sound]
'SixCone.rh_marked_dot6' depends on axioms: [propext, Quot.sound]
'DotRH.misaddressed_general' does not depend on any axioms
'LogicalForm.universal_posits_nothing' does not depend on any axioms
'LogicalForm.denial_posits_a_witness' does not depend on any axioms
'LogicalForm.fails_only_by_witness' does not depend on any axioms
'LogicalForm.witness_is_a_finite_check' does not depend on any axioms
'DirCant.cant_refute_seals' does not depend on any axioms
'DirCant.rh_iff_cant_refute' does not depend on any axioms
'DirCant.cant_prove_does_not_seal_false' does not depend on any axioms
'DirCant.asymmetry' does not depend on any axioms
\end{Verbatim}

\hypertarget{appendix-g-completed-the-seventeen-theorems-the-file-does-not-print}{%
\subsection{Appendix G, completed: the seventeen theorems the file does
not
print}\label{appendix-g-completed-the-seventeen-theorems-the-file-does-not-print}}

The file prints 83 dependency sets. A supplementary run of the same
file, with one \texttt{\#print\ axioms} line appended for each of the 17
theorems it does not print, exit 0: 9 axiom-free, 8 on \texttt{propext}
and \texttt{Quot.sound}, none on \texttt{Classical.choice}; 100 sets in
all, 70 axiom-free.

\begingroup\scriptsize

\begin{Verbatim}[numbers=left,numbersep=6pt,xleftmargin=2.2em,fontsize=\scriptsize,breaklines,breakanywhere,frame=lines,framesep=1.5mm,rulecolor=\color{black!35}]
'RALi.ra_holds_in_model' does not depend on any axioms
'RALi.mertens_identity' depends on axioms: [propext, Quot.sound]
'RALi.N1_sub_N0' depends on axioms: [propext, Quot.sound]
'RALi.grows_left' depends on axioms: [propext, Quot.sound]
'RALi.dies_right' depends on axioms: [propext, Quot.sound]
'RALi.real_at_Lam' depends on axioms: [propext, Quot.sound]
'RALi.iter_succ' depends on axioms: [propext, Quot.sound]
'RALi.leg_bridge' does not depend on any axioms
'RALi.leg_unity' does not depend on any axioms
'RALi.leg_monism' depends on axioms: [propext, Quot.sound]
'RALi.leg_transport' depends on axioms: [propext, Quot.sound]
'RALi.leg_upstream' does not depend on any axioms
'TimeLocus.test_same_locus' does not depend on any axioms
'TimeLocus.test_fold_inv' does not depend on any axioms
'TimeLocus.test_not_bound' does not depend on any axioms
'TimeLocus.witness_proves_bound' does not depend on any axioms
'DotRH.dot_outside_economy' does not depend on any axioms
\end{Verbatim}

\endgroup

\hypertarget{appendix-h-the-kernel-files-of-part-iii}{%
\section{Appendix H: The Kernel Files of Part
III}\label{appendix-h-the-kernel-files-of-part-iii}}

Nine files, each compiled under core Lean 4.19.0 at exit 0 with zero
warnings, no \texttt{sorry}, and no axiom declared: forty-five
dependency sets, forty-two axiom-free and three on \texttt{propext}
alone, none on \texttt{Classical.choice}; the tenth file of Part III,
the formal block with its nine sets, is Appendix I. Then the Turing run
on \(\zeta\) to height 100, in floating point, testimony and not a
certificate.

\hypertarget{rh_gold_and_lead.lean}{%
\subsection{RH\_Gold\_And\_Lead.lean}\label{rh_gold_and_lead.lean}}

SHA-256
\texttt{\seqsplit{f5a399cd986dd472d36dab7b802e97260d0a93c66d73c336f72a905fc12253de}},
70 lines.

\begin{Verbatim}[numbers=left,numbersep=6pt,xleftmargin=2.2em,fontsize=\scriptsize,breaklines,breakanywhere,frame=lines,framesep=1.5mm,rulecolor=\color{black!35}]
/-
RH_Gold_And_Lead.lean · why the Real part and the Unicorn part should not be mixed.
Core Lean 4.19.0, no library, no sorry. Drafted 2026-09-24.
Companion to RH_Real_Unicorn_Split.lean, where RH ⟺ Real part ∧ Unicorn part at every height.
A statement's standing is proved, open, or refuted. A conjunction stands only as high as its
weakest part. So the conventional RH, the two parts fused, inherits the Unicorn part's standing:
the gold of the Real part does not survive the alloy.
-/
namespace GoldLead

/-- The three standings of a statement, ranked. -/
inductive Standing | refuted | opn | proved
  deriving DecidableEq, Repr

def Standing.rank : Standing → Nat
  | .refuted => 0 | .opn => 1 | .proved => 2

/-- A conjunction stands at the weaker of its two parts. -/
def join (a b : Standing) : Standing := if a.rank ≤ b.rank then a else b

-- G1.
/-- The alloy never outranks either metal. -/
theorem join_below_both (a b : Standing) :
    (join a b).rank ≤ a.rank ∧ (join a b).rank ≤ b.rank := by
  cases a <;> cases b <;> decide

-- G2.
/-- THE FUSED RH. The Real part stands proved (certified to 3 × 10¹²), the Unicorn part stands
    open; fused, the conventional RH stands open. -/
def realPart : Standing := .proved
def unicornPart : Standing := .opn
def fusedRH : Standing := join realPart unicornPart

theorem fused_rh_is_open : fusedRH = .opn := by decide

-- G3.
/-- GOLD NO LONGER PURE. Fusion strictly lowers the Real part's standing: kept apart it is proved,
    carried inside the fused statement it is only open. -/
theorem fusion_lowers_the_gold : fusedRH.rank < realPart.rank := by decide

-- G4.
/-- Logically, every proof of the fused statement carries a proof of the Unicorn part, and every
    refutation of the Unicorn part refutes the fused statement, whatever the Real part holds. -/
theorem fused_carries_unicorn (Real Unicorn : Prop) :
    ((Real ∧ Unicorn) → Unicorn) ∧ (¬ Unicorn → ¬ (Real ∧ Unicorn)) :=
  ⟨fun h => h.2, fun hU h => hU h.2⟩

-- G5.
/-- Kept apart, each part keeps its own standing: the Real part is asserted exactly as proved,
    with nothing borrowed from the Unicorn part, and the Unicorn part is asserted exactly as open. -/
theorem kept_apart (Real Unicorn : Prop) (hReal : Real) :
    Real ∧ ((Real ∧ Unicorn) ↔ Unicorn) :=
  ⟨hReal, ⟨fun h => h.2, fun hU => ⟨hReal, hU⟩⟩⟩

-- G6.
/-- THE RULE, IN ONE THEOREM. Given the Real part, the fused RH is exactly the Unicorn part, and
    fusing lowers the Real part from proved to open. So state the Real part as gold, the Unicorn
    part as the open bit, and do not alloy them. -/
theorem do_not_alloy (Real Unicorn : Prop) (hReal : Real) :
    ((Real ∧ Unicorn) ↔ Unicorn) ∧ fusedRH = .opn ∧ fusedRH.rank < realPart.rank :=
  ⟨(kept_apart Real Unicorn hReal).2, fused_rh_is_open, fusion_lowers_the_gold⟩

end GoldLead

#print axioms GoldLead.join_below_both
#print axioms GoldLead.fused_rh_is_open
#print axioms GoldLead.fusion_lowers_the_gold
#print axioms GoldLead.fused_carries_unicorn
#print axioms GoldLead.kept_apart
#print axioms GoldLead.do_not_alloy
\end{Verbatim}

\begin{Verbatim}[numbers=left,numbersep=6pt,xleftmargin=2.2em,fontsize=\scriptsize,breaklines,breakanywhere,frame=lines,framesep=1.5mm,rulecolor=\color{black!35}]
'GoldLead.join_below_both' does not depend on any axioms
'GoldLead.fused_rh_is_open' does not depend on any axioms
'GoldLead.fusion_lowers_the_gold' does not depend on any axioms
'GoldLead.fused_carries_unicorn' does not depend on any axioms
'GoldLead.kept_apart' does not depend on any axioms
'GoldLead.do_not_alloy' does not depend on any axioms
\end{Verbatim}

\hypertarget{rh_unicorn_isomorphism.lean}{%
\subsection{RH\_Unicorn\_Isomorphism.lean}\label{rh_unicorn_isomorphism.lean}}

SHA-256
\texttt{\seqsplit{0c3796bb1f36a429f85824edee71c1e177280523650db313835c0b79672b0ef3}},
48 lines.

\begin{Verbatim}[numbers=left,numbersep=6pt,xleftmargin=2.2em,fontsize=\scriptsize,breaklines,breakanywhere,frame=lines,framesep=1.5mm,rulecolor=\color{black!35}]
/-
RH_Unicorn_Isomorphism.lean · the Riemann Hypothesis and the unicorn condition are one statement.
Core Lean 4.19.0, no library, no sorry. Harvested 2026-09-24.
Scope: an equivalence (RH ⟺ the off-line zeros are unicorns), not a derivation of RH.
"Every white unicorn has its white fur attached only to its body" is true because there are no
unicorns: a universal over an empty class holds of anything. Here the class is the off-line zeros.
-/
namespace Unicorn

variable {α : Type} (Z onL : α → Prop)

/-- The off-line zeros: zeros not on the line. -/
def OffLine (z : α) : Prop := Z z ∧ ¬ onL z

/-- The hypothesis: every zero is on the line. -/
def RH : Prop := ∀ z, Z z → onL z

/-- The off-line zeros are unicorns: EVERY property holds of all of them, vacuously. -/
def Unicorns : Prop := ∀ (Q : α → Prop) (z : α), OffLine Z onL z → Q z

-- U1.
/-- If the hypothesis holds, the off-line zeros are unicorns: anything at all is true of them. -/
theorem rh_makes_unicorns : RH Z onL → Unicorns Z onL :=
  fun h _ z ⟨hz, off⟩ => absurd (h z hz) off

-- U2.
/-- If the off-line zeros are unicorns, the hypothesis holds (with a decidable line predicate):
    take the property "false"; it holds of every off-line zero, so none exists. -/
theorem unicorns_make_rh [∀ z, Decidable (onL z)] : Unicorns Z onL → RH Z onL :=
  fun h z hz => Decidable.byContradiction (fun off => h (fun _ => False) z ⟨hz, off⟩)

-- U3.
/-- THE SEAL. The Riemann Hypothesis is exactly the statement that the off-line zeros are
    unicorns: a class of which every sentence is vacuously true, because it has no member. -/
theorem rh_iff_unicorns [∀ z, Decidable (onL z)] : RH Z onL ↔ Unicorns Z onL :=
  ⟨rh_makes_unicorns Z onL, unicorns_make_rh Z onL⟩

-- U4.
/-- And the unicorn class is empty exactly when no off-line zero can be exhibited. -/
theorem unicorns_iff_no_witness : Unicorns Z onL ↔ ¬ ∃ z, OffLine Z onL z :=
  ⟨fun h ⟨z, hz⟩ => h (fun _ => False) z hz, fun h _ z hz => absurd ⟨z, hz⟩ h⟩

end Unicorn

#print axioms Unicorn.rh_makes_unicorns
#print axioms Unicorn.unicorns_make_rh
#print axioms Unicorn.rh_iff_unicorns
#print axioms Unicorn.unicorns_iff_no_witness
\end{Verbatim}

\begin{Verbatim}[numbers=left,numbersep=6pt,xleftmargin=2.2em,fontsize=\scriptsize,breaklines,breakanywhere,frame=lines,framesep=1.5mm,rulecolor=\color{black!35}]
'Unicorn.rh_makes_unicorns' does not depend on any axioms
'Unicorn.unicorns_make_rh' does not depend on any axioms
'Unicorn.rh_iff_unicorns' does not depend on any axioms
'Unicorn.unicorns_iff_no_witness' does not depend on any axioms
\end{Verbatim}

\hypertarget{rh_real_unicorn_split.lean}{%
\subsection{RH\_Real\_Unicorn\_Split.lean}\label{rh_real_unicorn_split.lean}}

SHA-256
\texttt{\seqsplit{6979d9a5aaac8ef2651e4984d11cf1b03542e85993539dac513541adee285235}},
72 lines.

\begin{Verbatim}[numbers=left,numbersep=6pt,xleftmargin=2.2em,fontsize=\scriptsize,breaklines,breakanywhere,frame=lines,framesep=1.5mm,rulecolor=\color{black!35}]
/-
RH_Real_Unicorn_Split.lean · the Riemann Hypothesis divided into a Real part and a Unicorn part.
Core Lean 4.19.0, no library, no sorry. Drafted 2026-09-24.
Split at a height T. The Real part is the witnessed region: every zero at height ≤ T lies on the
line. The Unicorn part is the unwitnessed region: above T, the off-line zeros are unicorns, a
class of which every sentence is vacuously true because it has no member.
-/
namespace Split

variable {α : Type} (Z onL : α → Prop) (height : α → Int)

def RH : Prop := ∀ z, Z z → onL z

/-- THE REAL PART at height T: every zero up to T is on the line. -/
def RealPart (T : Int) : Prop := ∀ z, Z z → height z ≤ T → onL z

/-- THE UNICORN PART at height T: above T, every property holds of every off-line zero, because
    there are none. -/
def UnicornPart (T : Int) : Prop :=
  ∀ (Q : α → Prop) (z : α), T < height z → Z z → ¬ onL z → Q z

-- S1.
/-- THE SPLIT. At every height T, the hypothesis is exactly the Real part and the Unicorn part
    together. Nothing is lost and nothing is added. -/
theorem rh_iff_real_and_unicorn [∀ z, Decidable (onL z)] (T : Int) :
    RH Z onL ↔ RealPart Z onL height T ∧ UnicornPart Z onL height T := by
  constructor
  · intro h
    exact ⟨fun z hz _ => h z hz, fun _ z _ hz off => absurd (h z hz) off⟩
  · intro ⟨hR, hU⟩ z hz
    by_cases hle : height z ≤ T
    · exact hR z hz hle
    · exact Decidable.byContradiction
        (fun off => hU (fun _ => False) z (Int.lt_of_not_ge hle) hz off)

-- S2.
/-- The parts are separate: the Real part can hold while the Unicorn part fails. An off-line
    zero above T leaves every zero up to T on the line. -/
def twoZeros (z : Int × Bool) : Prop := z = (5, true) ∨ z = (20, false)
def onLine2 (z : Int × Bool) : Prop := z.2 = true
def height2 (z : Int × Bool) : Int := z.1

theorem parts_are_separate :
    RealPart twoZeros onLine2 height2 10 ∧ ¬ UnicornPart twoZeros onLine2 height2 10 := by
  refine ⟨?_, ?_⟩
  · intro z hz hle
    rcases hz with h | h <;> subst h
    · rfl
    · exact absurd hle (by decide)
  · intro hU
    exact hU (fun _ => False) (20, false) (by decide) (Or.inr rfl) (fun h => Bool.noConfusion h)

-- S3.
/-- The Unicorn part is empty exactly when no off-line zero above T can be exhibited. -/
theorem unicorn_part_iff_no_witness (T : Int) :
    UnicornPart Z onL height T ↔ ¬ ∃ z, T < height z ∧ Z z ∧ ¬ onL z :=
  ⟨fun hU ⟨z, hT, hz, off⟩ => hU (fun _ => False) z hT hz off,
   fun h _ z hT hz off => absurd ⟨z, hT, hz, off⟩ h⟩

-- S4.
/-- Raising T moves zeros from the Unicorn part into the Real part and never the other way: a
    Real part at a higher height contains the Real part at a lower one. -/
theorem real_part_monotone (T T' : Int) (hTT : T ≤ T') :
    RealPart Z onL height T' → RealPart Z onL height T :=
  fun h z hz hle => h z hz (Int.le_trans hle hTT)

end Split

#print axioms Split.rh_iff_real_and_unicorn
#print axioms Split.parts_are_separate
#print axioms Split.unicorn_part_iff_no_witness
#print axioms Split.real_part_monotone
\end{Verbatim}

\begin{Verbatim}[numbers=left,numbersep=6pt,xleftmargin=2.2em,fontsize=\scriptsize,breaklines,breakanywhere,frame=lines,framesep=1.5mm,rulecolor=\color{black!35}]
'Split.rh_iff_real_and_unicorn' depends on axioms: [propext]
'Split.parts_are_separate' does not depend on any axioms
'Split.unicorn_part_iff_no_witness' does not depend on any axioms
'Split.real_part_monotone' depends on axioms: [propext]
\end{Verbatim}

\hypertarget{ram_zfc_placement.lean}{%
\subsection{RAM\_ZFC\_Placement.lean}\label{ram_zfc_placement.lean}}

SHA-256
\texttt{\seqsplit{d7b0d3da00ca3f6e1340440a775beeb539599b66a9f189deecade5309fb3086b}},
85 lines.

\begin{Verbatim}[numbers=left,numbersep=6pt,xleftmargin=2.2em,fontsize=\scriptsize,breaklines,breakanywhere,frame=lines,framesep=1.5mm,rulecolor=\color{black!35}]
/-
RAM_ZFC_Placement.lean · ZFC subsumed by RAM, and ZFC's location under RAM.
Core Lean 4.19.0, no library, no sorry, no axiom declaration. Harvested 2026-09-24.

Location. ZFC is a ladder at L2m (provability); its intended world, the cumulative hierarchy and
the numbers inside it, is the Ground at L1m (formal being as grounding); its computable part is
the rung at L3m. The order rung ⊆ ladder ⊆ Ground holds under two named premises: the rung on
the ladder (for ZFC, Σ₁-completeness, a theorem of arithmetic) and soundness (which ZFC cannot
prove of itself, Gödel II, so it is carried openly).

Subsumption is by placement: RAM locates ZFC, as it locates NBG, PA, and every other ladder; it
does not derive ZFC's axioms and does not replace ZFC's proofs. Any theory, ZFC included, is
modelled abstractly: sentences, a provability predicate, a decidable fragment, and truth in the
intended world.
-/
namespace RAMZFC

structure Theory where
  Sent  : Type
  Prov  : Sent → Prop        -- L2m, the ladder
  Comp  : Sent → Prop        -- L3m, the rung: sentences settled by finite computation
  True_ : Sent → Prop        -- L1m, the Ground: truth in the intended world

/-- The rung is on the ladder: what finite computation settles, the theory proves.
    (For ZFC: Σ₁-completeness, a theorem of arithmetic, carried as a named premise.) -/
def RungOnLadder (T : Theory) : Prop := ∀ s, T.Comp s → T.Prov s

/-- The ladder reaches only the Ground: soundness. For ZFC this is exactly what ZFC cannot
    prove of itself (Gödel II); it is carried openly, never derived. -/
def Sound (T : Theory) : Prop := ∀ s, T.Prov s → T.True_ s

-- Z1.
/-- THE PLACEMENT. Under the two named premises, every theory sits in RAM's strata in order:
    rung ⊆ ladder ⊆ Ground. -/
theorem ram_places_theory (T : Theory) (hR : RungOnLadder T) (hS : Sound T) :
    (∀ s, T.Comp s → T.Prov s) ∧ (∀ s, T.Prov s → T.True_ s) ∧ (∀ s, T.Comp s → T.True_ s) :=
  ⟨hR, hS, fun s h => hS s (hR s h)⟩

-- Z2.
/-- The placement is not free: an unsound theory breaks ladder ⊆ Ground. -/
def unsound : Theory := ⟨Bool, fun _ => True, fun _ => False, fun b => b = true⟩

theorem soundness_is_load_bearing :
    RungOnLadder unsound ∧ ¬ Sound unsound := by
  refine ⟨fun _ h => h.elim, fun h => ?_⟩
  have := h false trivial
  cases this

-- Z3.
/-- The Ground can exceed the ladder: a sound theory can leave a truth unproved (Gödel's
    incompleteness region, L1m minus L2m), so the strata are genuinely three, not one. -/
def incomplete : Theory := ⟨Bool, fun b => b = true, fun _ => False, fun _ => True⟩

theorem ground_exceeds_ladder :
    Sound incomplete ∧ incomplete.True_ false ∧ ¬ incomplete.Prov false :=
  ⟨fun _ _ => trivial, trivial, fun h => Bool.noConfusion h⟩

-- Z4.
/-- Axioms as bits. Two independent axioms are two bits, not one composite paradox: all four
    combinations of their truth values are realized, so a theory with several independent
    axioms is a product of bits. (For ZFC: Choice is independent of ZF, Gödel 1938 and Cohen
    1963; this is cited, not proved here.) -/
theorem independent_axioms_are_separate_bits :
    ∀ a b : Bool, ∃ w : Bool × Bool, w.1 = a ∧ w.2 = b :=
  fun a b => ⟨(a, b), rfl, rfl⟩

-- Z5.
/-- THE VERDICT. RAM places every theory, ZFC included, as rung ⊆ ladder ⊆ Ground, conditional
    on the two named premises; soundness cannot be dropped; the Ground genuinely exceeds the
    ladder. Placement is structural subsumption: RAM does not derive ZFC's axioms, and does
    not replace ZFC's proofs. -/
theorem ram_subsumes_by_placement :
    (∀ T : Theory, RungOnLadder T → Sound T →
        (∀ s, T.Comp s → T.True_ s)) ∧
    (RungOnLadder unsound ∧ ¬ Sound unsound) ∧
    (Sound incomplete ∧ incomplete.True_ false ∧ ¬ incomplete.Prov false) :=
  ⟨fun T hR hS => (ram_places_theory T hR hS).2.2, soundness_is_load_bearing, ground_exceeds_ladder⟩

end RAMZFC

#print axioms RAMZFC.ram_places_theory
#print axioms RAMZFC.soundness_is_load_bearing
#print axioms RAMZFC.ground_exceeds_ladder
#print axioms RAMZFC.independent_axioms_are_separate_bits
#print axioms RAMZFC.ram_subsumes_by_placement
\end{Verbatim}

\begin{Verbatim}[numbers=left,numbersep=6pt,xleftmargin=2.2em,fontsize=\scriptsize,breaklines,breakanywhere,frame=lines,framesep=1.5mm,rulecolor=\color{black!35}]
'RAMZFC.ram_places_theory' does not depend on any axioms
'RAMZFC.soundness_is_load_bearing' does not depend on any axioms
'RAMZFC.ground_exceeds_ladder' does not depend on any axioms
'RAMZFC.independent_axioms_are_separate_bits' does not depend on any axioms
'RAMZFC.ram_subsumes_by_placement' does not depend on any axioms
\end{Verbatim}

\hypertarget{ra_to_zfc_chain.lean}{%
\subsection{RA\_to\_ZFC\_Chain.lean}\label{ra_to_zfc_chain.lean}}

SHA-256
\texttt{\seqsplit{5714fbd5dc59b7ec79e251daca934819f5663e5554bad94d5537ed9f6c618a3e}},
90 lines.

\begin{Verbatim}[numbers=left,numbersep=6pt,xleftmargin=2.2em,fontsize=\scriptsize,breaklines,breakanywhere,frame=lines,framesep=1.5mm,rulecolor=\color{black!35}]
/-
RA_to_ZFC_Chain.lean · RA mapped to ZFC through the RA–RAM bridge, link by link.
Core Lean 4.19.0, no library, no sorry, no axiom declaration. Harvested 2026-09-24.

The chain RA → Bridge → RAM → ZFC. The RA–RAM bridge is the junction's discriminator
(Master Codex Ch. 10B, `Bridge.discriminator`): a property whose every world has it is KEYLESS
and crosses from the root; a property some world lacks is KEYED, one bit, supplied and never
derived. ZFC's location under RAM is RAM_ZFC_Placement.lean. Worlds are ways a theory's content
could be, which of its independent sentences hold.

Verdict. Upward, every act of using ZFC grounds RA. Downward, RA crosses the bridge as presence
into every world of ZFC and carries no keyed content: no independent sentence of ZFC, Choice
over ZF the cited example (Gödel 1938, Cohen 1963), is decided by RA.
-/
namespace RAtoZFC

structure Substrate where
  U  : Type
  ΔE : U → Int

def RA (S : Substrate) : Prop := ∀ x : S.U, 0 < S.ΔE x

structure SelfGrounding (R : Prop) where
  Act       : Type
  anAct     : Act
  instances : Act → R

def unitSub : Substrate := ⟨Unit, fun _ => 1⟩
theorem ra_unit : RA unitSub := fun _ => show (0 : Int) < 1 by decide

/-- A theory, abstractly: sentences and a provability predicate (the RAM ladder). -/
structure Theory where
  Sent : Type
  Prov : Sent → Prop

-- LINK 1 · ZFC → RA: every act of using a theory grounds RA. This link holds.
theorem using_grounds_root (T : Theory) (s : T.Sent) (d : T.Prov s) :
    Nonempty (SelfGrounding (RA unitSub)) :=
  ⟨⟨PLift (T.Prov s), ⟨d⟩, fun _ => ra_unit⟩⟩

-- LINK 2 · RA → Bridge: the bridge carries only KEYLESS properties, those no world denies.
/-- A property of worlds is keyless when every world has it, keyed when some world lacks it. -/
def Keyless {W : Type} (P : W → Prop) : Prop := ∀ w, P w
def Keyed   {W : Type} (P : W → Prop) : Prop := ∃ w, ¬ P w

/-- RA crosses to every world whatever its content: RA is keyless over worlds. -/
theorem ra_is_keyless {W : Type} : Keyless (fun _ : W => RA unitSub) := fun _ => ra_unit

/-- A keyless property is carried by RA to every world. -/
theorem keyless_crosses {W : Type} (P : W → Prop) (hP : Keyless P) :
    ∀ w, RA unitSub → P w := fun w _ => hP w

-- LINK 3 · Bridge → RAM → ZFC: a keyed sentence, one the theory's worlds disagree on (for ZFC,
-- Choice over ZF: Gödel 1938, Cohen 1963, cited), is not carried. RA holds in both worlds.
/-- Two worlds of ZF: in one Choice holds, in the other it fails. -/
def choiceHolds : Bool → Prop := fun w => w = true

theorem choice_is_keyed : Keyed choiceHolds := ⟨false, fun h => Bool.noConfusion h⟩

/-- THE BREAK. RA holds in both worlds, so RA cannot carry Choice into ZFC, nor its denial. -/
theorem ra_does_not_cross_keyed :
    (∀ _w : Bool, RA unitSub) ∧ choiceHolds true ∧ ¬ choiceHolds false :=
  ⟨fun _ => ra_unit, rfl, fun h => Bool.noConfusion h⟩

/-- In general: RA decides no keyed property of worlds. -/
theorem ra_decides_no_keyed {W : Type} (P : W → Prop) (hK : Keyed P) :
    ¬ (∀ w, RA unitSub → P w) :=
  fun h => let ⟨w, hw⟩ := hK; hw (h w ra_unit)

-- THE CHAIN, IN ONE THEOREM.
/-- Upward, every act of using ZFC grounds RA. Downward, RA carries exactly the keyless and
    none of the keyed: it reaches every world of ZFC as presence, and no independent sentence of
    ZFC as content. -/
theorem chain_verdict :
    (∀ (T : Theory) (s : T.Sent), T.Prov s → Nonempty (SelfGrounding (RA unitSub))) ∧
    (∀ {W : Type} (P : W → Prop), Keyless P → ∀ w, RA unitSub → P w) ∧
    (∀ {W : Type} (P : W → Prop), Keyed P → ¬ ∀ w, RA unitSub → P w) ∧
    ((∀ _w : Bool, RA unitSub) ∧ choiceHolds true ∧ ¬ choiceHolds false) :=
  ⟨using_grounds_root, fun P hP => keyless_crosses P hP, fun P hK => ra_decides_no_keyed P hK,
   ra_does_not_cross_keyed⟩

end RAtoZFC

#print axioms RAtoZFC.using_grounds_root
#print axioms RAtoZFC.ra_is_keyless
#print axioms RAtoZFC.keyless_crosses
#print axioms RAtoZFC.choice_is_keyed
#print axioms RAtoZFC.ra_does_not_cross_keyed
#print axioms RAtoZFC.ra_decides_no_keyed
#print axioms RAtoZFC.chain_verdict
\end{Verbatim}

\begin{Verbatim}[numbers=left,numbersep=6pt,xleftmargin=2.2em,fontsize=\scriptsize,breaklines,breakanywhere,frame=lines,framesep=1.5mm,rulecolor=\color{black!35}]
'RAtoZFC.using_grounds_root' does not depend on any axioms
'RAtoZFC.ra_is_keyless' does not depend on any axioms
'RAtoZFC.keyless_crosses' does not depend on any axioms
'RAtoZFC.choice_is_keyed' does not depend on any axioms
'RAtoZFC.ra_does_not_cross_keyed' does not depend on any axioms
'RAtoZFC.ra_decides_no_keyed' does not depend on any axioms
'RAtoZFC.chain_verdict' does not depend on any axioms
\end{Verbatim}

\hypertarget{zfc_ra_roundtrip.lean}{%
\subsection{ZFC\_RA\_RoundTrip.lean}\label{zfc_ra_roundtrip.lean}}

SHA-256
\texttt{\seqsplit{72a7571509a019e4b5af3f3e7324c698ffa66e64a69c6332f91be8080059310a}},
64 lines.

\begin{Verbatim}[numbers=left,numbersep=6pt,xleftmargin=2.2em,fontsize=\scriptsize,breaklines,breakanywhere,frame=lines,framesep=1.5mm,rulecolor=\color{black!35}]
/-
ZFC_RA_RoundTrip.lean · the round trip ZFC → RAM → Bridge → RA → Bridge → RAM → ZFC.
Core Lean 4.19.0, no library, no sorry, no axiom declaration. Drafted 2026-09-24.
A sentence of ZFC is read in its worlds (ways ZFC's content could be). Up the chain it is placed
on RAM's Ground; across the RA–RAM bridge it meets RA; the arrow back returns through the bridge,
which carries keyless content, to RAM and then to ZFC.
-/
namespace RoundTrip

structure Substrate where
  U  : Type
  ΔE : U → Int

def RA (S : Substrate) : Prop := ∀ x : S.U, 0 < S.ΔE x
def unitSub : Substrate := ⟨Unit, fun _ => 1⟩
theorem ra_unit : RA unitSub := fun _ => show (0 : Int) < 1 by decide

/-- Up: a sentence's content, read on RAM's Ground in each world. -/
def up {W : Type} (P : W → Prop) : W → Prop := P

/-- Across and back: what the bridge returns to RAM from RA, in each world: RA together with the
    content it received. -/
def viaRA {W : Type} (P : W → Prop) : W → Prop := fun w => RA unitSub ∧ up P w

/-- Down: from RAM back to the sentence of ZFC. -/
def down {W : Type} (Q : W → Prop) : W → Prop := Q

-- R1.
/-- THE ROUND TRIP IS THE IDENTITY ON CONTENT. In every world, the sentence that returns to ZFC
    is equivalent to the sentence that left. -/
theorem round_trip_identity {W : Type} (P : W → Prop) (w : W) :
    down (viaRA P) w ↔ P w :=
  ⟨fun h => h.2, fun h => ⟨ra_unit, h⟩⟩

-- R2.
/-- An absolute sentence, one with the same truth value in every world (an arithmetic sentence
    such as RH, whose value is fixed in every model with the standard numbers): the round trip
    returns that same value, true if it left true and false if it left false. -/
theorem absolute_returns_unchanged {W : Type} (P : W → Prop) (hAbs : ∀ w v, P w ↔ P v)
    (w : W) : (down (viaRA P) w ↔ P w) ∧ (∀ v, down (viaRA P) v ↔ P w) :=
  ⟨round_trip_identity P w, fun v => (round_trip_identity P v).trans (hAbs v w)⟩

-- R3.
/-- What the trip adds: presence. RA holds in every world on return; it holds identically for a
    sentence and for its negation, so it tells them apart nowhere. -/
theorem trip_adds_presence_only {W : Type} (P : W → Prop) :
    (∀ _w : W, RA unitSub) ∧ (∀ w, (down (viaRA P) w ↔ P w) ∧ (down (viaRA (fun v => ¬ P v)) w ↔ ¬ P w)) :=
  ⟨fun _ => ra_unit, fun w => ⟨round_trip_identity P w, round_trip_identity (fun v => ¬ P v) w⟩⟩

-- R4.
/-- THE VERDICT. Whatever ZFC sends up returns to ZFC unchanged in content, with RA's presence
    attached in every world; a sentence ZFC does not settle leaves undecided and returns
    undecided. -/
theorem round_trip_verdict :
    (∀ {W : Type} (P : W → Prop) (w : W), down (viaRA P) w ↔ P w) ∧
    (∀ {W : Type} (P : W → Prop), (∀ w v, P w ↔ P v) → ∀ w v, down (viaRA P) v ↔ P w) :=
  ⟨fun P w => round_trip_identity P w, fun P h w v => (absolute_returns_unchanged P h w).2 v⟩

end RoundTrip

#print axioms RoundTrip.round_trip_identity
#print axioms RoundTrip.absolute_returns_unchanged
#print axioms RoundTrip.trip_adds_presence_only
#print axioms RoundTrip.round_trip_verdict
\end{Verbatim}

\begin{Verbatim}[numbers=left,numbersep=6pt,xleftmargin=2.2em,fontsize=\scriptsize,breaklines,breakanywhere,frame=lines,framesep=1.5mm,rulecolor=\color{black!35}]
'RoundTrip.round_trip_identity' does not depend on any axioms
'RoundTrip.absolute_returns_unchanged' does not depend on any axioms
'RoundTrip.trip_adds_presence_only' does not depend on any axioms
'RoundTrip.round_trip_verdict' does not depend on any axioms
\end{Verbatim}

\hypertarget{massless_arrow_test.lean}{%
\subsection{Massless\_Arrow\_Test.lean}\label{massless_arrow_test.lean}}

SHA-256
\texttt{\seqsplit{7822a00ba1a2540891cdfb5888aac7d533d3e4f84c623c08b879c0751d17e5ee}},
42 lines.

\begin{Verbatim}[numbers=left,numbersep=6pt,xleftmargin=2.2em,fontsize=\scriptsize,breaklines,breakanywhere,frame=lines,framesep=1.5mm,rulecolor=\color{black!35}]
/-
Massless_Arrow_Test.lean · can a zero-content arrow from RA deliver a proposition?
Core Lean 4.19.0, no library, no sorry. Drafted 2026-09-24.
An arrow sends a proposition to what it delivers. It is massless (zero content) when it delivers
the same for every proposition, and sound when what it delivers for P implies P.
-/
namespace Massless

def Massless (f : Prop → Prop) : Prop := ∀ P Q, f P ↔ f Q
def Sound (f : Prop → Prop) : Prop := ∀ P, f P → P

-- M1.
/-- A sound massless arrow delivers nothing: if it delivered P, it would deliver ¬P too, and
    both would hold. So its delivery is empty for every proposition. -/
theorem massless_sound_delivers_nothing (f : Prop → Prop) (hM : Massless f) (hS : Sound f) :
    ∀ P, ¬ f P :=
  fun P h => hS (¬ P) ((hM P (¬ P)).mp h) (hS P h)

-- M2.
/-- The RA-shaped arrow: it delivers RA's presence, the same for every proposition. -/
def raArrow (RA : Prop) : Prop → Prop := fun _ => RA

theorem raArrow_massless (RA : Prop) : Massless (raArrow RA) := fun _ _ => Iff.rfl

/-- If RA holds, the RA-shaped arrow is NOT sound: it delivers every proposition, true or false. -/
theorem raArrow_not_sound (RA : Prop) (hRA : RA) : ¬ Sound (raArrow RA) :=
  fun hS => hS False hRA

-- M3.
/-- THE VERDICT. A zero-content arrow either delivers nothing (if it is sound) or delivers every
    proposition and its negation alike (if it is not). Either way it cannot deliver RH rather than
    ¬RH. Content is what separates a proposition from its negation, and a massless arrow has none. -/
theorem massless_cannot_separate (f : Prop → Prop) (hM : Massless f) (RH : Prop) :
    (f RH ↔ f (¬ RH)) :=
  hM RH (¬ RH)

end Massless

#print axioms Massless.massless_sound_delivers_nothing
#print axioms Massless.raArrow_massless
#print axioms Massless.raArrow_not_sound
#print axioms Massless.massless_cannot_separate
\end{Verbatim}

\begin{Verbatim}[numbers=left,numbersep=6pt,xleftmargin=2.2em,fontsize=\scriptsize,breaklines,breakanywhere,frame=lines,framesep=1.5mm,rulecolor=\color{black!35}]
'Massless.massless_sound_delivers_nothing' does not depend on any axioms
'Massless.raArrow_massless' does not depend on any axioms
'Massless.raArrow_not_sound' does not depend on any axioms
'Massless.massless_cannot_separate' does not depend on any axioms
\end{Verbatim}

\hypertarget{arrow_translation.lean}{%
\subsection{Arrow\_Translation.lean}\label{arrow_translation.lean}}

SHA-256
\texttt{\seqsplit{f57924e397b9ec7e8b9f654547ff96c3f69d7aa7e41b4e32ee88afcb0a972ea9}},
74 lines.

\begin{Verbatim}[numbers=left,numbersep=6pt,xleftmargin=2.2em,fontsize=\scriptsize,breaklines,breakanywhere,frame=lines,framesep=1.5mm,rulecolor=\color{black!35}]
/-
Arrow_Translation.lean · Part Two translated into the purely formal domain: the Root Axiom
replaced by a bare computation arrow. Core Lean 4.19.0, no library, no sorry, no axiom declared.
The arrow is the identity computation on any type, `id : α → α`; its existence needs nothing.
Every place Part Two used the Root Axiom, it is replaced by the proposition that an arrow exists.
-/
namespace ArrowForm

/-- THE ARROW: a computation from a type to itself. It exists on every type, with no premise. -/
def Arrow (α : Type) : Prop := Nonempty (α → α)

theorem arrow_exists (α : Type) : Arrow α := ⟨id⟩

/-- The arrow used throughout, on the unit type. -/
def A : Prop := Arrow Unit
theorem a_holds : A := arrow_exists Unit

-- T1 · the upward link. Every act of proof yields an arrow: the proof itself, carried to itself.
theorem proof_yields_arrow (P : Prop) (h : P) : Arrow (PLift P) :=
  ⟨fun _ => ⟨h⟩⟩

-- T2 · keyless crossing. The arrow holds in every world and carries every keyless property.
def Keyless {W : Type} (P : W → Prop) : Prop := ∀ w, P w
def Keyed   {W : Type} (P : W → Prop) : Prop := ∃ w, ¬ P w

theorem arrow_is_keyless {W : Type} : Keyless (fun _ : W => A) := fun _ => a_holds
theorem keyless_crosses {W : Type} (P : W → Prop) (hP : Keyless P) : ∀ w, A → P w :=
  fun w _ => hP w

-- T3 · the break. The arrow decides no keyed property.
theorem arrow_decides_no_keyed {W : Type} (P : W → Prop) (hK : Keyed P) : ¬ ∀ w, A → P w :=
  fun h => let ⟨w, hw⟩ := hK; hw (h w a_holds)

-- T4 · the round trip. Carried through the arrow and back, every sentence returns unchanged.
def viaArrow {W : Type} (P : W → Prop) : W → Prop := fun w => A ∧ P w
theorem round_trip_identity {W : Type} (P : W → Prop) (w : W) : viaArrow P w ↔ P w :=
  ⟨fun h => h.2, fun h => ⟨a_holds, h⟩⟩

-- T5 · the massless arrow. An arrow that delivers the same for every proposition cannot
-- separate a proposition from its negation.
def Massless (f : Prop → Prop) : Prop := ∀ P Q, f P ↔ f Q
theorem arrowDelivery_massless : Massless (fun _ => A) := fun _ _ => Iff.rfl
theorem massless_cannot_separate (f : Prop → Prop) (hM : Massless f) (RH : Prop) :
    f RH ↔ f (¬ RH) := hM RH (¬ RH)

-- T6 · the Real part, closed by a certificate, with no premise at all.
structure Certificate {α : Type} (Z onL : α → Prop) (height : α → Int) (T : Int) where
  zeros    : List α
  complete : ∀ z, Z z → height z ≤ T → z ∈ zeros
  onLine   : ∀ z, z ∈ zeros → onL z

theorem arrow_real_part_of_certificate {α : Type} (Z onL : α → Prop) (height : α → Int) (T : Int)
    (c : Certificate Z onL height T) : ∀ z, Z z → height z ≤ T → onL z :=
  fun z hz hle => c.onLine z (c.complete z hz hle)

-- T7 · THE TRANSLATION IS FAITHFUL. Any inhabited proposition plays the root's part in every
-- result above: the road's theorems hold for an arbitrary true premise R, so the root is idle in
-- the formal content, and the bare arrow is enough.
theorem any_true_premise_serves (R : Prop) (hR : R) {W : Type} (P : W → Prop) :
    ((∀ w, P w) → ∀ w, R → P w) ∧
    ((∃ w, ¬ P w) → ¬ ∀ w, R → P w) ∧
    (∀ w, (R ∧ P w) ↔ P w) :=
  ⟨fun hP w _ => hP w, fun ⟨w, hw⟩ h => hw (h w hR), fun _ => ⟨fun h => h.2, fun h => ⟨hR, h⟩⟩⟩

end ArrowForm

#print axioms ArrowForm.arrow_exists
#print axioms ArrowForm.proof_yields_arrow
#print axioms ArrowForm.keyless_crosses
#print axioms ArrowForm.arrow_decides_no_keyed
#print axioms ArrowForm.round_trip_identity
#print axioms ArrowForm.massless_cannot_separate
#print axioms ArrowForm.arrow_real_part_of_certificate
#print axioms ArrowForm.any_true_premise_serves
\end{Verbatim}

\begin{Verbatim}[numbers=left,numbersep=6pt,xleftmargin=2.2em,fontsize=\scriptsize,breaklines,breakanywhere,frame=lines,framesep=1.5mm,rulecolor=\color{black!35}]
'ArrowForm.arrow_exists' does not depend on any axioms
'ArrowForm.proof_yields_arrow' does not depend on any axioms
'ArrowForm.keyless_crosses' does not depend on any axioms
'ArrowForm.arrow_decides_no_keyed' does not depend on any axioms
'ArrowForm.round_trip_identity' does not depend on any axioms
'ArrowForm.massless_cannot_separate' does not depend on any axioms
'ArrowForm.arrow_real_part_of_certificate' does not depend on any axioms
'ArrowForm.any_true_premise_serves' does not depend on any axioms
\end{Verbatim}

\hypertarget{rh_real_part_computed.lean}{%
\subsection{RH\_Real\_Part\_Computed.lean}\label{rh_real_part_computed.lean}}

SHA-256
\texttt{\seqsplit{a5c432c7eaf5f68ed1117368e71450cfad5a52b4d1cf331d7db9a98e21721645}},
54 lines.

\begin{Verbatim}[numbers=left,numbersep=6pt,xleftmargin=2.2em,fontsize=\scriptsize,breaklines,breakanywhere,frame=lines,framesep=1.5mm,rulecolor=\color{black!35}]
/-
RH_Real_Part_Computed.lean · the Real part of RH, closed by computation and carried up the ladder.
Core Lean 4.19.0, no library, no sorry. Drafted 2026-09-24.
The Real part at height T (RH_Real_Unicorn_Split.lean) is a bounded statement: finitely many
zeros lie below T. It is closed by a finite certificate, a complete list of the zeros below T and
a check that each is on the line. That is the rung (L3m); the rung is on the ladder (L2m), so a
certified Real part is provable; and the round trip returns it unchanged.
-/
namespace RealPart

variable {α : Type} (Z onL : α → Prop) (height : α → Int)

def RealPartAt (T : Int) : Prop := ∀ z, Z z → height z ≤ T → onL z

/-- A certificate for the Real part at T: a list of all zeros up to T, and each is on the line. -/
structure Certificate (T : Int) where
  zeros    : List α
  complete : ∀ z, Z z → height z ≤ T → z ∈ zeros
  onLine   : ∀ z, z ∈ zeros → onL z

-- C1.
/-- THE REAL PART, CLOSED. A certificate proves the Real part at T completely. -/
theorem real_part_of_certificate (T : Int) (c : Certificate Z onL height T) :
    RealPartAt Z onL height T :=
  fun z hz hle => c.onLine z (c.complete z hz hle)

-- C2.
/-- An executed instance: a zero set whose members up to T are listed and checked by the kernel's
    own computation (`decide`), the Real part following from the certificate. -/
def zs : List (Int × Bool) := [(14, true), (21, true), (25, true)]
def Zs (z : Int × Bool) : Prop := z ∈ zs
def onLs (z : Int × Bool) : Prop := z.2 = true

theorem real_part_executed : RealPartAt Zs onLs (fun z => z.1) 30 :=
  real_part_of_certificate Zs onLs (fun z => z.1) 30
    ⟨zs, fun _ hz _ => hz, fun z hz => by
      simp only [zs, List.mem_cons, List.mem_nil_iff, or_false] at hz
      rcases hz with h | h | h <;> subst h <;> rfl⟩

-- C3.
/-- The rung is on the ladder: a true Real part is a bounded statement, and a theory complete for
    true bounded statements proves it (for ZFC and PA, bounded completeness, a theorem of
    arithmetic, carried as the named premise `bounded`). A certificate therefore yields a proof. -/
theorem certified_is_provable (Prov : Prop → Prop)
    (bounded : ∀ T, RealPartAt Z onL height T → Prov (RealPartAt Z onL height T))
    (T : Int) (c : Certificate Z onL height T) :
    Prov (RealPartAt Z onL height T) :=
  bounded T (real_part_of_certificate Z onL height T c)

end RealPart

#print axioms RealPart.real_part_of_certificate
#print axioms RealPart.real_part_executed
#print axioms RealPart.certified_is_provable
\end{Verbatim}

\begin{Verbatim}[numbers=left,numbersep=6pt,xleftmargin=2.2em,fontsize=\scriptsize,breaklines,breakanywhere,frame=lines,framesep=1.5mm,rulecolor=\color{black!35}]
'RealPart.real_part_of_certificate' does not depend on any axioms
'RealPart.real_part_executed' depends on axioms: [propext]
'RealPart.certified_is_provable' does not depend on any axioms
\end{Verbatim}

\hypertarget{rh_real_part_turing_check.py}{%
\subsection{RH\_Real\_Part\_Turing\_Check.py}\label{rh_real_part_turing_check.py}}

\begin{Verbatim}[numbers=left,numbersep=6pt,xleftmargin=2.2em,fontsize=\scriptsize,breaklines,breakanywhere,frame=lines,framesep=1.5mm,rulecolor=\color{black!35}]
# Real part executed to height T, three independent counts, in floating point:
#   (a) zeros ON the line, as sign changes of Hardy's Z(t) on a fine grid;
#   (b) all zeros in the strip with 0 < Im < T, N(T), by mpmath's counting function nzeros(T),
#       computed independently of the grid; the earlier form theta(T)/pi + 1 + arg zeta(1/2 + iT)/pi
#       with a principal-branch endpoint argument is not the counting function S(T) and is not used;
#   (c) the ordinates of the first zeros themselves, from mpmath.zetazero, counted below T.
# Agreement of (a), (b), (c) at T means every zero up to T is on the line, in 30-digit floating
# point without interval arithmetic: testimony, not a certificate. The certificate of record for
# zeta to 3 x 10^12 is Platt and Trudgian (2021), cited in the paper, not this script.
import mpmath as mp
mp.mp.dps = 30
T = mp.mpf(100)
def Z(t): return mp.siegelz(t)
grid = [mp.mpf(k)/20 for k in range(1, int(T*20)+1)]
on_line, prev = 0, Z(grid[0])
for t in grid[1:]:
    z = Z(t)
    if z == 0 or mp.sign(z) != mp.sign(prev): on_line += 1
    prev = z
N = int(mp.nzeros(T))
ords = []
k = 1
while True:
    g = mp.zetazero(k).imag
    if g > T: break
    ords.append(g); k += 1
print("T =", T)
print("(a) zeros ON the line, sign changes of Z on a 1/20 grid :", on_line)
print("(b) N(T), all zeros in the strip, mpmath nzeros          :", N)
print("(c) ordinates of zeros below T, enumerated               :", len(ords),
      " first", mp.nstr(ords[0], 8), " last", mp.nstr(ords[-1], 8), " next", mp.nstr(mp.zetazero(k).imag, 8))
ok = (on_line == N == len(ords))
print("Real part at T =", int(T), ":", "HOLDS (three counts equal)" if ok else "count mismatch")
\end{Verbatim}

\begin{Verbatim}[numbers=left,numbersep=6pt,xleftmargin=2.2em,fontsize=\scriptsize,breaklines,breakanywhere,frame=lines,framesep=1.5mm,rulecolor=\color{black!35}]
T = 100.0
(a) zeros ON the line, sign changes of Z on a 1/20 grid : 29
(b) N(T), all zeros in the strip, mpmath nzeros          : 29
(c) ordinates of zeros below T, enumerated               : 29  first 14.134725  last 98.831194  next 101.31785
Real part at T = 100 : HOLDS (three counts equal)
\end{Verbatim}

\hypertarget{appendix-i-rh_unicorn_block.lean-and-its-audit}{%
\section{Appendix I: RH\_Unicorn\_Block.lean and Its
Audit}\label{appendix-i-rh_unicorn_block.lean-and-its-audit}}

SHA-256
\texttt{\seqsplit{222bd4c985b99df22435e5f3b04a66171304377177f042495caa5be2e5a4921b}},
Lean 4.19.0, core only, 142 lines; nine dependency sets, every one
reported as depending on no axiom.

\begingroup\scriptsize

\begin{Verbatim}[numbers=left,numbersep=6pt,xleftmargin=2.2em,fontsize=\scriptsize,breaklines,breakanywhere,frame=lines,framesep=1.5mm,rulecolor=\color{black!35}]
/-
RH_Unicorn_Block.lean · the Unicorn part of the Riemann Hypothesis, closed to derivation, its
aperture one bit wide, its supply side silent. Core Lean 4.19.0, no library, no sorry, no user-declared axiom. 2026-09-25.

Part III divided the hypothesis at a height T into a Real part and a Unicorn part and proved
the Real part by certificate. Parts I and II proved the wall, the crossing, the bound, and the
cure theorem, and did not say what they say of the Unicorn part. This file says it: closed to
derivation, the aperture one bit wide, the register silent beyond it. ΔM = 0.
-/
namespace RHUnicorn

variable {α β : Type}

/-- The line property, the Real part at T, the Unicorn part at T, as in Part III. -/
def RH (Z onL : α → Prop) : Prop := ∀ z, Z z → onL z
def RealPartAt (Z onL : α → Prop) (height : α → Nat) (T : Nat) : Prop :=
  ∀ z, Z z → height z ≤ T → onL z
def UnicornAt (Z onL : α → Prop) (height : α → Nat) (T : Nat) : Prop :=
  ∀ z, Z z → height z > T → onL z

/-- THE FORMAL BLOCK OF THE UNICORN PART. Let τ act on the zeros, let ρ be a record even at a
    seat z above the height, and let d decide the line property, odd at z. Then no reading g of
    the record agrees with d even on the region above the height alone. The Unicorn part is
    closed to every even register; the closure is a theorem and does not expire. -/
theorem unicorn_block (height : α → Nat) (T : Nat) (τ : α → α) (ρ : α → β) (d : α → Bool)
    (z : α) (habove : height z > T) (habove' : height (τ z) > T)
    (hρ : ρ (τ z) = ρ z) (hd : d (τ z) = !d z) :
    ¬ ∃ g : β → Bool, ∀ w, height w > T → g (ρ w) = d w := by
  intro ⟨g, hg⟩
  have h1 := hg z habove
  have h2 := hg (τ z) habove'
  rw [hρ, hd, h1] at h2
  revert h2
  cases d z <;> intro h2 <;> exact Bool.noConfusion h2

/-- THE APERTURE, ONE BIT WIDE. At a seat where d is odd, one odd bit s at the seat would decide d
    on the seat and its partner through a calibration that exists and is unique. The width of
    the aperture, a theorem; it says nothing of whether anything passes through it. -/
theorem aperture_one_bit_wide (τ : α → α) (s d : α → Bool) (z : α)
    (hs : s (τ z) = !s z) (hd : d (τ z) = !d z) :
    ∃ c : Bool, (d z = xor (s z) c ∧ d (τ z) = xor (s (τ z)) c) ∧
      ∀ c' : Bool, (d z = xor (s z) c' ∧ d (τ z) = xor (s (τ z)) c') → c' = c :=
  ⟨xor (d z) (s z),
    ⟨by cases s z <;> cases d z <;> rfl,
     by rw [hs, hd]; cases s z <;> cases d z <;> rfl⟩,
    by
      intro c' hc
      obtain ⟨h1, -⟩ := hc
      generalize hsz : s z = sv
      generalize hdz : d z = dv
      rw [hsz, hdz] at h1
      cases sv <;> cases dv <;> cases c' <;> first | rfl | exact absurd h1 (by decide)⟩

/-- EXISTENCE SUPPLIES NOTHING. Conditioning the Unicorn part on an inhabited premise, the root
    axiom for one, leaves it exactly where it was; and no class of frames on which the premise
    is to decide it decides more than already held on the class. Theorems E and I of Part I,
    read at the Unicorn part. -/
theorem unicorn_rule_exact (Z onL : α → Prop) (height : α → Nat) (T : Nat) (A : Prop) (ha : A) :
    (A → UnicornAt Z onL height T) ↔ UnicornAt Z onL height T :=
  ⟨fun h => h ha, fun hu _ => hu⟩

theorem unicorn_no_cure {Frame : Type} (A : Prop) (ha : A) (C U : Frame → Prop) :
    (∀ X, C X → A → U X) ↔ (∀ X, C X → U X) :=
  ⟨fun h X hc => h X hc ha, fun h X hc _ => h X hc⟩

/-- NO BYPASS. Given the Real part, the fused hypothesis is exactly the Unicorn part, so a proof
    of the hypothesis contains a proof of the Unicorn part and the block is not avoided by
    aiming at the whole. -/
theorem unicorn_no_bypass (Z onL : α → Prop) (height : α → Nat) (T : Nat)
    (hr : RealPartAt Z onL height T) : RH Z onL ↔ UnicornAt Z onL height T := by
  constructor
  · intro h z hz _; exact h z hz
  · intro hu z hz
    cases Nat.lt_or_ge T (height z) with
    | inl hgt => exact hu z hz hgt
    | inr hle => exact hr z hz hle

/-- THE SUPPLY SIDE IS SILENCE, AT THE REGISTER. Over one even record, both orientations of d at
    the seat are equally refused to every reading above the height: the register cannot say
    which way the seat lies, so it cannot say what a supply would bring, nor that one exists,
    nor that none does. -/
theorem supply_side_silent (height : α → Nat) (T : Nat) (τ : α → α) (ρ : α → β) (d : α → Bool)
    (z : α) (habove : height z > T) (habove' : height (τ z) > T)
    (hρ : ρ (τ z) = ρ z) (hd : d (τ z) = !d z) :
    (¬ ∃ g : β → Bool, ∀ w, height w > T → g (ρ w) = d w) ∧
    (¬ ∃ g : β → Bool, ∀ w, height w > T → g (ρ w) = !d w) ∧
    (!d z) ≠ d z :=
  ⟨unicorn_block height T τ ρ d z habove habove' hρ hd,
   unicorn_block height T τ ρ (fun w => !d w) z habove habove' hρ (by
     show (!d (τ z)) = !(!d z)
     rw [hd]),
   by cases d z <;> decide⟩

/-- The division at the height, on natural heights: the hypothesis is exactly its Real part and
    its Unicorn part. -/
theorem rh_iff_real_and_unicorn (Z onL : α → Prop) (height : α → Nat) (T : Nat) :
    RH Z onL ↔ RealPartAt Z onL height T ∧ UnicornAt Z onL height T := by
  constructor
  · intro h; exact ⟨fun z hz _ => h z hz, fun z hz _ => h z hz⟩
  · intro ⟨hr, hu⟩ z hz
    cases Nat.lt_or_ge T (height z) with
    | inl hgt => exact hu z hz hgt
    | inr hle => exact hr z hz hle

/-- A certificate at the height: a complete list of the zeros up to T, each checked on the line. -/
structure Certificate (Z onL : α → Prop) (height : α → Nat) (T : Nat) where
  items    : List α
  complete : ∀ z, Z z → height z ≤ T → z ∈ items
  checked  : ∀ z, z ∈ items → onL z

theorem real_part_of_certificate (Z onL : α → Prop) (height : α → Nat) (T : Nat)
    (c : Certificate Z onL height T) : RealPartAt Z onL height T :=
  fun z hz hle => c.checked z (c.complete z hz hle)

/-- THE THREE BITS AS ONE TERM · THE REGISTER'S PROOF OF THE HYPOTHESIS, COMPLETE. Under a
    certificate at T and the seat data above T, the hypothesis divides exactly (bit one), its
    Real part at T holds (bit two), every reading of the even register is refused above T (bit
    three, refused), and one supplied bit calibrates the seat uniquely (the door). The hypothesis
    is not thereby proved; nothing derivable about it is left underived. -/
theorem rh_register_proof_complete (Z onL : α → Prop) (height : α → Nat) (T : Nat)
    (c : Certificate Z onL height T) (τ : α → α) (ρ : α → β) (d s : α → Bool) (z : α)
    (habove : height z > T) (habove' : height (τ z) > T) (hρ : ρ (τ z) = ρ z)
    (hd : d (τ z) = !d z) (hs : s (τ z) = !s z) :
    (RH Z onL ↔ RealPartAt Z onL height T ∧ UnicornAt Z onL height T) ∧
    RealPartAt Z onL height T ∧
    (¬ ∃ g : β → Bool, ∀ w, height w > T → g (ρ w) = d w) ∧
    ∃ k : Bool, (d z = xor (s z) k ∧ d (τ z) = xor (s (τ z)) k) ∧
      ∀ k' : Bool, (d z = xor (s z) k' ∧ d (τ z) = xor (s (τ z)) k') → k' = k :=
  ⟨rh_iff_real_and_unicorn Z onL height T, real_part_of_certificate Z onL height T c,
   unicorn_block height T τ ρ d z habove habove' hρ hd, aperture_one_bit_wide τ s d z hs hd⟩

end RHUnicorn

#print axioms RHUnicorn.unicorn_block
#print axioms RHUnicorn.aperture_one_bit_wide
#print axioms RHUnicorn.unicorn_rule_exact
#print axioms RHUnicorn.unicorn_no_cure
#print axioms RHUnicorn.unicorn_no_bypass
#print axioms RHUnicorn.supply_side_silent
#print axioms RHUnicorn.rh_iff_real_and_unicorn
#print axioms RHUnicorn.real_part_of_certificate
#print axioms RHUnicorn.rh_register_proof_complete
\end{Verbatim}

\endgroup

Kernel audit, verbatim: \texttt{lean\ RH\_Unicorn\_Block.lean}, exit 0.

\begingroup\scriptsize

\begin{Verbatim}[numbers=left,numbersep=6pt,xleftmargin=2.2em,fontsize=\scriptsize,breaklines,breakanywhere,frame=lines,framesep=1.5mm,rulecolor=\color{black!35}]
'RHUnicorn.unicorn_block' does not depend on any axioms
'RHUnicorn.aperture_one_bit_wide' does not depend on any axioms
'RHUnicorn.unicorn_rule_exact' does not depend on any axioms
'RHUnicorn.unicorn_no_cure' does not depend on any axioms
'RHUnicorn.unicorn_no_bypass' does not depend on any axioms
'RHUnicorn.supply_side_silent' does not depend on any axioms
'RHUnicorn.rh_iff_real_and_unicorn' does not depend on any axioms
'RHUnicorn.real_part_of_certificate' does not depend on any axioms
'RHUnicorn.rh_register_proof_complete' does not depend on any axioms
\end{Verbatim}

\endgroup
\end{document}
